         \textbf{Franco-Prussian} 

          \underline{Onset:} 1870         

\underline{Reasons:}
\begin{itemize}
\item Nationalism
\item Power Transition or Security Dilemma
\item Revenge/Retribution
\end{itemize}

         \underline{Sources:} 
         \rowcolors{3}{white}{light-gray} 
         {\scriptsize 
         \begin{longtable}{p{0.2\textwidth}p{0.8\textwidth}} 
         \hiderowcolors% 
         \toprule 
         Source & Excerpt \\ 
         \midrule 
         \showrowcolors% 
         \endfirsthead% 
         \hiderowcolors% 
         \toprule 
         Source & Excerpt \\ 
         \midrule 
         \showrowcolors% 
         \endhead% 
         \midrule 
         \endfoot% 
         \hiderowcolors% 
         \bottomrule 
         \showrowcolors% 
         \endlastfoot% 
         
Gibler, Douglas M. 2018. International Conflicts, 1816–2010: Militarized Interstate Dispute Narratives. Vol. 2 Rowman \& Littlefield. & This dispute describes the Franco-Prussian War fought between July 1870 and February 1871. Prussia had been taking active means to unify the German states, and their efforts had resulted in war elsewhere, especially with Denmark over Schleswig-Holstein and later Austria over the same territory. This caused suspicions in France, which tried to keep a favorable balance of power with respect to central Europe. Various incidents brought the two to war. France neutralized Luxembourg (administered by the Dutch), effecting the withdrawal of Prussian troops from the territory. Prussia had attempted to maximize leverage over the south German states, which were firmly anti-French. Meanwhile, Isabella had been deposed from the Spanish throne in 1868, and Spain, wanting to maintain its monarchy, offered the position to a Prussian prince from the House of Hohenzollern. France forced Spain to withdraw the offer, which infuriated Otto von Bismarck. The Ems Dispatch was released by the Prussians on July 13, 1870, detailing (and exaggerating) a demand made by French diplomat Count Benedetti to Bismarck to never again support a Prussian for the throne in Spain. This ultimately made the French declare war on the Prussians on July 19, and Baden, Bavaria, and Wuerttemberg took arms with the Prussians. The war was a major victory for the Prussians. The allied forces greatly outnumbered and overwhelmed the French forces in battle. Napoleon III commanded the French forces that were routed at Sedan on August 31. Napoleon III surrendered on September 2 and was taken as prisoner of war, ending the Second French Empire and beginning the Third French Republic. However, negotiations between the new republic and the Prussians produced no agreement, and the French continued fighting. Continued embarrassments for the French led to a surrender on January 28, 1871. The embarrassments continued for France. The Prussians had a victory parade in Paris and the Prussian king was proclaimed the German emperor in Versailles. Two treaties end this war. The Treaty of Versailles, signed on February 26, 1871, ceded Alsace and Lorraine to Prussia and relented to an indemnity of five million francs. The Treaty of Frankfurt replaced the Versailles agreement on May 10 and ultimately concluded the war. (pp. 979-980)
\\
Britannica. 2023. “Franco-German War.” Encyclopædia Britannica. Accessed August 19, 2023. https://www.britannica.com/event/Franco-German-War. & Prussia's defeat of Austria in the Seven Weeks' War in 1866 had confirmed Prussian leadership of the German states and threatened France's position as the dominant power in Europe. The immediate cause of the Franco-German War, however, was the candidacy of Prince Leopold of Hohenzollern-Sigmaringen (who was related to the Prussian royal house) for the Spanish throne, which had been left vacant when Queen Isabella II had been deposed in 1868. The Prussian chancellor, Otto von Bismarck, and Spain's de facto leader, Juan Prim, persuaded the reluctant Leopold to accept the Spanish throne in June 1870. This move greatly alarmed France, who felt threatened by a possible combination of Prussia and Spain directed against it.
\\
Wawro, Geoffrey. 2011. “German Unification, Wars of (1864–1871).” In The Encyclopedia of War, edited by Gordon Martel. & (...). The Prussian did secure big indemnities from the South Germans, as well as the right to forge “national connections” with them, a fateful French concession that would permit Bismarck to sign military pacts with the South German states in 1867–1868 that bound them to the Prussian war with France in 1870. (...).
Many in Paris called for a preemptive attack on Prussia, before Berlin had time to absorb the new territories. The usually moderate Adolphe Thiers declared after Königgrätz that “the only way to save France is to declare war on Prussia immediately.” Eugène Rouher agreed: “Smash Prussia and take the Rhine,” he advised the emperor, but Napoleon III hung fire until 1870. By then, the French emperor had absorbed a series of provocations from Bismarck, who privately sought a war with the French as the fastest way to dismantle the North and South German Confederations formed after Nikolsburg and absorb their members into a unitary German Empire. In 1867, Napoleon IIII demanded Luxembourg—a key German fortress on the left bank of the Rhine—as compensation for Prussia's gains in 1866. Bismarck used the French demand to mobilize German opinion and win the signature of military treaties with the South German states. In 1868, Bismarck goaded the French again, this time extending the North German customs union to the south, provoking more saberrattling in Paris. More crises followed, each deliberately manufactured by Bismarck to goad the French into war. He suggested that King Wilhelm I might make himself Kaiser Wilhelm I of Germany to reflect his new responsibilities. “No more violations,” Napoleon III growled. Bismarck then invested in a railway linking the German states and Italy and hinted that the line through Switzerland might have the “strategic purpose” of encircling France. In 1870, he put forward a Catholic Hohenzollern for the vacant Spanish throne. That Spanish throne crisis improbably triggered the Franco Prussian War. When Bismarck rudely rejected Napoleon III's demand that King Wilhelm withdraw his nephew's candidacy forever, Napoleon III, wanting a showdown with Prussia as much as Prussia wanted one with France, declared war in July 1870.
\\
Wright, Oliver. 2025. “Franco-Prussian War (1870–71).” In The Encyclopedia of Diplomacy, edited by Gordon Martel. & The Franco-Prussian War was a conflict fought between France and Prussia in 1870–71. It resulted directly from a dispute between Paris and Berlin over the candidacy of a German prince for the succession to the Spanish throne, although it occurred in a climate of increasing rivalry between the two most powerful states on the European continent. As the third and final of a series of wars, including the Prusso-Danish War of 1864 and the Austro-Prussian War of 1866, it marked the final chapter in Otto von Bismarck's campaign to establish Prussian dominance over Germany. Prussia's quick and unexpected victory enabled Bismarck to unite Germany into a single empire under Prussian leadership in 1871, and tipped the balance of power in Europe very much in Berlin's favor. The Franco-Prussian War was, therefore, the conclusive stage in a sequence of events back to the outbreak of the Crimean War in 1853, which gradually destroyed the Vienna Settlement of 1815. The resentment that it caused in France can be considered a contributing factor to the outbreak of the two world wars of the twentieth century. (...). 
By the time that hostilities opened, however, both sides had been looking for a war for some time. The conflict arose principally from Franco-Prussian rivalry on the Rhine. In Prussia, Bismarck wished to strengthen the position of his country within the German confederation. He hoped also to exploit German nationalism to bolster his position at a time when rapid economic and social changes were putting pressure on the political system (Blackbourn 2003: 135–84).
\\
Morillo, Stephen, Jeremy Black, and Paul Lococo. 2008. War in World History: Society, Technology, and War from Ancient Times to the Present. New York: McGraw-Hill. & But these terms, which included the cession of Alsace and part of Lorraine, as well as the payment of an indemnity, were far harsher than necessary to achieve Bismarck's initial war aims, which largely had to do with incorporating the remaining independent, mostly southern, German states into the Prussian confederation that would after the war become the German Empire. (p. 458)
\\
Parker, Geoffrey, ed. 2020. The Cambridge History of Warfare. 2nd ed. Cambridge: Cambridge University Press. & Bismarck now sought to consolidate his gains. He felt no great desire to create a united Germany - after all, southern Germany was the bastion of two of his great hates: liberalism and Catholicism - but the French refused to accommodate him. (...) [He] edited the account of a minor confrontation between King Wilhelm and the French resident ambassador, creating a dispatch that made Prussians believe that their king had been insulted and Frenchmen consider that their honor had been impugned. (pp. 247-248)
\\
\end{longtable}}
\clearpage
         \textbf{First Central American} 

          \underline{Onset:} 1876         

\underline{Reasons:}
\begin{itemize}
\item Nationalism
\item Power Transition or Security Dilemma
\item Economic, Long-Run
\end{itemize}

         \underline{Sources:} 
         \rowcolors{3}{white}{light-gray} 
         {\scriptsize 
         \begin{longtable}{p{0.2\textwidth}p{0.8\textwidth}} 
         \hiderowcolors% 
         \toprule 
         Source & Excerpt \\ 
         \midrule 
         \showrowcolors% 
         \endfirsthead% 
         \hiderowcolors% 
         \toprule 
         Source & Excerpt \\ 
         \midrule 
         \showrowcolors% 
         \endhead% 
         \midrule 
         \endfoot% 
         \hiderowcolors% 
         \bottomrule 
         \showrowcolors% 
         \endlastfoot% 
         
Gibler, Douglas M. 2018. International Conflicts, 1816–2010: Militarized Interstate Dispute Narratives. Vol. 1. Rowman \& Littlefield. & Honduras, not an interstate system member at the time, was the focal point of this war between Guatemala and El Salvador. In 1872, the fledgling liberal regimes of Guatemala and El Salvador allied together to overthrow the conservative regime in Honduras and install a liberal head of state. They were successful, prompting all three sides to enter into an alliance to preserve their liberal regimes. The leader of the overthrown conservative regime in Honduras, Jose Maria Medina, returned in 1876 to attempt to reestablish his position in Honduras. Guatemala welcomed him, thinking that the situation in Honduras was not liberal enough. El Salvador changed sides on Guatemala, and attacked Medina and his troops. Guatemala eventually joined Honduras in supporting another contender, Marco Aurelio Soto, as he successfully assumed power in Honduras. This led to a rupture of the alliance and of diplomatic relations. Guatemala responded by declaring war on March 27, 1876, and sending troops across the border to fight in El Salvador. They were successful. El Salvador sued for peace on April 25. Guatemala replaced Salvadoran president Andres del Valle with Rafael Zaldivar, a crony of Guatemalan president Justo Rufino Barrios. With Guatemala having established a regime to its liking in El Salvador, both states concluded the conflict with a friendship treaty. (pp. 78-79)
\\
Bancroft, Hubert Howe. 1887. History of Central America. 1801–1887. Volume 8 of The Works of Hubert Howe Bancroft. San Francisco: The History Company. & Barrios contended that, though Valle was president, Gonzalez was the real power in Salvador, whom he accused in a public manifestor of hypocrisy and treachery. Angry words continued, the two nations being now armed for the conflict, till they agreed to disband their forces. [...] The probability is, that, distrusting one another, they merely pretended compliance, keeping their troops ready for action. Barrios sent 1,500 men into Honduras, and came himself with a force to threaten Salvador on the west, and actually invaded the latter without previous declaration of war.
\\
Holden, Robert H. 2004. Armies Without Nations: Public Violence and State Formation in Central America, 1821–1960. Oxford and New York: Oxford University Press. & No sooner had the Remincheros [revolt] been overcome than the governments of both Guatemala and El Salvador sent their troops back into Honduras in 1876 to intervene in a civil war. El Salvador's president promptly betrayed Barrios and joined a Honduran faction to overthrow Barrios himself. Now allying himself with a different Honduran faction, Barrios invaded El Salvador. As combat engulfed the whole of El Salvador and part of Honduras, El Salvador surrendered to Barrios on 25 April 1876. The Guatemalan seized the opportunity to install his own candidates as the presidents of both countries: Rafael Zandivar in El Salvador and Marco Aurelio Soto (...) in Honduras. (pp.53-54)
\\
Karnes, Thomas L. 1961. The Failure of Union: Central America, 1824–1960. Chapel Hill, NC: University of North Carolina Press. & [Barrios] distributed a circular to the other four states proposing a conference in 1876 to be held at Guatemala City for the purpose of planning a federal union. (…). While seeking friendly cooperation from the states, Barrios concluded that he could not get it from the incumbent officials next door. So his forces invaded El Salvador, in the usual manner of aiding the party out of power, and succeeded in installing a more cooperative president, while Guatemalan troops were still present. (p. 153)
\\
\end{longtable}}
\clearpage
         \textbf{Second Russo-Turkish} 

          \underline{Onset:} 1877         

\underline{Reasons:}
\begin{itemize}
\item Nationalism
\item Power Transition or Security Dilemma
\item Economic, Long-Run
\item Revenge/Retribution
\end{itemize}

         \underline{Sources:} 
         \rowcolors{3}{white}{light-gray} 
         {\scriptsize 
         \begin{longtable}{p{0.2\textwidth}p{0.8\textwidth}} 
         \hiderowcolors% 
         \toprule 
         Source & Excerpt \\ 
         \midrule 
         \showrowcolors% 
         \endfirsthead% 
         \hiderowcolors% 
         \toprule 
         Source & Excerpt \\ 
         \midrule 
         \showrowcolors% 
         \endhead% 
         \midrule 
         \endfoot% 
         \hiderowcolors% 
         \bottomrule 
         \showrowcolors% 
         \endlastfoot% 
         
Gibler, Douglas M. 2018. International Conflicts, 1816–2010: Militarized Interstate Dispute Narratives. Vol. 2 Rowman \& Littlefield. & The Russo-Ottoman War of 1877–1878 followed nationalist sentiments in the Balkans that encouraged the Russians to make a bold move in order to acquire territory lost during the Crimean War (...). The time was ripe for Russian advances on the Ottomans. The Bulgarian uprising in April 1876 had been aggressively put down by the Ottomans, who used irregular troops in the action, which led to rape, arson, and the massacre of 12,000 Bulgarians. International outrage forced Britain to withdraw its support for the Ottomans. Thus, after reaching an agreement with Romania, nominally under Ottoman control, regarding the moving of Russian troops through the territory, Russia issued an ultimatum and partially mobilized its forces on October 31, 1876. The ultimatum required the Ottomans to sign a peace agreement with Serbia, with which it had been at war since June 30, 1876. After further securing Austro-Hungary's benevolent neutrality, and agreements (the London Treaty) from the major powers to force changes from the Ottomans, Russia declared war on April 24, 1877, and Romania declared its independence a month later. The war was fought on two fronts—on the Danube and in Transcaucasia—with the Ottoman forces largely checking the Russians. However, the Russians had a breakthrough as they finally overwhelmed the fortress at Plevna. After achieving victory in December 1877, the Russian forces raced toward Constantinople. Under pressure of the Russian forces in the vicinity of the capital, an armistice was quickly signed. The Treaty of San Stefano on March 3, 1878, ended the war. The Treaty of Berlin was signed in July 1878 to modify the terms of that treaty. (pp. 575-576)
\\
Britannica. 2014. “Russo-Turkish Wars.” Accessed August 20, 2023. https://www.britannica.com/topic/Russo-Turkish-wars. & The last Russo-Turkish War (1877–78) was also the most important one. In 1877 Russia and its ally Serbia came to the aid of Bosnia and Herzegovina and Bulgaria in their rebellions against Turkish rule. The Russians attacked through Bulgaria, and after successfully concluding the Siege of Pleven they advanced into Thrace, taking Adrianople (now Edirne, Tur.) in January 1878. In March of that year Russia concluded the Treaty of San Stefano with Turkey. This treaty freed Romania, Serbia, and Montenegro from Turkish rule, gave autonomy to Bosnia and Herzegovina, and created a huge autonomous Bulgaria under Russian protection.
\\
Van Der Oye, David S. 2011. “Russo-Turkish War (1877–1878).” In The Encyclopedia of War, edited by Gordon Martel. & The immediate casus belli was the Ottoman refusal to yield to St. Petersburg's demands for reforms to ease the plight of Turkey's Orthodox Christian subjects in the Balkans. Nascent nationalist sentiments in recent years had made these Slavs increasingly restive, leading to revolts in Herzegovina in 1875 and Bulgaria the following year. Brutal repression of the latter by its Ottoman overlords led to considerable European outrage against the “Bulgarian horrors.” Russians felt these sentiments particularly keenly at a time of growing Pan-Slav solidarity with their co-religionists in southeastern Europe. Indeed, when Serbia and Montenegro launched a short but disastrous war against Turkey in 1876, some five thousand Russian volunteers joined them. Other factors included Russia's long-standing ambitions for the straits separating the Black Sea from the Mediterranean and a burning desire to expunge the humiliation of defeat in 1855 during the Crimean War.
\\
Morillo, Stephen, Jeremy Black, and Paul Lococo. 2008. War in World History: Society, Technology, and War from Ancient Times to the Present. New York: McGraw-Hill. & Russian Pan-Slavism led to the dispatch of many volunteers to help the Serbs when they declared war on the Turks in 1876, only to be defeated at Aleksinac and Djunis. The Russians formally entered the war later in 1876 (…). (…). Peace was dictated at San Stefano outside Constantinople in February 1878, with Bulgaria established as a major state including Macedonia and stretching to the Aegean Sea. (pp. 458-459)
\\
Parker, Geoffrey, ed. 2020. The Cambridge History of Warfare. 2nd ed. Cambridge: Cambridge University Press. & The Ottoman empire continued to decline. In 1876, with the enthusiastic help of local Muslims, the Turks brutally suppressed rebellions in Herzegovina, Bosnia, and Bulgaria; and, when the Serbs came to help their brothers, the Turks administered a thorough thrashing. The following year Russia declared war on the Ottomans and, assisted by Romania, won a series of victories that threatened to destroy Turkish control in the southern Balkans. (...) [The] Treaty of San Stefano (3 March 1878) forced the Turks to recognize the independence of Serbia, Montenegro, and Romania, and to grant Bulgaria a large measure of autonomy. Shortly afterwards, Austria received authorization to occupy Bosnia and Herzegovina, although for thirty years they would remain under Turkish suzerainty. Russia (...) had avenged the disgrace of the Crimean War (...). (p. 258)
\\
Uyar, Mesut, and Edward J. Erickson. 2009. A military history of the Ottomans: From Osman to Atatürk. Bloomsbury Publishing, USA.  & The heavy-handed tactics and techniques of the Ottoman army [in Bulgaria] (...) worsened the already tarnished image of the empire. (...). The fate of the empire or so-called ‘‘Eastern Question became a topic of everyday talk in parliaments and newspaper columns, thanks to efforts of several influential leaders, of which British politician and Prime Minister William Gladstone was the most famous. (...). [Neither the new sultan nor any of his advisors] did (...) comprehend the drastic changes in the balance of European politics (...). (...). The new administration not only ignored last-minute changes for peace, but also further strengthened Russia's position by paying no attention to Romanian requests for autonomy. Frustrated Romania would become a valuable ally of the Russians. War became inevitable and bad crisis management and publicity isolated the Ottoman Empire further. (...).
[On the Russian side] political pressures, as well as traditional misgivings about the validity of the army's staff planning, took precedence and war was declared on April 24, 1877, one day after the inspection of combat units by Tsar Alexander II in Kishinev. (pp. 183-185)
\\
\end{longtable}}
\clearpage
         \textbf{War of the Pacific} 

          \underline{Onset:} 1879         

\underline{Reasons:}
\begin{itemize}
\item Border Clashes
\item Economic, Long-Run
\end{itemize}

         \underline{Sources:} 
         \rowcolors{3}{white}{light-gray} 
         {\scriptsize 
         \begin{longtable}{p{0.2\textwidth}p{0.8\textwidth}} 
         \hiderowcolors% 
         \toprule 
         Source & Excerpt \\ 
         \midrule 
         \showrowcolors% 
         \endfirsthead% 
         \hiderowcolors% 
         \toprule 
         Source & Excerpt \\ 
         \midrule 
         \showrowcolors% 
         \endhead% 
         \midrule 
         \endfoot% 
         \hiderowcolors% 
         \bottomrule 
         \showrowcolors% 
         \endlastfoot% 
         
Gibler, Douglas M. 2018. International Conflicts, 1816–2010: Militarized Interstate Dispute Narratives. Vol. 1 Rowman \& Littlefield. & The War of the Pacific pitted Chile against the allied forces of Peru and Bolivia in a conflict that had profound consequences for all states involved. The borders of Chile, Peru, and Bolivia converged in an area that was minimally populated but rich in nitrates that had become a significant source of trade with European powers. Chile had already been mining the area for its natural resources by commission from Bolivia, which had territorial control over the region. However, these commissions soon became the pretext for making territorial claims on Bolivia regarding the area. Chile wanted the Antofagasta and Atacama regions as their own. Their interests were not just limited to the Bolivian provinces—Chilean interest went as far north as Tacna and Arica, then Peruvian territory. Chile struck Bolivia first on February 14, 1879. Peru went on alert and tried to mediate, but Chile's interest in Peruvian territory ultimately made Peru a target as well. Chile declared war on Peru on April 5. Chile emerged victorious over the course of four years of war and significantly altered the territorial composition of all three states involved in the conflict. The Treaty of Ancon, signed on October 20, 1883, imposed the terms of settlement on Peru. Chile acquired Tacna and Arica, holding them for over 40 years until mediation from Herbert Hoover returned Tacna to Peru. Chile kept Arica. Chile's imposed terms on Bolivia ultimately gave Chile all of Bolivia's coast, landlocking Bolivia ever since. (p. 176)
\\
Britannica. 2023. “War of the Pacific.” Accessed August 20, 2023. https://www.britannica.com/event/War-of-the-Pacific. & It grew out of a dispute between Chile and Bolivia over control of a part of the Atacama Desert that lies between the 23rd and 26th parallels on the Pacific coast of South America. The territory contained valuable mineral resources, particularly sodium nitrate.
\\
Farcau, Bruce. 2011. “Pacific, War of the (1879–1883).” In The Encyclopedia of War, edited by Gordon Martel. & Like most wars in the former colonial world, the genesis of the dispute was in a casually drawn border between what had been two provinces of the Spanish Empire, casual because it was placed in the midst of what appeared to be totally valueless land and because there would always be a higher authority to which to appeal should a disagreement ever arise. With the independence of both Bolivia and Chile, however, that higher authority had ceased to exist, and, more importantly, the Atacama Desert, possibly the most desolate piece of real estate on the planet, suddenly became a treasure house in the late nineteenth century as it happened to be rich in nitrates, now a valuable ingredient in the manufacture of both fertilizers and explosives. (...). 
[When] nitrates suddenly became a valuable commodity in the second half of the century, Chile began gradually encroaching on the territory through a series of “reinterpretations” of where the border lay, to which Bolivia meekly succumbed. On February 23, 1878, however, the Bolivian government finally took a stand. Always strapped for cash, it then attempted to impose a modest tax of ten centavos per hundredweight of nitrates exported from the area to finance modest maintenance expenses for the port, but the levy violated an agreement with Chilean entrepreneurs in Antofagasta. The Chilean government hit upon this as a pretext for invasion.
\\
Scheina, Robert L. 1987. Latin America: A Naval History, 1810–1987. Annapolis, MD: Naval Institute Press. & The War of the Pacific (1879-83) was a struggle to control the guano- and nitrate-producing Atacama Desert. The area and control of ist resources had long been disputed by Chile and Bolivia, the latter with the active support of Perú. Bolivia administered much of the disputed territory, while Chile had certain commercial rights guaranteed by treaty. (...). Matters came to a head when Bolivia, contrary to an agreement with Chile, increased taxes on nitrate exports. (p. 31)
\\
Kiernan, Victo G. 1955. “Foreign Interests in the War of the Pacific.” The Hispanic American Historical Review, 35(1): 14–36. & Exceedingly tortuous in detail, the causes of the war were ex- ceedingly simple in essence. Bolivia had nitrate deposits in her coastal province in the Atacama desert; Peru, her ally since 1873, had guano and nitrates in the Tarapaca province bordering it on the north. Chile, to the south, with few deposits of her own, had invested in the development work in Bolivia and to some extent in Tarapaca. All three countries were hard up, and run by oligarchies which dis- liked paying taxes and looked to revenue from these fertilizers as a substitute. Peru set up a state monopoly, taking over private enter- prises in exchange for certificates. Bolivia put an export tax on the Chilean coinpany at Antofogasta. Chile denounced this as a breach of agreement, and in February, 1879, seized the port. Bolivia de- clared war on Chile, and Peru supported her ally. (p. 1)
\\
\end{longtable}}
\clearpage
         \textbf{Conquest of Egypt} 

          \underline{Onset:} 1882         

\underline{Reasons:}
\begin{itemize}
\item Nationalism
\item Economic, Long-Run
\item Domestic Politics
\end{itemize}

         \underline{Sources:} 
         \rowcolors{3}{white}{light-gray} 
         {\scriptsize 
         \begin{longtable}{p{0.2\textwidth}p{0.8\textwidth}} 
         \hiderowcolors% 
         \toprule 
         Source & Excerpt \\ 
         \midrule 
         \showrowcolors% 
         \endfirsthead% 
         \hiderowcolors% 
         \toprule 
         Source & Excerpt \\ 
         \midrule 
         \showrowcolors% 
         \endhead% 
         \midrule 
         \endfoot% 
         \hiderowcolors% 
         \bottomrule 
         \showrowcolors% 
         \endlastfoot% 
         
Gibler, Douglas M. 2018. International Conflicts, 1816–2010: Militarized Interstate Dispute Narratives. Vol. 2 Rowman \& Littlefield. & This dispute is the Anglo-Egyptian War of 1882, which culminated in the British occupation of Egypt. Egypt was still in its infancy as a state independent of French colonial rule and Ottoman control and detested what it deemed to be foreign meddling in affairs, including the operation of the Suez Canal. Meanwhile, the British took interest in the canal. It provided a clear link to India, its prized colonial possession. When news of Europeans being targeted and massacred in Egyptian riots reached the governments of Paris and London, both states decided to send warships. Only the British opted for war. The war started with bombardments of Alexandria, the place where the Egyptians were constructing garrisons to defend themselves against European forces. The Egyptians were no match for the British in either numbers or expertise, and the British quickly pushed the Egyptians southeast. The Egyptian army was eventually annihilated on September 13, and the British assumed control of Egypt on September 15. (pp. 522-523)
\\
Hopkins, Antony. G. 1882. The Victorians and Africa: A Reconsideration of the Occupation of Egypt, 1882. The Journal of African History. & Britain, like France, had long-standing and extensive commitments in Egypt. But, whereas French connections covered a wide span of interests, Britain's were more narrowly focused on the economy. [...] In 1880 Britain took 80 per cent of Egypt's exports and supplied 44 per cent of her imports; and a notable feature of the composition of trade was the shipment of raw cotton and the sale of Manchester cotton goods. Commercial expansion was accompanied by railway and harbour construction, and by the installation of industrial machinery - all of which gave employment to British manufacturers and personnel. Export development was financed largely by external borrowing, and in this activity, too, British investors were dominant.108 The volume and distribution of overseas investment are subjects which have to be approached with caution. Very little is known about private investment except that it was substantial and substantially British. More than half the public debt was owed to British creditors in I873, before the rash of short-term loans raised the proportion held by French investors; even so, the funded debt remained largely in British hands. [...] To accept these objections, however, is not to deny the possibility that economic considerations played a part, perhaps even a crucial part, in Britain's decision to invade Egypt [...] By then Egypt had served its purpose: Britain's economic interests had been upheld; the Liberal party had been united; and the Conservatives had been rendered at least temporarily speechles
\\
Stapleton, T.J. (2011). Mahdist War (1881–1899). In The Encyclopedia of War, G. Martel (Ed.). & (...). The first major victory for the Mahdi occurred in 1882 when his supporters captured the Egyptian military base of El-Obeid west of the Nile. Although Mahdist forces initially tried to storm the fort and were cut down by breach-loading rifles, they eventually undertook a successful siege. As a result of this lesson, the Mahdi formed part of his army into a near-professional infantry corps equipped with firearms. The British occupied Egypt in 1882 to secure the Suez Canal as a vital choke point for shipping to and from India.
\\
Parker, Geoffrey, ed. 2020. The Cambridge History of Warfare. 2nd ed. Cambridge: Cambridge University Press. & As the greatest colonial power in the world, Britain took part in the largest number of conflicts with non-Europeans. The construction of the Suez Canal gave Egypt a central position in the British empire because of the importance of the lines of communication to India. In September 1882 anti-Western riots in Alexandria resulted in British intervention, and an army under Sir Garnet Wolseley launched a surprise night attack that crushed the Egyptian army at Tel-el-Kebir. An immediate pursuit finished off the war and placed Egypt under British domination for the next seventy-four years. (p. 254)
\\
Galbraith, John S., and Afaf Lutfi al-Sayyid-Marsot. 1978. “The British Occupation of Egypt: Another View.” International Journal of Middle East Studies, 9(4): 471–488. & After January 1882, when French policy changed drastically, the British line against Egyptian nationalism remained inflexible. Once committed to a har line against Urabi, the cabinet committed its prestige to removing him without regard to whether it was in the British national interest to do so. To retreat before Urabi would mean a loss of prestige, it was said, and the West could not lose face and maintain its mastery over the 'Oriental races'. (...) The 'canal in danger' (...) was an issue that united Britain behind the intervention that led to the occupation of Egypt. The Liberal party indeed had demostrated that it could be as aggressive as its Conservative opponents. (pp. 487-488)
\\
\end{longtable}}
\clearpage
         \textbf{Sino-French} 

          \underline{Onset:} 1884         

\underline{Reasons:}
\begin{itemize}
\item Nationalism
\item Power Transition or Security Dilemma
\item Economic, Long-Run
\end{itemize}

         \underline{Sources:} 
         \rowcolors{3}{white}{light-gray} 
         {\scriptsize 
         \begin{longtable}{p{0.2\textwidth}p{0.8\textwidth}} 
         \hiderowcolors% 
         \toprule 
         Source & Excerpt \\ 
         \midrule 
         \showrowcolors% 
         \endfirsthead% 
         \hiderowcolors% 
         \toprule 
         Source & Excerpt \\ 
         \midrule 
         \showrowcolors% 
         \endhead% 
         \midrule 
         \endfoot% 
         \hiderowcolors% 
         \bottomrule 
         \showrowcolors% 
         \endlastfoot% 
         
Gibler, Douglas M. 2018. International Conflicts, 1816–2010: Militarized Interstate Dispute Narratives. Vol. 2 Rowman \& Littlefield. & The Sino-French War of 1883 paved the way for the establishment of French Indochina. France desired to expand its empire into Southeast Asia, targeting the rest of present-day Vietnam to add to what it already had in central Vietnam, which posed a problem for China. Though Vietnam was nominally independent, it was heavily influenced by China. China's preeminence had been slipping through the 19th century with the decline of the Qing Dynasty, allowing European powers and Japan to make moves on areas where China had previously been preeminent (for example, in Korea). (...). The fighting assumed an interstate dimension when the Chinese military took arms along with the [Vietnamese] Black Flags, aware that the French would eventually overwhelm them if the Chinese did not act. Further, China acted under the faulty assumption that France would not commit to the war and that the joint forces with the Vietnamese would be sufficient. (...). China's optimism was not entirely borne out, though their contribution was able to change the course of what would have been certain disaster for the Vietnamese. Nowhere was this more pronounced than the Battle at Bang Bo, prompting French retreat. This was an embarrassment for the French, but the Chinese victory here was pyrrhic. They had to exhaust available personnel to drive back the French, limiting themselves should a conflict against expanding Japan be on the horizon. France sued for peace following battles in March 1885, and China welcomed negotiations. (p.731-732)
\\
Britannica. 2023. “Sino-French War.” Accessed August 20, 2023. https://www.britannica.com/event/Sino-French-War. & Sino-French War, conflict between China and France in 1883–85 over Vietnam, which disclosed the inadequacy of China's modernization efforts and aroused nationalistic sentiment in southern China. The French had already begun to encroach on Vietnam, China's major protectorate in the south, and by 1880 France controlled the three southern provinces, known as Cochinchina. In the 1880s the French began to expand northward in Vietnam, stationing troops in Hanoi and Haiphong. The Chinese responded by building up their forces in the area and engaging the French in a series of limited battles.
\\
Vandervort, Bruce. 2011. “French Foreign Legion.” In The Encyclopedia of War, edited by Gordon Martel. & The [French Foreign Legion] sent its first battalion to Indochina in 1882, to take part in the conquest of Tonkin or northern Vietnam, inaugurating a continual presence in that French colony until the defeat at Dien Bien Phu in 1954. French efforts to take over Tonkin led to war with China, which had traditionally viewed Vietnam as a Chinese tributary state, in 1884–1885. The Sino-French War saw French marines and legionnaires landing in Taiwan, as well as fighting Chinese troops and irregulars along the North Vietnamese border.
\\
Morillo, Stephen, Jeremy Black, and Paul Lococo. 2008. War in World History: Society, Technology, and War from Ancient Times to the Present. New York: McGraw-Hill. & But the final conquest of the rest of Vietnam in 1884–85 proved to be no walkover for the French. (...).  According to the terms of the resulting treaty, China accepted French control of Vietnam. Thus, in Vietnam, anti-Western reactions undermined early moves toward modernization, and Vietnam joined Cambodia and Laos as French Indochina.
\\
\end{longtable}}
\clearpage
         \textbf{Second Central American} 

          \underline{Onset:} 1885         

\underline{Reasons:}
\begin{itemize}
\item Nationalism
\item Power Transition or Security Dilemma
\item Religion or Ideology
\end{itemize}

         \underline{Sources:} 
         \rowcolors{3}{white}{light-gray} 
         {\scriptsize 
         \begin{longtable}{p{0.2\textwidth}p{0.8\textwidth}} 
         \hiderowcolors% 
         \toprule 
         Source & Excerpt \\ 
         \midrule 
         \showrowcolors% 
         \endfirsthead% 
         \hiderowcolors% 
         \toprule 
         Source & Excerpt \\ 
         \midrule 
         \showrowcolors% 
         \endhead% 
         \midrule 
         \endfoot% 
         \hiderowcolors% 
         \bottomrule 
         \showrowcolors% 
         \endlastfoot% 
         
Gibler, Douglas M. 2018. International Conflicts, 1816–2010: Militarized Interstate Dispute Narratives. Vol. 1 Rowman \& Littlefield. & Guatemala was one of the stronger Central American states by the time of the second Central American War, and its leader, Justo Rufino Barrios, was a strong proponent for unification of Central America. After a European trip, he saw the unification successes of Italy and Germany as signs to move forward with unification once more. On February 28, 1885, Barrios declared the Union of Central America and offered himself to lead this new union. Honduras ratified the decree on March 7 of that same year. However, El Salvador, which was a participant at the previous conference for union (convened on September 15, 1884) was reticent to move forward with ratifications. Public opinion was against the idea, and the Salvadoran Congress eventually killed ratification for El Salvador. Salvadoran head of state Rafael Zaldivar became vocal in opposing Barrios. In response, Barrios threatened to attack El Salvador and then Mexico. Honduras, which supported Barrios, would neutralize Nicaragua and Costa Rica, as both were vocal in opposing Barrios as well. Barrios followed through with his threats and marched into conflict with El Salvador, winning the first battle on March 30. Three days later, he was killed in battle. Guatemala's troops could not press on in his absence and the conflict ended abruptly. Barrios's death prevented a broader Central American conflict from forming, as it seemed likely to embroil the entire region in war, Mexico included. The dream for a unified Central America died with Barrios as well since the idea never gained serious traction ever again. (p. 79)
\\
Palmer, Steven. 1993. “Central American Union or Guatemalan Republic? The National Question in Liberal Guatemala, 1871–1885.” The Americas. & The 1885 war for Central American unification, and Guatemalan irredentism in general, has too easily been painted as simply the manifestation of a rather base imperialism.8 Elements of such an imperialist dream wer surely present in this last significant coercive attempt at creating a Centra American nation. Yet this explanation does not offer much insight into the persistent and central problems that always blocked a coherent nationa Guatemalan vision that could sit easily with the society inside its borders. Liberal intellectual life during the Barrios regime was influenced by a wide variety of Unionist currents that can scarcely be reduced to mere ulterior ambition, as Ralph Lee Woodward has characterized the bulk of attempts at creating a Central American nation-state. Indeed, Woodward's nation divided thesis is never really coherently defined or demonstrated, and appears to rest largely on a general and trans-historical logic of union: a rational perception of the material and geo-political benefits that might accrue to a united isthmian country.9 The Unionist culture of Guatemalan liberals during the Reform period, while willing to embrace such logic, was driven by more complex historical force [...] aftermath of Central Ame almost exclusively as a result of the liberal elements of the isthmian Republic to supercede [...]
\\
Holden, Robert H. 2004. Armies Without Nations: Public Violence and State Formation in Central America, 1821–1960. Oxford: Oxford University Press. & Despite the cautions against the use of force voice by the governments of the United States, Mexico, and Spain, and the even more significant obstacle posed by a defensive alliance of three of the five governments (Costa Rica, Nicaragua, and El Salvador) against their unification, General Barrios organized an army of 14,500 men and invaded El Salvador on 31 March 1885. Three days later, he took a Salvadoran bullet in the heart, inspiring his army to retreat to Guatemala. The grand campaign to unify the isthmus collapsed, tragically but ludicrously, with the death of the general. (...). Barrios may have been the born leader of a rising class of rural, mestizo, medium-sized entrepreneurs in whose interests liberal ideology was devised. But to succeed in Guatemala, Barrios would above all have to emulate Carrera, showing that he too was prepared to kill ruthlessly. (p. 53)
\\
Karnes, Thomas L. 1961. The Failure of Union: Central America, 1824–1960. Chapel Hill, NC: University of North Carolina Press. & On Febuary 28, 1885 (...) at the Teatro Nacional in Guatemala City (...) an officer interrupted the performance (...) and read to the (...) audience a proclamation written by [Barrio]. The paper began with a lenghty argument for Barrios' assuming supreme military command of all of Central America, but doing this as the tool of the people. (...). Doubtless there were enemies, but they would be subdued. Germany and Italy had united by force of arms, and Central America could follow their steps.
(...).
[Barrio] may have wanted to extend to all of Central America the progress which he had brought to Guatemala, but it was to be by the use of force, and this would have meant that the evils, too, would have been exported. It has been said that his real purpose was to acquire Nicaragua before a canal treaty could be concluded so that he might gain the expected revenues. This theory would help to account for his apparent haste, since such a treaty was then pending in the Senate of the United States. (pp. 158-163)
\\
\end{longtable}}
\clearpage
         \textbf{First Sino-Japanese} 

          \underline{Onset:} 1894         

\underline{Reasons:}
\begin{itemize}
\item Power Transition or Security Dilemma
\item Economic, Long-Run
\end{itemize}

         \underline{Sources:} 
         \rowcolors{3}{white}{light-gray} 
         {\scriptsize 
         \begin{longtable}{p{0.2\textwidth}p{0.8\textwidth}} 
         \hiderowcolors% 
         \toprule 
         Source & Excerpt \\ 
         \midrule 
         \showrowcolors% 
         \endfirsthead% 
         \hiderowcolors% 
         \toprule 
         Source & Excerpt \\ 
         \midrule 
         \showrowcolors% 
         \endhead% 
         \midrule 
         \endfoot% 
         \hiderowcolors% 
         \bottomrule 
         \showrowcolors% 
         \endlastfoot% 
         
Gibler, Douglas M. 2018. International Conflicts, 1816–2010: Militarized Interstate Dispute Narratives. Vol. 2 Rowman \& Littlefield. & Korea had remained successfully isolated from the rest of the international system up until the middle of the 19th century. China had kept Korea as a dependent but was struggling to keep its presence in Korea strong and deter outside meddling. Interest in Korea came from across Europe and also Japan. Japan had successfully opened up Korea to the international system following a militarized incident in 1875. Japan's growing presence and China's desire to maintain its position in Korea led to the Tientsin Convention in 1885 in which both China and Japan agreed to notify the other if their troops were going to be sent to Korea. Nine years later, the Donghak Peasant Uprising (1894) against the Korean government led to Korean pleas from China for assistance. China sent 2,000 soldiers and various elements of its fleet, and Japan responded under the auspices of the Tientsin Convention. The uncertain atmosphere and tension resulted in the First Sino-Japanese War when Japan sank a Chinese naval transport on July 25, 1894. The war occurred in a period when Japan was coming into its own as a military power while China was certainly on the decline for many years before this war. The Chinese lasted until November 21, 1894. Port Arthur fell victim to a Japanese siege, and the Chinese sued for peace, resulting in a treaty at Shimonoseki on April 17, 1895. The terms were severe. Japan acquired Formosa, an indemnity, and most-favored-nation status in China. China also formally renounced all claims to suzerainty over Korea, recognizing Korea's total independence. Growing Japanese presence in China soon led to intrigue from Russia, a pretext for the Russo-Japanese War 10 years later. (pp. 806-807)
\\
Britannica. 2023. “First Sino-Japanese War.” Encyclopædia Britannica. Accessed August 19, 2023. https://www.britannica.com/event/First-Sino-Japanese-War-1894-1895. & The war grew out of conflict between the two countries for supremacy in Korea. Korea had long been China's most important client state, but its strategic location opposite the Japanese islands and its natural resources of coal and iron attracted Japan's interest. In 1875 Japan, which had begun to adopt Western technology, forced Korea to open itself to foreign, especially Japanese, trade and to declare itself independent from China in its foreign relations. [...] In 1894, however, Japan, flushed with national pride in the wake of its successful modernization program and its growing influence upon young Koreans, was not so ready to compromise.
\\
Mackenzie, Simon P. 2011. “Sino-Japanese War (1894–1895).” In The Encyclopedia of War, edited by Gordon Martel. & A war fought between Meiji Japan and Manchu China over influence in Korea. Often better armed and always better organized and led than their opponents, Japanese forces on land and sea were almost always victorious. The results highlighted to the world the extent to which the two Asian states involved had either succeeded, in the case of Japan, or failed, in the case of China, to transform themselves along western lines sufficiently to fight and win an industrial-era war. (...).
Korea had long been a vassal state of imperial China, but by the 1880s a rapidly modernizing Japan was sponsoring those in Korea who challenged the traditional governmental forms that were still supported by the Manchus. (...) [The] murder of a pro-Japanese Korean leader (...) heightened (...) internal upheaval in Korea (...). Beijing had agreed to (...) send troops to help suppress a religious insurrection (...) which Tokyo [interpreted] as an effort to reassert the vassal status of Korea at the expense of Japanese interests.
\\
Morillo, Stephen, Jeremy Black, and Paul Lococo. 2008. War in World History: Society, Technology, and War from Ancient Times to the Present. New York: McGraw-Hill. & Fought for control of Korea, the Sino-Japanese War of 1894–95 reflects to China's disadvantage the relative success of two military modernization programs. Japanese forces swiftly forced the Chinese out of Korea and Manchuria and then, after securing a foothold on the Shandong Peninsula, prepared a two-pronged assault on the Qing capital of Beijing, whereupon the Qing court sued for peace.(...). Japan's favorable terms and territorial gains agreed to in the peace treaty were, however, largely reversed under pressure from Western imperial powers, which moved both to prop up the now-collapsing Qing dynasty and, even more, with Qing weakness suddenly exposed, to carve up the Chinese empire for themselves. (pp. 474-478)
\\
Parker, Geoffrey, ed. 2020. The Cambridge History of Warfare. 2nd ed. Cambridge: Cambridge University Press. & By 1894, the country had modernized so effectively that it was able to deploy its new-style armed forces – the navy trained by the British, the army by the French and Germans – to challenge China for control over the Korean peninsula. After six months of defeats, in April 1895 China ceded to Japan the island of Taiwan and the Liaodong peninsula, and agreed to pay a massive war indemnity, but diplomatic pressure by Russia, supported by France and Germany, induced Japan to sell Liaodong back to China. (p. 262)
\\
\end{longtable}}
\clearpage
         \textbf{Greco-Turkish} 

          \underline{Onset:} 1897         

\underline{Reasons:}
\begin{itemize}
\item Nationalism
\item Religion or Ideology
\item Border Clashes
\end{itemize}

         \underline{Sources:} 
         \rowcolors{3}{white}{light-gray} 
         {\scriptsize 
         \begin{longtable}{p{0.2\textwidth}p{0.8\textwidth}} 
         \hiderowcolors% 
         \toprule 
         Source & Excerpt \\ 
         \midrule 
         \showrowcolors% 
         \endfirsthead% 
         \hiderowcolors% 
         \toprule 
         Source & Excerpt \\ 
         \midrule 
         \showrowcolors% 
         \endhead% 
         \midrule 
         \endfoot% 
         \hiderowcolors% 
         \bottomrule 
         \showrowcolors% 
         \endlastfoot% 
         
Gibler, Douglas M. 2018. International Conflicts, 1816–2010: Militarized Interstate Dispute Narratives. Vol. 1 Rowman \& Littlefield. & The origins of the Greco-Ottoman War of 1897 lay in Greek irredentism and, specifically, the desire to own Crete outright. Crete was formally under Ottoman rule, but the Turks gave Crete substantial autonomy under the Halepa Pact of October 1878. After a revolt by the Cretans, a Greek boat was dispatched to prevent Ottoman forces from trying to maintain order of the island (December 1896). Great power intervention, which consistently blocked Greece from changing the status quo to preserve the status of the Ottoman Empire and the balance of power in Europe, proved ineffective. After the Turks secured the neutrality of Bulgaria and Serbia, they declared war on Greece in April 1897. The war effort was poorly designed and ultimately ill-fated for Greece. Greek forces in the north of Crete were routed within a month. Intervention from the czar of Russia was able to secure a ceasefire on May 20, 1897. Negotiations started in June, and the treaty of Constantinople followed on December 4. (pp. 355-356)
\\
Britannica. 2016. “Greco-Turkish Wars.” Encyclopædia Britannica. Accessed August 19, 2023. https://www.britannica.com/event/Greco-Turkish-wars. & The first war, also called the Thirty Days' War, took place against a background of growing Greek concern over conditions in Crete, which was under Turkish domination and where relations between the Christians and their Muslim rulers had been deteriorating steadily. The outbreak in 1896 of rebellion on Crete, fomented in part by the secret Greek nationalistic society called Ethniki Etairia, appeared to present Greece with an opportunity to annex the island.
\\
Pizzo, David. 2011. “Greco-Turkish War (1897).” In The Encyclopedia of War, edited by Gordon Martel. & The war began initially over the unresolved issue of the island of Crete, with which many Greeks desired enosis, or union. (...) [It] did result in special, autonomous status for the island under a Greek high commissioner. The Cretan revolt that provoked the conflict and the war itself also profoundly changed the ethnic composition of the island, which until that point had been a point of contention between Greece and the Ottomans (...). (...). The War of Greek Independence had left large portions of what is today Greece out-side of the kingdom's borders, giving rise to passionate, irredentist claims for enosis with Macedonia, Thessaly, and particularly “the Great Island,” Crete. Indeed, many ofthe most prominent irredentist politicians of independent Greece came from Crete,most notably Eleftherios Venizelos. These ambitions were embodied in the “Megali [Great] Idea (...). [The] concept of a Greater Greece was in many ways similar to the idea of Pan-Slavism (...). (...). The Greek population [of Crete] engaged in frequent rebellions (...) in an attempt to unify the island with the Greek Kingdom, which placed severe strain on the already problematic relations between Greece and the Ottoman Empire. 
\\
Shevill, Ferdinand. 1971. The History of the Balkan Peninsula, 428–429. & (...) [The] Christians (...) were disinclined to put an end to agitation until they had achieved thier one absorbing purpose, which was union with the motherland of Greece. (…) [An international guarantee for Crete guaranteed by the treaty of Berlin], among other benefits, established a local assembly with a Christian majority [opening] the prospect of a genuine amelioration of the lot of the harassed Greeks. (...). From a troubled picture of more or less chronic revolt the insurrection of February, 1897, stands out because, comet-wise, it carried in its wake a long trail of consequences. The new rebellion was caused (...) by continued Moslem misgovernment; its end (...) was union with the Greek Kingdom. 
\\
Uyar, Mesut, and Edward J. Erickson. 2009. A Military History of the Ottomans: From Osman to Atatürk. Bloomsbury Publishing USA. & The combat actions lasted barely a month. Only 10 Ottoman divisions, reinforced with partial mobilization, took part, and overall casualty figures were low. But it was large enough for an evaluation of the extent of Hamidian reforms. 
The Ottoman administration tried its best to stay away from war. However, the
over-confident Greek leadership saw the situation for annexing Crete and even expanding on the mainland further north as ripe for exploitation. This was partly due to miscalculation of the Great Powers' policy and an exaggerated view of the internal problems of the Ottomans, especially regarding the recent Armenian rebellions. (p.208)
\\
\end{longtable}}
\clearpage
         \textbf{Spanish-American} 

          \underline{Onset:} 1898         

\underline{Reasons:}
\begin{itemize}
\item Nationalism
\item Power Transition or Security Dilemma
\item Economic, Long-Run
\item Domestic Politics
\end{itemize}

         \underline{Sources:} 
         \rowcolors{3}{white}{light-gray} 
         {\scriptsize 
         \begin{longtable}{p{0.2\textwidth}p{0.8\textwidth}} 
         \hiderowcolors% 
         \toprule 
         Source & Excerpt \\ 
         \midrule 
         \showrowcolors% 
         \endfirsthead% 
         \hiderowcolors% 
         \toprule 
         Source & Excerpt \\ 
         \midrule 
         \showrowcolors% 
         \endhead% 
         \midrule 
         \endfoot% 
         \hiderowcolors% 
         \bottomrule 
         \showrowcolors% 
         \endlastfoot% 
         
Gibler, Douglas M. 2018. International Conflicts, 1816–2010: Militarized Interstate Dispute Narratives. Vol. 1 Rowman \& Littlefield. & The United States had long sought control over Cuba, the biggest reminder of European defiance of the Monroe Doctrine in its own backyard. Before the American Civil War, the interest in Cuba was particularly strong in the American South and, thus, not as strong of an issue among the northern sentiments that strongly suspected southern interests. As generations passed after the American Civil War, interest in Cuba grew precisely when Spain's continued presence in Cuba had defied the autonomy the Cubans were promised. The casus belli for this war was infamously instigated by newspaper publishers in the United States following the sinking of the USS Maine in Havana on February 15, 1898. The United States quickly blamed Spain (rightly or wrongly), and American newspapers fanned the flames of war. US President McKinley petitioned Congress for an invasion of Cuba; Congress responded by authorizing the recognition of Cuban independence on April 19. The United States dispatched its naval fleet to Cuba shortly thereafter, and both sides declared war on each other just days later. The Americans made short work of the Spaniards in a conflict where Americans simultaneously attacked Spanish positions in Cuba and the Philippines. The Spanish signed an armistice on August 12, 1898, that ended the war. Manila fell two days later. A final peace accord was reached in Paris on December 10. The United States gained Spain's colony in the Philippines, though their presence was short-lived. Cuba was made independent, but the agreement provided for considerable US leverage in Cuban affairs. Puerto Rico and Guam also became American territories. (p. 205)
\\
Britannica. 2023. “Spanish-American War.” Accessed August 20, 2023. https://www.britannica.com/event/Spanish-American-War. & The war originated in the Cuban struggle for independence from Spain, which began in February 1895. The Cuban conflict was injurious to U.S. investments in the island, which were estimated at USD50 million, and almost ended U.S. trade with Cuban ports, normally valued at USD100 million annually.
\\
Watson, Samuel. 2011. “Spanish–American War (1898).” In The Encyclopedia of War, edited by Gordon Martel. & The Spanish–American War was rooted in the growth of American power and the proximity of Spanish colonial possessions to the United States. Historians have assigned a variety of reasons for the US declaration of war against Spain, from humanitarian altruism to popular jingoism,a flight from the domestic divisions of the 1890s, or the influence of sugar planters, to a vision of national power that demanded stable pro-American regimes in nearby nations and along major trade routes. The answer is a mix of all these factors, but the decision to go to war, the United States' first international war in more than half a century, certainly signaled a new national assertiveness. (...) [Cuba was] one of the nation's principal international security concerns since 1810s due to its position at the entrance of the Gulf of Mexico, the gate to New Orleans and the Mississippi Valley. (...) US intervention [after the Cuban insurrection of 1895] may have been precipitated by the imminence of an insurgent victory, for the majority of the rebels were in fact revolutionaries, dedicated to refashioning the plantation system into a smallholding economy that would support the landless majority of Afro-Cubans. (...) Whether from narrow economic interests, capitalist economic ideology, racial fears (there was a wave of lynching, voter disenfranchisement, segregation, and localized ethnic cleansing in the American South during the 1890s), or national security concerns, US policy-makers had ample rationale for intervention. (...). [President] McKinley pushed US mediation more strongly [than Cleveland before him]; Spain responded by proposing internal self-government for Cuba, but both the rebels and local Spanish commanders refused. In January 1898 the US Navy sent the battleship Maine to demonstrate American interest; a month later the ship exploded, probably due to engine problems,with a loss of 253 men. 
\\
Gompert, David C., Hans Binnendijk, and Bonny Lin. 2014. “The American Decision to Go to War with Spain, 1898.” In Blinders, Blunders, and Wars: What America and China Can Learn, 53–62. RAND Corporation. & The premise of this case is that the United States decided to go to war with Spain, not the other way around. It is nearly certain that the USS Maine was not blown up in Havana Harbor by a Spanish naval mine but that it blew up on its own. Regardless, major players in the U.S. government, political elite, and journalism were maneuvering the United States toward hostilities with Spain before the Maine went to the bottom. For that matter, powerful Americans—notably, Theodore Roosevelt, Henry Cabot Lodge, and William Randolph Hearst—were itching for empire and war even prior to choosing Spain as the ideal enemy by virtue of its weakness. (p. 53)
\\
\end{longtable}}
\clearpage
         \textbf{Boxer Rebellion} 

          \underline{Onset:} 1900         

\underline{Reasons:}
\begin{itemize}
\item Religion or Ideology
\item Economic, Short-Run
\end{itemize}

         \underline{Sources:} 
         \rowcolors{3}{white}{light-gray} 
         {\scriptsize 
         \begin{longtable}{p{0.2\textwidth}p{0.8\textwidth}} 
         \hiderowcolors% 
         \toprule 
         Source & Excerpt \\ 
         \midrule 
         \showrowcolors% 
         \endfirsthead% 
         \hiderowcolors% 
         \toprule 
         Source & Excerpt \\ 
         \midrule 
         \showrowcolors% 
         \endhead% 
         \midrule 
         \endfoot% 
         \hiderowcolors% 
         \bottomrule 
         \showrowcolors% 
         \endlastfoot% 
         
Gibler, Douglas M. 2018. International Conflicts, 1816–2010: Militarized Interstate Dispute Narratives. Vol. 2. Rowman \& Littlefield. & The Boxer Rebellion, pitting China against a grand alliance of foreign powers, had multiple causes. Frustrations in China mounted over unequal treaties with foreign powers, the growing presence of foreign nationals, and the influences and ideas that came with the new population. A resistance group, known in English as the “Boxers,” emerged. They engaged in riots and violence against Chinese Christians and British missionaries, who bore the brunt of their rage. A May 30, 1900, riot resulted in the death of two British missionaries in Pao Ting Fu, prompting an ultimatum by Western diplomats. The Chinese had 24 hours to restore peace to the area or Western forces would enter. The Chinese lacked an adequate time to respond. The Boxers rioted throughout Beijing, threatening the Legation Quarter, and the powers ultimately took action. The ensuing war had the effect of quelling the Boxers and restoring order to the besieged Legation Quarter. The issue dragged on (see MID\#2314) until the Qing government in China was compelled to pay a large indemnity for the troubles of the eight-nation alliance. (p. 723)
\\
Britannica. 2023. “Boxer Rebellion.” Encyclopædia Britannica. Accessed August 19, 2023. https://www.britannica.com/event/Boxer-Rebellion. & The Boxer Rebellion targeted foreigners first and foremost, Western missionaries in particular. It also targeted Chinese converts to Christianity, who drew ire for flouting traditional Chinese ceremonies and family relations. [...] In the late 19th century, because of growing economic impoverishment, a series of unfortunate natural calamities, and unbridled foreign aggression in the area, the Boxers began to increase their strength in the provinces of North China. In 1898 conservative, antiforeign forces won control of the Chinese government and persuaded the Boxers to drop their opposition to the Qing dynasty and unite with it in destroying the foreigners. [...] An international force of some 19,000 troops was assembled, most of the soldiers coming from Japan and Russia but many also from Britain, the United States, France, Austria-Hungary, and Italy. On August 14, 1900, that force finally captured Beijing, relieving the foreigners and Christians besieged there since June 20.
\\
Reist, Katherine. 2011. “Boxer Uprising (1899–1901).” In The Encyclopedia of War, edited by Gordon Martel. & The Boxer Uprising has been variously interpreted as a rebellion, an uprising, or an anti-foreign expression of nationalism. (...). [Boxer] groups began to focus on the foreign presence in the area following the Second Opium War (1858–1860) and the first Sino-Japanese War (1894–1895), in both of which the Chinese suffered defeat. The Boxer explanation was that the foreign presence had upset the gods, whose displeasure was evidenced in the droughts and floods that plagued the area. (...).
The Boxers, who referred to themselves as the Righteous and Harmonious Fists, and later the Righteous and Harmonious Militia (therefore invoking a measure of legall status), began to use their slogan Support the Qing. Destroy the foreign.
\\
\end{longtable}}
\clearpage
         \textbf{Sino-Russian} 

          \underline{Onset:} 1900         

\underline{Reasons:}
\begin{itemize}
\item Nationalism
\item Religion or Ideology
\end{itemize}

         \underline{Sources:} 
         \rowcolors{3}{white}{light-gray} 
         {\scriptsize 
         \begin{longtable}{p{0.2\textwidth}p{0.8\textwidth}} 
         \hiderowcolors% 
         \toprule 
         Source & Excerpt \\ 
         \midrule 
         \showrowcolors% 
         \endfirsthead% 
         \hiderowcolors% 
         \toprule 
         Source & Excerpt \\ 
         \midrule 
         \showrowcolors% 
         \endhead% 
         \midrule 
         \endfoot% 
         \hiderowcolors% 
         \bottomrule 
         \showrowcolors% 
         \endlastfoot% 
         
Gibler, Douglas M. 2018. International Conflicts, 1816–2010: Militarized Interstate Dispute Narratives. Vol. 2 Rowman \& Littlefield. & Russia had a variety of interests in China, including the presence of Russian missionaries in China involved in the Boxer Rebellion. Russia also had keen economic interest in China, especially involving the establishment of a Manchurian railway. In 1898, Russia successfully used gunboat diplomacy to force China to accept a railway linking Port Arthur and the Manchurian city of Harbin. In the process, Harbin became a de facto Russian colony in China. The Boxer Rebellion pitted China against an eight-nation alliance that included Russia. Harbin became involved because the Boxer Rebellion eventually consumed the whole of that region. Russia responded with a full-scale occupation of Manchuria, which was complete by November 1900. In spite of a successful Russian occupation of the region, Chinese rebels and bandits destroyed approximately 560 miles of the 800 miles of railway owned by the Russian administered Chinese Eastern Railway. Losses were estimated at approximately USUSD35,783,000 and Russia demanded compensation. The Boxer Protocol, the treaty after the “disturbances” of 1900, addressed the indemnity issue. Russia gained most of the sum to be paid by China and distributed among the victors. That protocol was signed on September 1901, but an earlier protocol between Russia (through the Chinese Eastern Railway) and China was signed in July 1901. This treaty reaffirmed the previous one that gave control of Harbin to Russia and specifically addressed issues of punishment for bandits. (p. 753)
\\
Britannica. 2023. “Boxer Rebellion.” Accessed August 19, 2023. https://www.britannica.com/event/Boxer-Rebellion. & The Boxer Rebellion targeted foreigners first and foremost, Western missionaries in particular. It also targeted Chinese converts to Christianity, who drew ire for flouting traditional Chinese ceremonies and family relations.
\\
Glebov, Sergey. 2023. “11 Blagoveshchensk Massacre and Beyond: The Landscape of Violence in the Amur Province in the Spring and Summer of 1900.” In Russia's North Pacific, 211. Heidelberg University Publishing. & Mass violence erupted along the Amur River in the summer of 1900, an
event that occurred in the context of the Boxer Rebellion in China and
the subsequent invasion of Manchuria by the Russian imperial army.
\\
\end{longtable}}
\clearpage
         \textbf{Russo-Japanese} 

          \underline{Onset:} 1904         

\underline{Reasons:}
\begin{itemize}
\item Power Transition or Security Dilemma
\item Economic, Long-Run
\end{itemize}

         \underline{Sources:} 
         \rowcolors{3}{white}{light-gray} 
         {\scriptsize 
         \begin{longtable}{p{0.2\textwidth}p{0.8\textwidth}} 
         \hiderowcolors% 
         \toprule 
         Source & Excerpt \\ 
         \midrule 
         \showrowcolors% 
         \endfirsthead% 
         \hiderowcolors% 
         \toprule 
         Source & Excerpt \\ 
         \midrule 
         \showrowcolors% 
         \endhead% 
         \midrule 
         \endfoot% 
         \hiderowcolors% 
         \bottomrule 
         \showrowcolors% 
         \endlastfoot% 
         
Gibler, Douglas M. 2018. International Conflicts, 1816–2010: Militarized Interstate Dispute Narratives. Vol. 2 Rowman \& Littlefield. & The Russo-Japanese War occurs largely as a result of the total vulnerability of China at the end of the 19th century and the desire for Russia and Japan to fulfill their imperial ambitions as a result (...). Russia focused on Manchuria while Japan was monopolizing its influence in Korea. During the Boxer Rebellion and immediate fallout, Russia began to expand its influence into Korea, much to the distress of Japan. Throughout 1903, Russia and Japan tried to negotiate an end to the issue, but to no avail. Russia was ultimately served a demand by Japan to evacuate Manchuria. Russia refused. Japan responded by severing ties with Russia on February 6, 1904. Four days later, Japan declared war. The war proved disastrous for Russia. The Japanese surprise attack on Port Arthur blockaded Russian ships there. The Russians on site ultimately surrendered after the long siege. In a desperate move, the Russians rerouted its Baltic fleet to the Pacific. Once they arrived, they were routed in the Battle of Tshushima on May 27, 1905, in what was the most decisive naval rout since Trafalgar in 1805. Russia sought peace. The Treaty of Portsmouth was signed on September 5, with mediation from US President Teddy Roosevelt. Russia lost most of its holdings in the Pacific as a result. Port Arthur was ceded to Japan, as was the southern half of the Sakhalin Island, and Japanese dominance of Korea was recognized. Manchuria was returned to China while the Chinese Eastern Railway (see MID\#3250) was ceded to Japan as well (p. 1003)
\\
Britannica. 2023. “Russo-Japanese War.” Accessed August 20, 2023. https://www.britannica.com/event/Russo-Japanese-War. & Russo-Japanese War, (1904–05), military conflict in which a victorious Japan forced Russia to abandon its expansionist policy in East Asia, thereby becoming the first Asian power in modern times to defeat a European power.
\\
Van Der Oye, David S. 2011. “Russo-Japanese War (1904–1905).” In The Encyclopedia of War, edited by Gordon Martel. & This was a major war in Manchuria and the surrounding waters fought for dominance in northeast Asia. Although Japan was unable decisively to defeat Russia on land, its naval victories together with revolutionary unrest in Russia enabled the island empire to prevail. Japan's victory marked the first significant defeat of a leading European power by a non-European in the modern age. The conflict's origins lie in the decline of imperial China, which whetted the territorial appetites of more vigorous powers. Russia had already prized large swaths of land on the Amur and Ussuri rivers from the Middle Kingdom in 1860 and considered the rest of Manchuria an area of interest. Japan's easy victory in the Sino-Japanese War (1894–1895), however, encouraged that nation's expansive designs as well.  Russian diplomatic pressure during the peace talks at Shimonoseki in 1895, which forced Japan to retrocede to China the Liaodong Peninsula with its naval base of Port Arthur at the southern tip, thoroughly aggravated its rival. Tensions were hardly eased by Russia's occupation of the same naval base two years later. A clash might nevertheless have been averted had St.Petersburg responded to Tokyo's repeated efforts to negotiate respective spheres of influence in the region. However, the former's endless prevarications ultimately convinced Japan to go to war by early 1904.
\\
Orgill, Nathan N. 2025. “Russo-Japanese War (1904–5).” In The Encyclopedia of Diplomacy, edited by Gordon Martel. & The Russo-Japanese War (1904–5) was a momentous event in the diplomatic history the twentieth century. Originating in a competition for influence and territory in Manchuria and Korea a decade earlier, the conflict led to a significant victory for Japan, as Russia recognized her preponderance in the contested regions that sparked the war. The long-term origins of the Russo-Japanese War stretch back to the late nineteenth century, when both Russia and Japan set themselves on trajectories that led to the ultimate collision of 1904. (...). Russia's conflicting interests with Japan in the Far East originated out of her own designs in the region, particularly in Manchuria and Korea. While she had focused her attention on expanding her influence at the expense of the Ottoman Empire and in competition with Austria-Hungary during the 1870s and 1880s, a shift occurred after 1897 (...). (...) shifting attention away from her European interests to the Far East. The shift allowed influential advisers like Sergey Witte (1849–1915) (...) to gain ascendency with Tsar Nicholas II (r. 1894–1917). The result was Russian railroad development in Siberia and particularly the construction of the Trans-Siberian Railway – underway after 1891 – with its terminus at the Pacific port of Vladivostok in the Far East and all-important implications for Russo-Japanese relations. (...). The Russians made the Boxer uprising a pretext to intervene militarily in southern Manchuria, deploying troops in September and October 1900 (...). The Japanese desired to work towards a peaceful solution and Russian evacuation of Manchuria, while defending their interests in Korea from being undermined further (Gooch and Temperley 1926–38: 3.207). But Witte's dismissal in 1903 marked a defeat for moderation and the triumph of the so-called “Koreans” around the tsar, who desired aggressive measures to establish a position of predominance in Manchuria and Korea (Lee 1974: 85). The outward signal of these intentions came in the form of an administrative reorganization in Russia's Far Eastern holdings in August 1903. The two Russian provinces bordering on Manchuria and northern Korea were consolidated into one massive viceroyalty under Vice Admiral Yevgeny I. Alexeyev (1843–1918), a proponent of expansion into Korea.


\\
Yanagihara, Masaharu. 2025. “Treaty of Portsmouth (1905).” In The Encyclopedia of Diplomacy, edited by Gordon Martel. & The Treaty of Portsmouth concluded the 1904–5 Russo-Japanese War – a war which is sometimes referred to as “World War Zero,” as it was the first major international war of the twentieth century. The climactic land battle of the war took place at Mukden, which resulted in its occupation by the Japanese army on March 10, 1905. This was certainly a victory for Japan, but the main Russian forces succeeded in retreating to the northern part of Manchuria, while the Japanese army had difficulty in continuing to provision its troops. At sea, the Japanese navy achieved a remarkable and one-sided victory at the battle of Tsushima (May 27–28, 1905).
\\
\end{longtable}}
\clearpage
         \textbf{Third Central American} 

          \underline{Onset:} 1906         

\underline{Reasons:}
\begin{itemize}
\item Nationalism
\item Religion or Ideology
\item Revenge/Retribution
\end{itemize}

         \underline{Sources:} 
         \rowcolors{3}{white}{light-gray} 
         {\scriptsize 
         \begin{longtable}{p{0.2\textwidth}p{0.8\textwidth}} 
         \hiderowcolors% 
         \toprule 
         Source & Excerpt \\ 
         \midrule 
         \showrowcolors% 
         \endfirsthead% 
         \hiderowcolors% 
         \toprule 
         Source & Excerpt \\ 
         \midrule 
         \showrowcolors% 
         \endhead% 
         \midrule 
         \endfoot% 
         \hiderowcolors% 
         \bottomrule 
         \showrowcolors% 
         \endlastfoot% 
         
Gibler, Douglas M. 2018. International Conflicts, 1816–2010: Militarized Interstate Dispute Narratives. Vol. 1 Rowman \& Littlefield. & This war began as a two-pronged assault on Guatemala by emigres from El Salvador and Honduras, aimed at replacing Guatemalan dictator Estrada Cabrera with Manuel Lisandro Barillas. Guatemala blamed the Salvadoran and Honduran governments for fostering the attacks and began an offensive against both states. The US minister to Guatemala tried to intervene under the Roosevelt Corollary and the Corinto Pact of 1904, but the latter did not involve Guatemala and was thus rejected. When the conflict had reached the threshold for war, Roosevelt personally took charge of the negotiation process. Mexico joined the diplomatic effort with the United States since both were interested in promoting business initiatives in Central America and required peace in the area to do so. With approval of the Mexican government, the United States dispatched the warship USS Marblehead to Central American waters as a show of force and a signal to resolve the dispute among the three states. The Marblehead became the de facto good offices of the United States. The three belligerents met on board on July 20, 1906, and concluded a preliminary peace treaty the same day. The treaty called for disarmament of all sides within three days and a negotiation of a more comprehensive peace within two months. This more comprehensive peace came on September 25 and took effect on September 14, 1907. (p. 80)
\\
Slade, William F. 1917. The Journal of Race Development. The Federation of Central America & Guatemala invaded Honduras and Salvador. The latter republic was anxious for peace, but the war spirit was strong in Guatemala. 
\\
\end{longtable}}
\clearpage
         \textbf{Fourth Central American} 

          \underline{Onset:} 1907         

\underline{Reasons:}
\begin{itemize}
\item Nationalism
\item Power Transition or Security Dilemma
\item Economic, Long-Run
\end{itemize}

         \underline{Sources:} 
         \rowcolors{3}{white}{light-gray} 
         {\scriptsize 
         \begin{longtable}{p{0.2\textwidth}p{0.8\textwidth}} 
         \hiderowcolors% 
         \toprule 
         Source & Excerpt \\ 
         \midrule 
         \showrowcolors% 
         \endfirsthead% 
         \hiderowcolors% 
         \toprule 
         Source & Excerpt \\ 
         \midrule 
         \showrowcolors% 
         \endhead% 
         \midrule 
         \endfoot% 
         \hiderowcolors% 
         \bottomrule 
         \showrowcolors% 
         \endlastfoot% 
         
Gibler, Douglas M. 2018. International Conflicts, 1816–2010: Militarized Interstate Dispute Narratives. Vol. 1. Rowman \& Littlefield. & This conflict follows shortly after MID\#1205—the Third Central American War in 1906. That war ended with American intervention that frustrated Nicaraguan President Jose Santos Zelaya's interest to bring back and rule a united confederation in Central America. Zelaya tried again seven months later and first targeted Manuel Bonilla in Honduras. Zelaya started by supporting rebel groups in Honduras. When these rebel groups fled into Nicaragua, Bonilla asked Zelaya for permission to pursue them across the border. Zelaya refused and Bonilla's forces ransacked the relevant villages in Nicaragua anyway. Zelaya wanted reparations, but Bonilla refused. Rather than waste time with the American-Mexican requested mediation, Zelaya invaded Honduras in February 1907, and El Salvador entered the fray in defense of Honduras. Zelaya's Nicaraguan forces ultimately overwhelmed Honduras and El Salvador, winning a pivotal battle at Namasigue on March 18, 1907. Zelaya tried to go further, wanting to instigate rebellion in El Salvador. This would have brought regional heavyweight Guatemala into the fray. Here, as before, Zelaya was frustrated by the intervention of the Americans and the Mexicans. The Americans in particular took an active role in ending this dispute short of what Zelaya desired, much to Zelaya's chagrin. Bonilla was expelled from Honduras, ultimately taking refuge on the USS Chicago before living in exile in the United States. He was replaced by Manuel Davila, who was distrusted by Zelaya and shortly signed an agreement with El Salvador to try to oust Zelaya. On June 11, a Nicaraguan warship attacked and captured the Salvadoran port of Acajutla. Honduras then joined in the dispute on the side of Nicaragua. On September 24 Guatemala announced it would support an uprising in Honduras, and the ouster of puppet President Davila, effectively taking sides with El Salvador. The United States sponsored a Washington Conference of 1907 to try to enact some kind of settlement to the issues that would bring some semblance of peace to the region. (...). During the course of the negotiations, it became clear to the Americans that Zelaya was a menace in the region and was working toward a situation where he was the hegemon. His presence made Nicaragua more worrisome to peace in the area than Guatemala. Shortly after the conclusion of the Washington Conference, Zelaya abrogated the terms of the peace treaty by leading rebel groups against Davila in Honduras. The United States ultimately removed him from Nicaragua in 1909. (pp. 87-88)
\\
Slade, William F. 1917. “The Federation of Central America.” The Journal of Race Development. & On January 9, 1907, that republic let it be known that it refused to be bound by the terms of the Treaty of Marblehead. The President of Honduras found it necessary to order troops near the Nicaraguan frontier to suppress revolutionary movements directed against his government; these he believed were aided by President Zelaya of Nicaragua. The Hondurean troops finally invaded Nicaragua, on the ground that in order to suppress the revolutionists who sought refuge in Nicaragua, it was necessary to cross the boundary. 
\\
Martin, Percy F. 1911. Salvador of the Twentieth Century, 72–74. & In the troubles which afflicted the country in the years 1907-08, the whole cause was the incitement which was offered to them by their turbulent and troublesome neighbours the Nicaraguans and the Honduraneans. As I have shown very conclusively, it was the long-established policy of Santos Zelaya to foster an outbreak in Salvador which should broaden into a revolution, in the course of which Salvadorean troops would be compelled innocently to commit some overt act which could give Honduras or Nicaragua a cause for the initiation of a movement against the Republic. This, it was hoped, would ultimately result in the election to Presidency of Salvador of Dr. Prudencio Alfaro, who was aleways a creature of Santos Zelaya [...]
\\
\end{longtable}}
\clearpage
         \textbf{Second Spanish-Moroccan} 

          \underline{Onset:} 1909         

\underline{Reasons:}
\begin{itemize}
\item Nationalism
\item Power Transition or Security Dilemma
\item Economic, Long-Run
\end{itemize}

         \underline{Sources:} 
         \rowcolors{3}{white}{light-gray} 
         {\scriptsize 
         \begin{longtable}{p{0.2\textwidth}p{0.8\textwidth}} 
         \hiderowcolors% 
         \toprule 
         Source & Excerpt \\ 
         \midrule 
         \showrowcolors% 
         \endfirsthead% 
         \hiderowcolors% 
         \toprule 
         Source & Excerpt \\ 
         \midrule 
         \showrowcolors% 
         \endhead% 
         \midrule 
         \endfoot% 
         \hiderowcolors% 
         \bottomrule 
         \showrowcolors% 
         \endlastfoot% 
         
Gibler, Douglas M. 2018. International Conflicts, 1816–2010: Militarized Interstate Dispute Narratives. Vol. 2 Rowman \& Littlefield. & The Spanish-Moroccan war of 1909 is commonly known as one of the Rif Wars. Spain was limited in its colonial ambitions, certainly in Africa. A weak state relative to the rest of Europe that had been dealt a damaging blow by the United States in Cuba and the Philippines (the Spanish-American War, see MID\#1557), the north of Morocco was the extent to which Spain could be involved in the African scramble. Spain already had establishments in Morocco acquired in the past, and additional inquiries into Morocco came only after an agreement was reached with France regarding a division of territory. Even this brought intrigue from the Germans, though mostly directed at the French. This war is commonly known as one of the Rif Wars because it is a conflict between Spain and Morocco regarding the Riffian tribesmen that lived in the north of Morocco, where the Spanish settlements were. Spanish nationals working on railways were killed by some Riffians near Melilla, which was one of the Spanish settlements. After some early setbacks, in part because of Spain's dilapidated infrastructure and discontent in the Spanish population, the Spanish military eventually overwhelmed Morocco in battle. The fighting ultimately ended on March 23, 1910, and was formally concluded with an agreement signed on November 16. An indemnity was fixed for the conflict, and an outline was provided for customs houses in Ceuta and the maintenance of order. Morocco would not remain independent much longer as France soon coerced Morocco into a protectorate, eliminating it from the interstate system until 1956. (pp. 541-542)
\\
Chandler, James A. 1975. “Spain and Her Moroccan Protectorate 1898–1927.” Journal of Contemporary History. & The first interventions by Spain in Morocco were largely pre-emptions of possible French initiatives. They were also excursions into what was for them largely unknown territory. In 1906 a trading station at Restinga, near Melilla, had been set up by a French adventurer in order to supply El Rogui, a local pretender to the Sultan's throne. In 1907 the region was occupied by the Sultan's troops largely at Spanish insistence but this force was defeated by El Rogui at the end of the year. Spain then moved into the area in February 1908 in order to prevent any possible re-establishment of the French statio
\\
Alvarez, José E. 2011. “Rif Wars (1893, 1909, 1920).” In The Encyclopedia of War, edited by Gordon Martel. & There were three significant colonial wars fought between the Spanish army and Moroccan tribesmen starting in the late nineteenth century and continuing into the early twentieth century. (...). The situation in Morocco once again heated up on July 9, 1909, when a force of Riffian tribesmen attacked a military outpost protecting Spanish workers building a railway to serve Spanish-owned iron ore mines outside Melilla. Four workers and one sentry were killed. (...). This episode visibly demonstrated the government's commitment not only to Spanish commercial interests in northwestern Morocco by reinforcing the number of troops posted there, but also to the expansion of its Melillan hinterlands. Moreover, it brought to light the deep divisions that existed in Spain between those who advocated overseas colonization and those who opposed it. (...). Moreover, this campaign had gained new territory for Spain, which had been losing its overseas colonies around the world for the last two hundred years. In addition, the Melillan campaign of 1909–1910 provided the Spanish army an opportunity to gain glory, respect, and promotions. (...). Militarily speaking, the 1909 campaign served as a wake-up call for the army.
\\
\end{longtable}}
\clearpage
         \textbf{Italian-Turkish} 

          \underline{Onset:} 1911         

\underline{Reasons:}
\begin{itemize}
\item Nationalism
\item Economic, Long-Run
\item Domestic Politics
\item Economic, Short-Run
\end{itemize}

         \underline{Sources:} 
         \rowcolors{3}{white}{light-gray} 
         {\scriptsize 
         \begin{longtable}{p{0.2\textwidth}p{0.8\textwidth}} 
         \hiderowcolors% 
         \toprule 
         Source & Excerpt \\ 
         \midrule 
         \showrowcolors% 
         \endfirsthead% 
         \hiderowcolors% 
         \toprule 
         Source & Excerpt \\ 
         \midrule 
         \showrowcolors% 
         \endhead% 
         \midrule 
         \endfoot% 
         \hiderowcolors% 
         \bottomrule 
         \showrowcolors% 
         \endlastfoot% 
         
Gibler, Douglas M. 2018. International Conflicts, 1816–2010: Militarized Interstate Dispute Narratives. Vol. 2 Rowman \& Littlefield. & The Italian-Turkish War of 1911 to 1912 resulted in the establishment of an Italian colony over the entirety of Libya. Tensions ran deep between the two states (...), and the Ottoman Empire had lost its tight grip on its North African territories. While France and Great Britain scrambled for various parts of the Ottoman Empire's administrative districts in North Africa, Italy had its intentions locked on Libya. Acquiring Libya would exponentially improve Italy's position in the Mediterranean as well as balance the power between France and Britain elsewhere in North Africa—the French in Morocco, Algeria and Tunis, and the British in Egypt. The occasion for war came as a result of the Agadir Crisis involving the Great Powers. Italy used the distraction to attack the Ottoman Empire and declared war on September 29, 1911. The ensuing battles between Turkey and Italy were evenly matched, owing mostly to Italian incompetence and poor planning. Unfortunately for the Ottoman Empire, the fallout of the Agadir Crisis and the Italian surprise attack on Ottoman controlled Libya enabled the Balkan uprisings in 1912. Distracted with conflicts closer to home, the Ottoman Empire sued for peace. Italy used this to their advantage, gaining favorable terms in the Treaty of Ouchy on October 15 and the Treaty of Lausanne on October 18, 1912. The Treaty of Ouchy announced the Ottoman evacuation of Libya, though the sultan retained the right to appoint religious officials. Italy gained complete sovereignty over Libya. The Treaty of Lausanne called for the Italian evacuation of the Dodecanese upon completion of the Turkish evacuation of Libya. This never happened. World War I started around the time the Italians were scheduled to make the transition, and the Italians ended up possessing these islands until after World War II. Greece inherited them thereafter. (pp. 558-559)
\\
Clark, Christopher M. 2012. The Sleepwalkers: How Europe Went to War in 1914. London: Allen Lane.  & From 1909 onwards, the nationalist Enrico Corradini, supported by the nationalist organ L'Idea Nazionale, campaigned energetically for an imperialist enterprise focused on Libya; by the spring of 1911 he was openly demanding a policy of invasion and seizure [...] Until the summer of 1911, however, Italy's leading statesmen remained faithful to the country's ancient axiom that Italy must not provoke the break-up of the Ottoman Empire. As late as the summer of 1911, Prime Minister Giovanni Giolitti was still firmly rejecting calls to adopt a more aggressive position vis-à-vis Constantinople on a range of issues relating to the governance of Ottoman Albania. It was the French intervention in Morocco that changed everything. The Italian foreign ministry believed it had excellent grounds for demanding a quid pro quo in Libya. In view of France's ‘radical modification' of the situation in the Mediterranean, it would now be impossible, an Italian foreign ministry senior official pointed out, to ‘justify' a policy of continuing inaction ‘before public opinion'
\\
Vandervort, Bruce. 2011. “Turco-Italian War (1911–1912).” In The Encyclopedia of War, edited by Gordon Martel. & The Italian invasion of Ottoman-ruled Libya, which comprised the provinces of Tripolitania, with its capital at Tripoli, and Cyrenaica, with its capital at Benghazi,in October 1911 was, in the words of Italian historian Angelo del Boca, “a project nurtured for 30 years.” (...) [Actual] planning for the conquest of the region only began in earnest in November 1884, when the Italian government came to believe that the French, who had just established a protectorate over Tunisia, were about to seize Morocco, which would have given them control over all of North Africa, save for Egypt and Libya. Rome feared both for Italy's maritime security, as the Mediterranean risked becoming a French lake, and for Italy's future as an imperial power, a future already darkened by the French occupation of Tunisia, where many Italians lived and Italy had played a leading economic role for decades. Some Italians also cast covetous eyes on what they believed, quite erroneously as it turned out, was a rich caravan trade across the desert from sub-Saharan Africa to the Libyan port of Tripoli and feared its diversion to entrepots in Tunisia or Algeria controlled by France. (...). Italians liked to portray their attempt to seize Libya as the reconquest of a lost colony. Also, the Italian government saw Libya – and Eritrea and Ethiopia – as places in which to settle the country's surplus population. The loss of so many citizens to new allegiances in the Americas was a burning issue in Italian politics in the early twentieth century. This led some Italian politicians and journalists to make exaggerated claims about the agricultural and commercial potential of Libya. Some Italian politicians and statesmen coveted Libya for strategic reasons. They believed Italy needed portsand naval bases in North Africa in order to keep control over the central Mediterranean in the hands of Italy's new and formidable navy.
\\
Die Zeit. 2003. “Libyen, verheißenes Land.” Die Zeit, May 15, 2003. & Italien war ein bitterarmes, industriell zuruckgebliebenes Land. [...] Aus schierer Not kehrten zwischen 1901 und 1911 nicht weniger als 1,6 Millionen Italiener der Großen Proletarierin den Rucken, um sich in den USA und in Argentinien ein besseres Leben aufzubauen. [...] Nationalistische Intellektuelle im Umfeld der jungen Associazione nazionalista italiana lancierten die Idee, dass nur ein weiteres koloniales Ausgreifen die sozialen Probleme Italiens lösen könne. Die großen Tageszeitungen griffen den Gedanken auf. Tripolitanien und die Cyrenaika, die beiden letzten nordafrikanischen Provinzen des Osmanischen Reiches (auf dem Gebiet des heutigen Libyen), wurden die naturliche Verlängerung der Halbinsel in Afrika bilden. In den menschenleeren Weiten gebe es fur die Söhne Italiens fruchtbare Ländereien im Überfluss. Das Gebiet um den Syrte-Bogen sei die terra
promessa, die verheißene Erde, und fur die Auswanderer eine heimatnahe Alternative zu Amerika. Nichts davon stimmte: Das Land glich eher einer Sandschachtel, und die fruchtbaren Kustenstriche am Meer waren dicht besiedelt.
\\
\end{longtable}}
\clearpage
         \textbf{First Balkan} 

          \underline{Onset:} 1912         

\underline{Reasons:}
\begin{itemize}
\item Nationalism
\item Power Transition or Security Dilemma
\item Domestic Politics
\end{itemize}

         \underline{Sources:} 
         \rowcolors{3}{white}{light-gray} 
         {\scriptsize 
         \begin{longtable}{p{0.2\textwidth}p{0.8\textwidth}} 
         \hiderowcolors% 
         \toprule 
         Source & Excerpt \\ 
         \midrule 
         \showrowcolors% 
         \endfirsthead% 
         \hiderowcolors% 
         \toprule 
         Source & Excerpt \\ 
         \midrule 
         \showrowcolors% 
         \endhead% 
         \midrule 
         \endfoot% 
         \hiderowcolors% 
         \bottomrule 
         \showrowcolors% 
         \endlastfoot% 
         
Gibler, Douglas M. 2018. International Conflicts, 1816–2010: Militarized Interstate Dispute Narratives. Vol. 1. Rowman \& Littlefield. & The First Balkan War (in 1912 and 1913) was fought by the Ottoman Empire against the allied forces of Serbia, Greece, and Bulgaria. Once an international force that threatened to conquer all of Europe, the Ottoman Empire of the 19th century and onward was unable to prevent independence movements within its confines and was unable to prevent European states—notably, the Russian Empire—from ripping away territory from its borders. The issue was only getting worse in the beginning of the 20th century. (...). After conclusion of the Italian-Turkish War, the Balkan League declared war on the Ottoman Empire. Their aims were to take away then-Turkish territories remaining in the Balkan Peninsula—the areas of Thrace, Macedonia, and Eastern Rumelia. The Ottoman Empire, the Sick Man of Europe, was routed. Exhausted from the conflict with Italy, the quick turnaround and strategic disadvantage ended the war quickly. Serbia and Bulgaria concluded an armistice with the Turks on December 3, 1912, shortly after the war began. However, a conference in London with the warring parties and the great powers in December 1912 produced no satisfactory outcome. (pp. 372-373)
\\
Britannica. 2023. “Balkan Wars.” Encyclopædia Britannica. Accessed August 19, 2023. https://www.britannica.com/topic/Balkan-Wars. & The Balkan Wars had their origin in the discontent produced in Serbia, Bulgaria, and Greece by disorder in Macedonia. The Young Turk Revolution of 1908 brought into power in Constantinople (now Istanbul) a ministry determined on reform but insisting on the principle of centralized control. There were, therefore, no concessions to the Christian nationalities of Macedonia, which consisted not only of Macedonians but also of Serbs, Bulgarians, Greeks, and Vlachs. The Albanians, whose growing sense of nationalism had been awakened by the Albanian League, were likewise discontented with the Young Turks' centralist policy.
\\
Hall, Richard C. 2011. “Balkan Wars (1912–1913).” In The Encyclopedia of War, edited by Gordon Martel. & In the second half of the nineteenth century,the Balkan Peninsula increasingly became an area of nationalist restiveness, directed mainly against the waning power of the Ottoman Empire. (...). 
The Bulgarians,Greeks, Montenegrins, and Serbs all sought additional territories from the Ottoman Empire to realize their national unity. (...). 
The Young Turk coup of 1908 aroused fears among these Balkan states that resulting reforms could strengthen the Ottoman Empire and thus deny the Bulgarians,Greeks, Montenegrins, and Serbs Ottoman territories containing their co-nationals.These concerns enabled the Balkan states to overcome their national rivalries and to cooperate against the Ottomans before the Young Turk reforms had a chance to succeed. By the summer of 1912, with Russian encouragement, these efforts resulted in the establishment of a loose Balkan League directed against the Ottoman Empire. (...).
The Balkan League decided to act in the fall of 1912, before the anticipated end of the Italo-Turkish War could bring additional Ottoman forces to the Balkans.Montenegro began the war on October 8,1912. The war became general ten days later.
\\
Brennan, Christopher. 2025. “Balkan Wars (1912–13).” In The Encyclopedia of Diplomacy, edited by Gordon Martel. & The Balkan Wars (1912–13) arose from the desire of former Ottoman territories in the Balkans – now independent, nationalistic, land-hungry countries: Bulgaria, Greece, Serbia, Montenegro, and Romania – to carve up the remainder of the empire's European possessions among themselves. (...).
The Treaty of Berlin in 1878 had already begun dismantling the [Ottoman] empire, granting de facto independence to Romania, Serbia, and Montenegro, along with autonomy for Bulgaria (which had previously made far greater gains during the Preliminary Treaty of San Stefano early that year, to general consternation). Nevertheless, an independent Greater Bulgaria was wiped off the map, and one of its vital constituent parts, Macedonia, was returned to the Ottomans. This was the first of many unsatisfactory outcomes of the Eastern Question. Despite the reformist Young Turk revolution of 1908 (...) [the Ottoman Empire] continued to suffer from internal turmoil and to lose its grip on its European possessions. (...). Italy lit the fire for the First Balkan War by successfully attacking and occupying Tripolitania (now in Libya) and the Dodecanese (in the Aegean Sea, off the Anatolian coast), both under Ottoman rule. This whet the appetite of the remaining Balkan countries.
\\
Parker, Geoffrey, ed. 2020. The Cambridge History of Warfare. 2nd ed. Cambridge: Cambridge University Press. & A series of crises in the Balkans provided tinder for the impending European conflagration. The weakness of the Ottoman Empire combined with Austrian and Russian ambitions to exacerbate the situation. Moreover, all three imperial powers felt threatened by internal political developments and therefore sought to escape their domestic dilemmas through foreign policy successes - starting in the Balkans. (...). 
The next round of Balkan troubles began in September 1911 when the Italians seized the Turkish province of Libya in north Africa, using aerial bombardment for the first time in war, and then the Dodecanese islands off the Turkish coast. Italy's easy success encouraged four Balkan states - Serbia, Bulgaria, Greece, and Montenegro - to jump on the prostrate Turks in October 1912; but the thieves soon fell out over the spoils. In June 1913 Serbia, Greece, and Montenegro joined Romania to pummel Bulgaria. Even the Turks joined in. While Serbia gained much from the war, almost doubling its size, Austria issued an ultimatum threatening war unless Serbia withdrew its troops from some contested lands. When Germany once again deterred the tsar from intervening, Serbia immediately capitulated. (pp. 271-272)
\\
Morillo, Stephen, Jeremy Black, and Paul Lococo. 2008. War in World History: Society, Technology, and War from Ancient Times to the Present. New York: McGraw-Hill Professional. & The Balkans were a major area of European conflict in the 1870s, 1910s, 1940s, and 1990s. As
such, they registered both changes in the nature of war and the continued potency of ethnic animosities. The desire for territorial control was the key objective, seen as the way to ensure ethnic survival and success. (...). As with the conflicts in the 1990s that caused, and arose from, the collapse of Yugoslavia, so the key in 1912–13 was the fall of the Turkish Empire in Europe. The First Balkan War (1912–13) saw Turkey attacked by its neighbors—Bulgaria, Greece, Montenegro, and Serbia—and most of its European empire conquered. At the tactical, strategic, and operational level, this war demonstrated the risks associated with the absence of mass formations, as Turkish strategic dispersal—designed to prevent territorial loss in the face of attack from a number of directions—enabled their opponents to achieve key superiorities in particular areas of attack.
\\
\end{longtable}}
\clearpage
         \textbf{Second Balkan} 

          \underline{Onset:} 1913         

\underline{Reasons:}
\begin{itemize}
\item Nationalism
\item Border Clashes
\item Revenge/Retribution
\end{itemize}

         \underline{Sources:} 
         \rowcolors{3}{white}{light-gray} 
         {\scriptsize 
         \begin{longtable}{p{0.2\textwidth}p{0.8\textwidth}} 
         \hiderowcolors% 
         \toprule 
         Source & Excerpt \\ 
         \midrule 
         \showrowcolors% 
         \endfirsthead% 
         \hiderowcolors% 
         \toprule 
         Source & Excerpt \\ 
         \midrule 
         \showrowcolors% 
         \endhead% 
         \midrule 
         \endfoot% 
         \hiderowcolors% 
         \bottomrule 
         \showrowcolors% 
         \endlastfoot% 
         
Gibler, Douglas M. 2018. International Conflicts, 1816–2010: Militarized Interstate Dispute Narratives. Vol. 1 Rowman \& Littlefield. & The Second Balkan War directly succeeded and was a consequence of the First Balkan War. The Balkan League scored a decisive victory on the battlefield against the Ottoman Empire, but interference by the Great Powers leading up to the Treaty of London left no one satisfied with the conflict's resolution. Serbia was denied Albania, and the treaty gave no provisions for the divisions of the spoils in Macedonia, Thrace and Eastern Rumelia. That was left to the allies, whose understandings were shaped by the alliances that brought them into concert with one another. Bulgaria also had strong intentions of becoming a hegemon in the Balkans and was intent on maximizing its control over Macedonia. Serbia, which unexpectedly lost any claim to Albania by great power interference, asked the Bulgarians to cede more of Macedonia as a result. Bulgaria refused, and Romania, which did not participate in the First Balkan War, demanded some of the war spoils as well. Finally, Greece and Bulgaria were having difficulties over who would gain Thessaloniki. The Greeks had occupied it first and were not going to give it to the Bulgarians. Under pressure and seeing no better way to maximize its share of the Balkan Peninsula, Bulgaria surprised its former allies by attacking them on June 30, 1913. Bulgaria's surprise attack was foolish. Serbia and Greece's combined forces outmatched Bulgaria's personnel. This became even more apparent when Romania and the Ottoman Empire threw their lot in with the Serbians and Greeks in July. Bulgaria sued for peace on July 30, 1913, ending the fighting, and the Treaty of Bucharest and Treaty of Constantinople concluded the war on August 10, 1913. Bulgaria was forced out of Macedonia, which was ultimately shared between the Serbians and Greeks. The Greeks were big winners for these two Balkan conflicts, ultimately doubling their territory in a matter of months. Romania gained the Dobruja region on the lower Danube. The Ottoman Empire recovered Edirne, part of Thrace that it had previously lost. (pp. 337-338)
\\
Britannica. 2023. “Balkan Wars.” Accessed August 19, 2023. https://www.britannica.com/topic/Balkan-Wars. & The second conflict erupted when the Balkan allies Serbia, Greece, and Bulgaria quarreled over the partitioning of their conquests. The result was a resumption of hostilities in 1913 between Bulgaria on the one hand and Serbia and Greece, which were joined by Romania, on the other.
\\
Hall, Richard C. 2011. “Balkan Wars (1912–1913).” In The Encyclopedia of War, edited by Gordon Martel. & In the second half of the nineteenth century,the Balkan Peninsula increasingly became an area of nationalist restiveness, directed mainly against the waning power of the Ottoman Empire. (...). The Bulgarians,Greeks, Montenegrins, and Serbs all sought additional territories from the Ottoman Empire to realize their national unity. (...).
[By the Spring of 1913] tensions were rising among the Balkan allies. The failure of the Serbs to retain northern Albania increased their determination to retain Macedonia in the face of growing Bulgarian opposition.The Bulgarians and Greeks failed to reach any agreement for the disposition of conquered Ottoman territories and soon fell to skirmishing over northern Macedonia. By May 5, 1913, the Greeks and Serbs had concluded an alliance directed against Bulgaria. To complicate the situation, the Romanians, who wanted compensation for any Bulgarian gains in the war, began to make demands on Bulgarian (southern) Dobrudzha. A Great Powers ambassadors' conference held in St. Petersburg in April 1913 failed to resolve the issue to the satisfaction of either the Bulgarians or Romanians. The conclusion of the London treaty enabled the Bulgarians to transfer the bulk of their army from the Chataldzha lines to the southwestern part of their country in order to enforce their claims to Macedonia. Before the Russians could act upon their promise to mediate the dispute,an explosion occurred.
\\
Brennan, Christopher 2025. “Balkan Wars (1912–13).” In The Encyclopedia of Diplomacy, edited by Gordon Martel. & The Balkan Wars (1912–13) arose from the desire of former Ottoman territories in the Balkans – now independent, nationalistic, land-hungry countries: Bulgaria, Greece, Serbia, Montenegro, and Romania – to carve up the remainder of the empire's European possessions among themselves. (...).
By the Treaty of London on May 30, 1913, the allies retained their military gains while the Ottoman Empire kept only a minute strip of European land west of Constantinople (...). (...). Motives of dissatisfaction quickly arose. For instance, the treaty placed the fate of Albania in the hands of the Great Powers (according to the wishes of Austria-Hungary and Italy, but against those of Russia), to the chagrin of Greece and Serbia, who occupied the land at the time and had expected to partition it. The Serbs in particular did not gain the access to the sea they had been promised. To compensate for this loss, both looked towards Macedonia, which had been partly annexed by Bulgaria but partly unassigned. Greeks and Serbs thus reached a secret agreement on the partition of the territory and on mutual assistance in case of war. (...).
The Bulgarians – who had mobilized 25 percent of their male population (500,000 combatants), had done most of the fighting, in Thrace, and lost more men than any of their coalition partners – were aggrieved by the fact that they had, proportionally, gained far less territory than the Greeks and the Serbs (though Bulgaria was now the largest country in the Balkans). What is more, the overwhelming majority of Bulgarians considered Macedonia rightfully theirs and its acquisition had been the principal topic of Bulgarian politics in preceding years. Accordingly, and despite the lack of Great Power support and the huge sacrifices it had already consented to during the First Balkan War, a war-weary, weakened, and sapped Bulgaria (...) launched an attack on its former allies Greece and Serbia on the night of June 29–30, 1913.
\\
Morillo, Stephen, Jeremy Black, and Paul Lococo. 2008. War in World History: Society, Technology, and War from Ancient Times to the Present. New York: McGraw-Hill. & The Balkans were a major area of European conflict in the 1870s, 1910s, 1940s, and 1990s. As
such, they registered both changes in the nature of war and the continued potency of ethnic animosities. The desire for territorial control was the key objective, seen as the way to ensure ethnic survival and success. (...). As with the conflicts in the 1990s that caused, and arose from, the collapse of Yugoslavia, so the key in 1912–13 was the fall of the Turkish Empire in Europe. (...). The dynamics of power politics quickly led the victors to fall out, and in the Second Balkan War (1913), Bulgaria unsuccessfully fought its recent allies, as well as Romania and Turkey. The Bulgarians fought well in Macedonia, where they concentrated their forces, but this allowed the Turks to recapture Adrianople and the Romanians to advance on the capital, Sofi a. Bulgaria capitulated and surrendered territory. Renewed confl ict in the Balkans—what became the First World War—was to begin a year later.

\\
Parker, Geoffrey, ed. 2020. The Cambridge History of Warfare. 2nd ed. Cambridge: Cambridge University Press. & A series of crises in the Balkans provided tinder for the impending European conflagration. The weakness of the Ottoman Empire combined with Austrian and Russian ambitions to exacerbate the situation. Moreover, all three imperial powers felt threatened by internal political developments and therefore sought to escape their domestic dilemmas through foreign policy successes - starting in the Balkans. (...). 
The next round of Balkan troubles began in September 1911 when the Italians seized the Turkish province of Libya in north Africa, using aerial bombardment for the first time in war, and then the Dodecanese islands off the Turkish coast. Italy's easy success encouraged four Balkan states - Serbia, Bulgaria, Greece, and Montenegro - to jump on the prostrate Turks in October 1912; but the thieves soon fell out over the spoils. In June 1913 Serbia, Greece, and Montenegro joined Romania to pummel Bulgaria. Even the Turks joined in. While Serbia gained much from the war, almost doubling its size, Austria issued an ultimatum threatening war unless Serbia withdrew its troops from some contested lands. When Germany once again deterred the tsar from intervening, Serbia immediately capitulated. (pp. 271-272)
\\
\end{longtable}}
\clearpage
         \textbf{World War I} 

          \underline{Onset:} 1914         

\underline{Reasons:}
\begin{itemize}
\item Nationalism
\item Power Transition or Security Dilemma
\item Economic, Long-Run
\end{itemize}

         \underline{Sources:} 
         \rowcolors{3}{white}{light-gray} 
         {\scriptsize 
         \begin{longtable}{p{0.2\textwidth}p{0.8\textwidth}} 
         \hiderowcolors% 
         \toprule 
         Source & Excerpt \\ 
         \midrule 
         \showrowcolors% 
         \endfirsthead% 
         \hiderowcolors% 
         \toprule 
         Source & Excerpt \\ 
         \midrule 
         \showrowcolors% 
         \endhead% 
         \midrule 
         \endfoot% 
         \hiderowcolors% 
         \bottomrule 
         \showrowcolors% 
         \endlastfoot% 
         
Norwich University. 2017. “Six Causes of World War I.” Accessed August 20, 2023. https://online.norwich.edu/academic-programs/resources/six-causes-of-world-war-i. & The expansion of European nations as empires (also known as imperialism) can be seen as a key cause of World War I, because as countries like Britain and France expanded their empires, it resulted in increased tensions among European countries. [...] Nationalism was one of many political forces at play in the time leading up to World War I, with Serbian nationalism in particular, playing a key role. [...] The alliance, between France, Britain and Russia, formed in 1907 and called the Triple Entente, caused the most friction among nations. Germany felt that this alliance surrounding them was a threat to their power and existence. As tensions continued to rise over alliances, the preexisting alliances fed into other countries declaring war against one another in the face of conflict.
\\
(1) Woodward, David R. 2011. “World War I: Western Front.” In The Encyclopedia of War, edited by Gordon Martel.

(2) Vandervort, Bruce. 2011. “World War I: Southern Front.” In The Encyclopedia of War, edited by Gordon Martel.

(3) Dinardo, Richard L. 2011. “World War I: Eastern Front.” In The Encyclopedia of War, edited by Gordon Martel.

(4) Hughes, Matthew. 2011. “World War I: Afro-Asian Theaters.” In The Encyclopedia of War, edited by Gordon Martel. & (1) The strengthening of the nation-state prior to 1914, bolstered by a strong sense of national identity, gave Europe's leaders unprecedented popular support. Patriotic young men from all classes were prepared to make the ultimate sacrifice in defense of their country.

(2) [Italian prime minister and foreign minister] had come to support intervention, originally if for no other reason than that they feared the international and domestic consequences of neutrality. If the Central Powers won, they reasoned, an Austrian war of revenge, probably abetted by the Germans, was sure to follow. If, on the other hand, the Entente won, they believed Italy would be likely to face isolation, hostility, and contempt. Domestically, they thought that by standing aside the government would risk condemnation by all those who had been weaned on the notion of Italy as one of Europe's Great Powers.

(3) The Austro-Hungarians continued to believe that the Germans would launch an offensive into Poland even after Moltke had told Conrad in a 1913 letter that Austria's fate would be settled on the Seine River, not the Bug. Russian war planning was based on the expectation that a war against Germany would be difficult at best, but that Austria-Hungary could be dealt with successfully. In any case, the Russians regarded it as imperative that an offensive be launched to aid France in a war with Germany.

(4) During World War I, fighting spread to Africa and Asia as Entente forces attacked the Ottoman (Turkish) Empire in the Middle East, and German colonies in Africa and the Far East (in China and the Pacific islands). (...) British-led forces (...) occupied (...) Palestine, Trans-Jordan, Syria, Lebanon, and Iraq. (...). In Africa, the Belgians, British, and French invaded and occupied German colonies (...). In 1916 Portugal joined the war (...) [by] defending Portugal's African colonies. (...) Japan took German territory in China (...) and Australian, Japanese and New Zealand forces did the same in the Pacific islands of Micronesia (...), northern Papua New Guinea, Nauru, the northern Solomon islands, and Samoa. The conquering powers wanted to establish traditional colonies in captured territory but under pressure—not least from the United States, opposed to European imperial expansion—they were forced to accept a new system of imperial control known as the League of Nations Mandates system.
\\
Otte, Thomas G. 2025. “July Crisis (1914).” In The Encyclopedia of Diplomacy, edited by Gordon Martel. & The July Crisis of 1914 remains the single most studied diplomatic prelude to any war. Scholarly assessments are diverse, and many aspects of Europe's final crisis remain controversial. The “Long Debate” on the origins of the First World War began almost as soon the fighting itself commenced. The crisis itself, meanwhile, underscores the potential limits of diplomacy as a tool for solving international disputes.
A vast array of longer-term causes has been adduced to explain the outbreak of the war. Among them are imperial competition and arms races between the Great Powers, the rise of nationalism, and the structure of the European alliance system, as well as the fragile domestic politics of Europe at the end of the “long” nineteenth century (1800–1914). The trigger for the chain of events that culminated in the descent into the first general war in Europe since 1815 was the assassination of the Archduke Franz Ferdinand of Austria-Este, the nephew of the emperor of Austria-Hungary, Franz Joseph, and the future ruler of that empire, in Sarajevo, capital of the Austrian-administered province of Bosnia, on June 28, 1914. This act of terrorism was perpetrated by a group of self-radicalized Bosnian-Serb students. They had connections with, and had been equipped, financed, and trained by, elements in Serbian military intelligence. That connection – suspected at the time but not proven – was the vital consideration for the key decision-makers of the Habsburg Empire. (...). [The murder of Franz Joseph] was seen at Vienna as the latest and most violent manifestation of Southern Slav nationalism. (...). There was also a sense that international mediation during the two previous Balkan wars of 1912 and 1913 had done little to safeguard its vital interests. (...). The experience of 1912–13 suggested that Belgrade would yield in the face of a determined Austria-Hungary, and it was assumed that Russia, too, would ultimately desist given her financial and domestic weaknesses. There was also a strong sense that a more assertive policy in the Balkans would check the recent spread of Russian influence in the region (...). (...). A further motivating factor was the perceived decline of the Triple Alliance with Germany and Italy. (...). The fragile condition of the Habsburg economy made another mobilization unaffordable, unless it guaranteed to enhance the monarchy's strategic position in the Balkans in the longer term; and only the crushing of Serbia could achieve this. (...). The established consensus favors the view that the German leadership took the “calculated risk” of a continental war, and that Berlin's aim was either a diplomatic triumph or a preventive war. Such interpretation assumes more coherence in German policy than was actually the case because, if anything, Emperor Wilhelm II's promise of unconditional support reflected a lack of policy coordination rather than political intent. Neither Wilhelm's entourage nor his ministers had prepared him for the meeting with the Austro-Hungarian ambassador. His assurances were made almost casually. Moral outrage and a sense of monarchical solidarity prevailed over cold calculation. (...). [Chancellor] Theobald von Bethmann Hollweg, provided no check. (...). Spurred on by a sense of impending change, the two Great Powers had embarked on a dangerous course – Austria-Hungary out of blasé indifference to the wider risks and Germany owing to her chaotic governance.
\\
\end{longtable}}
\clearpage
         \textbf{Latvian Liberation} 

          \underline{Onset:} 1918         

\underline{Reasons:}
\begin{itemize}
\item Nationalism
\item Religion or Ideology
\end{itemize}

         \underline{Sources:} 
         \rowcolors{3}{white}{light-gray} 
         {\scriptsize 
         \begin{longtable}{p{0.2\textwidth}p{0.8\textwidth}} 
         \hiderowcolors% 
         \toprule 
         Source & Excerpt \\ 
         \midrule 
         \showrowcolors% 
         \endfirsthead% 
         \hiderowcolors% 
         \toprule 
         Source & Excerpt \\ 
         \midrule 
         \showrowcolors% 
         \endhead% 
         \midrule 
         \endfoot% 
         \hiderowcolors% 
         \bottomrule 
         \showrowcolors% 
         \endlastfoot% 
         
Gibler, Douglas M. 2018. International Conflicts, 1816–2010: Militarized Interstate Dispute Narratives. Vol. 1 Rowman \& Littlefield. & Latvian territory was under Russian control, but in World War I Latvian officials pushed for separate Latvian units in the Russian army. The Germans captured Riga, Latvia's capital, and the Russian civil war broke out in 1917. The Latvians took the opportunity to move for independence. They first declared Latvia to be an autonomous republic within Russia, but in November 1918 they declared the independent state of Latvia. Russia sought to reassert control quickly. On January 3, 1919, Bolshevik forces invaded Latvia, seized the capital, announced the Latvian Soviet Republic, and called for its union with Soviet Russia. The Soviets controlled nearly all of Latvia's claimed territory by the end of January. It did not take long for other states to get involved. Three days later the British announced its intention to support Latvia in the struggle for independence, and British warships were soon off the coast. German troops also joined the Latvians in March, with the approval of the Allies. The Germans and Latvians then pushed the Soviets out of Riga. The Germans managed to recapture most of Latvia's territory, but in October and November 1919 Russian forces pushed back. Poland, which was involved in its own dispute with the Russians, sent a force in January 1920, pushing the Soviets back toward the border. On February 1, 1920, Latvia and Soviet Russia signed a ceasefire. On August 11, they signed the Treaty of Riga, which granted Soviet recognition of the new Latvian state. (pp. 382-383)
\\
Britannica. 2023. “Baltic War of Liberation.” Accessed August 20, 2023. https://www.britannica.com/event/Baltic-War-of-Liberation. & After World War I ended, however, Soviet Russia, hoping to advance through the Baltic states in order to bring about a Socialist revolution in Germany, attacked in November 1918 and conquered three-quarters of Estonia's territory by the end of the year. In January the Red Army seized the capitals of Latvia and Lithuania, advanced to the Venta River in Latvia, and occupied northern and eastern Lithuania. [...] The commander of the German troops in Latvia, Gen. Rüdiger, Graf von der Goltz, sought to transform Latvia into a base for a new anti-Communist German–Russian force and to form Baltic regimes loyal to imperial Germany and pre-revolutionary Russia. 
\\
Stone, David R. 2011. “Russian Civil War (1917–1920).” In The Encyclopedia of War, edited by Gordon Martel. & Over half the population of the Russian Empire was ethnically non-Russian and took advantage of the breakdown of control in the center to secede. (...). In the Treaty of Brest-Litovsk of March 3, 1918, Lenin surrendered to German control territory in the Baltics, Belorussia, and Ukraine. (...). [Around mid-October 1919] another White force under N. N. Iudenich attempted to take Petrograd, pushing east from the Baltics into the city's eastern suburbs. (...). Both Denikin and Iudenich stalled well short of victory, and with the loss of momentum their seeming triumphs quickly became ignominious retreat. The remnants of Iudenich's force fled into the Baltic states (...). (...). Only on the western frontier, where Poland, Finland, and the Baltic states established effective governments and enjoyed western support, did the Bolsheviks fail to regain large territories that had been part of the Russian Empire.
\\
\end{longtable}}
\clearpage
         \textbf{Estonian Liberation} 

          \underline{Onset:} 1918         

\underline{Reasons:}
\begin{itemize}
\item Nationalism
\item Religion or Ideology
\item Revenge/Retribution
\end{itemize}

         \underline{Sources:} 
         \rowcolors{3}{white}{light-gray} 
         {\scriptsize 
         \begin{longtable}{p{0.2\textwidth}p{0.8\textwidth}} 
         \hiderowcolors% 
         \toprule 
         Source & Excerpt \\ 
         \midrule 
         \showrowcolors% 
         \endfirsthead% 
         \hiderowcolors% 
         \toprule 
         Source & Excerpt \\ 
         \midrule 
         \showrowcolors% 
         \endhead% 
         \midrule 
         \endfoot% 
         \hiderowcolors% 
         \bottomrule 
         \showrowcolors% 
         \endlastfoot% 
         
Gibler, Douglas M. 2018. International Conflicts, 1816–2010: Militarized Interstate Dispute Narratives. Vol. 1. Rowman \& Littlefield. & This dispute is the Estonian War for Independence. Estonia declared independence on November 28, 1917, in the midst of Russia's civil war. The Soviets attempted to reassert control in Estonia, but the Germans moved in to occupy. Under German protection, Estonia reasserted its claim to independence on February 24, 1918. Russia signed the Treaty of Brest-Litovsk to exit World War I on March 3 and began focusing on its civil war. Estonia repeated its claim to independence on November 19th following the defeat of Germany in World War I, but the Soviets invaded Estonia on November 22 and within months controlled most of the country. However, Estonia continued to fight, and the British Royal Navy joined the Estonian effort, driving the Soviets out of Estonia. Peace talks began in December 1919. The two sides agreed to an armistice on January 3, 1920. The Soviet Union and Estonia signed the Tartu Peace Treaty recognizing Estonian independence on February 2. (p. 380)
\\
Minnik, Taavi. 2015. “The Cycle of Terror in Estonia, 1917–1919: On Its Preconditions and Major Stages.” Journal of Baltic Studies 46 (1): 35–47.  & The temporary return of the Bolsheviks at the beginning of the Estonian War of Independence was often seen as a pretext to avenge the injustices suffered under German occupation.
\\
Stone, David R. 2011. “Russian Civil War (1917–1920).” In The Encyclopedia of War, edited by Gordon Martel. & Over half the population of the Russian Empire was ethnically non-Russian and took advantage of the breakdown of control in the center to secede. (...). In the Treaty of Brest-Litovsk of March 3, 1918, Lenin surrendered to German control territory in the Baltics, Belorussia, and Ukraine. (...). [Around mid-October 1919] another White force under N. N. Iudenich attempted to take Petrograd, pushing east from the Baltics into the city's eastern suburbs. (...). Both Denikin and Iudenich stalled well short of victory, and with the loss of momentum their seeming triumphs quickly became ignominious retreat. The remnants of Iudenich's force fled into the Baltic states (...). (...). Only on the western frontier, where Poland, Finland, and the Baltic states established effective governments and enjoyed western support, did the Bolsheviks fail to regain large territories that had been part of the Russian Empire.
\\
Republic of Estonia, Ministry of Foreign Affairs. 2019. Estonian War of Independence 1918–1920: Estonia’s Allies. Accessed August 20, 2023. \url{https://copenhagen.mfa.ee/wp-content/uploads/sites/10/2019/11/Vabadussõja_NÄITUS_2019_DEMO.pdf}. & The Soviets were looking at it as part of the future worldwide workers' revolution. For this reason, a Soviet puppet government, the Commune of the Working People of Estonia, in the border city of Narva was established as a counter to the Republic of Estonia's government.
\\
\end{longtable}}
\clearpage
         \textbf{Second Greco-Turkish} 

          \underline{Onset:} 1919         

\underline{Reasons:}
\begin{itemize}
\item Nationalism
\item Border Clashes
\end{itemize}

         \underline{Sources:} 
         \rowcolors{3}{white}{light-gray} 
         {\scriptsize 
         \begin{longtable}{p{0.2\textwidth}p{0.8\textwidth}} 
         \hiderowcolors% 
         \toprule 
         Source & Excerpt \\ 
         \midrule 
         \showrowcolors% 
         \endfirsthead% 
         \hiderowcolors% 
         \toprule 
         Source & Excerpt \\ 
         \midrule 
         \showrowcolors% 
         \endhead% 
         \midrule 
         \endfoot% 
         \hiderowcolors% 
         \bottomrule 
         \showrowcolors% 
         \endlastfoot% 
         
Gibler, Douglas M. 2018. International Conflicts, 1816–2010: Militarized Interstate Dispute Narratives. Vol. 1 Rowman \& Littlefield. & The Second Greco-Turkish War, fought between 1919 and 1922, followed the dissolution of the Ottoman Empire and long-standing Greek territorial claims against the successor state. Greece used the occasion of occupation by the victorious Allies to move on centuries' old Turkish territory that had large Greek populations. The Greek move was met by resistance and ultimately a war, which the Turks won. A Greek retreat prompted the Armistice of Mudanya on October 11, 1922. The Treaty of Lausanne, signed in July 1923, updated the previous Sevres treaty and gave Turkey a more favorable division of the former Ottoman Empire than it was previously afforded under the Sevres treaty. The Lausanne agreement mostly, but not perfectly, mirrors the existing territorial boundaries between Greece and Turkey. (p. 357)
\\
Britannica. 2016. “Greco-Turkish Wars.” Accessed August 19, 2023. https://www.britannica.com/event/Greco-Turkish-wars. & The second war occurred after World War I, when the Greeks attempted to extend their territory beyond eastern Thrace (in Europe) and the district of Smyrna (İzmir; in Anatolia). These territories had been assigned to them by the Treaty of Sèvres, August 10, 1920, which was imposed upon the weak Ottoman government. In January 1921 the Greek army, despite its lack of equipment and its unprotected supply lines, launched an offensive in Anatolia against the nationalist Turks, who had defied the Ottoman government and would not recognize its treaty. Although repulsed in April, the Greeks renewed their attack in July and advanced beyond the Afyonkarahisar-Eskişehir railway line toward Ankara. The Turks, however, commanded by the nationalist leader Mustafa Kemal (Kemal Atatürk), defeated them at the Sakarya River (August 24–September 16, 1921).
\\
Pizzo, David 2011. “Greco-Turkish War (1919–1922).” In The Encyclopedia of War, edited by Gordon Martel. & The [Greco-Turkish War] represented both the final stage of disintegration of the Ottoman Empire and the culmination ofthe Greek “Megali [Great] Idea” of uniting all Greeks in the eastern Mediterranean under a single Greek state. Early Greek successes seemed to offer the prospect of a pan-Hellenic Greek state on both sides of the Aegean, but the Turkish revolutionaries' military successes of 1921–1922 turned victory into catastrophe, resulting in the collapse of Greek irredentist dreams, large refugee flows, and the destruction of both the Greek communities in Anatolia and Turkish communities in Greece. For the Turkish national movement, on the other hand, the war represented a crucial phase of their war of independence. (...).
[Greece] had been party to discussions among the Allies about the division of the post-war Ottoman Empire, as the Entente powers sought to balance their various and competing claims to Ottoman territory. [The Greek Prime Minister] (...) pushed very hard (...) for a Greek military occupation of western Anatolia. (...). The British soon came to view this as a preferable outcome (...). Britain (...) hoped to impose a harsh settlement on the Ottomans and prevent the victory of the nationalists without directly committing its own forces (...). (...). (...) [The] Turks (...) saw the conflict with the Greeks as a struggle for their very existence.
\\
\end{longtable}}
\clearpage
         \textbf{Russo-Polish} 

          \underline{Onset:} 1919         

\underline{Reasons:}
\begin{itemize}
\item Nationalism
\item Border Clashes
\item Economic, Long-Run
\end{itemize}

         \underline{Sources:} 
         \rowcolors{3}{white}{light-gray} 
         {\scriptsize 
         \begin{longtable}{p{0.2\textwidth}p{0.8\textwidth}} 
         \hiderowcolors% 
         \toprule 
         Source & Excerpt \\ 
         \midrule 
         \showrowcolors% 
         \endfirsthead% 
         \hiderowcolors% 
         \toprule 
         Source & Excerpt \\ 
         \midrule 
         \showrowcolors% 
         \endhead% 
         \midrule 
         \endfoot% 
         \hiderowcolors% 
         \bottomrule 
         \showrowcolors% 
         \endlastfoot% 
         
Gibler, Douglas M. 2018. International Conflicts, 1816–2010: Militarized Interstate Dispute Narratives. Vol. 1 Rowman \& Littlefield. & Following World War I, Poland, again independent after more than a century of occupation, sought the territory it originally possessed in 1772. The Soviets were distracted with the Russian Civil War, and this prompted Poland to begin military operations at the beginning of January 1919. They established a 300-mile front by March and took their first city in April. By the end of the year Polish armies had recovered all the territory Russia gained by partition in 1795, and Poland and the Soviets secretly agreed to a truce. By April 1920 the Soviets had removed Lieutenant General Denikin's armies as a threat to Poland. However, diplomatic negotiations soon broke down, and Poland launched additional military operations on April 25. Polish troops then captured Kiev on May 7, but a Soviet counterattack compelled Poland to evacuate Ukraine on June 11. By the middle of August, the Soviets were 20 miles from Warsaw, then a counterattack by Poland pushed Soviet troops out of ethnic Poland toward the end of the month. On September 18, the two sides agreed to a ceasefire, and on March 18, 1921, both settled the dispute with the Treaty of Riga, which defined the boundary between them. (p. 287)
\\
Britannica. 2023. “Russo-Polish War.” Accessed August 20, 2023. https://www.britannica.com/event/Russo-Polish-War-1919-1920. & Russo-Polish War, also called Polish-Soviet War, (1919–20), military conflict between Soviet Russia and Poland. It was the result of the German defeat in World War I, Polish nationalism sparked by the re-creation of an independent Polish state, and the Bolsheviks' determination to carry the gains they had achieved during the Russian Civil War to central Europe. The decisive Polish victory resulted in the establishment of the Russo-Polish border that existed until 1939.
\\
Stone, David R. 2011. “Russian Civil War (1917–1920).” In The Encyclopedia of War, edited by Gordon Martel. & The Russo-Polish War was unleashed by the end of World War I for control of the east European borderlands between Polish and Soviet territory. The war grew out of the collapse of the three empires of eastern Europe at the end of World War I: Russia in 1917, then Germany and Austria–Hungary in 1918. The resulting power vacuum created an opportunity for smaller nations of eastern Europe to create independent states. The new Polish republic emerged from territories and populations previously part of all three empires, with an army cobbled together from multiple sources, including Polish formations of the now-defunct empires,under the leadership of former socialist,now Polish nationalist, Jozef Klemens Piłsudski (1867–1935). In Russia, the Bolshevik Party under Lenin seized power in late 1917 and then fought a civil war to establish its own authority over what had been the Russian Empire.
\\
Stone, David R. 2011. “Russian Civil War (1917–1920).” In The Encyclopedia of War, edited by Gordon Martel. & Over half the population of the Russian Empire was ethnically non-Russian and took advantage of the breakdown of control in the center to secede. (...). A final reckoning with the White forces in the Crimea, now commanded by Baron Peter Wrangel in place of the defeated and disgraced Denikin, was delayed by the outbreak of the Russo-Polish War. Once the chance of an actual White victory had passed, with its danger of restoration of the Russian Empire, the newly created Polish state attacked east in spring 1920 to bring as much as possible of the Polish–Russian borderlands in Belorussia and Ukraine under its control. This Russo-Polish War pulled manpower and supplies temporarily away from the Red fight against Wrangel. (...). Only on the western frontier, where Poland, Finland, and the Baltic states established effective governments and enjoyed western support, did the Bolsheviks fail to regain large territories that had been part of the Russian Empire.
\\
\end{longtable}}
\clearpage
         \textbf{Hungarian Adversaries} 

          \underline{Onset:} 1919         

\underline{Reasons:}
\begin{itemize}
\item Nationalism
\item Power Transition or Security Dilemma
\item Religion or Ideology
\end{itemize}

         \underline{Sources:} 
         \rowcolors{3}{white}{light-gray} 
         {\scriptsize 
         \begin{longtable}{p{0.2\textwidth}p{0.8\textwidth}} 
         \hiderowcolors% 
         \toprule 
         Source & Excerpt \\ 
         \midrule 
         \showrowcolors% 
         \endfirsthead% 
         \hiderowcolors% 
         \toprule 
         Source & Excerpt \\ 
         \midrule 
         \showrowcolors% 
         \endhead% 
         \midrule 
         \endfoot% 
         \hiderowcolors% 
         \bottomrule 
         \showrowcolors% 
         \endlastfoot% 
         
Gibler, Douglas M. 2018. International Conflicts, 1816–2010: Militarized Interstate Dispute Narratives. Vol. 1 Rowman \& Littlefield. & Hungary was a new interstate system member after World War I, which eliminated Austria-Hungary (to which Hungary was a junior partner) from the international system. However, Hungary's new independence was not without peril. Hungary was under pressure by the Allied powers, and they expected proper reparation for Hungary's previous association with the Central Powers. Conservative sentiments and Communists alike made the transition difficult for new prime minister Mihaly Karolyi. Frustrated, Karolyi resigned his post, which opened the door for the Communists to take power and declare the Hungarian Soviet Republic in March 1919. Romania, Czechoslovakia, and Yugoslavia used the opportunity to press their demands for Hungarian territory, though the World War I's conclusion already gave these states de facto occupation of what they were demanding. With the blessing of the French and Italians, Romania attacked Hungary on April 16, 1919. Czechoslovakia joined them on April 26, though a Hungarian counteroffensive against Czechoslovakia led to the occupation of eastern Slovakia and the temporary emergence of the Slovak Socialist Republic on June 16. The Allies pressured all sides to end the conflict, which ultimately came when the Romanians and Czechs occupied Budapest on August 14. As a result, Hungary capitulated to the ongoing Treaty of Trianon discussions. The treaty was concluded on June 4, 1920, effectively determining what Hungary's punishment would be for World War I. The treaty provided a strong punishment, too. Beyond the war reparations paid to the Allies, Hungary lost approximately three-fourths its prewar territory as the Hungarian part of Austria-Hungary. It also lost well over half of its population. These were divvied up among the allies in the Balkans. (p. 218)
\\
University of Central Arkansas. Accessed August 19, 2023. https://uca.edu/politicalscience/home/research-projects/dadm-project/europerussiacentral-asia-region/hungary-1918-present/ & Hungary declared its independence from the Austria-Hungary Empire on October 17, 1918. Prime Minister Sandor Wekerle resigned on October 24, 1918, and the Hungarian National Council (HNC) headed by Count Mihaly Karolyi was established on October 25, 1918. Government troops and demonstrators clashed in Budapest on October 28, 1918, resulting in the deaths of three individuals. [...] The HNC proclaimed the Hungarian People's Republic on November 16, 1918. Bela Kun established the Communist Party of Hungary (CPH) on November 24, 1918. Count Karolyi was elected provisional president by the HNC on January 11, 1919, and Dezso Berinkey was appointed prime minister on January 18, 1919. [...] Romanian troops intervened against the government of Bela Kun in southern Hungary beginning on April 16, 1919. Opponents of the government of Bela Kun, including Count Gyula Karolyi, Count Istvan Bethlen, and Admiral Miklos Horthy de Nagybanya, established a rival government in Arad, Romania (and later Szeged, Hungary) on May 5, 1919. The RGC collapsed on August 1, 1919, and Bela Kun fled into exile to Vienna
\\
Lojko, Miklos. 2025. “Treaty of Trianon (1920).” In The Encyclopedia of Diplomacy, edited by G. Martel. & The historic kingdom of Hungary had been part of the dual monarchy of Austria-Hungary since the so-called Compromise of 1867. (...). Once the world war had broken out, aims and opportunities changed for the nationalities of Hungary, who made up roughly 45 percent of the population. Leaders of the Entente powers identified the liberation of oppressed nationalities in Austria-Hungary and Ottoman Turkey, the multinational empires who fought against the Entente, as one of their war aims. The idea of self-determination was not only a liberal commitment, but also a strategic tool in the circumstances: encouragement to millions of subjects of Austria-Hungary and Ottoman Turkey to break ranks, desert their political masters, and join the Entente in the war. (...).
While no “gray zones” corresponding to Seton-Watson's scheme were set up systematically, the military occupation zone between Hungary proper and Transylvania communicated in the so-called Vix note to Mihály Károlyi's interim government in March 1919 was of a similar character. The Vix note, foreshadowing the final borders in the east, caused the Károlyi administration to resign and hand over executive power to Béla Kun's communists on March 21, 1919. Merging with mainstream Social Democrats, the communists declared a Soviet Republic which lasted 133 days, during which time the Hungarian Red Army, a well-organized force, managed to temporarily reverse the occupation of Hungarian territory by armies of the neighboring countries, especially in the north. (...).
Due to a combination of diplomatic maneuvering by the British military representative in Vienna, Sir Thomas Montgomery-Cuninghame, and the operations of the Romanian army, implicitly encouraged by the peace conference, Béla Kun's communist dictatorship collapsed on August 1, 1919. Following intensive mediation conducted by emissaries of the peace conference, the Romanian army was made to withdraw from Budapest, Admiral Miklós Horthy was accepted as regent, and democratic elections to the national assembly were scheduled.
\\
\end{longtable}}
\clearpage
         \textbf{Franco-Turkish} 

          \underline{Onset:} 1919         

\underline{Reasons:}
\begin{itemize}
\item Nationalism
\item Power Transition or Security Dilemma
\end{itemize}

         \underline{Sources:} 
         \rowcolors{3}{white}{light-gray} 
         {\scriptsize 
         \begin{longtable}{p{0.2\textwidth}p{0.8\textwidth}} 
         \hiderowcolors% 
         \toprule 
         Source & Excerpt \\ 
         \midrule 
         \showrowcolors% 
         \endfirsthead% 
         \hiderowcolors% 
         \toprule 
         Source & Excerpt \\ 
         \midrule 
         \showrowcolors% 
         \endhead% 
         \midrule 
         \endfoot% 
         \hiderowcolors% 
         \bottomrule 
         \showrowcolors% 
         \endlastfoot% 
         
Gibler, Douglas M. 2018. International Conflicts, 1816–2010: Militarized Interstate Dispute Narratives. Vol. 2 Rowman \& Littlefield. & At the center of the Franco-Turkish War was control of Cilicia, which was in southeastern Anatolia bordering on Syria. During World War I the Great Powers signed secret agreements which, among other things, gave France a sphere of influence in Cilicia, and France moved into Cilicia in November 1919. Turkey did not recognize French control of Cilicia and in January 1920 attacked the French and their Armenian supporters in Maras Province, pushing the French out within three weeks. On May 23, Turkey and France signed an armistice, but the French broke it with attacks on June 10. Fighting continued until October 20, 1921, when France and Turkey signed the Ankara Accord. (pp. 536-537)
\\
Britannica. 2023. “The Nationalist Movement and the War for Independence.” Encyclopædia Britannica. Accessed August 19, 2023. https://www.britannica.com/biography/Kemal-Ataturk/The-nationalist-movement-and-the-war-for-independence. & On February 8, 1919, the French general Franchet d'Espèrey entered the city in a spectacle compared to the entrance of Mehmed the Conqueror in 1453—but this time signifying that Ottoman sovereignty over the imperial city was over. The Allies made plans to incorporate the provinces of eastern Anatolia into an independent Armenian state.
\\
Compton, Thomas E. 1922. “The French Campaign of 1920–21 in Cilicia.” Royal United Services Institution Journal 67 (465): 68–79. & Why, with all their other difficulties, the French were in Cilicia at all is somewhat surprising; but it would seem to have been partly due to the demands of the Armenians, who furnished three out of the four battalions occupying the province, when, after the armistice, Lord Allenby (in accordance with the Sykes-Picot agreement of 1916) divided the whole country, west of the Euphrates and Syrian Desert, into three spheres of influence (...).
\\
Güçlü, Yücel. 2001. “The Struggle for Mastery in Cilicia: Turkey, France, and the Ankara Agreement of 1921.” The International History Review 23 (3): 580–603.
 & The armistice of Mudros left the victorious Entente powers poised to partition almost all of die Ottoman Empire with die aim of extinguishing it as an independent international actor. The Allies envisaged that the future Turkey, confined to Istanbul and central and northern Anatolia, would have few resources and little freedom of action in the economic sphere. Aldiough the Allies' plans were not formalized until die treaty of Sèvres, made with the Imperial Ottoman government on 10 August 1920, several of the Allied states had tried to arrange faits accomplis. In die spring of 1919, Greek troops landed at Izmir and Italian troops at Antalya, while the French occupied Cilicia, and British, French, and Italian troops were stationed at the Straits.
\\
\end{longtable}}
\clearpage
         \textbf{Lithuanian-Polish} 

          \underline{Onset:} 1920         

\underline{Reasons:}
\begin{itemize}
\item Nationalism
\item Power Transition or Security Dilemma
\item Religion or Ideology
\item Border Clashes
\end{itemize}

         \underline{Sources:} 
         \rowcolors{3}{white}{light-gray} 
         {\scriptsize 
         \begin{longtable}{p{0.2\textwidth}p{0.8\textwidth}} 
         \hiderowcolors% 
         \toprule 
         Source & Excerpt \\ 
         \midrule 
         \showrowcolors% 
         \endfirsthead% 
         \hiderowcolors% 
         \toprule 
         Source & Excerpt \\ 
         \midrule 
         \showrowcolors% 
         \endhead% 
         \midrule 
         \endfoot% 
         \hiderowcolors% 
         \bottomrule 
         \showrowcolors% 
         \endlastfoot% 
         
Gibler, Douglas M. 2018. International Conflicts, 1816–2010: Militarized Interstate Dispute Narratives. Vol. 1 Rowman \& Littlefield. & This Polish-Lithuanian conflict occurred shortly after both states became independent in the aftermath of World War I. Poland's new independence led Jozef Pilsudski, then leader of the fledgling republic, to attempt restoration of the former status of the Polish-Lithuanian Commonwealth. Lithuania, now independent as well, felt any union with Poland as it had previously would be a simple subjugation and loss of cultural autonomy and refused all overtures. Poland then opted to restore its prepartition territories by force. Vilnius was the heart of the conflict. The predominantly Polish city in Lithuania was under siege by Russia's Red Army, who moved west after Germany retreated from the area. Poland, who despised the Russians for their role in Poland's elimination from the interstate system, moved against Russia. Lithuania, for whom Vilnius was a capital, was caught in the middle of this power struggle. Lithuania was forced to withdraw from Vilnius on January 3, 1919, under threat from the Soviets (see MID\#2603). The Lithuanian government relocated to Kaunas. Russian occupation was tenuous and both the Poles and Lithuanians seized the opportunity to expel the Bolsheviks from Vilnius. The Poles got to Vilnius first. Arriving on April 18, the Poles occupied Vilnius, effecting the call to occupy the city and unify it with Poland that the Sejm passed on April 4. Pilsudski had complete control of Vilnius by April 21, prompting the Lithuanians to declare Poland as an invading force that did not cooperate with the government in Kaunas. (pp. 290-291)
\\
Britannica. 2023. “Vilnius Dispute.” Accessed August 20, 2023. https://www.britannica.com/event/Vilnius-dispute. & Although the new Lithuanian government established itself at Vilnius in late 1918, it evacuated the city when Soviet forces moved in on January 5, 1919. A few months later Polish forces drove the Red Army out of Vilnius and occupied it themselves (April 20, 1919). The Lithuanians rejected the demands of the Polish chief of state, Józef Piłsudski, for union with Poland, and hostilities were avoided only by the Allies' creation of a demarcation line (the Foch Line) to separate the armies of the two countries. Vilnius was left on the Polish side of the line.
\\
Balkelis, Tomas. 2018. War, Revolution, and Nation-Making in Lithuania, 1914–1923. New York: Oxford University Press. & After the German defeat, the Polish–Lithuanian clash over Vilnius became immediately evident. Lithuania became a territory of violence through many different forces: the German army and later Freikorps, Bolshevik troops, Polish soldiers and the first units of the Lithuanian army. Balkelis gives an overview of the ethnic groups that were taking part in the clash. There were ‘Two visions of Lithuania' (78) that had a Socialist-Bolshevik and a national concept as their core. The different aims often intermingled: ethnic conflicts leading to international disputes and confrontations. But there also was a connection between national and social upheavals among the Lithuanian population.
\\
\end{longtable}}
\clearpage
         \textbf{Manchurian} 

          \underline{Onset:} 1929         

\underline{Reasons:}
\begin{itemize}
\item Nationalism
\item Economic, Long-Run
\item Domestic Politics
\item Revenge/Retribution
\end{itemize}

         \underline{Sources:} 
         \rowcolors{3}{white}{light-gray} 
         {\scriptsize 
         \begin{longtable}{p{0.2\textwidth}p{0.8\textwidth}} 
         \hiderowcolors% 
         \toprule 
         Source & Excerpt \\ 
         \midrule 
         \showrowcolors% 
         \endfirsthead% 
         \hiderowcolors% 
         \toprule 
         Source & Excerpt \\ 
         \midrule 
         \showrowcolors% 
         \endhead% 
         \midrule 
         \endfoot% 
         \hiderowcolors% 
         \bottomrule 
         \showrowcolors% 
         \endlastfoot% 
         
Gibler, Douglas M. 2018. International Conflicts, 1816–2010: Militarized Interstate Dispute Narratives. Vol. 2 Rowman \& Littlefield. & The Sino-Soviet conflict of 1929 was fought over ownership of the Chinese Eastern Railway in Manchuria. The dispute followed from the Chinese attempt under the aegis of its new nationalist government to rectify its previous treaties that favored the Soviets, especially the latest agreements in 1924 and 1925 that reestablished Soviet control of the rails. When negotiations failed, China attempted to seize control of the rails in a raid. The attempt failed and only prompted further Soviet preparedness in the region. Subsequent raids and attacks failed as well. Soviet pressure ultimately led to a split between the republican government in Nanjing and the Manchurian government in Mukden. The nationalist government in China conceded to the status quo ante bellum on December 22, 1929. (p. 759)
\\
Siegelbaum, Lewis. “Chinese Railway Incident.” Michigan State University. Accessed August 20, 2023. https://soviethistory.msu.edu/1929-2/chinese-railway-incident/ & The year 1929 found the Soviet Union's fortunes in the Far East at a low ebb. Two years earlier, the Kuomintang under Chiang Kai-shek had turned against and crushed their erstwhile allies, the Chinese Communists, and severed diplomatic relations with Moscow. Consolidating their position as claimants to rule China, the Kuomintang sought to extend their authority to Manchuria, presenting a real threat to Soviet interests in the region. The Soviet presence in Manchuria derived from the Chinese Eastern Railroad (CER), the ownership of which had passed into Soviet hands with the overthrow of the tsarist government and the October Revolution. Recognition of the Soviet Union's title to the railroad had been secured via treaties in 1924 with both the Chinese government in Peking and the government of the warlord, Chang Tso-lin, in Mukden. Yet, the existence of a substantial “White” Russian community in the Manchurian city of Harbin, the activities of Soviet consular officials and commercial agents, and the fact that as of 1929, 75 per cent of the railroad's employees were Russians who held all the controlling posts constituted an affront to the Kuomintang government in Nanking and its claim to represent Chinese sovereignty.
\\
Lee, Chong-Sik. 1983. “Sino-Soviet Conflict of 1929 and the Li Li-San Line.” In Revolutionary Struggle in Manchuria: Chinese Communism and Soviet Interest, 1922–1945. 1st ed. Berkeley: University of California Press. & The polemic between the Manchurian Provincial Committee and the Chinese Communist Party (CCP) headquarters in 1927 and 1928 clearly showed that, in order to protect the Soviet interest over the Chinese Eastern Railway, the leaders of the CCP had decided to hold the revolutionary movement in Manchuria in abeyance. Its organizational movements being in the doldrums, the party also decided to tone down its attack against imperialism. Provocations from imperialism were to be evaded until the party could receive the support of the international revolutionary proletariat (i.e., the Soviet Union) under conditions most advantageous to the revolution. It was natural, therefore, that the CCP's policies toward Manchuria would change when the Soviet perspective on Manchuria changed. The Soviet policy toward Manchuria did change in 1929 because of Chang Hsiieh-liang's attempt to take over the Chinese Eastern Railway (CER). Indeed, the conflict over the CER and the corollary shift in the Soviet policy toward China was to affect the CCP's development throughout China between 1929 and 1931. The events in Manchuria precipitated the emergence of radical strategies under Li Li-san, which came to be known as the Li Li-san line. (p. 95)
\\
Walker, Michael. 2021. The 1929 Sino-Soviet War: The War Nobody Knew. Lawrence, KS: University Press of Kansas. & The 1929 Sino-Soviet conflict was a short and bloody war fought over the jointly operated Chinese Eastern Railroad in China's Northeast between two powers mostly relegated to the dustbin of history, the Republic of China and the Union of Soviet Socialist Republics. A modern limited war, it proved to be the largest military clash between China and a Western power ever fought on Chinese soil. Over 300,000 soldiers, sailors, and aviators served in the war, although only a part participated in the heavy fighting. (...). The conflict was the first major combat test of the reformed Soviet Red Army (...) and ended with the mobilization and deployment of 156,000 troops to the Manchurian border. (...). A path to war was created when Chiang Kai-shek and Chang Hsueh-liang miscalculated, both diplomatically and militarily, as they viewed the Soviets as politically isolated and militarily weak and were convinced that the time was right to reassert full authority over the CER. For the Soviets, Stalin dominated the action, and he saw war, not negotiations, as the preferred option. Once Stalin approved the large-scale offensive, the Soviet Red Army unexpectedly scored a decisive victory, disproving the assumption that it was incapable of fighting a modern war. (pp. 1-2)
\\
\end{longtable}}
\clearpage
         \textbf{Second Sino-Japanese} 

          \underline{Onset:} 1931         

\underline{Reasons:}
\begin{itemize}
\item Nationalism
\item Economic, Long-Run
\item Economic, Short-Run
\end{itemize}

         \underline{Sources:} 
         \rowcolors{3}{white}{light-gray} 
         {\scriptsize 
         \begin{longtable}{p{0.2\textwidth}p{0.8\textwidth}} 
         \hiderowcolors% 
         \toprule 
         Source & Excerpt \\ 
         \midrule 
         \showrowcolors% 
         \endfirsthead% 
         \hiderowcolors% 
         \toprule 
         Source & Excerpt \\ 
         \midrule 
         \showrowcolors% 
         \endhead% 
         \midrule 
         \endfoot% 
         \hiderowcolors% 
         \bottomrule 
         \showrowcolors% 
         \endlastfoot% 
         
Gibler, Douglas M. 2018. International Conflicts, 1816–2010: Militarized Interstate Dispute Narratives. Vol. 2 Rowman \& Littlefield. & The Second Sino-Japanese War fought between 1931 and 1933 largely concerned Japan's push west for hegemony in northeast China (Manchuria). The opportunity to make further incursions into Manchuria presented itself to Japan around 1929 and 1930, when China became involved in conflicts with Russia (see MID\#041) and civil conflicts with Chinese warlords in the west and the growing Communist faction. The casus belli was the Mukden incident of September 18, 1931, when a section of the Japanese-owned South Manchuria Railway was dynamited under suspicious circumstances. While potential Japanese planning of this explosion remains a debate, Japan certainly used the event as a pretext to begin a conflict against China. The conflict quickly escalated to war. Protests ensued from observing states, like Britain, France, and the United States. Two years into the conflict, Japan was able to overwhelm the Chinese troops and coerce them into an armistice that conceded defeat. The fighting stopped on May 22, 1933, and the armistice was signed on May 31 of the same year. Conflicts between China and Japan would later resume in 1937, ultimately getting absorbed into World War II. Japan left the League of Nations amid the protests of this war. (p. 811)
\\
Britannica. 2022. “Mukden Incident.” Accessed August 20, 2023. https://www.britannica.com/event/Mukden-Incident. & Mukden Incident, (September 18, 1931), also called Manchurian Incident, seizure of the Manchurian city of Mukden (now Shenyang, Liaoning province, China) by Japanese troops in 1931, which was followed by the Japanese invasion of all of Manchuria (now Northeast China) and the establishment of the Japanese-dominated state of Manchukuo (Manzhouguo) in the area. Most observers believe the incident was contrived by the Japanese army, without authorization of the Japanese government, to justify the Japanese invasion and occupation that followed.
\\
Bjorge, Gary J. 2011. “China, Invasion of (1931, 1937–1945).” In The Encyclopedia of War, edited by Gordon Martel. & Japan's invasions of China during the 1930s initiated and then expanded the largest war fought between two countries during the twentieth century. For Japan, this was a war fought for resources and geopolitical position. For China, it was a war for national survival. (...). These invasions were an extension of Japan's long-standing national security goal of establishing control over adjacent areas on the Asian continent. Such control was considered necessary in order to counter possible military threats and ensure access to the natural resources needed to guarantee Japan's economic independence. Manchuria (...) became the focal point of Japanese expansionist aspirations because of its mineral wealth, rich farmland, and potential value as a buffer shielding Korea from both China and Russia, later the Soviet Union. (...). 
A bomb was to be exploded along the railroad south of Mukden to provide an excuse for Kwantung army units to seize the nearby Manchurian army barracks and the arsenal in the city. Once those objectives were taken, the area under Kwantung army control was to be gradually expanded until all of Manchuria was occupied. (...). China's response to the Kwantung army's attacks in and around Mukden was a policy of non-resistance and an appeal to the League of Nations for support. At this time, Jiang Jieshi, the leader of the Guomindang (Nationalist Party) (...)  did notwant a war with Japan. Jiang was determined to destroy the communists and achieve internal unity before confronting the Japanese.
\\
Doenecke, Justus D. 2025. “Manchurian Crisis (1931–33).” In The Encyclopedia of Diplomacy, edited by Gordon Martel. & On September 19, 1931, Japan's Guandong (Kwantung) Army, a semi-autonomous force, seized Manchuria's capital city of Shenyang (Mukden), claiming that three companies of Chinese soldiers had destroyed tracks of the Japanese-owned South Manchuria Railway at a village three miles away. Japan's embarrassed foreign office pledged that its country would soon withdraw its troops, but it quickly became apparent that the promise would not be honored. As the Japanese attack was the first military action by a Great Power since the 1918 Armistice, it caused global consternation.(...). By 1931, Manchuria held over half Japan's foreign investment, some 1.2 billion yen, as well as the South Manchurian Railroad. In addition to possessing 690 miles of track, the railway owned ironworks, coal mines, harbors, police forces, agricultural stations, public utilities, schools, and libraries. Furthermore, the Japanese economy had become increasingly dependent upon such Manchurian resources as lumber. Fifty percent of Japan's food and pig iron came from Manchuria; a third of Manchuria's coal crossed the Yellow Sea to Japan. In addition, a quarter of a million Japanese subjects, mostly Koreans, lived there, and it absorbed about 40 percent of Japan's China trade. In time total control could lead the Japanese Empire to complete economic independence. Because of the worldwide depression, by summer 1931 Japan's economic situation was becoming increasingly desperate. Compounding the problem was continued population growth: 65 million Japanese were squeezed into an area smaller than the American state of Texas. A land mass the size of France and Germany combined, therefore, appeared most promising. Yet, although Japan had long treated Manchuria as if it were its colony, culturally and demographically 28 of its 30 million people were Chinese; all Japanese privileges were acknowledged in formal treaties with the Chinese government. The regional warlord Zhang Zoulin (Chang Tso-lin), whom the Japanese assassinated in 1928, had originally been Japan's client, but his son-in-law and heir, Zhang Xueliang (Chang Hsueh-liang), allied himself with Nationalist leader Jiang Jieshi (Chiang Kai-shek). To counter the Japanese, Jiang's government persuaded several million Chinese to emigrate to Manchuria, there to spur economic development. In addition, their presence would emphasize China's legal, if always tenuous, ownership of the province.
\\
Parker, Geoffrey, ed. 2020. The Cambridge History of Warfare. 2nd ed. Cambridge: Cambridge University Press. & By closing their markets during the Great Depression, the Western powers encouraged aggressive Japanese policies towards the Asian mainland and helped Japanese militarists to gain power. In 1931, Japanese army units seized Manchuria without Tokyo's authorization. Six years later that army initiated an undeclared war against China, and Japanese troops soon controlled China's coastal regions and most of the important Chinese cities, leaving a trail of atrocities in their wake.  (p. 356)
\\
\end{longtable}}
\clearpage
         \textbf{Chaco} 

          \underline{Onset:} 1932         

\underline{Reasons:}
\begin{itemize}
\item Border Clashes
\item Economic, Long-Run
\item Revenge/Retribution
\end{itemize}

         \underline{Sources:} 
         \rowcolors{3}{white}{light-gray} 
         {\scriptsize 
         \begin{longtable}{p{0.2\textwidth}p{0.8\textwidth}} 
         \hiderowcolors% 
         \toprule 
         Source & Excerpt \\ 
         \midrule 
         \showrowcolors% 
         \endfirsthead% 
         \hiderowcolors% 
         \toprule 
         Source & Excerpt \\ 
         \midrule 
         \showrowcolors% 
         \endhead% 
         \midrule 
         \endfoot% 
         \hiderowcolors% 
         \bottomrule 
         \showrowcolors% 
         \endlastfoot% 
         
Gibler, Douglas M. 2018. International Conflicts, 1816–2010: Militarized Interstate Dispute Narratives. Vol. 1. Rowman \& Littlefield. & The Chaco War between Bolivia and Paraguay contested 116,000 square miles of territory, rumored to have large oil deposits and giving access to the Paraguay River, Bolivia's only access to the sea. Bolivia and Paraguay mobilized to fight over the Chaco Boreal in 1928 and 1929. The sides never resolved the issue. On June 15, 1931, Paraguayan troops advanced into Bolivian-controlled Chaco, and Paraguayan gunboats advanced up the river. On June 18, 1932, Bolivia seized a Paraguayan fort. Paraguay responded less than a month later with its own seizure of a Bolivian fort Bolivian forces countered with the capture of Paraguayan-controlled territory. Paraguayan forces then pushed the Bolivian troops back past their original lines. Fighting continued nearly through 1933. A truce took effect in December 1933, but fighting broke out again on January 6, 1934. (...).
(pp. 173-174)
\\
Britannica. 2023. “Chaco War.” Encyclopædia Britannica. Accessed August 19, 2023. https://www.britannica.com/event/Chaco-War. & The conflict stemmed from the outcome of the War of the Pacific (1879–84), in which Chile defeated Bolivia and annexed that country's entire coastal region. Thereafter, Bolivia attempted to break out of its landlocked situation through the Río de La Plata system to the Atlantic coast; athwart that route lay the Gran Chaco, which the Bolivians thought had large oil reserves. [...] Bolivia attempted to break out of its landlocked situation through the Río de La Plata system to the Atlantic coast; athwart that route lay the Gran Chaco, which the Bolivians thought had large oil reserves.
\\
Farcau, Bruce W. 2011. “Chaco War (1932–1935).” In The Encyclopedia of War, edited by Gordon Martel. & As in many conflicts in the formerly colonial world, the roots of the Chaco War lay in the casual demarcation of borders between what had once been provinces of a single empire. The boundary between what would become Bolivia and Paraguay was rather arbitrarily set at the confluence of the Paraguay and Pilcomayo rivers, which caused no trouble since the area comprises some of the most inhospitable terrain onthe planet (...). Most of the Chaco remained both uninhabited and unexplored well into the twentieth century. However, nationstates have “patrimony” to defend, and when outposts of the two rivals began to bump up against each other in the wilderness, conflict ensued. Bolivia had vague hopes of establishing a connection (...) to the sea, having lost its Pacific Coast to Chile half a century earlier (...). Paraguay had begun to settle the eastern portion of the Chaco, which began right across the river from the capital of Asuncion, and to sell the marginal agricultural land to pay off the massive debts incurred during the War of the Triple Alliance in which Paraguay took on (and lost to) Brazil, Argentina, and Uruguay. Consequently, both nations began pushing forward their outposts in an effort to lay claim to as much real estate as possible. (...) [For] Bolivia this would essencially be a colonial war (...). (...). For Paraguay, conversely, the Chaco War was seen as an existential threat.
\\
\end{longtable}}
\clearpage
         \textbf{Saudi-Yemeni} 

          \underline{Onset:} 1934         

\underline{Reasons:}
\begin{itemize}
\item Power Transition or Security Dilemma
\item Border Clashes
\end{itemize}

         \underline{Sources:} 
         \rowcolors{3}{white}{light-gray} 
         {\scriptsize 
         \begin{longtable}{p{0.2\textwidth}p{0.8\textwidth}} 
         \hiderowcolors% 
         \toprule 
         Source & Excerpt \\ 
         \midrule 
         \showrowcolors% 
         \endfirsthead% 
         \hiderowcolors% 
         \toprule 
         Source & Excerpt \\ 
         \midrule 
         \showrowcolors% 
         \endhead% 
         \midrule 
         \endfoot% 
         \hiderowcolors% 
         \bottomrule 
         \showrowcolors% 
         \endlastfoot% 
         
Gibler, Douglas M. 2018. International Conflicts, 1816–2010: Militarized Interstate Dispute Narratives. Vol. 2 Rowman \& Littlefield. & The Saudi-Yemeni War of 1934 was fought over the extent of the boundaries claimed by Ibn Saud and provided for the emergence of the modern Saudi Arabian kingdom. The local Yemeni noble, Imam Yahya Muhammad Hamid ed-Din, tried to formalize the extent of his boundaries after the dissolution of the Ottoman Empire as the Saudis began encroaching on Yemeni-controlled territories. Several autonomous Yemeni regions joined the Yemeni King, but the Saudis still won convincingly. The Treaty of Taif formalized the borders between Saudi Arabia and Yemen, giving the Saudi king the disputed territory of Najran. (p. 702)
\\
Britannica. 2023. “The Kingdom of Saudi Arabia.” Accessed August 20, 2023. https://www.britannica.com/place/Saudi-Arabia/The-Kingdom-of-Saudi-Arabia. & In 1934 Ibn Saud was involved in war with Yemen over a boundary dispute. An additional cause of the war was Yemen's support of an uprising by an Asiri prince against Ibn Saud. In a seven-week campaign, the Saudis were generally victorious.
\\
Schofield, Richard. 1996. “The Last Missing Fence in the Desert: The Saudi‐Yemeni Boundary.” Geopolitics and International Boundaries, 1(3): 247–299. & During the inter-war years Britain was responsible for the foreign affairs and defence of all Saudi Arabia's neighbours in the peninsula except for one, the Imamate of Yemen. As a consequence, Britain had virtually nothing to do with the evolution of the boundary between Saudi Arabia and the Imamate, which was defined in its western reaches (from near Midi on the Red Sea to Najran) by the May 1934 Treaty of 'Islamic Friendship and Brotherhood' (Taif treaty). That the modern state territories of Saudi Arabia and Yemen met at all on the eastern shores of the Red Sea was only as a result of a series of agreements reached between Ibn Saud and the Idrisi of Asir during the 1920s. An agreement of 1920 had seen the north-eastern half of Asir, formerly an Ottoman administrative unit linked indirectly to the wilayat of Yemen, incorporated into the expanding Najdi state. This would later become the Saudi province of Asir Surat. Saudi-Idrisi agreements of 1926 and 1930 saw Ibn Saud extend protectorate facilities over and then formally annex the remaining portions of the Idrisi's territory, including the Farasan archipelago. This would later become the Saudi province of Tihamat Asir. In the early 1920s the Idrisi had administered the coastal Tihama plain as far south as the Yemeni port of Hudaida. The Zaidi Imam of San'a then recaptured the coastal plain and advanced as far north as Ibn Saud's forces would permit. By 1927 a territorial equilibrium of sorts had been reached. Essentially, it was this de facto 1927 line that was formally recognised as the Saudi-Yemen boundary by the treaty of May 1934, but not before Saudi forces had overrun the Tihama, again as far south as Hudaida, during the brief Saudi-Yemeni war fought out in the spring of that year. The Saudi-Yemeni treaty of 'Islamic Friendship and Brotherhood' was signed on 20 May 1934 at the small Yemeni village of Taif (lying just below Hudaida), the most southerly point to which the Saudi campaign had advanced during the hostilities of the previous weeks. The treaty formalised the Imam of Yemen's earlier undertaking to comply with Ibn Saud's terms for peace. The text of the treaty was ratified by Saudi Arabia on 16 July 1934 and by the Imamate only two days later. Exchange of ratifications took place on 22 June 1934.
\\
\end{longtable}}
\clearpage
         \textbf{Conquest of Ethiopia} 

          \underline{Onset:} 1935         

\underline{Reasons:}
\begin{itemize}
\item Nationalism
\item Domestic Politics
\item Revenge/Retribution
\item Economic, Short-Run
\end{itemize}

         \underline{Sources:} 
         \rowcolors{3}{white}{light-gray} 
         {\scriptsize 
         \begin{longtable}{p{0.2\textwidth}p{0.8\textwidth}} 
         \hiderowcolors% 
         \toprule 
         Source & Excerpt \\ 
         \midrule 
         \showrowcolors% 
         \endfirsthead% 
         \hiderowcolors% 
         \toprule 
         Source & Excerpt \\ 
         \midrule 
         \showrowcolors% 
         \endhead% 
         \midrule 
         \endfoot% 
         \hiderowcolors% 
         \bottomrule 
         \showrowcolors% 
         \endlastfoot% 
         
Gibler, Douglas M. 2018. International Conflicts, 1816–2010: Militarized Interstate Dispute Narratives. Vol. 1. Rowman \& Littlefield. & This dispute concerns Italian conquest of Ethiopia following the Italo-Ethiopian War. Italy rounded into its modern form much later than the other colonizing powers in Europe, though Italy too desired an empire in Africa. However, as a latecomer Italy had only two possible routes for empire in Africa, and neither was particularly desirable. Italy could either try to detach territory where one of the other European powers (France, Great Britain, Ottoman Empire) was the prevailing influence or attempt to subdue Ethiopia, a very well-armed, defended and self-governing state in Africa. Italy tried this in 1896 and was humiliated. Italy then returned in the mid-1930s and were moderately successful in this campaign. The Ethiopian capital of Addis Ababa was captured in early May 1936, and the King of Italy was declared ruler of Ethiopia shortly thereafter. Ethiopia joined with Eritrea and Italian Somaliland as constitutive elements of Italian East Africa. The Italians held Ethiopia until 1941, when they were expelled by the British as part of World War II. (p. 412)
\\
Britannica. 2023. “Italo-Ethiopian War.” Encyclopædia Britannica. Accessed August 19, 2023. https://www.britannica.com/event/Italo-Ethiopian-War-1935-1936. & The war, by giving substance to Italian imperialist claims, contributed to international tensions between the fascist states and the Western democracies.
\\
Finaldi, Guiseppe M. 2011. “Italo-Abyssinian Wars (1887–1896, 1935–1936).” In The Encyclopedia of War, edited by Gordon Martel. & Accepting the defeat of Adowa as final was a bitter pill for some Italians to swallow and an undercurrent of dissatisfaction with what was considered to be an abdicating and humiliatingly weak Italian state provided some of the impetus behind the forces that would converge in Fascism in the years after World War I. (...). Mussolini's decision to return Italy to a forward policy in Abyssinia was related to the need to recast his regime at home after the difficulties brought on by the Great Depression and by the fact that with Hitler's dynamism in full view, the “Fascist Revolution” in Italy appeared to be tired and lacking in direction The prestige of the regime both at home and abroad is the key to understanding Mussolini's reignition of the Scramble for Africa in the mid-1930s. (...). Haile Selassie trusted in Europe but was perhaps unaware that with the rise to power of Hitler, the League (and indeed the Versailles settlement) had become a dead letter. It was paramount for Mussolini's regime to distinguish itself from the bungling of liberal Italy that had supposedly been the cause of the Adowa disaster and no expense was spared to ensure success.
\\
Wright, Oliver. 2025. “Abyssinian Crisis (1935).” In The Encyclopedia of Diplomacy, edited by Gordon Martel. & The Abyssinian Crisis was sparked by Benito Mussolini leading fascist Italy into its first foreign war. After thirteen years in power, the dictator decided to add to Italy's existing overseas territories. (...). The Abyssinian Empire of Haile Selassie was a vast territory conveniently located between (...) existing Italian colonies (...). It was already effectively a part of Italy's informal empire, having been subjected to considerable economic penetration for some decades. It was the one place in Africa where the other European powers had no particular interests, and where Italy's predominant concerns were already recognized by them (Mack Smith 1976: 59). (...). [The] Italians conducted their relations with Abyssinia in an increasingly provocative manner. (...). Fascist propaganda justified war, by exaggerating the likely economic benefits of the conquest of Abyssinia, and presenting the territory as large enough to provide an outlet for Italian migration and rich enough in raw materials to solve the country's energy problems (Mack Smith 1976: 64). (...). The victory undoubtedly enhanced the Duce's prestige, especially considering that he had forged his entire career through the use of warlike rhetoric, yet until this point had precious little evidence to show that there was any substance to his words. (...). Certainly, by the mid-1930s, a dictatorship that was entirely built around the dictator and which relied largely upon a personality cult invented in order to strengthen his position and prolong the survival of the regime was in need of ever-greater successes, in order to shore up Mussolini's flagging popularity as fascism became “middle-aged” (Melograni 1976). There is also the possibility that the action served to distract Italy's domestic population from the economic crisis developing from the shortcomings of Mussolini's corporative state (Aquarone 1969). The rise of Hitler in Germany can also be considered to have been an impetus. Until 1933 Mussolini had dominated the international limelight as Europe's “leading” dictator. He had even blocked the upstart Hitler's first attempt at Anschluss in 1934. By 1935, however, Hitler's achievement on the domestic scene in Germany and his first tentative violations of the Versailles settlement were beginning to attract international attention, and Mussolini was in danger of being overshadowed. In addition, he believed that Italy would need to act sooner rather than later to achieve its aims in Africa, so that Italy's armed forces could be repatriated and positioned strategically on the Austrian frontier in order to counter any future attempts by Hitler to achieve German expansion in that sphere.
\\
Morillo, Stephen, Jeremy Black, and Paul Lococo. 2008. War in World History: Society, Technology, and War from Ancient Times to the Present. New York: McGraw-Hill. & After the mid-1920s, the extent and severity of resistance to imperial power declined. In part, this was a consequence of the use of brutal methods of reprisal. This was particularly true of the Italians (...). (...) in conquering Ethiopia, in 1935–36,  the Italians employed poison gas and bacteriological weapons, as well as motorized columns. However, rather than seeing this as a triumph for modern weaponry, it is necessary to note other factors in Italian success, including the deployment of large numbers of troops—nearly 600,000 men. Furthermore, fortunately for the Italians, the Ethiopians chose to engage in battle rather than to avoid engagements and rely on guerrilla tactics. This misguided native strategy permitted the Italians to focus on their strengths rather than to face the problems arising from the logistical handicaps of campaigning in the difficult mountainous terrain. (p. 526)
\\
\end{longtable}}
\clearpage
         \textbf{Third Sino-Japanese} 

          \underline{Onset:} 1937         

\underline{Reasons:}
\begin{itemize}
\item Nationalism
\item Economic, Long-Run
\end{itemize}

         \underline{Sources:} 
         \rowcolors{3}{white}{light-gray} 
         {\scriptsize 
         \begin{longtable}{p{0.2\textwidth}p{0.8\textwidth}} 
         \hiderowcolors% 
         \toprule 
         Source & Excerpt \\ 
         \midrule 
         \showrowcolors% 
         \endfirsthead% 
         \hiderowcolors% 
         \toprule 
         Source & Excerpt \\ 
         \midrule 
         \showrowcolors% 
         \endhead% 
         \midrule 
         \endfoot% 
         \hiderowcolors% 
         \bottomrule 
         \showrowcolors% 
         \endlastfoot% 
         
Gibler, Douglas M. 2018. International Conflicts, 1816–2010: Militarized Interstate Dispute Narratives. Vol. 2 Rowman \& Littlefield. & The third Sino-Japanese War started in 1937 and was ultimately absorbed into the pacific front of the broader world war that ended in 1945. The armistice that ended second Sino-Japanese War (see MID\#129) only accomplished some of the Japanese objectives. The Marco Polo Bridge Incident, like the Mukden Incident for the previous conflict, served as the casus belli. Here, a July 7, 1937, clash between Japanese and Chinese forces at the Marco Polo Bridge near Beijing served as the pretext for a Japanese declaration of war. Fighting followed for several years, to considerable opposition from Western powers, especially the United States (see also, simultaneous occupation of Thailand in 1941, MID\#613). Annoyed with American objections and mindful the United States might buttress their objections with force, the Japanese bombed the American naval forces at Pearl Harbor on December 7, 1941. The United States declared war, following its previous threat regarding the occupation of Thailand. (p. 815)
\\
Britannica. 2023. “Second Sino-Japanese War.” Accessed August 20, 2023. https://www.britannica.com/event/Second-Sino-Japanese-War. & Second Sino-Japanese War, (1937–45), conflict that broke out when China began a full-scale resistance to the expansion of Japanese influence in its territory (which had begun in 1931). The war, which remained undeclared until December 9, 1941, may be divided into three phases: a period of rapid Japanese advance until the end of 1938, a period of virtual stalemate until 1944, and the final period when Allied counterattacks, principally in the Pacific and on Japan's home islands, brought about Japan's surrender.
\\
Bjorge, Gary J. 2011. “China, Invasion of (1931, 1937–1945).” In The Encyclopedia of War, edited by Gordon Martel. & Japan's invasions of China during the 1930s initiated and then expanded the largest war fought between two countries during the twentieth century. For Japan, this was a war fought for resources and geopolitical position. For China, it was a war for national survival. (...). These invasions were an extension of Japan's long-standing national security goal of establishing control over adjacent areas on the Asian continent. Such control was considered necessary in order to counter possible military threats and ensure access to the natural resources needed to guarantee Japan's economic independence. (...).
(...) [In December 1936] Jiang had to agree to end anti-communist operations and form a united front with them to resist Japan. The formation of the united front to resist Japanese aggression put Chinese nationalism on a collision course with Japanese nationalism and expansionism. The collision finally came as the result of an incident that occurred on the night of July 7, 1937, near the Marco Polo Bridge (Lugouqiao) southwest of Beiping. A Tianjin garrison unit that was conducting maneuvers in a sensitive area containing the railroads linking Beiping to central China and the sea encountered unanticipated Chinese interference and shots were fired.
\\
Parker, Geoffrey, ed. 2020. The Cambridge History of Warfare. 2nd ed. Cambridge: Cambridge University Press. & By closing their markets during the Great Depression, the Western powers encouraged aggressive Japanese policies towards the Asian mainland and helped Japanese militarists to gain power. In 1931, Japanese army units seized Manchuria without Tokyo's authorization. Six years later that army initiated an undeclared war against China, and Japanese troops soon controlled China's coastal regions and most of the important Chinese cities, leaving a trail of atrocities in their wake.  (p. 356)
\\
\end{longtable}}
\clearpage
         \textbf{Changkufeng} 

          \underline{Onset:} 1938         

\underline{Reasons:}
\begin{itemize}
\item Power Transition or Security Dilemma
\end{itemize}

         \underline{Sources:} 
         \rowcolors{3}{white}{light-gray} 
         {\scriptsize 
         \begin{longtable}{p{0.2\textwidth}p{0.8\textwidth}} 
         \hiderowcolors% 
         \toprule 
         Source & Excerpt \\ 
         \midrule 
         \showrowcolors% 
         \endfirsthead% 
         \hiderowcolors% 
         \toprule 
         Source & Excerpt \\ 
         \midrule 
         \showrowcolors% 
         \endhead% 
         \midrule 
         \endfoot% 
         \hiderowcolors% 
         \bottomrule 
         \showrowcolors% 
         \endlastfoot% 
         
Gibler, Douglas M. 2018. International Conflicts, 1816–2010: Militarized Interstate Dispute Narratives. Vol. 2 Rowman \& Littlefield. & The border around Changkufeng was poorly delineated, leading to several clashes between Soviet Russia and Japan throughout the 1930s. On July 9, 1938, Soviet forces occupied Changkufeng and began to construct fortified positions. Japanese commanders on the scene attacked the new Soviet positions. Heavy fighting began July 27 with a Japanese offensive. On July 29, the Soviets began construction of new positions nearby at Shatsofeng, and on August 2, they launched a counterattack. On August 7, the Soviets employed heavy artillery. The Japanese wanted a quick settlement to the dispute because they were involved in major operations in their campaign in China. On August 10, both sides agreed to a ceasefire that went into effect the next day, terminating the dispute. (p. 1005)
\\
Blumenson, Martin. 1960. “The Soviet Power Play at Changkufeng.” World Politics 12, no. 2: 249–63. & The immediate background of the Changkufeng incident was the growing hostility between Japan and Soviet Russia over the Japanese invasion of China-Manchuria in 1931. The Soviets strengthened their border defenses in Siberia. The Japanese joined Germany in the Anti-Comintern Pact and thereby placed the Soviet Union in the classic military dilemma of facing potential ene mies on two separate fronts. Japan thus insured herself against direct Soviet action and intervention in support of China
\\
Glantz, David M. 2011. “Khalkhin Gol, Battle of (1939).” In The Encyclopedia of War, edited by Gordon Martel. & (...). Both [border] clashes resulted from conflicting Soviet and Japanese territorial interests in northeastern Asia. The Japanese Empire (...) began pursuing an aggressive policy of expanding its territory on the Asian continent in the 1930s (...) ostensibly to protect its economic interest in the region (...). (...).
The Kwantung Army, which espoused a continental military strategy in opposition to the Japanese navy's maritime (Pacific) strategy, then “flexed its muscles” in the late 1930s by provoking aseries of border incidents with Manchukuo's neighbors, two of which produced significant armed clashes. The first of these clashes took place in July and August 1938 in the Lake Khasan region situated along the border between the USSR's Far East province and Japanese-occupied Korea, when a small Japanese force crossed the border and seized sev-eral small hills in territory claimed by the Soviet Union. After both sides escalated the brief armed clash, the incident ended in stalemate, when the two sides signed an armistice.
\\
Parker, Geoffrey, ed. 2020. The Cambridge History of Warfare. 2nd ed. Cambridge: Cambridge University Press. & Now, even as the war in China sucked [Japan's] forces in ever more deeply, the Japanese army also initiated a series of incidents along the Manchurian frontier to test the reaction of the Red Army, until in August 1939 Soviet troops commanded by Georgy Zhukov annihilated a Japanese reinforced division at Kalkhin Gol ('Nomonhan' in Japanese sources) on the Mongolian border. Zhukov kept his adversaries tied down in the centre, while building up an undetected mass in the rear that he later launched in a pincer attack to trap his opponents - the 'Cannae' tactic he would later use, on a grander scale and with even greater success, at Stalingrad (now Volgograd). Kalkhin Gol persuaded the Japanese government that the Red Army was not an easy mark, and on 15 September it signed a ceasefire with the Soviet Union. (pp. 356-357)
\\
\end{longtable}}
\clearpage
         \textbf{World War II} 

          \underline{Onset:} 1939         

\underline{Reasons:}
\begin{itemize}
\item Nationalism
\item Power Transition or Security Dilemma
\item Religion or Ideology
\item Economic, Long-Run
\item Domestic Politics
\end{itemize}

         \underline{Sources:} 
         \rowcolors{3}{white}{light-gray} 
         {\scriptsize 
         \begin{longtable}{p{0.2\textwidth}p{0.8\textwidth}} 
         \hiderowcolors% 
         \toprule 
         Source & Excerpt \\ 
         \midrule 
         \showrowcolors% 
         \endfirsthead% 
         \hiderowcolors% 
         \toprule 
         Source & Excerpt \\ 
         \midrule 
         \showrowcolors% 
         \endhead% 
         \midrule 
         \endfoot% 
         \hiderowcolors% 
         \bottomrule 
         \showrowcolors% 
         \endlastfoot% 
         
Vasquez, John A. 1996. “The Causes of the Second World War in Europe: A New Scientific Explanation.” & The war was initiated because of specific territorial demands made by Germany on Poland, reflecting a string of claims coming out of the Versailles settlement  sought to bring all Germans under the sovereignty of the Third Reich [...] According to this explanation, such a war might be expected to expand, since there were ongoing rivalries and alliances. In fact, the war expanded along a path that is often typical of world war, namely, a rival of the major state came to the aid of the minor state ally [...] A host of other territorial claims by other states-some longstanding disputes, others the product of new demands-were present prior to the outbreak of the war. They added potential fuel to the conflagration and worked to disrupt the earlier norm against the use of violence to transfer territory or revise the Versailles settlement [...] While territorial demands explain what brought about the initiation of the attacks in 1939 and 1941, they do not directly explain British and French intervention. Here one must look at the steps to war, for while territory is an underlying cause of war, the proximate causes of wars between rivals lie in how the territorial issues are handled (see Vasquez, 1993, Chap. 9). British and French intervention is generally thought to have resulted from a series of repeated crises in which they learned that the only way to deal with Hitler was to fight him; war was necessary, from their perspective, because the territorial demands became relentless and began to spill over to non-German areas, such as Bohemia and Moravia. [...] Hitler had long-term plans, most clearly evinced in his program for self-sufficiency, such as the production of synthetics for oil and rubber (Bell, 1986: 153). The German rearmament program was so ambitious that it failed to meet its goals in 1937 and 1938 (Bell, 1986: 154) because of a lack of raw materials; it placed such a strain on the German economy that it provided a separate impetus to conquer new territory that had these resources.
\\
(1) Nolan, Cathal J. 2011. “World War II: German Invasion of Poland (September 1–October 5, 1939).” In The Encyclopedia of War, edited by Gordon Martel.

(2) Glantz, David M. 2011. “World War II: Eastern Front.” In The Encyclopedia of War, edited by Gordon Martel.

(3) Baxter, Colin F. 2011. “World War II: Mediterranean Campaign.” In The Encyclopedia of War, edited by Gordon Martel.

(4) Kuehn, John T. 2011. “World War II: War in Asia.” In The Encyclopedia of War, edited by Gordon Martel. & (1) In older accounts of the start of World War II, what supposedly opened the door to the German invasion of Poland was the Nazi–Soviet Pact announced on August 23, 1939. In fact, Adolf Hitler had ordered the Oberkommando der Wehrmacht (OKW, or German High Command) to draw up operational plans for invading Poland four months earlier, in April 1939. These plans, code-named Fall Weiss (Case White), were complete before the ink was dry on the Nazi–Soviet Pact. Moreover, his intentions to eliminate Poland as an independent state, and even as a nation, long predated any operational plans. (...). (...) [The] operational achievement [that] Hitler and the OKW most desired: a quick and decisive “battle of annihilation” of the entire Polish army. (...). 
[On] September 17 Joseph Stalin declared that the state of Poland had de facto and de jure “ceased to exist,” broke an existing non-aggression treaty with Warsaw, and sent Red Army divisions slashing into eastern Poland.

(2) The Eastern Front of World War II—called the Great Patriotic War by Russians—refers to the Soviet–German War during which Hitler's Third German Reich attempted to defeat and subjugate Stalin's Soviet Union. The war began on June 22, 1941, when Germany's Wehrmacht (armed forces) invaded the Soviet Union in accordance with Plan Barbarossa, a military operation designed to defeat the Soviet Union's Red Army, overthrow its communist government, and dismember and exploit its territories for the benefit of Germany. The war resulted from political and ideological competition between Nazi Germany and the communist Soviet Union.

(3) World War II came to the Mediterranean on June 10, 1940, when Italy's Benito Mussolini declared war on Britain and France. For the next five years, the Allied and Axis powers would struggle for control of that ancient sea and its shores: a maritime highway, and a geographical link between continents and oceans. The Mediterranean campaign was and is the subject of intense debate and controversy.

(4) Relations between the United States and the Empire of Japan had been strained for a long time, principally over the issue of Japan's policies in China. At the Washington Naval Conference in 1922 the two nations signed a naval treaty in order to end an incipient arms race that had developed after World War I. This treaty established capital ship ratio between the United States, Great Britain, and Japan as 5–5–3. Japan's “inferior” position with respect to the United States caused considerable resentment by some officers in the Imperial Japanese Navy (IJN). As the prosperous 1920s gave way to worldwide depression, militant elements inside Japan began to undermine the basis for peace in the Far East. After the signing of a second naval treaty in London in 1930, insubordinate Imperial Japanese Army (IJA) officers stationed in southern Manchuria invaded the rest of that Chinese province and then established a Japanese puppet state (Manchukuo). Japan subsequently withdrew from the League of Nations in 1932. After another armyengineered “incident” in 1937, Japan became involved in open conflict in China against Chiang Kai-shek's Nationalists and Communist forces under Mao Zedong. Although victorious on the battlefield, Japan found China simply too big to conquer. Incidents such as the Rape of Nanking further alienated Japan from world opinion. (...). 
In 1939 the United States terminated its commercial treaty with Japan and in the following year President Roosevelt ordered the Pacific Fleet to remain permanently at Pearl Harbor instead of California. The United States followed these actions in the summer of 1940 with further economic sanctions. The Japanese responded by moving into northern Indochina at the “invitation” of the new Vichy French government. The United States countered by expanding its embargo to include scrap metal. Japan then formally joined the Axis Powers of Europe in the Tripartite Pact. The lack of military progress in China resulted in formation of a new government under Map 97 World War II: War in Asia. the militarist general Tojo Hideki. Japan's economic situation worsened and she bullied the Vichy French into accepting fifty thousand Japanese troops into southern Indochina. This caused the United States, Great Britain, and the Netherlands to cut off all her access to fuel oil and resulted in the Japanese decision for war against the United States. Admiral Yamamoto Isoruku, an opponent of the Tripartite Pact, was assigned to command the Combined Fleet and reluctantly began planning for war. He conceived the bold idea for a surprise air and mini-submarine attack on the US fleet in Hawaii. On November 26, 1941 the six big aircraft carriers of Dai Ichi Kido Butai (First Mobile Strike Force) secretly departed from Takan Bay in the Kurile Islands under the command of Vice Admiral Nagumo Chuichi for Pearl Harbor.
\\
Folly, Martin H. 2025. “Second World War.” In The Encyclopedia of Diplomacy, edited by Gordon Martel. & The Second World War was a global conflict from December 1941 to August 1945, though the war in Asia began with the start of the Sino-Japanese War in July 1937, and in Europe with the German invasion of Poland on September 1, 1939. The first of the two conflicts that eventually merged into the World War arose out of Japan's expansionism on mainland Asia. Badly affected by the Great Depression, and with its politics increasingly dominated by nationalistic ideas and powerful armed forces, Japan occupied Manchuria, in northern China, from 1931. China was deeply divided by power-struggles between local leaders, the communist party, and the nationalist government, and the opportunity to seize important resources was too tempting for the Japanese army. Ideas developed in Japan of expelling European imperialists from Asia and creating a “Great East Asia Co-Prosperity Sphere” of liberated peoples, which in practice would be a thinly veiled Japanese hegemony. In July 1937, local incidents near Beijing flared up into full-scale war, and Japan in the next two years seized most of the Chinese coastal ports and much of the eastern part of the country, including Beijing, Shanghai, and the capital Nanjing. (...).

\\
\end{longtable}}
\clearpage
         \textbf{Russo-Finnish} 

          \underline{Onset:} 1939         

\underline{Reasons:}
\begin{itemize}
\item Power Transition or Security Dilemma
\item Religion or Ideology
\item Economic, Long-Run
\end{itemize}

         \underline{Sources:} 
         \rowcolors{3}{white}{light-gray} 
         {\scriptsize 
         \begin{longtable}{p{0.2\textwidth}p{0.8\textwidth}} 
         \hiderowcolors% 
         \toprule 
         Source & Excerpt \\ 
         \midrule 
         \showrowcolors% 
         \endfirsthead% 
         \hiderowcolors% 
         \toprule 
         Source & Excerpt \\ 
         \midrule 
         \showrowcolors% 
         \endhead% 
         \midrule 
         \endfoot% 
         \hiderowcolors% 
         \bottomrule 
         \showrowcolors% 
         \endlastfoot% 
         
Gibler, Douglas M. 2018. International Conflicts, 1816–2010: Militarized Interstate Dispute Narratives. Vol. 1 Rowman \& Littlefield. & The Winter War between Russia and Finland took its name from when it was fought in late 1939 and early 1940. Russia was very deliberate in reacquiring territories lost following the revolutions in 1917. An important step was signing the nonaggression pact with a German state that Stalin deeply distrusted. With nonaggression seemingly guaranteed between the two states, a sphere of influence was given that allowed Russia leeway in reacquiring territories in the Baltic (...). Russia had previously acquired Finland from Sweden as a result of the Finnish War in 1809 but lost the territory to independence in the wake of World War I. Beyond satisfying whatever irredentist claims it had in readjusting its border to resemble the pre-1917 arrangement, Finland would be an important buffer between Moscow and Germany. The nonaggression pact was still in effect, but Stalin still did not trust Hitler. When overtures for alliance were rejected by the Finnish, the Soviets invaded on November 30, 1939. The Finnish military fought the best they could and were able to prolong the conflict longer than the Soviets had anticipated. However, ultimately, the numbers caught up with the Finns; the Soviet personnel committed to the fight dwarfed anything the Finns could mobilize. The Finns relented to the terms given by the Soviets in a treaty on March 12, 1940. Finland maintained independence, but lost the Karelian Isthmus that the Soviets strongly desired. The Hanko Peninsula was leased to the Soviets for a period of 30 years, though was reclaimed by the joint forces of the Finns and Nazi Germans as part of World War II. The territory of Salla was ceded to Russia, along with the Rybachy Peninsula in the Barents Sea. Most of Salla was returned to Finland as part of World War II. The losses of Rybachy and the Karelian Isthmus remain for Finland. (p. 392)
\\
Britannica. 2023. “Russo-Finnish War.” Accessed August 20, 2023. https://www.britannica.com/event/Russo-Finnish-War. & Russo-Finnish War, also called Winter War, (November 30, 1939–March 12, 1940), war waged by the Soviet Union against Finland at the beginning of World War II, following the conclusion of the German-Soviet Nonaggression Pact (August 23, 1939). During the 1920s the Finnish government, wary of the threat posed by the Soviet Union, pursued a defense alliance with Estonia, Latvia, and Poland. However, that effort was squelched when the Finnish parliament chose not to ratify the agreement. The Finnish-Soviet nonaggression pact of 1932 was directed at the same concern but failed to quell Finnish fears of Soviet expansionism. Following the invasion, defeat, and partitioning of Poland by Germany and the Soviets in 1939, the Soviet Union sought to push its border with Finland on the Karelian Isthmus westward in an attempt to buttress the security of Leningrad (St. Petersburg) from potential German attack.
\\
Ailes, Mary E. 2011. “Russo-Finnish War (1939–1940).” In The Encyclopedia of War, edited by Gordon Martel. & The Russo-Finnish War, also known as the Winter War, began on November 30,1939 when Soviet forces invaded Finland. Beginning in 1938, the Soviet government pressured Finland to allow it to build a military base on the Finnish coast to guard the maritime approaches to Leningrad. The Soviet leadership feared that Germany would invade the Soviet Union through Finland and that the Finnish military could not withstand such an attack. The Finnish government denied this request. It believed that such an action would make it difficult for Finland to resist further Soviet expansionin the future. Throughout 1938 and 1939 the Soviet government continued to make territorial demands of Finland. The signing of the Nazi–Soviet Non-Aggression Pact in August 1939 led the Soviet government to increase its diplomatic pressure upon Finland as it believed that Germany would not react to a Soviet invasion of its western neighbor.On November 26, 1939 the Soviets broke off negotiations after alleging an artillery attack along the Finnish–Russian border.
\\
\end{longtable}}
\clearpage
         \textbf{Nomonhan} 

          \underline{Onset:} 1939         

\underline{Reasons:}
\begin{itemize}
\item Nationalism
\item Border Clashes
\item Economic, Long-Run
\end{itemize}

         \underline{Sources:} 
         \rowcolors{3}{white}{light-gray} 
         {\scriptsize 
         \begin{longtable}{p{0.2\textwidth}p{0.8\textwidth}} 
         \hiderowcolors% 
         \toprule 
         Source & Excerpt \\ 
         \midrule 
         \showrowcolors% 
         \endfirsthead% 
         \hiderowcolors% 
         \toprule 
         Source & Excerpt \\ 
         \midrule 
         \showrowcolors% 
         \endhead% 
         \midrule 
         \endfoot% 
         \hiderowcolors% 
         \bottomrule 
         \showrowcolors% 
         \endlastfoot% 
         
Gibler, Douglas M. 2018. International Conflicts, 1816–2010: Militarized Interstate Dispute Narratives. Vol. 2 Rowman \& Littlefield. & The border between Mongolia and Japanese-controlled Manchukuo was both poorly defined and proximate to the Soviet border. Japan claimed the border lay at the Khalkha River, while Mongolia's claim extended 15 kilometers east to Nomonhan. On May 12, 1939, Japanese soldiers attacked 700 Mongolian cavalry who had entered the disputed territory. Clashes continued for weeks, each time leading to defeat for Japanese forces. At the end of May, Soviet forces joined the battles with planes, tanks, and troops. The Chief of Staff of the Kwantung Army mobilized all available forces to the fight, despite orders from Tokyo to avoid action that might spread the war. On June 27, 1939, 140 Japanese planes attacked inside Mongolia. The Soviets and the Japanese then reinforced. On August 20, the Soviets launched an offensive that inflicted 17,000 casualties on the Japanese. Tokyo initiated discussions with Moscow on August 28, while the Kwantung Army sent fresh divisions to the front. Tokyo and Moscow reached a truce on September 15, and fighting ceased the next day. On November 19, the Soviets and Japanese agreed to establish a commission to demarcate Manchukuo-Mongol border. The commission met several times through January 20, 1940, but never came to a settlement. On June 9, Molotov and Togo finally reached agreement on the general shape of the boundary, leaving the details to a commission operating on the scene. Border demarcation finished May 15, 1942. Mongolia received its desired boundary in the basin of Holsten while Manchukuo received territory southeast of Nomonhan. (pp. 838-839)
\\
Britannica. 2023. “Mongolia - Counterrevolution and Japan.” Accessed August 20, 2023. https://www.britannica.com/place/Mongolia/Reform-and-the-birth-of-democracy & The MPRP congress held in October 1934 drew attention to a new threat to Mongolia: that from Japan. In 1927 Tanaka Giichi, then the Japanese prime minister, had called for a policy of Japanese expansion in East Asia. [...] The uncertainty, both internal and external, posed by the Japanese threat engendered political hysteria that developed into wholesale arrests and executions on charges of counterrevolutionary activity and spying for the Japanese. [...] Mongolia's slow collision with Japan escalated into a series of battles after the Japanese and Manchukuo troops invaded the northeastern corner of Mongolia in 1939. The fighting reached its climax in August, when Mongolian and Soviet troops under the command of Soviet Gen. Georgy Zhukov annihilated a large Japanese force near the Khalkhyn (Halhïn) River on the Mongolia-Manchukuo border.
\\
Glantz, David M. 2011. “Khalkhin Gol, Battle of (1939).” In The Encyclopedia of War, edited by Gordon Martel. & (...). Both [border] clashes resulted from conflicting Soviet and Japanese territorial interests in northeastern Asia. The Japanese Empire (...) began pursuing an aggressive policy of expanding its territory on the Asian continent in the 1930s (...) ostensibly to protect its economic interest in the region (...). (...).
The Kwantung Army, which espoused a continental military strategy in opposition to the Japanese navy's maritime (Pacific) strategy, then “flexed its muscles” in the late 1930s by provoking a series of border incidents with Manchukuo's neighbors, two of which produced significant armed clashes. (...).
[A] major border clash took place in the vicinity of the Khalkhin River and Nomonhan village, which were located in a disputed border region between Manchukuo and Mongolia, a Soviet client state with close military ties to the Soviet Union. Although both sides asserted that the other provoked the incident by violating the border, in mid-May 1939 a force ofsix hundred cavalrymen from the 6th Mongolian Cavalry Division crossed the Khalkhin Gol, which the Japanese claimed demarcated the border, and occupied disputed terrain near the Japanese village of Nomonhan.
\\
Parker, Geoffrey, ed. 2020. The Cambridge History of Warfare. 2nd ed. Cambridge: Cambridge University Press. & Now, even as the war in China sucked [Japan's] forces in ever more deeply, the Japanese army also initiated a series of incidents along the Manchurian frontier to test the reaction of the Red Army, until in August 1939 Soviet troops commanded by Georgy Zhukov annihilated a Japanese reinforced division at Kalkhin Gol ('Nomonhan' in Japanese sources) on the Mongolian border. Zhukov kept his adversaries tied down in the centre, while building up an undetected mass in the rear that he later launched in a pincer attack to trap his opponents - the 'Cannae' tactic he would later use, on a grander scale and with even greater success, at Stalingrad (now Volgograd). Kalkhin Gol persuaded the Japanese government that the Red Army was not an easy mark, and on 15 September it signed a ceasefire with the Soviet Union. (pp. 356-357)
\\
Otterstedt, Charles. 2000. “The Kwantun Army and the Nomonhan Incident: Its Impact on National Security.” USAWC Strategy Research Project. & From May to September 1939 Japan and the Soviet Union engaged in what started out as a small border clash but quickly escalated into a large undeclared war on the Mongolian plains near the city of Nomonhan. Both countries committed tens of thousands of troops, hundreds of tanks and airplanes. 
\\
\end{longtable}}
\clearpage
         \textbf{Franco-Thai} 

          \underline{Onset:} 1940         

\underline{Reasons:}
\begin{itemize}
\item Nationalism
\item Border Clashes
\item Economic, Long-Run
\item Revenge/Retribution
\end{itemize}

         \underline{Sources:} 
         \rowcolors{3}{white}{light-gray} 
         {\scriptsize 
         \begin{longtable}{p{0.2\textwidth}p{0.8\textwidth}} 
         \hiderowcolors% 
         \toprule 
         Source & Excerpt \\ 
         \midrule 
         \showrowcolors% 
         \endfirsthead% 
         \hiderowcolors% 
         \toprule 
         Source & Excerpt \\ 
         \midrule 
         \showrowcolors% 
         \endhead% 
         \midrule 
         \endfoot% 
         \hiderowcolors% 
         \bottomrule 
         \showrowcolors% 
         \endlastfoot% 
         
Gibler, Douglas M. 2018. International Conflicts, 1816–2010: Militarized Interstate Dispute Narratives. Vol. 2 Rowman \& Littlefield. & Thailand, unlike the rest of southeast Asia, was independent and able to adequately resist pressures from foreign powers for empire. France was the biggest nemesis in the region via French Indochina. The two had clashes through the end of the 19th century and beginning of the 20th century, resulting in Siamese (Thai) concessions to France. World War II changed that relationship. Thailand was led by a nationalist government that was seeking to avenge its losses and France had just been routed by Germany. The partition of France on July 10, 1940, presented the perfect opportunity for Thailand to take advantage of a situation in southeast Asia that France could no longer administer. The Japanese, with whom the Thai government was sympathetic and friendly, had already occupied most of France's colonies in Southeast Asia. The Thai government demanded the return of provinces in Cambodia and Laos, territories that France had wrested from Thailand. On August 16, 1940, five divisions of Thai forces were reported to be fortifying the Thai border with French Indochina. On September 17, France rejected Thai claims to territory within the borders of French Indochina and stated that it was “resolved to defend the territorial integrity of Indochina in all circumstances and against all foreign enterprises.” Border clashes started in November, and war began in earnest in December 1940. Thai victory was incomplete. Thailand had the upper hand after the January 9, 1941 offensive, but this changed in a matter of days. Though severely limited in southeast Asia, France won an important naval battle—the Battle of Ko Chang—on January 17. Japan then intervened and mediated a ceasefire signed on January 28. Japan forced an agreement on March 11, 1941, in Tokyo that gave Thailand the desired territory. France ceded three Cambodian and two Laotian provinces to Thailand, roughly 42,000 square miles in total. Thailand and Japan's relationship grew much closer, but not close enough for Japan which occupied Thailand (see MID\#1785). (pp. 734-735)
\\
Flood, Thadeus. 1969. “The 1940 Franco-Thai Border Dispute and Phibuun Sonkhraam's Commitment to Japan.” Journal of Southeast Asian History. & The border problems between Thailand and French Indochina (...) require brief elaboration. The points of conflict between the Thai and French Indochina in the Mekhong dispute hark back to the basic Franco-Thai treaty of October 3, 1893. (...) (1) Thailand renounced all claims on (...) [Laos] as well as on islands in the river [Mekhong] itself, and (2) Thailand agreed not to construct any military fortifications [25km] of the (...) [Thai] bank of the Mekhong. This frontier situation was complicated in 1904 when the Thai were forced to cede one enclave in the north (on their side of the river opposite Luang Prabang) and one further south near the Cambodian border (on the Thai side of the river opposite Pakse) to France. (...). This made for difficulties with the French administration, particularly in the case of Thai nationals desiring to use the river in these regions. (...). This situation continued to embitter Thailand's relations with French Indochina until World War I, when Thailand entered the war as an ally of France against the Germans. (...) Despite the French concessions in the 1926 convention, it did not, in the Thai view, bring the situation on the Mekhong into accord with international laws and usages. (...). To sum up these background factors: (1) The [border dispute along the Mekhong was not a new issue in 1949]. (2) The non-aggression talks of 1939-40 were initiated by the French (...). (3) The Thai had no particular need for a non-aggression pact but agreed to it on the understanding that France would agree to redress some long-standing injustices on the Franco-Thai border (...). (4) The French minister at Bangkok and Paris officials were sympathetic to the Thai requests (...). (5) The [French] Indochina colonial regime was adamantly opposed to the Thai request (...). (...). (...) [It] is here that the French side would default in the summer of 1940. The encouraging attitude of the French side (...) was radically transformed after the French defeat at the hands of Germany and the governmental changes that followed in France. (...). (...) France was obliged to shift the burden of the negotiations to the Indochina regime (...). (...). Since 1938 the [Imperial Japanese] army had concrete plans for using Thai territory to attack British possessions in Malaya and Burma, but the army was also under strict orders directly from the Emperor not to violate Thai neutrality. (...). Since the military coup d'etat of June 20, 1933, Thailand's military leaders had maintained a rather special relationship with the Japanese, nurtured by a tenuously shared antipathy towards the Western powers in Asia.
\\
Murashima, Eiji. 2005. “Opposing French Colonialism: Thailand and the Independence Movements in Indo-China in the Early 1940s.” South East Asia Research 13 (3): 333–383. & France's surrender to Germany in June 1940, followed by Japan's demands that it should be allowed to station troops in the northern part of Vietnam completely altered the circumstances of Indo-China. On 17 September, the Thai government contacted the new Vichy French government to demand a revision of the border that would make the Mekong River the line between Thailand and Indo-China, and requested the Vichy government to guarantee that in the event of a future change of sovereign authority in Indo-China, all of Thailand's lost territories would be restored. In demanding that the Mekong should be made the border, the Thais were trying to get back at least a portion of the land they had lost. But the French rejected the demand twice, and by the end of September the Phibun government had decided to pursue return of the lost territories resolutely.
\\
Strate, Shane. 2011. “An Uncivil State of Affairs: Fascism and Anti-Catholicism in Thailand, 1940–1944.” Journal of Southeast Asian Studies 42 (1): 59–87. & During the border crisis of 1940, nationalists referred back to the loss of territory in 1893 at every opportunity in order to remind people how Thailand had been humiliated by France in the past. These handbills also conjured up long-held suspicions that the Church was merely an imperialist tool employed to facilitate French expansion into Thailand.
\\
\end{longtable}}
\clearpage
         \textbf{First Kashmir} 

          \underline{Onset:} 1947         

\underline{Reasons:}
\begin{itemize}
\item Nationalism
\item Religion or Ideology
\item Border Clashes
\end{itemize}

         \underline{Sources:} 
         \rowcolors{3}{white}{light-gray} 
         {\scriptsize 
         \begin{longtable}{p{0.2\textwidth}p{0.8\textwidth}} 
         \hiderowcolors% 
         \toprule 
         Source & Excerpt \\ 
         \midrule 
         \showrowcolors% 
         \endfirsthead% 
         \hiderowcolors% 
         \toprule 
         Source & Excerpt \\ 
         \midrule 
         \showrowcolors% 
         \endhead% 
         \midrule 
         \endfoot% 
         \hiderowcolors% 
         \bottomrule 
         \showrowcolors% 
         \endlastfoot% 
         
Gibler, Douglas M. 2018. International Conflicts, 1816–2010: Militarized Interstate Dispute Narratives. Vol. 2 Rowman \& Littlefield. & This territorial dispute is the first of several between India and Pakistan over the contested region of Kashmir. As India split into two at the end of British occupation the Kashmiri maharaja, a Hindu ruler in a Muslim-dominated territory waited to make a decision about whether his kingdom should join Pakistan or India. An armed revolt erupted. In October 1947, 2,000 Pashtuns from Pakistan invaded Kashmir. The maharaja requested Indian help in exchange for temporary accession to India, after which a plebiscite would be held. Pakistan objected; at first individual Pakistani soldiers joined the fight but by May 1948 the Pakistani military had become involved. In response to the violence the United Nations created the Commission for India and Pakistan (UNCIP). Pakistan and India accepted a UNCIP ceasefire resolution on January 1, 1949, and signed a formal delimitation of the ceasefire line in July 1949. (pp. 866-867)
\\
Britannica. 2023. “Kashmir.” Encyclopædia Britannica. Accessed August 19, 2023. https://www.britannica.com/place/Kashmir-region-Indian-subcontinent. & Hari Singh, the maharaja of Kashmir, initially believed that by delaying his decision he could maintain the independence of Kashmir, but, caught up in a train of events that included a revolution among his Muslim subjects along the western borders of the state and the intervention of Pashtun tribesmen, he signed an Instrument of Accession to the Indian union in October 1947. This was the signal for intervention both by Pakistan, which considered the state to be a natural extension of Pakistan, and by India, which intended to confirm the act of accession.
\\
Barua, Pradeep P. 2011. “Indo-Pakistani Wars (1947–1948, 1965, 1971).” In The Encyclopedia of War, edited by Gordon Martel. & India and Pakistan became independent nations in 1947. Almost at their very incep-tion the two nations went to war over the disputed territory of Kashmir.(...).
The Hindu ruler of Jammu and Kashmir, Maharaja Hari Singh, wanted his princely state to be independent. However, the Pakistani president, Mohammed Ali Jinnah, was determined that the Muslim majority state be a part of Pakistan. On October 22, Jinnah launched Operation Gulmarg, a full-scale guerilla invasion of Kashmir by tribal raiders recruited and supported by the Pakistani army. On October 26 Hari Singh signed the instrument of accession handing over his kingdom to India. The next day Indiaflew the 1st Sikhs (infantry battalion) into the capital, Srinagar.
\\
Constable, Philip. 2025. “Kashmir Dispute since 1947.” In The Encyclopedia of Diplomacy, edited by Gordon Martel. & Both India and Pakistan claimed sovereignty over Kashmir leading to war in 1947–48 (...). Within this dispute, Kashmir pursued its own autonomy and independence of regional power interference from the late 1940s. (...). Kashmir reflects the difficulties that India and Pakistan have experienced since 1947 in terms of nationalist secession and the increasingly militant measures that nationalist insurgencies have taken for autonomy and independence. The Kashmir dispute has its roots in the British colonial period and partition of British-India between India and Pakistan in 1947. British rule had utilized Hindu-Muslim communal tensions to divide nationalist opposition into a majority Hindu Congress Party and minority Muslim League from the 1920s. These divisions became magnified by wider electorates after 1935 leading to Muslim League concern over the influence of a majority Hindu Congress in an independent India, and consequently demands for a Pakistani state. A logic of partition became the strategic consolidation (...). Kashmir became caught up in this logic of partition.
\\
Morillo, Stephen, Jeremy Black, and Paul Lococo. 2008. War in World History: Society, Technology, and War from Ancient Times to the Present. New York: McGraw-Hill. & In South Asia, for example, the end of British rule led to wars between India and Pakistan in 1947, 1965, and 1971. Pakistan's defeat in the last ensured success for the rebellion in East Pakistan, which became independent as Bangladesh.
\\
\end{longtable}}
\clearpage
         \textbf{Arab-Israeli} 

          \underline{Onset:} 1948         

\underline{Reasons:}
\begin{itemize}
\item Nationalism
\item Religion or Ideology
\item Domestic Politics
\item Revenge/Retribution
\end{itemize}

         \underline{Sources:} 
         \rowcolors{3}{white}{light-gray} 
         {\scriptsize 
         \begin{longtable}{p{0.2\textwidth}p{0.8\textwidth}} 
         \hiderowcolors% 
         \toprule 
         Source & Excerpt \\ 
         \midrule 
         \showrowcolors% 
         \endfirsthead% 
         \hiderowcolors% 
         \toprule 
         Source & Excerpt \\ 
         \midrule 
         \showrowcolors% 
         \endhead% 
         \midrule 
         \endfoot% 
         \hiderowcolors% 
         \bottomrule 
         \showrowcolors% 
         \endlastfoot% 
         
Gibler, Douglas M. 2018. International Conflicts, 1816–2010: Militarized Interstate Dispute Narratives. Vol. 2. Rowman \& Littlefield. & From mid-1947 through early 1948, Palestinians and Israelis fought over who would control the area once the British Mandate expired on May 15, 1947. As each side vied for control, the Israelis declared independence on May 14, 1948, prompting a coalition of Arab nations—Iraq, Egypt, Syria, Lebanon, and Jordan—to launch an attack that very day. The Arab coalition wanted to install the United State of Palestine, foregoing the two-state solution that had been proposed by the 1947 UN Partition Plan. Jordan began the war by invading the Corpus separatum area around Jerusalem. Fighting intensified. The Israeli Defense Force (IDF), founded on May 26, managed to eventually field more troops than the Arab countries and was able to gain ground and air superiority over the Arab nations fairly quickly. From June to October, the Israelis had the advantage, pushing back the Arab armies and gaining territory. Several attempts at truces did not hold, and in late October, the IDF initiated operations in Egypt, Lebanon, Jordan, and Syria, pushing out the remaining Arab troops and holding 5,000 square kilometers more than the territory allotted Israel by the United Nations. In early November, the United Nations proposed the creation of an armistice, which Israel accepted on November 18. In 1949, the four primary Arab states in the conflict also signed armistices with Israel, officially ending the conflict. (p. 657)
\\
Cashman, Greg, and Leonard C. Robinson. 2007. An Introduction to the Causes of War: Patterns of Interstate Conflict from World War I to Iraq. Rowman \& Littlefield Publishers, Inc & The tension caused by the proposed partition of Palestine boiled over into full-scale civil war between Palestinians and Jews [...]
The Zionist finally realized their dream of establishing a Jewish state in Palestine on May 14, 1948. [...] Arab states of Egypt, Syria, Jordan, and Iraq launched an offensive [...] stated intention of their attack being to protect 45 percent of the territory that had been set aside for the Palestinian state. [...] Arab leaders harbored conflicting objectives in fighting the war, with each seeking to use the conflict to bolster his domestic legitimacy as well as his claim to regional leadership.
\\
Bregman, Ahron. 2011. “Arab–Israeli Conflict.” In The Encyclopedia of War, edited by Gordon Martel. & [Before 1948] there was a bloody civil war between Arabs and Jews in Palestine, an area that was under British control from 1917 to 1948. What sparked it was the changing demography of Palestine brought about by the influx of Jewish immigrants who came to Palestine in search of shelter from pogroms and persecution in their native countries. (...). Jews (...) comprised only 4 percent of the total Palestinian population in 1882 (...). (...). By 1947 there were 608,230 Jews in Palestine, compared with about 1,364,330 Arabs (Bregman 2000:4–5). (...). (...) [The] United Nations proposed to partition Palestine between the two peoples (...); it offered the Jews 55 percent of Palestine and the Arabs (still the majority) 45 percent. The Jews accepted (...) Arabs objected and threatened (...) war. The UN proceeded anyway (...) on the next day a civil war broke out (...).
Friday, May 14, 1948 was the day the British departed Palestine and the Jews declared independence. In response to the Israeli declaration of independence, the Arab armies of Egypt, Syria, Lebanon, Iraq, and Transjordan, supported by units from Saudi Arabia and Yemen, invaded. Their aim was to destroy Israel, help Arab Palestinians, or perhaps, as some scholars claim, to grab some land for themselves in the absence of the “British policeman. (...). For the Israelis, this became their “War of Liberation,” or “War of Independence,” while for the Arab Palestinians, some 750,000 of whom became refugees as a result of the war, it became al Nakba, “The Catastrophe.”
\\
Siniver, Asaf. 2018. “Arab–Israeli Conflict.” In The Encyclopedia of Diplomacy, edited by Gordon Martel. & By 1947, with Jewish and Arab violence reaching unprecedented levels, Britain decided to refer the question of the future of Palestine to the nascent United Nations. On November 29, 1947, the UN General Assembly passed Resolution 181 which called for the partition of mandatory Palestine into a Jewish state and Arab state, with Jerusalem designated as an international city. The Jewish community in Palestine (the Yishuv) accepted the resolution but the Arabs rejected it as an unjust solution to the problem. Following the termination of the British mandate on May 14, 1948 the state of Israel was established, leading to its invasion by the armies of Egypt, Jordan, Syria, Lebanon, and Iraq the following day.
\\
Morillo, Stephen, Jeremy Black, and Paul Lococo. 2008. War in World History: Society, Technology, and War from Ancient Times to the Present. New York: McGraw-Hill. & The Arabs proved unwilling to accept the culmination of the Zionist movement in the form of an independent Israel, and this ensured a high level of tension in the region. Israel was able to establish its independence in the face of attacks by neighboring Arab states in 1948–49. (p. 581)
\\
Townshend, Charles, ed. 2005. The Oxford History of Modern War. No. 2. Oxford University Press, USA. & In 1948 the Israelis fought to carve a state out of the old
British Mandate of Palestine. (p.167)
\\
Bregman, Ahron. 2000. Israel's Wars, 1947–93. In Warfare and History. Routledge. & Yet, rather than easing tension, the resolution to partition the land and the subsequent British decision, made on 4 December 1947, to depart on Friday 14 May 1948, had increased tensions between the peoples of Palestine and ‘it was as if on a signal Arabs and Jews squeezed the trigger and exchanged fire'. (...) The principal aim of the Jews in Palestine in the period immediately after the UN resolution to partition Palestine, was to gain effective control over the territory allotted to them by the UN and to secure communication with thirty-three Jewish settlements which, according to the UN plan, fell outside the proposed Jewish state. (...) The Palestinians' strategic aim during the civil war was negative in nature, namely to prevent the implementation of the partition plan by disrupting and strangling Jewish lines of communication, and by cutting off Jewish settlements from localities and positions that were already occupied. (pp. 9-10)
\\
\end{longtable}}
\clearpage
         \textbf{Korean} 

          \underline{Onset:} 1950         

\underline{Reasons:}
\begin{itemize}
\item Nationalism
\item Power Transition or Security Dilemma
\item Religion or Ideology
\end{itemize}

         \underline{Sources:} 
         \rowcolors{3}{white}{light-gray} 
         {\scriptsize 
         \begin{longtable}{p{0.2\textwidth}p{0.8\textwidth}} 
         \hiderowcolors% 
         \toprule 
         Source & Excerpt \\ 
         \midrule 
         \showrowcolors% 
         \endfirsthead% 
         \hiderowcolors% 
         \toprule 
         Source & Excerpt \\ 
         \midrule 
         \showrowcolors% 
         \endhead% 
         \midrule 
         \endfoot% 
         \hiderowcolors% 
         \bottomrule 
         \showrowcolors% 
         \endlastfoot% 
         
Gibler, Douglas M. 2018. International Conflicts, 1816–2010: Militarized Interstate Dispute Narratives. Vol. 2 Rowman \& Littlefield. & The Korean War was fought between North Korea, South Korea, UN forces (under US leadership with American troops), China, and the Soviet Union, for control of the Korean peninsula. The Soviet Union and United States had divided the Korea peninsula along the 38th parallel at the end of World War II, with the Soviets supporting a new regime in the north led by Kim Il-Sung and the United States supporting one in the south led by Syngman Rhee. These competing regimes each claimed sovereignty over the peninsula. On June 25, 1950, North Korea crossed the 38th parallel, quickly captured Seoul, and pushed to the south. US troops landed with the assistance of other Western countries and pushed the North Koreans back across the 38th parallel, toward the Chinese border. The Chinese then intervened in October and pushed the troops from South Korea and its allies back toward the 38th parallel and a rough stalemate followed. Fighting continued through the summer of 1953, but by the end North Korea and South Korea still controlled nearly the same territory as before the conflict. The parties signed an armistice on July 27, 1953. (pp. 845)
\\
Britannica. 2023. “Korean War.” Accessed August 20, 2023. https://www.britannica.com/event/Korean-War. & The Korean War had its immediate origins in the collapse of the Japanese empire at the end of World War II in September 1945. Unlike China, Manchuria, and the former Western colonies seized by Japan in 1941–42, Korea, annexed to Japan since 1910, did not have a native government or a colonial regime waiting to return after hostilities ceased. Most claimants to power were harried exiles in China, Manchuria, Japan, the U.S.S.R., and the United States. They fell into two broad categories. The first was made up of committed Marxist revolutionaries who had fought the Japanese as part of the Chinese-dominated guerrilla armies in Manchuria and China. [...] The other Korean nationalist movement, no less revolutionary, drew its inspiration from the best of science, education, and industrialism in Europe, Japan, and America. These “ultranationalists” were split into rival factions [...] 
\\
Lee, Steven H. 2011. “Korean War (1949–1953).” In The Encyclopedia of War, edited by Gordon Martel. & The ideological foundations of the Korean War lie in the pre-1945 era, in the evolution of diverse and divisive strands of Korean nationalism, efforts by Koreans to eliminate Japanese colonial rule from the Korean Peninsula, the impact of the Russian Revolution on Asia, and the influence of Christianity and ideas of constitutional governance on Koreans. The conjunction of forces that led to the Korean conflict, however, were rooted in the immediate post-World War II era, in Soviet and American objectives toward Korea, the global beginnings of the Cold War, and the creation and consolidation of two rival Korean regimes.
\\
Qing, Simei. 2025. “Korean War (1950–53).” In The Encyclopedia of Diplomacy, edited by Gordon Martel. & On June 25, 1950, North Korea launched an offensive campaign against the south, a civil war took place between North and South Korea. On October 1, 1950, South Korean troops crossed the 38th parallel. On October 2, Mao Zedong informed Stalin that Beijing was not to dispatch troops for now. On October 3, Chinese premier Zhou Enlai informed the Truman administration that should American, rather than South Korean, troops cross the 38th parallel, China would join the war. On October 7, American troops crossed the 38th parallel and on October 19, Chinese troops crossed the Yalu River. The Korean War was transformed into a military confrontation between China and the United States in Korea. (...). Throughout 1949 both North and South Korea attempted to reunify their nation by means of military power (...). (...). On January 12, 1950, secretary of state Dean Acheson announced the new military strategy of a “defense perimeter” in the Western Pacific.
\\
\end{longtable}}
\clearpage
         \textbf{Off-shore Islands} 

          \underline{Onset:} 1954         

\underline{Reasons:}
\begin{itemize}
\item Nationalism
\item Power Transition or Security Dilemma
\item Religion or Ideology
\end{itemize}

         \underline{Sources:} 
         \rowcolors{3}{white}{light-gray} 
         {\scriptsize 
         \begin{longtable}{p{0.2\textwidth}p{0.8\textwidth}} 
         \hiderowcolors% 
         \toprule 
         Source & Excerpt \\ 
         \midrule 
         \showrowcolors% 
         \endfirsthead% 
         \hiderowcolors% 
         \toprule 
         Source & Excerpt \\ 
         \midrule 
         \showrowcolors% 
         \endhead% 
         \midrule 
         \endfoot% 
         \hiderowcolors% 
         \bottomrule 
         \showrowcolors% 
         \endlastfoot% 
         
Gibler, Douglas M. 2018. International Conflicts, 1816–2010: Militarized Interstate Dispute Narratives. Vol. 2 Rowman \& Littlefield. & On October 1, 1949, the Chinese Civil War ended and an interstate conflict began with the Nationalists who fled to Formosa, when the Chinese People's Republic was formally established on the mainland. Clashes followed frequently, and the Nationalists implemented a blockade of Communist ports. On June 27, 1950, the United States entered the dispute when President Truman announced his order to the
Seventh Fleet of the US navy to enter the Taiwan Strait and prevent any attack on Formosa. (...).
This dispute involves the attempts of Chinese Nationalists under the leadership of General Chiang to reclaim mainland China from the Communist government. The Communist leadership was in turn attempting to exercise what it believed was its territorial right to control the island of Formosa (Taiwan) and other areas in which exiled anticommunists had fled in 1949. As part of its anti-Communist foreign
policy, the United States became involved (...) while at the same time ensuring that the Nationalists would not try to invade the mainland. In early February 1953, President Eisenhower declared that the United States would no longer prevent the Nationalists from attacking China's mainland, while US forces would continue to protect Formosa from Communist attack. This change in foreign policy was seen as overt support for the Nationalist forces to attempt to take back the mainland. On February 9, the United States (in support of Nationalist forces) threatened to blockade China and use force if necessary in defense of Formosa. The Nationalist leader, Chiang, also began making claims that the time of his attack was drawing near. Through the summer of 1953, Chiang continued making threats (...) calling for the joint invasion of Communist China by the US and Nationalist forces. While the United States did not make plans for war, it offered military training to Nationalist forces. In January 1954, Chinese prisoners of war, who were released by South Korea on the grounds that they would fight the Communists, arrived in Formosa and were added to the Nationalist army. In May 1954, China issued a warning to the United States to end its involvement in Formosa and withdraw its advisers. Russia became involved in the dispute in June 1954 after one of its tankers was seized. The Soviets accused the United States of capturing the tanker, but the Nationalist forces took responsibility. Russia remained militarily and diplomatically involved from June to August 1954. Continuous action between the United States, Taiwan, and China was reported through the summer of 1954, and large-scale fighting took place during most of September 1954. Attacks, raids, threats, and blockades continued through 1955 until February 1956.
(pp. 790-792)
\\
Office of the Historian, Foreign Service Institute, United States Department of State. n.d. “The Taiwan Straits Crises: 1954–55 and 1958.” & Tensions between the People's Republic of China (PRC) and the Republic of China (ROC) in the 1950s resulted in armed conflict over strategic islands in the Taiwan Strait. On two separate occasions during the 1950s, the PRC bombed islands controlled by the ROC. The United States responded by actively intervening on behalf of the ROC. [...] The importance of the islands in the Taiwan Strait was rooted in their geographic proximity to China and Taiwan and their role in the Chinese Civil War. Jinmen (Quemoy), two miles from the mainland Chinese city of Xiamen, and Mazu, ten miles from the city of Fuzhou, are located approximately one hundred miles west of Taiwan. When the Nationalist Government of the ROC under Chiang Kai-shek recognized that it had lost control of mainland China during the Chinese Civil War, the officials and part of the Nationalist Army fled to the island of Taiwan, establishing troops on these two islands and the Dachen Islands further north [...]  The U.S. Government then announced its determination to defend Taiwan against communist attack, although it did not specify the territory included within its defensive perimeter
\\
Pruessen, Ronald W. 2001. “Over the Volcano: The United States and the Taiwan Strait Crisis, 1954–1955.” In Re-examining the Cold War: US-China Diplomacy, 1954–1973, Harvard University Asia Center. & In 1954. all were still controlled by Chiang Kai-shek, in spite of the fact that his Nationalist government's bastion was at least one hundred miles from any of them. Chiang had been eager to retain the offshore islands as jumping-off points for a return to the mainland-and as useful bases for harassment and intelligence gathering in the interim. Many Washington leaders in the later Truman and early Eisenhower years had been dubious about the real value of the tiny outposts, but a general sense of Chiang's value to the United States had ledto tolerance toward his continuing military buildup on these islands. During the more frustrating moments of the Korean War, in fact, there had even been American encouragement-and aid-for raids and other operations that diverted some Chinese Communist resources. 3 When provocative acts continued, even in the aftermath of the 1954 Geneva Conference, Beijing moved to correct at least some of the leftover anomalies of the civil war. Shelling of Jinmen, just four miles distant from the mainland and actually within the confines of Amoy harbor, began on September 3, 1954. (pp. 77-78)
\\
Stolper, Thomas E. 1985. “Climax of the 1954–55 Taiwan Affair.” In China, Taiwan and the Offshore Islands, 1st ed. Routledge. & The People's Republic of China (PRC) responded to the submission of the treaty by putting heavy military pressure on the northernmost islands held by the Nationalists. The most important of these were the Tachens, some 220 miles from Taiwan, which were garrisoned by 10,000 regular troops on the main islands and about 5,000 irregulars on associated islands. Though virtually everyone in Congress was prepared to defend Taiwan, there was great reluctance to become embroiled in the offshore islands. The logic of Peking's action during the 1954–55 Taiwan affair was to make the world realize that any international move tending to block the recovery of Taiwan entailed a risk of war with the PRC. The PRC would have made a great compromise by agreeing to any sort of formal cease-fire, and the world would have given it Taiwan. In the Soviet draft resolution, one side is the PRC, and the other is the United States.  (p. 1)
\\
\end{longtable}}
\clearpage
         \textbf{Soviet Invasion of Hungary} 

          \underline{Onset:} 1956         

\underline{Reasons:}
\begin{itemize}
\item Nationalism
\item Religion or Ideology
\end{itemize}

         \underline{Sources:} 
         \rowcolors{3}{white}{light-gray} 
         {\scriptsize 
         \begin{longtable}{p{0.2\textwidth}p{0.8\textwidth}} 
         \hiderowcolors% 
         \toprule 
         Source & Excerpt \\ 
         \midrule 
         \showrowcolors% 
         \endfirsthead% 
         \hiderowcolors% 
         \toprule 
         Source & Excerpt \\ 
         \midrule 
         \showrowcolors% 
         \endhead% 
         \midrule 
         \endfoot% 
         \hiderowcolors% 
         \bottomrule 
         \showrowcolors% 
         \endlastfoot% 
         
Gibler, Douglas M. 2018. International Conflicts, 1816–2010: Militarized Interstate Dispute Narratives. Vol. 1 Rowman \& Littlefield. & Imre Nagy, premier of Hungary from 1953 to 1955, fell in disfavor and was removed from his post and expelled from the Soviet-led Communist Party in April 1956. Rather than accept this result, the Hungarians revolted against Soviet policies and, on October 23, 1956, clashed with Soviet troops stationed in Hungary and soon restored Imre Nagy to his post. The Soviets invaded Hungary on November 4 with little response from the Western states or the United Nations. The fighting ended 10 days later on November 14, though the Soviets had effective control of Hungary days before that. Many of the protesters and supporters of Nagy were arrested, tried, and jailed. The Communist Party in Hungary was purged further of anti-Soviet elements. Nagy was given preliminary assurances by the Yugoslavian embassy that he could be exiled to another country, though he was eventually abducted by the Soviets. He was given a trial and executed two years later. (p. 312)
\\
Britannica. 2023. “Hungarian Revolution.” Accessed August 20, 2023. https://www.britannica.com/event/Hungarian-Revolution-1956. & Rebels won the first phase of the revolution, and Imre Nagy became premier, agreeing to establish a multiparty system. On November 1, 1956, he declared Hungarian neutrality and appealed to the United Nations for support, but Western powers were reluctant to risk a global confrontation. On November 4 the Soviet Union invaded Hungary to stop the revolution, and Nagy was executed for treason in 1958.
\\
\end{longtable}}
\clearpage
         \textbf{Sinai War} 

          \underline{Onset:} 1956         

\underline{Reasons:}
\begin{itemize}
\item Religion or Ideology
\item Economic, Long-Run
\item Revenge/Retribution
\end{itemize}

         \underline{Sources:} 
         \rowcolors{3}{white}{light-gray} 
         {\scriptsize 
         \begin{longtable}{p{0.2\textwidth}p{0.8\textwidth}} 
         \hiderowcolors% 
         \toprule 
         Source & Excerpt \\ 
         \midrule 
         \showrowcolors% 
         \endfirsthead% 
         \hiderowcolors% 
         \toprule 
         Source & Excerpt \\ 
         \midrule 
         \showrowcolors% 
         \endhead% 
         \midrule 
         \endfoot% 
         \hiderowcolors% 
         \bottomrule 
         \showrowcolors% 
         \endlastfoot% 
         
Gibler, Douglas M. 2018. International Conflicts, 1816–2010: Militarized Interstate Dispute Narratives. Vol. 2 Rowman \& Littlefield. & This dispute describes the Suez Crisis of 1956 that led to the Sinai War of the same year. Egypt had become a conundrum among the Western powers. Formerly a British protectorate that only gained system membership when the protectorate formally ended in 1936, the Egyptians sat out World War II (though Egyptian facilities were used by the Allied powers) and tried to pursue a middle way with respect to the ongoing Cold War. That changed with the emergence of Israel and when Nasser came to power in Egypt. Egypt was Israel's main enemy the day it gained statehood. Further, Soviet disposition toward Israel considered its emergence to be a result of bourgeois nationalism. Egypt found a sympathetic ally with the Soviets and fell into the Soviet camp when Nasser arrived in power and pursued a policy hostile to the West. The removal of all British personnel from the Suez Canal was part of this policy. When the Suez Canal was nationalized, the British personnel on site were forcibly removed and Israel was then shut out of an important part of its trade. The Soviets used its veto power in the United Nations to thwart any resolution on the matter. Britain, France, and Israel opted instead for punitive action on Egypt. The allies routed the Egyptians, leaving the Americans and Soviets as the principal contestants of the Cold War to plea for a ceasefire to minimize the chances of a broader war. The United States in particular played a very strong hand in forcing Britain out of Egypt. A UN-led ceasefire came into effect on November 6, 1956, which was administered by a UN peacekeeping force. Britain and France were finally coerced to give up the territory it had acquired during the conflict, vacating Egypt on December 22. Israel finally gave up the Gaza Strip on March 1957. (p. 660)
\\
Wright, William M., Michael C. Shupe, Niall M. Fraser, and Keith W. Hipel. 1980. “A Conflict Analysis of the Suez Canal Invasion of 1956.” Conflict Management and Peace Science. & On July 26, 1956, in a speech in Alexandria, President Nasser of Egypt announced the nationalization of the Suez Canal. This action was the climax of international tensions that have previously been studied using a game theoretic approach by Shupe et al. (1980). By controlling the Suez Canal, Nasser anticipated that its revenue could be utilized to assist in financing the Aswan Dam (Bowie, 1974). Britain and France were deeply shocked by the nationalization of the canal, as was the rest of the world. Britain was vitally concerned over the flow of oil from the Middle East, and because of its high investments in the canal and its declining position in the area, it viewed Nasser's action as a critical threat. France considered the Suez canal as a French undertaking since its construction was originally organized by a Frenchman, and France already viewed Nasser with hostility because of his support of the Algerian rebels. For Britain and France regaining control of the canal would guarantee a vital supply route and would serve to humiliate Nasser. Thus, Britain and France immediately wanted urgent and decisive action against Egypt. An integrated British-French military command was soon established with the British in charge.
\\
Bregman, Ahron. 2011. “Arab—Israeli Conflict.” In The Encyclopedia of War, edited by Gordon Martel. & Unlike 1948, when war was imposed on Israel, in 1956 it was Israel, in collusion with Britain and France, who went on the offensive. On July 26, 1956, President Gamal Abdel Nasser of Egypt nationalized the Suez Canal Company, of which France and Britain had been the majority shareholders. The two colonial powers resented Nasser's unilateral decision, as they would lose control over an important international waterway through which vital supplies came to Europe. France and Britain began considering the use of force to regain control of the Suez Canal. Israel was secretly invited to join the coalition against Nasser, which provided an opportunity to achieve some of her own aims – mainly to gain control of the Straits of Tiran. The Straits, at the foot of the Gulf of Aqaba, were Israel's primary trade route to East Africa and Asia, but had for several years been blockaded by Egypt. Now, Israel conditioned that if she was to join the planned war she should be allowed to move her troops south to the Straits and remove the blockade.
\\
Witty, David M. 2011. “Suez Crisis (1956).” In The Encyclopedia of War, edited by Gordon Martel. & On July 26, 1956, Egyptian President Gamal Abdel Nasser nationalized the internationally controlled Suez Canal. The resulting war was an attempt by Anglo-French and Israeli forces to take control of the canal and remove Nasser from power. Pressure from the United States and the United Nations (UN) forced a ceasefire in the conflict and the withdrawal of the attacking forces.The crisis was an important milestone in the anti-colonial movement. Great Britain suffered a diplomatic defeat and was permanently weakened, while Nasser emerged as the leader of the Arab world. Nasser came to power (...) with the goal of ending the British occupation of Egypt (...). The canal (...) was a vital security concern for Britain; avoiding an interruption of its traffic was a national priority. Nasser also aimed to revenge the defeat that the Arabs had suffered in their 1948–1949 war with Israel. (...) following a major Israeli raid into the Egyptian-controlled Gaza Strip in 1955, Egypt tightened its blockade of the Tiran Straits, preventing Israeli access to the Red Sea.
\\
Siniver, Asaf. 2018. “Arab–Israeli Conflict.” In The Encyclopedia of Diplomacy, edited by Gordon Martel. & The Second Arab–Israeli War – the Suez War of 1956 – came following years of cross-border infiltrations of guerrilla groups (Fedayeen) into Israel, and Egyptian president Gamal Abdel Nasser's decision to nationalize the Suez Canal and to close the Straits of Tiran at its southern tip to Israeli and Israel-bound shipping. At the same time, Israel's approach to asserting its borders, primarily via military retributions against Arab villages (most famously the October 1953 Qibya raid which resulted in dozens of civilian deaths) further contributed to the escalation of hostile relations between Israel and its neighbors. In October 1956 Israel colluded with France and Britain, who wished to maintain their strategic interests in the region, to attack Egypt and force it to reopen the Suez Canal. However, despite the successful military campaign the plan backfired and the three allies were forced to withdraw their forces amidst the condemnation of the United Nations and unprecedented cooperation between the United States and Soviet Union to bring the crisis to an end. 
\\
Jönsson, Christer. 2025. “Suez Crisis (1956).” In The Encyclopedia of Diplomacy, edited by Gordon Martel. & The Suez Crisis refers to the chain of events following the Egyptian nationalization of the Suez Canal in July 1956. In late October Israel, followed by Britain and France, attacked Egypt in an effort to regain Western control of the Suez Canal. After diplomatic interventions by the United States, the Soviet Union, and the United Nations, the invaders were forced to withdraw. Diplomatic communication throughout the crisis was complicated by vastly different perceptions among the main actors. (...). The Suez Crisis refers to the chain of events following the Egyptian nationalization of the Suez Canal and its operating Suez Canal Company in July 1956 in retaliation for the reneging of an agreement by the US and British governments to finance the construction of the Aswan Dam. (...).
In July 1954, Britain and Egypt signed an agreement requiring all British forces to withdraw from Egypt by June 1956. (...). This was a triumph for the Egyptian nationalists (...). In September 1955, Egypt revealed that it had concluded an agreement (...) to buy Soviet arms. US secretary of state John Foster Dulles labeled the Soviet entry into the Middle East “the most serious development since Korea, if not since World War II” (Finer 1964: 28). (...). Britain and France, the major shareholders in the Suez Canal Company, viewed Nasser's challenge as an act of aggression and launched plans for a military response. (...).
The main actors in the Suez drama viewed one another and unfolding events from markedly different vantage points. First, their image of the enemy diverged. The United States and the Soviet Union defined the conflict primarily in East–West terms, with “international communism” and “the West” as the principal villains. For Egypt, Britain and Israel were the main enemies; for Israel, the Arabs.
\\
\end{longtable}}
\clearpage
         \textbf{IfniWar} 

          \underline{Onset:} 1957         

\underline{Reasons:}
\begin{itemize}
\item Nationalism
\item Power Transition or Security Dilemma
\item Economic, Long-Run
\end{itemize}

         \underline{Sources:} 
         \rowcolors{3}{white}{light-gray} 
         {\scriptsize 
         \begin{longtable}{p{0.2\textwidth}p{0.8\textwidth}} 
         \hiderowcolors% 
         \toprule 
         Source & Excerpt \\ 
         \midrule 
         \showrowcolors% 
         \endfirsthead% 
         \hiderowcolors% 
         \toprule 
         Source & Excerpt \\ 
         \midrule 
         \showrowcolors% 
         \endhead% 
         \midrule 
         \endfoot% 
         \hiderowcolors% 
         \bottomrule 
         \showrowcolors% 
         \endlastfoot% 
         
Gibler, Douglas M. 2018. International Conflicts, 1816–2010: Militarized Interstate Dispute Narratives. Vol. 2 Rowman \& Littlefield. & This dispute describes the Ifni War, which was primarily a struggle over rich oil and mineral resources found in the Southern Moroccan Protectorate. Underscoring this issue was the Moroccan effort to claim this southern territory from Spain following its independence in April 1956. New fighting reportedly broke out in November 1957, with Moroccan Liberation Army forces attacking Spain's Ifni and Rio de Oro colonies. Spain sent reinforcements, and clashes ensued. On November 28, Moroccan Crown Prince Moulay Hassan accused Spanish forces of attacking Moroccan territory, and ordered the royal army to return fire. On December 8, Spain threatened Morocco regarding its aid to Ifni rebels. Though the Spanish government claimed all organized resistance had ended, it maintained a naval presence in the region for the duration of the dispute. Spain renewed its threats on December 31. Then Spain bombed the rebel-claimed Ifni region on February 19, 1958, and dropped about 500 paratroopers. In response, Morocco expelled the Spanish consul on February 21. Spanish threats and reinforcements on both sides continued into March. The dispute ended on April 2, 1958, with the Treaty of Angra de Cintra, in which Morocco would take over administration of the Southern Moroccan Protectorate. 
(p. 543)
\\
Britannica. 2023. “Ifni.” Accessed August 19, 2023. https://www.britannica.com/place/Ifni. & The enclave became part of Spanish West Africa in 1946. Spain retained its claim to Ifni following Moroccan independence in 1956 and repelled a series of Moroccan attempts to take the territory in the Ifni War (1957–58).
\\
US Army War College, Studies Institute. 2013. War and Insurgency in the Western Sahara. & When Morocco attained independence in 1956, Madrid first hesitated to relinquish any of Spanish West Africa. The Moroccan-supported Liberation Army (LA), later called the Sahrawi Liberation Army by the Spaniards, began to attack French outposts in Algeria and Mauritania, using Spanish territory as a safe haven. The Spanish military, lacking sufficient forces and clear instructions from Madrid, at first let guerrilla bands move across Spanish territory with surprising freedom, although the Spaniards provided the French with information about their movements. [...] The LA forces, however, found the French to be more than they could handle, and they shifted their efforts to the Spanish-controlled north, sparking the outbreak of the Ifni War (1957-58). The Spaniards fought back hard but eventually withdrew to a defensive parameter around the town of Sidi Ifni.
\\
\end{longtable}}
\clearpage
         \textbf{Taiwan Straits} 

          \underline{Onset:} 1958         

\underline{Reasons:}
\begin{itemize}
\item Nationalism
\item Power Transition or Security Dilemma
\item Religion or Ideology
\end{itemize}

         \underline{Sources:} 
         \rowcolors{3}{white}{light-gray} 
         {\scriptsize 
         \begin{longtable}{p{0.2\textwidth}p{0.8\textwidth}} 
         \hiderowcolors% 
         \toprule 
         Source & Excerpt \\ 
         \midrule 
         \showrowcolors% 
         \endfirsthead% 
         \hiderowcolors% 
         \toprule 
         Source & Excerpt \\ 
         \midrule 
         \showrowcolors% 
         \endhead% 
         \midrule 
         \endfoot% 
         \hiderowcolors% 
         \bottomrule 
         \showrowcolors% 
         \endlastfoot% 
         
Gibler, Douglas M. 2018. International Conflicts, 1816–2010: Militarized Interstate Dispute Narratives. Vol. 2 Rowman \& Littlefield. & This is the Second Taiwan Strait Crisis that began with China bombing the islands of Kinmen and the Matsu Islands along the coast of China. Communist Chinese forces followed with an attempted invasion of the Matsu but were repulsed by Nationalist forces. Most contemporaries contended that the Chinese actions were an attempt to probe the willingness and resolve of the United States to defend its Nationalist Chinese allies. The United States responded by quickly moving its Seventh Fleet into the area and warning the Communist Chinese. The Soviets pledged support for China, but also tried to caution against escalation. US forces were instrumental in reinforcing air defenses for the Nationalist forces. (p. 793)
\\
Office of the Historian, Foreign Service Institute United States Department of State. “The Taiwan Straits Crises: 1954–55 and 1958.” https://history.state.gov/milestones/1953-1960/taiwan-strait-crises. & Tensions between the People's Republic of China (PRC) and the Republic of China (ROC) in the 1950s resulted in armed conflict over strategic islands in the Taiwan Strait. On two separate occasions during the 1950s, the PRC bombed islands controlled by the ROC. The United States responded by actively intervening on behalf of the ROC. [...] The importance of the islands in the Taiwan Strait was rooted in their geographic proximity to China and Taiwan and their role in the Chinese Civil War. Jinmen (Quemoy), two miles from the mainland Chinese city of Xiamen, and Mazu, ten miles from the city of Fuzhou, are located approximately one hundred miles west of Taiwan. When the Nationalist Government of the ROC under Chiang Kai-shek recognized that it had lost control of mainland China during the Chinese Civil War, the officials and part of the Nationalist Army fled to the island of Taiwan, establishing troops on these two islands and the Dachen Islands further north [...]  The U.S. Government then announced its determination to defend Taiwan against communist attack, although it did not specify the territory included within its defensive perimeter [...] Although there were good reasons for the PRC to stand down in 1955, it resumed its bombardment of Jinmen and Mazu in 1958. This time, the PRC took advantage of the fact that international attention was focused on U.S. intervention in Lebanon and barred ROC efforts to re-supply garrisons on the off-shore islands.
\\
Pruessen, Ronald W. 2001. “‘A Thorn in the Side of Peace’: The Eisenhower Administration and the 1958 Offshore Islands Crisis.” In Re-examining the Cold War: US–China Diplomacy, 1954–1973. Harvard University Asia Center. & During both episodes [from September 1954 to April 1955 and from late August to late October 1958], senior officials in the Republican administration of Dwight D. Eisenhower, led by the president and Secretary of State John Foster Dulles, determined that the United States could not remain indifferent to the Chinese Communist challenge in the Taiwan Strait, and they made preparations to intervene militarily, possibly with nuclear weapons, to protect the Nationalist government's coastal possessions if circumstances so required. At the same time, they tried to deter the Chinese and find a satisfactory peaceful exit from the crisis while coping with a difficult Nationalist ally and the troubling domestic and international repercussions of their actions. (...). As in 1954-55, so in 1958, American policymaking was shaped in fundamental ways by the perceived U.S. stake in Taiwan and the preservation of a viable and friendly Nationalist government as well as by the necessity not just to support but to constrain the exiled regime of Chiang Kai-shek. Eisenhower and Dulles, together with other high-level civilian and military officials, were convinced that the loss of Jinmen (Quemoy) or the other larger offshore islands to the Chinese Communists by the threat or use of force would have dire consequences for the Nationalist government and thus for the security of Taiwan and the noncommunist Far East. This was the dominant consideration behind the Eisenhower administration's decision to chance a military collision with the PRC over what were otherwise considered to be militarily insignificant and vulnerable parcels of territory along the south China coast. American leaders believed that the foremost goal of Mao Zedong and his comrades, when they attempted to put a military stranglehold on Jinmen in late August, was the eventual takeover ofT aiwan and the ultimate undoing of the U.S. strategic position in the western Pacific. (pp. 106-107)
\\
Halperin, Morton H. 1988. The 1958 Taiwan Straits Crisis: A Documented History. Rand Corporation. & The first sign of a possible crisis in the Taiwan Straits came on June 30, 1958, when the Chinese Communists demanded a resumption of the Sino-American ambassadorial talks. The first military action came in lace July in the form of air clashes dver the Taiwan Straits and the Chinese mainland. During July the Chinese Nationalists began to anticipate a Connunist move against the Offshore Islands. Urging the United States to commit itself publicly to the defense of the Offshore Islands, they also sought modern equipment for their armed forces, including the delivery of American Sidewinder missiles. While the United States refused to issue a public statement indicating that it would defend Quemoy, it did increase its military assistance to the Government of the Republic of China (GRC) and began intensive contingency planning for a crisis in the Taiwan Straits. The basic policy of the American government was that it would help defend the Offshore Islands only if necessary for the defense of Taiwan. American officials in the field, how ever, were authorized to assist the GRC in planning for the defense of the Islands, and assumed that nuclear weapons would be used to counter anything but very light probing by the Chinese Communists. In early August, officials in Washington became concerned with the possibility of a crisis, although they did not expect the Chinese Communists to launch a major military attack. During that same month, a consensus developed that a high-level decision should be made as to what the American reaction would be to an air-sea interdiction campaign again,the Offshore Islands. There was also strong pressure for a diplomatic warning to the Chinese Communists that the United States would not tolerate the fall of Quemoy. On August 22 it was decided, just below the presidential level, that the United States would participate in the defense of the Offshore Islands if they came under attack. It was agreed that, as an attempt to deter a Chinese Communist move, a public statement clarifying the American position would be issued in the form of an exchange of letters between Secretary of State Dulles and Representative Thomas Morgan. On August 23, 1958, at 6:30 p.m. Taiwan time, the Chinese Communists launched a heavy artillery attack against the Quemoy Islands. (pp. v-vi)
\\
\end{longtable}}
\clearpage
         \textbf{Assam} 

          \underline{Onset:} 1962         

\underline{Reasons:}
\begin{itemize}
\item Nationalism
\item Religion or Ideology
\item Border Clashes
\item Revenge/Retribution
\end{itemize}

         \underline{Sources:} 
         \rowcolors{3}{white}{light-gray} 
         {\scriptsize 
         \begin{longtable}{p{0.2\textwidth}p{0.8\textwidth}} 
         \hiderowcolors% 
         \toprule 
         Source & Excerpt \\ 
         \midrule 
         \showrowcolors% 
         \endfirsthead% 
         \hiderowcolors% 
         \toprule 
         Source & Excerpt \\ 
         \midrule 
         \showrowcolors% 
         \endhead% 
         \midrule 
         \endfoot% 
         \hiderowcolors% 
         \bottomrule 
         \showrowcolors% 
         \endlastfoot% 
         
Gibler, Douglas M. 2018. International Conflicts, 1816–2010: Militarized Interstate Dispute Narratives. Vol. 2. Rowman \& Littlefield. & This dispute describes the Sino-Indian War fought in 1962, largely in the shadow of the ongoing Cuban Missile Crisis. India and China disputed where their borders met in the Himalayas, and India had been forceful in normalizing its borders after it had gained independence from the British following World War II. China had become aggressive too, occupying Tibet in 1950. This angered India, who had claims in the area as well. India seemingly began the dispute by proclaiming a Forward Policy on November 2, 1961, regarding its northern border with China. India followed by moving troops north to establish forts and police the area. China responded in kind. Diplomatic tensions and minor incidents between Chinese and Indian military personnel finally devolved into open war on October 20, 1962. China was much better equipped to fight the war and quickly routed the Indians. A ceasefire was reached on November 21, 1962, that was effectively an imposed settlement for China but was actually a pretty good deal for India. China occupied Aksai Chin but allowed India to resume control of Assam. Ten days later, on December 1, 1962, the Chinese government accused Indian troops of participating in an “armed provocation” on the eastern sector of the border between the two countries. Accusations between the two countries continued into 1963 as India accused China of border violations in the Spanggur Lake and the Ladakh area while China accused an Indian military aircraft of an airspace violation, “reconnoitering over the Yatung.” On October 14, 1963, China accused India of an earlier border violation into its Sinkiang territory on September 22, declaring it an “obvious intention to create tension and broaden the boundary dispute.” Finally, on December 25, 1963, a thinning of Chinese forces along the Tibetan border began in order to fortify the Sinkiang border with the Soviet Union. This was a signal to the Indian government that the Tibetan border dispute was coming to a close, with the belief that no additional Chinese incursions would occur in the manner of the 1962 movement. (pp. 818-819)
\\
Britannica. 2023. “Sino-Indian War.” Encyclopædia Britannica. Accessed August 19, 2023. https://www.britannica.com/topic/Sino-Indian-War. & The regime in Beijing, after suppressing the buffer state of Tibet in 1950, began disputing the border with India at several points between the tiny Himalayan states of Nepal, Bhutan, and Sikkim. [...] Aksai Chin in particular had been a long-ignored corner of the subcontinent because of its remoteness and isolation. However, this changed when the Chinese tried to connect Tibet with Xinjiang by building a military road through the region. India objected to the Chinese presence in the sector, which it claimed as part of the Ladakh region under Indian administration. [...] After a number of border skirmishes between 1959 and 1962, which began initially as a by-product of the uprising in Tibet, the People's Liberation Army (PLA) of China forcefully attacked across the disputed boundaries on October 20, 1962.
\\
Sinha, P. B., and A. A. Athale. 1992. “Chapter III: Towards Armed Conflict.” In History of the Conflict with China, 1962, edited by S. N. Prasad. Government of India, Ministry of Defense, History Division. & Therein seems to lie the basic cause of the debacle of 1962. India failed to avoid a war during the transition period. It became a typical example of changing horses in mid-stream. Lulled by faulty political assessment and wrong intelligence forecasts, the country got caught in a war when least prepared. (...).  [The] Chinese Government (...) converted [the] frontier (...) into an area of tension and conflict, thereby forcing the Government of India to adopt (...) military measures (...). (...). The basic cause for [the failure to resolve the border tangle peacefully] lay in the dual policy of China whereby it professed a desire for peaceful settlement of the border question while pursuing (...) the path of flagrant aggression. (...). China had recognised McMahon Line in the case of Burma but when it came to India [it] became  'so-called' or illegal.

In a fundamental sense, the origin of 1962 Sino-Indian conflict lay in Chinese expansion and occupation of Tibet. The issue got further aggravated due to failure of Chinese to win over the Tibetans. Indian asylum to Dalai Lama raised Chinese suspicions about ultimate Indian intentions.
\\
Chervin, Reed H. 2024.“The 1962 War and Domestic Reactions in China and India.” In The Cold War in the Himalayas: Multinational Perspectives on the Sino-Indian Border Conflict, 1950–1970, 33–70. Cold War in Asia and Beyond. Amsterdam University Press. & There are many questions as regards the causes of the War (...). (...). Overall it was a failure to assess the Chinese threat in correct perspective. The Indian government did not expect China to fight a War against India. (...). The independence of India resulted in boundaries which were disputable. (...). (...) [Possible] causes for the War can be classified as: Tibet and development in and around Tibet since 1950. Border dispute. Mao's leadership and personality, ambitions as also policies and actions. Indian policies and actions. (...). 
[The Chinese invasion of] Tibet is an important cause in the conflict between the two countries. (...). One of the reasons for the War was the belief that the Chinese perceived that India has been supporting the Khampas with assistance from the United States. (...). [The] Dalai Lama escaped on 17 March 1959 and sought asylum in India. This issue further intensified the disputes between the two countries. Thus, the occupation of Tibet led to Chinese assertiveness over the region, thus causing disputes with India. 
The Sino-Indian Border has a total length of 4056 km. [The] dispute stems on the western and the eastern portion. (...). The demarcation of the west- ern boundary commenced in the nineteenth century. (...). The eastern portion begins in Sikkim and extends up to the Lohit sub-division of Arunachal. In 1826, Britain annexed Assam. Subsequent annexations in further AngloBurmese wars expanded China's borders with British India eastwards, to include the border with what is now Burma. (...). The McMahon line [was] a proposed boundary between Tibet for the eastern sector (...). China soon refused to sign the agreement on a more detailed map. (...). 
While the boundary issue remained unsettled, there were other geostrategic factors (...). To add to China's discomfiture a full-blown revolt started in 1959 in the Kham and Amdo regions of mainland China inhabited by Tibetans. (...). During this period, there was acute tension between China and the Soviet Union. Mao felt that Khrushchev was giving in to the West and China deserved to lead the Socialist bloc rather than the Soviet Union. (...). Internationally, India was being seen as a major power. (...). The United States also sought closer relations with India. (...). The Soviet Union had continued to supply heavy machinery used in road construction [to India] (...). Khrushchev (...) visited India in February 1960 and stated that they were prepared to help if the need arises. This symbolism would obviously have been noticed in Beijing (...).
\\
Garver, John. 2004. “India, China, the United States, Tibet, and the Origins of the 1962 War.” India Review 3 (2): 9–20. & Mao Zedong became convinced (...) that India was actively and deliberately colluding with the United States to weaken and undermine China's rule in Tibet. (...) India–US collusion never developed into a major, large-scale effort toward Tibet, but I suspect that when scholars gain access to China's foreign policy archives for the 1960s they will find that Mao and other Chinese leaders were deeply apprehensive of a US–Indian effort to detach Tibet from China.
\\
Chervin, Reed H. 2024. “The 1962 War and Domestic Reactions in China and India.” In The Cold War in the Himalayas: Multinational Perspectives on the Sino-Indian Border Conflict, 1950–1970, 33–70. Cold War in Asia and Beyond. Amsterdam University Press. & [On] November 21 [1962], Zhou Enlai announced that a ceasefire would begin at midnight on November 22 and that on December 1 the People's Liberation Army would withdraw north of the line of actual control “which existed between China and India on November 7, 1959.” This unilateral decision elucidated Beijing's intent for the war to serve as a disciplinary action. According to General Lei Yingfu, Chairman Mao Zedong had proclaimed at an October 18 Politburo meeting that “our counterattack is only to warn and punish, only to tell Nehru and the Indian government that they cannot use military means to resolve the border problem.” Beijing would behave similarly during its intervention in Vietnam in 1979. The People's Republic attacked Vietnam in response to its invasion of Cambodia in 1978. Both the Sino-Indian and Sino-Vietnamese Wars served a punitive purpose, lasted approximately one month, and were justified by China as righteous.
\\
Devereux, David R. 2009. “The Sino-Indian War of 1962 in Anglo-American Relations.” Journal of Contemporary History 44 (1): 71–87. & Although disputed in two places along their common frontier (Aksai Chin inthe west, the Northeast Frontier Agency in the east, at opposite ends of the Himalayas), the boundary war of 1962 was fought mainly on the eastern border, between Tibet and India. The dispute dates to the nineteenth centuryand the efforts of British Indian officials to delineate a boundary in the largely unmapped high Himalayas. In 1914, an India–Tibet conference at Simla accepted the so-called McMahon Line as the official boundary, which followed (roughly) the crest of the mountains. Citing its claim to Tibet, neither the Chinese Republic nor the later communist regime accepted the boundary, but insisted on sovereignty all the way to the Indian side of the mountains. The dispute was purely theoretical until Indian independence and partition in 1947, which made the newIndian government acutely conscious of establishing its writ to the maximum extent. The Chinese communist invasion of Tibet in 1950 and suppression of Tibetan rebels in 1959 showed a similar effort by China to assert national dominion to the greatest possible degree. Although India recognized China's claim to Tibet in 1954, ushering in a relatively friendly period, the Dalai Lama's flight to India in 1959 and a ‘forward policy' of establishing Indian military outposts along the McMahon line led to border clashes in both the eastern and western zones and a steady deterioration of relations. (...).  
According to Neville Maxwell, Indian politicians and scholars remainwedded to the idea of India as an innocent victim of Chinese aggression. It was accepted at the time, and remains largely accepted, that ‘China took advantageof the then ongoing Cuban missile crisis, when the American attention wasdiverted, to teach India a lesson by launching a long-prepared surprise attack'.Nehru's decision to adopt a ‘forward policy' along the disputed McMahon lineand his refusal to negotiate was a source of grievance to China at a time whenthere were growing Chinese concerns about Tibetan resistance and the security of the entire border area.
\\
\end{longtable}}
\clearpage
         \textbf{Vietnam War, Phase 2} 

          \underline{Onset:} 1965         

\underline{Reasons:}
\begin{itemize}
\item Nationalism
\item Religion or Ideology
\end{itemize}

         \underline{Sources:} 
         \rowcolors{3}{white}{light-gray} 
         {\scriptsize 
         \begin{longtable}{p{0.2\textwidth}p{0.8\textwidth}} 
         \hiderowcolors% 
         \toprule 
         Source & Excerpt \\ 
         \midrule 
         \showrowcolors% 
         \endfirsthead% 
         \hiderowcolors% 
         \toprule 
         Source & Excerpt \\ 
         \midrule 
         \showrowcolors% 
         \endhead% 
         \midrule 
         \endfoot% 
         \hiderowcolors% 
         \bottomrule 
         \showrowcolors% 
         \endlastfoot% 
         
Gibler, Douglas M. 2018. International Conflicts, 1816–2010: Militarized Interstate Dispute Narratives. Vol. 2 Rowman \& Littlefield. & This dispute, commonly known as the Vietnam War in the United States, is best described as an internationalized civil war between North Vietnam and South Vietnam over the composition of government. The war was certainly a legacy of French colonial rule. French Indochina ended with the French embarrassment at Dien Bien Phu in 1954, and the French agreed to a de facto partition of Vietnam at the 17th parallel. The communists, who had been active in fighting the French colonial rule, occupied the north of the 17th parallel with a capital at Hanoi. The pro-French forces stayed south, with a capital of Saigon. The agreement called for elections to be held to determine a unified Vietnam, but this temporary expedient outlined in the Geneva Accords crystallized. Interstate system members after the final French withdrawal, this civil war became internationalized. It became a multilateral dispute as a result of the broader international context of containment and the specter of international communism. The United States intervened diplomatically to support the South, later mobilizing and committing a significant portion of its armed forces to combat. Australia, New Zealand, Philippines, South Korea, Thailand, and Cambodia assisted South Vietnam. Meanwhile, communist powers China and the Soviet Union provided support to the communists in Hanoi. The conflict lasted for almost a decade after the United States became involved. Unfortunately for South Vietnam, the Americans faced strong domestic pressure that greatly hindered American ability to commit to the conflict. The retirement of US President Lyndon Johnson after his one elected term and the election of Richard Nixon marked the beginning of a gradual withdrawal from the war. While American forays into Cambodia and Laos followed as a result of this withdrawal, it was clear the American government wanted out. The United States completely withdrew on January 27, 1973. South Korea, Thailand, and the Philippines withdrew with them. Secret negotiations involving the Americans agreed to the 17th parallel as a border in order for the United States to withdraw, but this does not end the conflict. Saigon fell on April 30, 1975. Vietnam was reunified and administered by the communists, ending the conflict. (p. 720)
\\
Britannica. 2023. “Vietnam War.” Accessed August 20, 2023. https://www.britannica.com/event/Vietnam-War. & Vietnam War, (1954–75), a protracted conflict that pitted the communist government of North Vietnam and its allies in South Vietnam, known as the Viet Cong, against the government of South Vietnam and its principal ally, the United States. Called the “American War” in Vietnam (or, in full, the “War Against the Americans to Save the Nation”), the war was also part of a larger regional conflict (see Indochina wars) and a manifestation of the Cold War between the United States and the Soviet Union and their respective allies. At the heart of the conflict was the desire of North Vietnam, which had defeated the French colonial administration of Vietnam in 1954, to unify the entire country under a single communist regime modeled after those of the Soviet Union and China. The South Vietnamese government, on the other hand, fought to preserve a Vietnam more closely aligned with the West.
\\
Möise, Edwin. 2011. “Vietnam War (1959–1975).” In The Encyclopedia of War, edited by Gordon Martel. & The Vietnam War was primarily a struggle for control of South Vietnam (...). (...). It was in many ways an outgrowth of the First Indochina War (1945-1954), in which the Vietminh, a Vietnamese nationalist organisation under communist leadership, had fought against France and a French puppet government, the State of Vietnam. (...). 
Lyndon Johnson, who became president of the United States after the assassination of John Kennedy in November 1963,wanted to reduce American military spending to make more money available for hisdomestic programs. But he feared the consequences, both strategic and political, of losing South Vietnam. Many Americans believed, and President Johnson at least took seriously, the “domino theory,” which stated that if South Vietnam fell to communists, then numerous other countries would quickly fall also. By early 1965 it seemed obvious that South Vietnam would fall quite soon if the United States did not intervene with direct military force.
\\
Dumbrell, John. 2018. “Vietnam War.” In The Encyclopedia of Diplomacy, edited by Gordon Martel. & The Vietnam War was a highly complex and long-sustained confrontation between Western forces, primarily France and the United States on the one hand, and, on the other, the forces of Vietnamese nationalism. (...). The “American war” began in earnest with President Lyndon Johnson's decision to escalate American involvement in 1965 and ended with the Paris Peace Agreement, concluded in January 1973. Following the exit of US forces, the war between North Vietnamese forces and the non-communist government in South Vietnam continued until April 1975, when communist forces from North Vietnam took the southern capital, Saigon. (...). 
The Vietnam-related diplomacy of the John F. Kennedy presidential years (1961–63) involved various proposals to achieve a “Cold War neutral” Vietnam on the model of the 1962 settlement in Laos. With the revival of conflict in the South between the communist-led National Liberation Front and South Vietnamese forces (assisted by US “advisers”), the NLF itself began to float the idea of a neutral South. Kennedy considered various proposals for neutralization, notably from under secretary of state Chester Bowles. Tentative approaches were made to Hanoi and Beijing via Indian and Burmese intermediaries. In the event, however, the failure of the Laotian neutralization – by 1962 communist forces were protecting Laotian connecting routes from North to South Vietnam along the so-called “Ho Chi Minh Trail” – persuaded JFK against further diplomatic initiatives. 
Lyndon Johnson (president, 1963–69) made his escalation decisions against a background of international peace diplomacy. Such efforts, including major mediation efforts by the United Nations in 1964–65, floundered in the face of mutual distrust between the warring parties, and a mutual assumption that each side was simply trying to woo international opinion rather than actually making a peace deal. LBJ has frequently been accused of virtual sabotage of such diplomacy, thus effectively “choosing war.” Johnson's own view seems to have been that he would give peace feelers what he called “the old college try” (Dallek 1996: 153), exploring possibilities at least partly for national and international public effect. LBJ's promise, in a major address at Johns Hopkins University in April 1965, of massive investment in a Vietnam stripped of “foreign interference” elicited no positive response. Hanoi was wary of being seen by Beijing to accept peace gestures which the latter interpreted as straightforward efforts to sow divisions between North Vietnam and China. At one level, LBJ's 1965 escalation of US involvement was an effort to augment American leverage in peace negotiations. However, LBJ's failure to think through exactly how such negotiations could proceed (especially in the context of the asymmetrical Hanoi–Beijing relationship) effectively doomed the strategy to failure.
\\
Parker, Geoffrey, ed. 2020. The Cambridge History of Warfare. 2nd ed. Cambridge: Cambridge University Press. & Only in autumn 1962 did the Americans finally recognize that Diem was a loser. Threats of US withdrawal eventually led to a coup by the Vietnamese military that overthrew the regime and assassinated the dictator and his brother. The coterie of generals who succeeded Diem proved even more inept than their predecessor, and by summer 1964 resistance in the South was collapsing. The new US president, Lyndon Johnson, refused to admit defeat at the hands of what he termed a 'piss-ant country', and in summer 1964 he launched a series of air raids against the North Vietnamese navy, supposedly in response to attacks on American destroyers in the Gulf of Tonkin. Johnson and his advisers hoped that such strikes would persuade the North Vietnamese to desist. Ho and his colleagues, however, had no intention of desisting. (...). In early 1965 Johnson authorized a bombing campaign against the North, code-named 'Rolling Thunder', which severely limited the targets that US aircraft could attack.  
\\
\end{longtable}}
\clearpage
         \textbf{Second Kashmir} 

          \underline{Onset:} 1965         

\underline{Reasons:}
\begin{itemize}
\item Nationalism
\item Religion or Ideology
\item Border Clashes
\end{itemize}

         \underline{Sources:} 
         \rowcolors{3}{white}{light-gray} 
         {\scriptsize 
         \begin{longtable}{p{0.2\textwidth}p{0.8\textwidth}} 
         \hiderowcolors% 
         \toprule 
         Source & Excerpt \\ 
         \midrule 
         \showrowcolors% 
         \endfirsthead% 
         \hiderowcolors% 
         \toprule 
         Source & Excerpt \\ 
         \midrule 
         \showrowcolors% 
         \endhead% 
         \midrule 
         \endfoot% 
         \hiderowcolors% 
         \bottomrule 
         \showrowcolors% 
         \endlastfoot% 
         
Gibler, Douglas M. 2018. International Conflicts, 1816–2010: Militarized Interstate Dispute Narratives. Vol. 2 Rowman \& Littlefield. & In January 1964, Pakistan brought the unsettled Kashmir question before the UN Security Council again. Pakistan wanted to have Kashmiris vote to resolve the issue of state allegiance, as the United Nations had argued several times but which India had rejected each time by arguing that Kashmir had acceded to Indian control in 1947. Various clashes began on February 21. Pakistan wanted to mediate the dispute, but India rejected the idea and called for negotiations without third-party intervention. The Security Council called for more Pakistani and Indian talks on the issue in May 1964. Shortly after this announcement, on May 19, Indian troops reportedly killed four civilians across the Pakistani border. Pakistan appealed this incident before the United Nations. Top state leaders were scheduled to meet in late May, but these talks were canceled upon the death of Prime Minister Nehru of India. Indian troops reportedly killed numerous civilians during June 1964, and Pakistan officially denounced India on these border violations in late July 1964. Deadly clashes took place from the summer of 1964 through the spring of 1965. 
\\
Britannica. 2023. “Kashmir.” Accessed August 20, 2023. https://www.britannica.com/place/Kashmir-region-Indian-subcontinent. & Although there was a clear Muslim majority in Kashmir before the 1947 partition, and its economic, cultural, and geographic contiguity with the Muslim-majority area of the Punjab could be convincingly demonstrated, the political developments during and after the partition resulted in a division of the region. Pakistan was left with territory that, although basically Muslim in character, was thinly populated, relatively inaccessible, and economically underdeveloped. The largest Muslim group, situated in the Vale of Kashmir and estimated to number more than half the population of the entire region, lay in Indian-administered territory, with its former outlets via the Jhelum valley route blocked. Many proposals were subsequently made to end the dispute over Kashmir, but tensions mounted between the two countries following the Chinese incursion into Ladakh in 1962, and warfare broke out between India and Pakistan in 1965.
\\
Barua, Pradeep P. 2011. “Indo-Pakistani Wars (1947–1948, 1965, 1971).” In The Encyclopedia of War, edited by Gordon Martel. & India and Pakistan became independent nations in 1947. (...). The Kashmir conflict would lead to war again in 1965 (...). (...). 
The 1965 or Second Kashmir War as it is sometimes called began with yet another ambitious Pakistani plan to use guerillas to destabilize Indian Kashmir. Code named Operation Gibraltar, the plan envisaged the use of seven thousand guerillas, who would infiltrate the Kashmir Valley and instigate an uprising by the local Kashmiris.The assault commenced on August 7, 1965.India's response was swift and decisive. Indian forces crossed the ceasefire line, seized a number of Pakistani posts, including the Haji Pir Pass, and crushed the guerillas by August.
\\
Constable, Philip. 2025. “Kashmir Dispute since 1947.” In The Encyclopedia of Diplomacy, edited by Gordon Martel. & In 1964 the PRC supported Pakistan's claim to Kashmir and the USSR withdrew support for India over Kashmir in the UN. Ayub Khan sought to use Soviet concerns over Chinese influence in Asia and US assistance to India in the early 1960s to generate support from both communist countries for Pakistan over Kashmir as a counter-balance to US ambivalence. Hindu–Muslim communal tension which had turned into regional conflict over consolidation of national territory and identity was thereby involved in Cold War rivalries. By the mid-1960s, on the one hand, India was less stable as a result of the Sino-Indian War in 1962 and Nehru's death in 1964. On the other hand, Pakistan was stable after a referendum on Ayub Khan's presidency in 1965, foreign economic development of Pakistan, and multilateral support from the PRC and USSR over Kashmir. In this context, Pakistan's prime minister Zulfikar Ali Bhutto re-invoked the issue of a plebiscite over Kashmir in a climate of communal tension within Kashmir itself.
\\
Morillo, Stephen, Jeremy Black, and Paul Lococo. 2008. War in World History: Society, Technology, and War from Ancient Times to the Present. New York: McGraw-Hill. & In South Asia, for example, the end of British rule led to wars between India and Pakistan in 1947, 1965, and 1971. Pakistan's defeat in the last ensured success for the rebellion in East Pakistan, which became independent as Bangladesh.
\\
\end{longtable}}
\clearpage
         \textbf{Six Day War} 

          \underline{Onset:} 1967         

\underline{Reasons:}
\begin{itemize}
\item Power Transition or Security Dilemma
\item Religion or Ideology
\item Border Clashes
\item Economic, Long-Run
\item Domestic Politics
\end{itemize}

         \underline{Sources:} 
         \rowcolors{3}{white}{light-gray} 
         {\scriptsize 
         \begin{longtable}{p{0.2\textwidth}p{0.8\textwidth}} 
         \hiderowcolors% 
         \toprule 
         Source & Excerpt \\ 
         \midrule 
         \showrowcolors% 
         \endfirsthead% 
         \hiderowcolors% 
         \toprule 
         Source & Excerpt \\ 
         \midrule 
         \showrowcolors% 
         \endhead% 
         \midrule 
         \endfoot% 
         \hiderowcolors% 
         \bottomrule 
         \showrowcolors% 
         \endlastfoot% 
         
Gibler, Douglas M. 2018. International Conflicts, 1816–2010: Militarized Interstate Dispute Narratives. Vol. 2 Rowman \& Littlefield. & This dispute describes the events that ultimately led to the Six-Day War pitting Israel against the allied Arab world (mostly Egypt, Syria, and Jordan in combat). An Israeli attack on Jordan led to a Jordanian-Egyptian defense pact on May 30, 1967. When Syria was thought to be the next target for an Israeli assault, Egypt intervened quickly. Egyptian troops moved into the Sinai, removed the UN peacekeepers stationed there since the Sinai War (...), blockaded the Tiran Straits and gave every indication that the conflict was likely to escalate. Israel decided to take the initiative, launching an attack on June 5, 1967. By demolishing the Egyptian air force on the first day, the Israelis set the path for a victory after six days of fighting. Israel concluded that it had met its military objectives and willingly agreed to a ceasefire agreement promulgated by the major powers. Israel's success was resounding and drastically rearranged its borders. Israel occupied the entirety of the Sinai Peninsula, which was eventually returned to Egypt following the Camp David Accords. Israel added buffers between Jerusalem and its various neighbors, and the Gaza Strip, West Bank, and Golan Heights were taken as well. UN Security Council Resolution 242 set the ceasefire and confirmed the new status quo. (pp. 686-687)
\\
Britannica. 2023. “Six-Day War.” Accessed August 20, 2023. https://www.britannica.com/event/Six-Day-War. & Prior to the start of the war, attacks conducted against Israel by fledgling Palestinian guerrilla groups based in Syria, Lebanon, and Jordan had increased, leading to costly Israeli reprisals. In November 1966 an Israeli strike on the village of Al-Samū' in the Jordanian West Bank left 18 dead and 54 wounded, and, during an air battle with Syria in April 1967, the Israeli Air Force shot down six Syrian MiG fighter jets. [...] In response to the apparent mobilization of its Arab neighbours, early on the morning of June 5, Israel staged a sudden preemptive air assault that destroyed more than 90 percent Egypt's air force on the tarmac. A similar air assault incapacitated the Syrian air force. Without cover from the air, the Egyptian army was left vulnerable to attack.
\\
Bregman, Ahron. 2011. “Arab—Israeli Conflict.” In The Encyclopedia of War, edited by Gordon Martel. & Imposing a blockade on the Straits of Tiran to all Israel-bound ships was precisely what President Nasser did ten years later on May 23, 1967. The debate continues as to why Nasser took this action, knowing full well that it amounted to a declaration of war. Perhaps Nasser – a self-declared leader of the Arab world – did it in response to growing pressure on him to stand up to Israel, or maybe he felt it was too good an opportunity to miss as Israel's old premier and defense minister, Levi Eshkol, who lacked any military experience, might not respond to the challenge. What ever the explanation, the blockade, along with other warlike actions such as removing UN observers from the Sinai, combined with bellicose rhetoric from Syria and Jordan, led to a significant escalation of tension in the Middle East. Feeling cornered, Israel decided to preempt any Arab attack and strike first.
\\
Siniver, Asaf. 2018. “Arab–Israeli Conflict.” In The Encyclopedia of Diplomacy, edited by Gordon Martel. & The June 1967 War, or Six-Day War, came as a result of a series of miscalculations by the Arabs, Israelis, and the Soviets. With the Egyptian and Syrian armies massed on its borders (though not assuming offensive postures), Israel launched pre-emptive strikes and within six days it captured the Gaza Strip and Sinai Peninsula from Egypt, the Golan Heights from Syria, and the West Bank and East Jerusalem from Jordan, which joined the war in the first day despite warnings from Israel to stay out of the fighting. The Six-Day War had dramatically changed the political and military balance of power in the Middle East. The Arabs and their Soviet backers were humiliated, with Nasser's pan-Arabism suffering a deadly blow. Arab attitudes towards Israel became even more intransigent and belligerent (...). The war also saw the resurrection of Palestinian national consciousness following the plight of 250,000 refugees from Gaza and the West Bank, while the Palestine Liberation Organization (PLO) gradually became a key non-state actor in the conflict by inflicting heavy casualties on Israeli civilians by means of terror.
\\
\end{longtable}}
\clearpage
         \textbf{Second Laotian, Phase 2} 

          \underline{Onset:} 1968         

\underline{Reasons:}
\begin{itemize}
\item Religion or Ideology
\end{itemize}

         \underline{Sources:} 
         \rowcolors{3}{white}{light-gray} 
         {\scriptsize 
         \begin{longtable}{p{0.2\textwidth}p{0.8\textwidth}} 
         \hiderowcolors% 
         \toprule 
         Source & Excerpt \\ 
         \midrule 
         \showrowcolors% 
         \endfirsthead% 
         \hiderowcolors% 
         \toprule 
         Source & Excerpt \\ 
         \midrule 
         \showrowcolors% 
         \endhead% 
         \midrule 
         \endfoot% 
         \hiderowcolors% 
         \bottomrule 
         \showrowcolors% 
         \endlastfoot% 
         
Gibler, Douglas M. 2018. International Conflicts, 1816–2010: Militarized Interstate Dispute Narratives. Vol. 2 Rowman \& Littlefield. & This dispute began in early 1962 when Pathet Lao communist rebels, with support from North Vietnam, laid siege to Nam Tha, near the Thai border. On February 13, Thai troops were deployed to the Lao border. On May 6, Pathet Lao forces attacked Nam Tha while Royal Lao forces ceded territory without a fight. The United States responded on May 12 by sending the Seventh Fleet to the Gulf of Siam and placing troops on alert. The 1,000 US marines already in Thailand for bilateral training were also moved to the border. Civil war fighting continued throughout the decade, with occasional transnational incidents affecting numerous neighbors and those countries fighting in Vietnam. Thailand took a direct role by training Laotian pilots, and South Vietnam trained Laotian officers. The United States and Soviets supplied the fighting forces and guarded their interests in the region. The overall fight was largely static until 1969, when Pathet Lao began to achieve battlefield victories. Then, by 1972, Pathet Lao controlled two-thirds of the country. Discussions between the Laotian government and the Pathet Lao began in 1970, and formal negotiations took place in Vientiane from October 1972. These talks did not advance, though, until the United States applied pressure to the Laotian government. The United States wanted to enter the Paris Peace talks for Vietnam with the Laotian war already concluded. On February 20, 1973, representatives from the Laotian government and Pathet Lao signed the Vientiane Agreement. (p. 921)
\\
Britannica. 2023. “History of Laos.” Accessed August 20, 2023. https://www.britannica.com/topic/history-of-Laos. & The uneasy peace in Laos was short-lived, as hostilities broke out between leftist and rightist factions in 1959. Another conference in Geneva in May 1961 culminated in an agreement in July 1962 that called for the country to become neutral and for a tripartite government to be formed. The new government consisted of factions from the left (the Pathet Lao, who were linked to North Vietnam), the right (linked to Thailand and the United States), and neutrals (led by Prince Souvanna Phouma). Once again, however, the cease-fire was brief. The coalition had split apart by 1964, and the larger war centred in Vietnam subsequently engulfed Laos. In that expanded war Laos, like Cambodia, was viewed by the major protagonists as a sideshow.
\\
Möise, Edwin. 2011. “Vietnam War (1959–1975).” In The Encyclopedia of War, edited by Gordon Martel. & The government of Laos, under Prince Souvanna Phouma, was nominally neutral but de facto allied with the United States.The United States had supported a genuinely neutral government under Souvanna Phouma until 1958, but then turned against him, preferring a rightist, anti-communist government. At the end of 1960, Souvanna Phouma was driven into an alliance with the Pathet Lao (the Laotian communist movement) and the People's Army of Vietnam (PAVN). Agreements signed in Geneva in 1962 stated that Laos was to be neutral once more, and returned Souvanna Phouma to office as prime minister. But in the months that followed, the PAVN significantly violated Laotian neutrality,especially along the Ho Chi Minh Trail. Souvanna Phouma drifted into an alliance with the United States, but he defined its terms. He allowed massive use of US airpower against communist forces in Laos from 1965 onward, but he did not want US ground troops.
There were two wars in Laos. In the North, PAVN and Pathet Lao forces fought the Royal Laotian Army and the “Secret Army” of Hmong tribesmen, armed and supported by the CIA. In the South, American aircraft attacked the PAVN forces along the Ho Chi Minh Trail. Air strikes were supplemented only by very small covert operations on the ground, small enough that American officials could reasonably hope Souvanna Phouma would not find out about them. (...).
In August 1973, the [US] Congress actually passed a law forbidding US military personnel from operating in Vietnam, Laos, or Cambodia (...). (...).
Pathet Lao and PAVN forces brought Laos under full communist control toward the end of 1975.

\\
Dumbrell, John. 2018. “Vietnam War.” In The Encyclopedia of Diplomacy, edited by Gordon Martel. & Kennedy considered various proposals for neutralization, notably from under secretary of state Chester Bowles. Tentative approaches were made to Hanoi and Beijing via Indian and Burmese intermediaries. In the event, however, the failure of the Laotian neutralization – by 1962 communist forces were protecting Laotian connecting routes from North to South Vietnam along the so-called “Ho Chi Minh Trail” – persuaded JFK against further diplomatic initiatives.
\\
\end{longtable}}
\clearpage
         \textbf{Football War} 

          \underline{Onset:} 1969         

\underline{Reasons:}
\begin{itemize}
\item Nationalism
\item Border Clashes
\item Economic, Short-Run
\end{itemize}

         \underline{Sources:} 
         \rowcolors{3}{white}{light-gray} 
         {\scriptsize 
         \begin{longtable}{p{0.2\textwidth}p{0.8\textwidth}} 
         \hiderowcolors% 
         \toprule 
         Source & Excerpt \\ 
         \midrule 
         \showrowcolors% 
         \endfirsthead% 
         \hiderowcolors% 
         \toprule 
         Source & Excerpt \\ 
         \midrule 
         \showrowcolors% 
         \endhead% 
         \midrule 
         \endfoot% 
         \hiderowcolors% 
         \bottomrule 
         \showrowcolors% 
         \endlastfoot% 
         
Gibler, Douglas M. 2018. International Conflicts, 1816–2010: Militarized Interstate Dispute Narratives. Vol. 1. Rowman \& Littlefield. & This dispute describes the Football War, which was fought over several issues. First, Honduran resentment built toward 300,000 El Salvadoran migrants who worked or squatted in agricultural plots and worked in factories. Second, Honduras and El Salvador contested portions of their border. Third, Honduras began to enforce ownership laws that restricted property rights for foreigners. These conflicts became palpable on June 8 when El Salvador lost a soccer match to Honduras 1-0 in overtime. Many El Salvadorans thought their team had been cheated. Rioting preceded the rematch in San Salvador on June 15, which El Salvador won 3-0. El Salvadorans threw rocks at Honduran team vehicles and broke windshields as the Hondurans made their way home. In response, Hondurans attacked Salvadoran shops in Tegucigalpa and San Pedro Sula, prompting 10,000 El Salvadorans to flee back to El Salvador. El Salvador broke diplomatic relations on June 26 and Honduras broke off relations on June 27. Early on July 14 three Honduran fighter jets entered Salvadoran territory near El Poy. The same day Salvadoran military planes attacked the Honduran capital and several other cities and ground troops entered Honduras. Although El Salvador managed to occupy parts of Honduras, neither side had the resources necessary to sustain an attack. Ammunition stocks were low and reportedly each side only possessed eight working military aircraft at the beginning of the conflict, all from World War II, and too few parts to keep them flying. On July 18 the Organization of American States (OAS), with pressure from the United States, called for a ceasefire, which both parties accepted. (The UN had also offered its good offices; and Costa Rica, Guatemala, and Nicaragua had offered mediation.) Fighting ended, but El Salvador refused to withdraw from Honduran territory. On July 30 the OAS passed additional resolutions calling on El Salvador to withdraw, which the Salvadorans accepted, and on October 30, 1980, El Salvador and Honduras signed the General Treaty of Peace. The treaty delimited the border where there was no disagreement and agreed to International Court of Justice (ICJ) resolution of the rest. The ICJ issued decisions in 1990 and 1992. In 1998 El Salvador and Honduras signed a border-demarcation treaty. (p. 85)
\\
Britannica. 2023. “El Salvador – Military Dictatorships.” Encyclopædia Britannica. Accessed August 19, 2023. https://www.britannica.com/place/El-Salvador/Military-dictatorships\#ref468021. & Most important was the problem raised by some 300,000 Salvadorans who had migrated to Honduras in search of land or jobs and who now found themselves threatened by an involuntary repatriation program begun by the Honduran government. Spurred by reports of the mistreatment of these refugees, the Salvadoran government opened hostilities on July 14, 1969.
\\
Cable, Vincent. 1969. “The ‘Football War’ and the Central American Common Market.” International Affairs (Royal Institute of International Affairs 1944–) 45 (4): 658–671. & Indeed, a war between the original 'banana republic', Honduras, and what looked like another obscure Central American military dictatorship, over issues that could not be easily labelled as racial, religious, ideological, linguistic or anticolonial and which had their immediate origin in a soccer match, appeared to deserve this lack of interest. (...). (...) [The] two protagonists are participants in one of the very few examples of successful economic integration amongst developing countries, which as such, has been held up as a prototype of the kind of arrangement which could well be emulated by the large and growing number of very small states. Ironically the origins of the war can be dated back to a phenomenon which in a common market of advanced market economies, such as the EEC, would be regarded as generally beneficial-the mobility of labour between partner states. For almost half a century, overcrowded El Salvador (...) has been exporting labour to Honduras, though attempts have been made to control the intake in recent years. Three hundred thousand have settled there-over 12.5 per cent of the total population of Honduras. While Salvadorian immigrants may have had a beneficial effect in a macro-economic sense (...) their presence has aroused resentment among the Hondurans who have faced competition on their own labour market. Animosity between El Salvador and Honduras is longstanding and certainly predates the formation of the Common Market in 1958. (...). (...) Honduras has begun to experience unemployment (...). The plantation sector has stagnated, and because of disease affecting the banana in Honduras the major banana companies (...) have turned their attention to Ecuador. The Salvadorians in Honduras are consequently held responsible for depressing the wage level and limiting opportunities for employment-grievances voiced by the powerful trade unions. (...). 
In addition to these economic and demographic factors, both governments have, in their different ways, done little to help. The Salvadorian Government has been criticised for exporting its social problems rather than solving them, and Honduran propagandists have drawn attention (...) to a feudal oligarchy of 'fourteen families' who supposedly own most private property in El Salvador. (...). The Honduran tradition is rather different: one in which the large concentration of wealth has been in foreign hands-notably the U.S. banana producers-rather than in those of a local élite, a situation perhaps less conductive to social divisiveness but more favorable to the development of nationalism. (...). El Salvador has complained about the manner in which this nationalistic fervour was carefully encouraged to help the Government overcome its immediate political problems-a strike of teachers, labour troubles generally and student disorders (...), Certainly the highly intemperate language used immediately after the first rumours trickled in from El Salvador of the defeat of Honduras on the football field was directly responsible for the anti-Salvadorian riots (...) which were the main pretext for the invasion by El Salvador.
\\
Rouquié, Alain, and Michel Vale. 1973. “Honduras – El Salvador, The War of One Hundred Hours: A Case of Regional 'Disintegration'.” International Journal of Politics 3 (3): 17–51. & The war did indeed originate in a soccer match. (...). Certain deplorable incidents occurred between supporters of the two teams. The return match, which took place in San Salvador on June 15, resulted in Salvadorian victory, which was followed by an unexpected outbreak of violence: hundreds of Hondurans who had travelled to support their team were attacked, and in some cases seriously wounded, by the crowd of Salvadorian fans. The same day in various parts of Honduras, particularly in the countryside, armed bands attacked Salvadorian residents, pillaging their property and forcing them to flee or hide. (...). (...) [The] Honduran press and radio (...) launched a virulent anti-Salvadorian campaign, amounting to a call for a pogrom. Also, on the eve of the first match, under a new agrarian reform program, about fifty Salavdorian families had been evicted from public lands they had been cultivating: only Hondurans by birth were to be allowed the benefit of this allocation of government lands. (...). However, if their ethnic and political history draws Honduras and El Salvador together, demographically and economically they remain far apart. (...). The common factor in such situations [of conflict] is the presence in one country of a large colony of immigrants from a neighboring country, the two being unequal in wealth and development, and between which long standing border disputes exist. In no other case in Latin America does the country of emigration enjoy a relatively more advanced development than the country receiving its surplus labor.  
\\
Tillema, Herbert K. International armed conflict since 1945: A bibliographic handbook of wars and military interventions. Routledge, 2019. & EI Salvador and Honduras disputed the location of border from the time of independence in the early nineteenth century. Negotiations eventually produced the Bonilla-Velasco agreement signed in 1895. The treaty was never ratified by either EI Salvador or Honduras, however, and later negotiations failed to produce a permanent understanding. Beginning in the 1920s, thousands of land-poor Salvadorans immigrated to work or squat within Honduran territory. In January 1969 the Honduran government announced social reforms that, if fully implemented, would result in deportation of as many as 300,000 Salvadorans who had failed to establish legal residence. In June 1970, Honduran and Salvadoran soccer teams met duringelimination rounds of the World Cup competition. Matches in Tegucigalpa, capital of Honduras, June 8, 1969, and in San Salvador, capital of EI Salvador, June 15 led to disputes, stadium violence, and rioting. EI Salvador dedared a state of alert. Formal diplomatic relations between the two countries were severed June 27. EI Salvador moved troops to the border. These conducted small raids upon Honduran border positions beginning July 8. On July 14, 1969, EI Salvador mounted a full-scale invasion that involved a combination of regular soldiers, designated militia and hastily recruited civilians. The Honduran Air Force undertook retaliatory attacks upon Salvadoran cities July 15 as well as attacked Salvadoran forces within Honduras. Fighting continued within Honduras until cease-fire was arranged July 18 with help of representatives of the Organization of American States. El Salvador withdrew from most of the territory she had conquered after a further agreement was reached July 30, 1969. (p. 27)
\\
\end{longtable}}
\clearpage
         \textbf{War of Attrition} 

          \underline{Onset:} 1969         

\underline{Reasons:}
\begin{itemize}
\item Power Transition or Security Dilemma
\item Religion or Ideology
\item Border Clashes
\end{itemize}

         \underline{Sources:} 
         \rowcolors{3}{white}{light-gray} 
         {\scriptsize 
         \begin{longtable}{p{0.2\textwidth}p{0.8\textwidth}} 
         \hiderowcolors% 
         \toprule 
         Source & Excerpt \\ 
         \midrule 
         \showrowcolors% 
         \endfirsthead% 
         \hiderowcolors% 
         \toprule 
         Source & Excerpt \\ 
         \midrule 
         \showrowcolors% 
         \endhead% 
         \midrule 
         \endfoot% 
         \hiderowcolors% 
         \bottomrule 
         \showrowcolors% 
         \endlastfoot% 
         
Gibler, Douglas M. 2018. International Conflicts, 1816–2010: Militarized Interstate Dispute Narratives. Vol. 2 Rowman \& Littlefield. & Following the loss in the devastating Six-Day War (MID\#345), Egypt was determined to reassert itself and reestablish control of the Sinai. Egypt maintained a constant state of military activity on one side of the Suez Canal, and Israel matched that activity across the new border. Low-level hostilities in the War of Attrition broke out into active fighting between February 1969 and August 1970. The Soviet Union supplied Egypt most of its arms and worked quickly to replace the United Arab Republic air force's planes and missiles that were destroyed during the war. However, the technology significantly trailed that available to the Israelis, and Israel retaliated on a far greater scale following each Egyptian attack. The conflict came to a head in July 1970 when Israel shot down four Egyptian planes flown by Soviet pilots. Fearing Soviet retaliation and unsure of American support, Israel accepted a ceasefire agreement on August 7, 1970. (p. 665)
\\
Britannica. 2020. “War of Attrition.” Accessed August 20, 2023. https://www.britannica.com/event/War-of-Attrition-1969-1970. & War of Attrition, inconclusive war (1969–70) chiefly between Egypt and Israel. The conflict, launched by Egypt, was meant to wear down Israel by means of a long engagement and so provide Egypt with the opportunity to dislodge Israeli forces from the Sinai Peninsula, which Israel had seized from Egypt in the Six-Day (June) War of 1967. Shortly after the end of the 1967 war, Egyptian Pres. Gamal Abdel Nasser made clear his intention to retake by force territory captured by Israel in the conflict. Egyptian losses in the war had been significant, but the support and material investment of the Soviet Union sped Egypt's recovery.
\\
Bregman, Ahron. 2011. “Arab—Israeli Conflict.” In The Encyclopedia of War, edited by Gordon Martel. & The Egyptian army, though badly beaten, had not been destroyed in the 1967 war, and reequipped by the Soviets with new arms, it attacked the IDF, which was now deployed along the eastern bank of the Suez Canal. The first major incident between Egypt and Israel after the June 1967 war took place on October 21, 1967, when an Egyptian destroyer torpedoed and sank the Israeli destroyer Eilat not far from Port Said. (...).
These [clashes] (...) were (...) part of a well-planned Egyptian military program which envisaged a total war against Israel in three main phases. (...).
President Nasser explained, The first priority (...) is the military front, for we must realize that the [Israeli] enemy will not withdraw [from land it occupied in 1967] unless we force him to withdraw through fighting.
\\
Britannica. 2023. “Six-Day War.” Accessed August 20, 2023. https://www.britannica.com/event/Six-Day-War. & In response to the apparent mobilization of its Arab neighbours, early on the morning of June 5, Israel staged a sudden preemptive air assault that destroyed more than 90 percent Egypt's air force on the tarmac. A similar air assault incapacitated the Syrian air force. Without cover from the air, the Egyptian army was left vulnerable to attack. Within three days the Israelis had achieved an overwhelming victory on the ground, capturing the Gaza Strip and all of the Sinai Peninsula up to the east bank of the Suez Canal.


\\
\end{longtable}}
\clearpage
         \textbf{Communist Coalition} 

          \underline{Onset:} 1970         

\underline{Reasons:}
\begin{itemize}
\item Nationalism
\item Power Transition or Security Dilemma
\item Religion or Ideology
\end{itemize}

         \underline{Sources:} 
         \rowcolors{3}{white}{light-gray} 
         {\scriptsize 
         \begin{longtable}{p{0.2\textwidth}p{0.8\textwidth}} 
         \hiderowcolors% 
         \toprule 
         Source & Excerpt \\ 
         \midrule 
         \showrowcolors% 
         \endfirsthead% 
         \hiderowcolors% 
         \toprule 
         Source & Excerpt \\ 
         \midrule 
         \showrowcolors% 
         \endhead% 
         \midrule 
         \endfoot% 
         \hiderowcolors% 
         \bottomrule 
         \showrowcolors% 
         \endlastfoot% 
         
Gibler, Douglas M. 2018. International Conflicts, 1816–2010: Militarized Interstate Dispute Narratives. Vol. 2 Rowman \& Littlefield. & On March 28, 1969, Cambodian Head of State Sihanouk acknowledged and denounced the presence of Vietnamese Communists in Cambodia in an effort to persuade them to leave the country. North Vietnamese soldiers were encamped in the northeastern provinces of Ratanakiri, Mondulkiri, and Stung Treng, and in the central provinces of Kompong Cham, Prey Vang and Svay Rieng. The North Vietnamese continued to build up their bases in Cambodia, and clashes became inevitable. The Cambodian prime minister, after June, reported that North Vietnamese soldiers had attacked Cambodian troops and army equipment. On October 16, 1969, Prime Minister Lon Nol released a statistical summary of why there had been so many clashes between Cambodian and North Vietnamese troops. Cambodia then demanded that North Vietnam withdraw its troops and Viet Cong from Cambodian territory. In March 1970, a coup deposed the Cambodian prince Sihanouk. Sihanouk formed an opposition in exile to the new government supported by North Vietnam, which invaded Cambodia with its own troops. Cambodia declared war, and the dispute ends with their participation in the Vietnam War (see MID\#611) (pp. 913-914)
\\
Pradhan, Prakash. C. “Cambodian Crisis of 1970.” Proceedings of the Indian History Congress. & The origin of this crisis goes back to the year 1970 when three important events occurred and changed the course of Cambodian history. Firstly , Norodom Sihanouk who had guided the destiny of millions of Cambodians for years together was ousted from power by a coup d'etat staged b General Lon Noi and Sirik Matak. Second , Cambodia which had been successfully keeping away from the flames of the Vietnam war, burnt itself in the same war when direct intervention was launched by the U.S. and South Vietnamese troops against Cambodia. Finally , monarchy which had been a vital factor in Cambodian unity was given a terrible blow when Cambodia was declared a Republic. 
\\
Chandler, David. 2018. A History of Cambodia. 4th ed. Routledge. & (...) Sisowath Sirik Matak, a career civil servant who had become deputy prime minister under Lon Nol. Matak had grown impatient with Sihanouk's mismanagement of the economy, and he was dismayed by the presence of Vietnamese bases on Cambodian soil and by his cousin's impulsive, contradictory foreign policy. Pro-Western himself, and with links to Phnom Penh's commercial elite, Matak was tired of playing a supporting role in Sihanouk's never-ending opera. Lon Nol, the prime minister, was a more enigmatic figure. He is not known to have objected to the prince about the cutoff of U.S. aid, the alliance with the Communists, or the shipments of arms through Cambodian territory. Indeed, many of his officers became rich by dealing in arms, medicines, and supplies for which the Vietnamese paid generously. By 1969, however, Cambodian troops were under fire from Communist insurgents, and the Vietnamese administration of many base areas had drawn complaints from local people. Lon Nol was also under pressure from some of his officers who saw Cambodia's isolation from U.S. aid and, by implication, from the Vietnam War as impediments to their financial ambitions. Moreover, Lon Nol was not immune to flattery, and he came to see himself as uniquely capable of saving Cambodia from the thmil, or “unbelievers,” as he called the Vietnamese hammering at his country's gates. (...). [Sihanouk's] decision to open a casino in Phnom Penh to raise revenue had disastrous results. In the last six months of 1969, thousands of Cambodians lost millions of dollars at its tables. Several prominent people, and dozens of impoverished ones, committed suicide after sustaining losses, and hundreds of families went bankrupt. Sihanouk, no gambler himself, was indifferent to the chaos he had caused, and his own expenditures continued to mount. A climax of sorts occurred in November, when an international film festival, stage-managed by the prince, ended with one of his own films, “Twilight,” being awarded a solid-gold statue sculpted from ingots donated for the purpose by Cambodia's national bank. When Sihanouk left the country for his annual holiday in January 1970, many people who remained behind interpreted his departure as a flight. (...). The National Assembly, as it was entitled to do, voted 86–3 to remove its confidence from the prince and to replace him as chief of state, pending elections, with the relatively colorless president of the assembly, Cheng Heng. Lon Nol remained prime minister with Matak as his assistant. The coup was popular among educated people in Phnom Penh and in the army, but rural Cambodians were totally unprepared for it. (...). [Sihanouk's] first thought was to seek political asylum in France, but after talks with Zhou Enlai and the Vietnamese premier, Pham Van Dong, he agreed to take command of a united front government allied to North Vietnam, whose Cambodian forces would consist largely of the Communists his army had been struggling to destroy only a month before.
\\
\end{longtable}}
\clearpage
         \textbf{Bangladesh} 

          \underline{Onset:} 1971         

\underline{Reasons:}
\begin{itemize}
\item Nationalism
\item Economic, Long-Run
\end{itemize}

         \underline{Sources:} 
         \rowcolors{3}{white}{light-gray} 
         {\scriptsize 
         \begin{longtable}{p{0.2\textwidth}p{0.8\textwidth}} 
         \hiderowcolors% 
         \toprule 
         Source & Excerpt \\ 
         \midrule 
         \showrowcolors% 
         \endfirsthead% 
         \hiderowcolors% 
         \toprule 
         Source & Excerpt \\ 
         \midrule 
         \showrowcolors% 
         \endhead% 
         \midrule 
         \endfoot% 
         \hiderowcolors% 
         \bottomrule 
         \showrowcolors% 
         \endlastfoot% 
         
Gibler, Douglas M. 2018. International Conflicts, 1816–2010: Militarized Interstate Dispute Narratives. Vol. 2. Rowman \& Littlefield. & As Pakistan considered a new constitution, unrest grew in East Pakistan, including demonstrations, strikes, civil disobedience, and refusals to pay taxes. On March 25, 1971, the Pakistani military struck Dacca, and in the following months the military burned villages and killed civilians. Over 10 million refugees fled to India, stressing the Indo-Pakistani relationship. On December 3, 1971, India and Pakistan went to war. The UN Security Council began deliberations on December 4, but because the council could not agree on a resolution it referred the matter to the General Assembly. On December 7, the General Assembly passed Resolution 2793 calling for an immediate ceasefire and a mutual troop withdrawal. India refused to comply with the resolution, so the council once again took up the issue. On December 16, Pakistani troops in East Pakistan surrendered, and on December 17, the Indian forces implemented a unilateral ceasefire. Four days later the Security Council passed Resolution 307 calling for a ceasefire while both sides withdrew troops. Afraid the loss of Kashmir in addition to the loss of East Pakistan would destabilize Pakistan, India did not attempt to seize Kashmir. On July 3, 1972, India and Pakistan signed the Simla Agreement, formally ending the dispute. The Simla Agreement cemented the Indian and Pakistani positions as the new line of control in Kashmir, granting more territory to India. (p. 876)
\\
The National Archives. 2023. “The Independence of Bangladesh in 1971.” Accessed August 19, 2023. https://www.nationalarchives.gov.uk/education/resources/the-independence-of-bangladesh-in-1971. & The separate north-eastern and north-western areas of the country, which were mostly Muslim, became a ‘united' Pakistan. The rest of the country, mostly non-Muslim, became known as India. [...] West and East Pakistan shared a religion, but not much else. For decades after Partition, the East Pakistanis (present-day Bangladeshis) were treated unfairly by the West Pakistani government over 1,000 miles away. [...] Bengalis have a rich and proud history and culture focused around language, art, food, fashion, community, family and religion. The 1971 Liberation War can be seen as a struggle to preserve and protect this heritage. [...] Maybe the most important event leading up to the war was the 1970 election. The winning Awami League Party was led by Sheikh Mujibur Rahman, who wanted more freedom and independence for East Pakistan. He won the national elections by a clear majority, as East Pakistan had a much larger population than West Pakistan. However, the centralised government in West Pakistan did not accept Rahman as leader. They declared the election results void. 
\\
Barua, Pradeep P. 2011. “Indo-Pakistani Wars (1947–1948, 1965, 1971).” In The Encyclopedia of War, edited by Gordon Martel. & The Third Indo-Pakistan War was caused not by events in Kashmir, but by the outbreak of civil war in East Pakistan. In December 1970, Bengali Muslim nationalists won an overwhelming victory in the East Pakistan elections. General Yahya Khan, the military dictator, refused to accept this result and ordered the Pakistani military garrison in East Pakistan to seize control, triggering a violent civil war. Millions of refugees flooded into India, creating immense pressure on the Indian government. In April 1971 the prime minister of India ordered General Sam Manekshaw, the army chief and the chairman of the joint chiefs, to prepare for military intervention in East Pakistan in support of the Bengali Muslims who were fighting to create the independent state of Bangladesh.
\\
Constable, Philip. 2025. “Kashmir Dispute since 1947.” In The Encyclopedia of Diplomacy, edited by Gordon Martel. & The repercussions of this second Indo-Pakistan War accelerated demands for autonomy by the Awami League led by Sheik Mujibur Rahman in East Pakistan. With East Pakistan separated geographically from West Pakistan by India, the Awami League realized that they had much to lose in economic, political, and security terms by involvement in West Pakistan's conflict with India over Kashmir. India's prime minister Indira Gandhi also saw benefit in fostering East Pakistan's secession from West Pakistan because it would reduce the two-sided military threat to India from Pakistan and undermine Pakistan's two-nation theory as a claim on Kashmir. When the Awami League won a majority in East Pakistan in the 1970 elections and Pakistan's president General Yahya Khan annulled the elections through martial law, India supported the Awami League's guerrilla war against the West Pakistan army leading to a third regional war on December 3, 1971. (...). The applicability of Pakistan's two-nation theory as justification for its right to Jammu-Kashmir was undermined, while Bangladesh's secession contributed to continuing Kashmiri hopes of independence from both India and Pakistan.
\\
Morillo, Stephen, Jeremy Black, and Paul Lococo. 2008. War in World History: Society, Technology, and War from Ancient Times to the Present. New York: McGraw-Hill. & In South Asia, for example, the end of British rule led to wars between India and Pakistan in 1947, 1965, and 1971. Pakistan's defeat in the last ensured success for the rebellion in East Pakistan, which became independent as Bangladesh.
\\
\end{longtable}}
\clearpage
         \textbf{Yom Kippur War} 

          \underline{Onset:} 1973         

\underline{Reasons:}
\begin{itemize}
\item Power Transition or Security Dilemma
\item Religion or Ideology
\item Border Clashes
\item Economic, Long-Run
\end{itemize}

         \underline{Sources:} 
         \rowcolors{3}{white}{light-gray} 
         {\scriptsize 
         \begin{longtable}{p{0.2\textwidth}p{0.8\textwidth}} 
         \hiderowcolors% 
         \toprule 
         Source & Excerpt \\ 
         \midrule 
         \showrowcolors% 
         \endfirsthead% 
         \hiderowcolors% 
         \toprule 
         Source & Excerpt \\ 
         \midrule 
         \showrowcolors% 
         \endhead% 
         \midrule 
         \endfoot% 
         \hiderowcolors% 
         \bottomrule 
         \showrowcolors% 
         \endlastfoot% 
         
Gibler, Douglas M. 2018. International Conflicts, 1816–2010: Militarized Interstate Dispute Narratives. Vol. 2 Rowman \& Littlefield. & From 1971 until approximate September 1973, there were some small border conflicts on both the Syrian and Egyptian borders; however, most of the conflict took place in the Golan Heights on the Syrian border and in the Sinai area of Egypt. In September 1973, Syria and Egypt, with a coalition of Arab countries that included Iraq, Jordan, and Saudi Arabia, began planning a coordinated surprise attack on Israel. They decided that Yom Kippur, one of the holiest days on the Jewish calendar, would be the best day to launch such a war and on October 6, 1973, all sides attacked. Egypt focused primarily on quickly breaking the Israeli defenses in the Sinai and holding that territory. In the Golan, where most of the fighting was concentrated due to its proximity to a majority of the Israeli population, Syria attempted to launch a quick offensive but was not as successful as the Egyptians. It became quickly apparent that Israel had the military upper hand as the Israeli Defense Forces took control of the Golan Heights and then advanced into Egypt, coming within 100 kilometers of Cairo. At this point, the UN Security Council passed a resolution on October 22 requesting that all parties cease hostilities, which they agreed to the next day. The official separation of forces agreements proposed by Israel was agreed to January 18, 1974, for Egypt, and May 31, 1974, for Syria. (p. 688)
\\
Britannica. 2023. “Yom Kippur War.” Accessed August 20, 2023. https://www.britannica.com/event/Yom-Kippur-War. & Yom Kippur War, also called the October War, the Ramadan War, the Arab-Israeli war of October 1973, or the Fourth Arab-Israeli War, fourth of the Arab-Israeli wars, which was initiated by Egypt and Syria on October 6, 1973, on the Jewish holy day of Yom Kippur. It also occurred during Ramadan, the sacred month of fasting in Islam, and it lasted until October 26, 1973. The war, which eventually drew both the United States and the Soviet Union into indirect confrontation in defense of their respective allies, was launched with the diplomatic aim of persuading a chastened—if still undefeated—Israel to negotiate on terms more favourable to the Arab countries.
\\
Bregman, Ahron. 2011. “Arab—Israeli Conflict.” In The Encyclopedia of War, edited by Gordon Martel. & After the 1967 war, Israel made it clear that she was reluctant to return the captured lands. She embarked on a creeping annexation, building settlements on the seized territories and exploiting resources such as oil in the Sinai and water in the Gaza Strip and West Bank. This deeply upset the Arabs. What's more, it seemed that the superpowers of the time – the United States and the Soviet Union – were enjoying an unusual period of détente and were reluctant to have their Middle Eastern clients ruining the improved atmosphere. Both, therefore, seemed to accept the new status quo and ignore Israel's gradual annexation of the seized lands. To break the deadlock and prevent the annexation becoming permanent, Egypt and Syria decided to launch a military attack on Israel, to liberate at least some of their lost land and perhaps force Israel into diplomatic negotiations over withdrawal from the rest. Egypt and Syria decided to attack on October 6, 1973, which was Yom Kippur, the holiest day in the Jewish calendar, thus catching Israel by surprise and unprepared. (...).
[The Israeli's] first priority was to contain the Syrian invasion of the Golan, where there was no strategic depth and Jewish settlements were close to the frontline (...).
\\
Siniver, Asaf. 2018. “Arab–Israeli Conflict.” In The Encyclopedia of Diplomacy, edited by Gordon Martel. & The tension between [the Soviet Union and the United States] over the conflict reached unprecedented levels in the last days of the Yom Kippur/October War of 1973. The war began with a combined surprise attack by Egypt and Syria on Israeli positions in the Sinai Peninsula and Golan Heights respectively, on Yom Kippur, the holiest day in the Jewish calendar. During the three-week war both superpowers supported their respective allies with arms; by the war's end the high stakes had brought the superpowers as close as they had been to a direct confrontation since the 1962 Cuban Missile Crisis, with American forces around the world placed on the highest level of alert short of nuclear war.
\\
Britannica. 2023. Six-Day War Accessed August 20, 2023. https://www.britannica.com/event/Six-Day-War & Prior to the start of the war, attacks conducted against Israel by fledgling Palestinian guerrilla groups based in Syria, Lebanon, and Jordan had increased, leading to costly Israeli reprisals. In November 1966 an Israeli strike on the village of Al-Samū' in the Jordanian West Bank left 18 dead and 54 wounded, and, during an air battle with Syria in April 1967, the Israeli Air Force shot down six Syrian MiG fighter jets. [...] In response to the apparent mobilization of its Arab neighbours, early on the morning of June 5, Israel staged a sudden preemptive air assault that destroyed more than 90 percent Egypt's air force on the tarmac. A similar air assault incapacitated the Syrian air force. Without cover from the air, the Egyptian army was left vulnerable to attack.
\\
\end{longtable}}
\clearpage
         \textbf{Turco-Cypriot} 

          \underline{Onset:} 1974         

\underline{Reasons:}
\begin{itemize}
\item Nationalism
\item Religion or Ideology
\item Economic, Long-Run
\end{itemize}

         \underline{Sources:} 
         \rowcolors{3}{white}{light-gray} 
         {\scriptsize 
         \begin{longtable}{p{0.2\textwidth}p{0.8\textwidth}} 
         \hiderowcolors% 
         \toprule 
         Source & Excerpt \\ 
         \midrule 
         \showrowcolors% 
         \endfirsthead% 
         \hiderowcolors% 
         \toprule 
         Source & Excerpt \\ 
         \midrule 
         \showrowcolors% 
         \endhead% 
         \midrule 
         \endfoot% 
         \hiderowcolors% 
         \bottomrule 
         \showrowcolors% 
         \endlastfoot% 
         
Gibler, Douglas M. 2018. International Conflicts, 1816–2010: Militarized Interstate Dispute Narratives. Vol. 1 Rowman \& Littlefield. & This dispute is the Turkish-Cypriot War fought in the summer of 1974. Cyprus was always a difficult issue for Turkey and Greece once Greece gained independence from the Ottoman Empire in the 19th century. Britain effectively quelled Greek and Turkish efforts to incorporate the territory by occupying it themselves and eventually reorganizing it as a colony. However, the 1959 London-Zurich Agreements made Cyprus independent and, by the terms of the agreement, indivisible. Though predominantly Greek, a small but active Turkish minority persistently resulted in the government in Ankara using its power to exercise leverage over Cyprus and Greece. In 1974, the colonels' regime in Greece initiated a coup against sitting head of state Makarios III. Though Greek himself, Makarios III was no friend of the Greek colonels' regime and wanted them to vacate their military bases. Rather than do so, Greece removed him on July 15, via the sympathetic Cypriot national guard. Turkey suspected that the ultimate end game was incorporation of Cyprus into Greece and acted fast. It already had a military presence on the island but mobilized an additional 8,000 troops and dispatched them to Cyprus. More rapid mobilization put the Turkish numbers at 40,000 on the island. An ultimatum was delivered to the Greek regime on July 18: Greece was to withdraw or suffer the consequences. Greece refused, and the Turks responded with an invasion of Cyprus on July 20. The Turks were superior in all aspects to the Greek forces, and the embarrassment here was the leading cause of the end of the colonels' regime in Greece. After failed negotiations from August 8 to 14, the Turks resumed military activities. A second ceasefire on August 16 ended the conflict. The result of the invasion was a new de facto partition of Cyprus. The de facto independent and Turkey-supported Northern Cyprus covers the northeast of the island and accounts for almost 40 percent of the island's territory. A UN-administered buffer zone separates Turkish Cyprus and Greek Cyprus while the remainder that was not the UN Green Line or territory allocated to Britain via the 1959 agreement became Greek Cypriot. (pp. 366-367)
\\
Bishku, Michael B. 1991. “Turkey, Greece and the Cyprus Conflict.” Journal of Third World Studies. & Cyprus' approximately 650,000 inhabitants, about eighty percent are Greekspeaking Christians while roughly twenty percent are Turkishspeaking Muslims; [...] Cyprus has been aptly described as an ethnic omelette that has been unscrambled. Such has been the case since 1963, just three years after Cyprus' independence from Great Britain, when a constitutional crisis and intercommunal violence brought an end to governmental power sharing and led to the development of fortified Turkish enclaves. This separation was solidified in 1974 when the Greek junta's attempt to annex Cyprus led the Turks to intervene militarily and to occupy more than one-third of the island.
\\
Johnson, Edward. 2025. “Cyprus Crisis of 1974.” In The Encyclopedia of Diplomacy, edited by Gordon Martel. & Cyprus occupies a strategic position in the eastern Mediterranean (...). Culturally and demographically, the island was predominantly, but not exclu-sively, Greek (...). (...) [Cyprus] was leased to Britain in 1878 as part of a defense agreement with Russia. (...). The British annexed Cyprus in 1914 and took formal, legal control of the island in 1923 (...) faced with a longstanding demand from the Greek Cypriots for énosis, or union with Greece. The British refused to accede to these demands and, in 1955, the Greek Cypriots waged a terrorist campaign (...) against them (...) to force (...) énosis. The British government was conscious (...) any ceding of union (...) would (...) undermine the Turkish minority (...) threaten a Turkish military response. (...) British were faced with two competing demands: énosis from the Greek Cypriots and Takism or partition, from the Turkish minority. (...) British granted neither, but in 1960 imposed independence on Cyprus with conditions. (...). (...) intercommunal fighting broke out in late 1963 (...). [In] Greece, in 1967 members of the army staged a coup d'état, overthrew the monarchy and instituted a right-wing military junta which embarked on a regime of torture and brutality towards its political opponents. On Cyprus, Makarios, himself a Greek Cypriot, appeared by the mid-1970s to have abandoned any idea of énosis, and was intent upon maintaining the island as an independent and non-aligned state outside NATO, even being prepared to court the Soviet Union as well as seeking the support of AKEL (...), the Cypriot Communist Party. It was against this background of political inertia over the delivery of énosis that on July 15, 1974, the Greek junta in Athens, using the Cypriot National Guard, engineered a coupd'état and removed Makarios from power, even at one point trying to assassinate him and replacing him with Nikos Sampson, a former member of EOKA and a supporter of énosis. (...). It fell then to Britain and Turkey to defend the Cypriot constitutional position and by extension that of Makarios.
\\
\end{longtable}}
\clearpage
         \textbf{War over Angola} 

          \underline{Onset:} 1975         

\underline{Reasons:}
\begin{itemize}
\item Nationalism
\item Power Transition or Security Dilemma
\item Religion or Ideology
\end{itemize}

         \underline{Sources:} 
         \rowcolors{3}{white}{light-gray} 
         {\scriptsize 
         \begin{longtable}{p{0.2\textwidth}p{0.8\textwidth}} 
         \hiderowcolors% 
         \toprule 
         Source & Excerpt \\ 
         \midrule 
         \showrowcolors% 
         \endfirsthead% 
         \hiderowcolors% 
         \toprule 
         Source & Excerpt \\ 
         \midrule 
         \showrowcolors% 
         \endhead% 
         \midrule 
         \endfoot% 
         \hiderowcolors% 
         \bottomrule 
         \showrowcolors% 
         \endlastfoot% 
         
Gibler, Douglas M. 2018. International Conflicts, 1816–2010: Militarized Interstate Dispute Narratives. Vol. 1 Rowman \& Littlefield. & Warring parties in Angola were divided between pro-Western and Communist factions. In August 1975, Cuba began to send troops and equipment to assist the People's Movement for the Liberation of Angola (MPLA) faction. The South African army then began intervening in the conflict shortly after Cuba, on September 24, when an officer was sent to help halt an MPLA advance on Huambo. Zambia made a brief entrance into the dispute in 1976, when President Kaunda stated to his people on January 28, 1976, that the country was “at war.” Zambia's presence in the conflict on the side of South Africa was short-lived, however, as the country restored relations with Angola the following September. Zaire's involvement in the conflict was announced in December when reports confirmed that the country had 11,200 troops in Angola, assisting the pro-Western factions. Zaire and Angola began to improve relations in early 1978. The civil conflict continued, as did the dispute between Angola and Cuba on one side and South Africa on the other. Peace talks began in July 1988, and South Africa began withdrawing its troops on August 10 and completed the withdrawal by the end of the month. (p. 480)
\\
Britannica. 2023. “Angola – Independence and Civil War.” Accessed August 20, 2023. https://www.britannica.com/place/Angola/Independence-and-civil-war. & The three liberation movements proved unable to constitute a united front after the Portuguese coup. The FNLA's internal support had dwindled to a few Kongo groups, but it had strong links with the regime in Zaire and was well armed; it thus made a bid to seize Luanda by force. The MPLA, with growing backing from the Portuguese Communist Party, Cuba, and the Soviet Union, defeated this onslaught and then turned on UNITA, chasing its representatives out of Luanda. UNITA was militarily the weakest movement, but it had the greatest potential electoral support, given the predominance of the Ovimbundu within the population, and it thus held out most strongly for elections. But the Portuguese army was tired of war and refused to impose peace and supervise elections. The Portuguese therefore withdrew from Angola in November 1975 without formally handing power to any movement, and nearly all the European settlers fled the country.
\\
Van der Waag, Ian. 2011. “Angolan Civil Wars (1975–2002).” In The Encyclopedia of War, edited by Gordon Martel. & The civil war was a continuation of the liberation war (1961–1975). As the Portuguese colonial empire collapsed, the nationalist movements, fractured along ethnic and ideological lines, intensified their struggles with each other. The Movimento Popular de Libertc¸a˜o de Angola (MPLA), essentially a Marxist party founded by “mesticos,” proclaimed a government in Luanda. A rival government was proclaimed in Ambriz by the Frente Nacional de Libertac¸a˜o de Angola (FNLA) and Unia˜o Nacional para a Independeˆncia Total de Angola (UNITA), groups that enjoyed strong support among the Bakongo and Ovimbundu. The MPLA, occupying the colonial capital and the major ports, drew Soviet and Cuban support. Cuba, acting independently of Moscow, soon took the lead (Gleijeses 2002). The FNLA looked to the United States and Zaire, and UNITA to South Africa, a country that was soon drawn to more direct involvement as a result of a growing refugee crisis and the presence of Soviet-bloc troops.
\\
\end{longtable}}
\clearpage
         \textbf{Second Ogaden War, Phase 2} 

          \underline{Onset:} 1977         

\underline{Reasons:}
\begin{itemize}
\item Nationalism
\item Religion or Ideology
\item Economic, Long-Run
\end{itemize}

         \underline{Sources:} 
         \rowcolors{3}{white}{light-gray} 
         {\scriptsize 
         \begin{longtable}{p{0.2\textwidth}p{0.8\textwidth}} 
         \hiderowcolors% 
         \toprule 
         Source & Excerpt \\ 
         \midrule 
         \showrowcolors% 
         \endfirsthead% 
         \hiderowcolors% 
         \toprule 
         Source & Excerpt \\ 
         \midrule 
         \showrowcolors% 
         \endhead% 
         \midrule 
         \endfoot% 
         \hiderowcolors% 
         \bottomrule 
         \showrowcolors% 
         \endlastfoot% 
         
Gibler, Douglas M. 2018. International Conflicts, 1816–2010: Militarized Interstate Dispute Narratives. Vol. 1 Rowman \& Littlefield. & Somalia fought a year-long war with Ethiopia over the Ogaden, a predominantly Somali and Muslim enclave in internationally recognized Ethiopian territory. Somali interest was largely irredentist and supportive of the Western Somali Liberation Front operating within Ethiopia. Cuba's foray into the conflict reversed their traditional role. Rather than support the Somali Communists against the Westernbacked Ethiopian government, Cuban troops came to the assistance of the Ethiopians as the Soviets tried to displace the United States as the most significant foreign ally of Ethiopian interests. The joint forces of the Cubans and Ethiopians were too much for the Somali forces, who conceded defeat on March 9, 1978, and started their withdrawal. (...). 
On April 14, 1978, the Ethiopian ambassador to Kenya threatened to invade Somalia if it did not end its support for the guerrillas within Ethiopia. The Somali presence did not cease, and reports continued of Ethiopians being killed in small numbers on a regular basis. Two months later on June 22, the Somali government reported air raids on the towns of Borama and Gebile and villages between the city of Hargeisa and the Ethiopian border. Six months later, on December 21, there were Somali raids in the Ogaden area. Then, on February 7, 1979, a Somali cabinet minister traveled to the United States to request assistance against what he said had been 151 air raids by planes based out of Ethiopia in the last nine months. (pp. 472-473)
\\
Lewis, Ioan M. 1989. “The Ogaden and the Fragility of Somali Segmentary Nationalism.” African Affairs. & The Ogaden, who traditionally live as herdsmen in the region named after them in eastern Ethiopia, bordering the Somali Republic, are part of the wider Darod family of Somali clans. They attracted world attention a decade ago during their abortive attempt to secure independence from Ethiopia in the 1977-78 war between Somalia and Ethiopia. In retrospect, it is easy to see that this represented the high point of Somali nationalist fervour in recent Somali history. It was also a testimony to the strenuous efforts of the present Somali head of state in seeking to transform Somali nationalism from its old segmentary style to a modern organic mode
\\
Stapleton, Timothy J. 2011. “Ethiopia–Somalia Conflicts (1960s–2000s).” In The Encyclopedia of War, edited by Gordon Martel. & During decolonization in the 1950s, British and Italian Somaliland united to form the Republic of Somalia, which became independent in 1960. In a 1958 referendum, French Somaliland had voted against joining Somalia and gained independence as Djibouti in 1977. Pan-Somalism, the desire to unite all Somalis in one nation-state, gained popularity and prompted armed insurgency (...). Staging a military coup in Somalia in 1969, Siad Barre aspired to create a “Greater Somalia.” (...). To gain a naval base at Berbera that would threaten western oil shipments, the Soviet Union militarily assisted the Barre regime. A 1974 military coup overthrew the Ethiopian Emperor Haile Selassie, a long-standing western ally, and created a period of instability in which separatist movements gained strength. This presented Barre with an opportunity to take the Ogaden. (...). 
In early 1977 Mengistu Haile Mariam emerged as the military ruler of Ethiopia and launched offensives against separatists in Eritrea, which would deprive Ethiopia of access to the sea if independent, and the Ogaden. Mengistu, who was a Marxist-Leninist, also expelled American military advisors and by mid-1977 imported Soviet weapons. In June 1977, five thousand Somali soldiers, posing as WSLF, attacked towns in eastern Ethiopia and were repulsed. This probing operation signaled the shift from guerilla to conventional warfare. In mid-July Somalia launched an all-out invasion of Ethiopia, utilizing tanks and mechanized infantry supported by air power, and quickly gained control of the Ogaden lowlands. (...). Initially supplying both sides, the Soviets threw all their support behind Ethiopia in late 1977. There were a number of reasons for this. Soviet arbitration of the conflict had failed, Barre expelled around six thousand Soviet advisors, Mengistu was ideologically suitable, and Ethiopia a potentially stronger ally.

\\
\end{longtable}}
\clearpage
         \textbf{Vietnamese-Cambodian} 

          \underline{Onset:} 1977         

\underline{Reasons:}
\begin{itemize}
\item Nationalism
\item Power Transition or Security Dilemma
\item Border Clashes
\item Domestic Politics
\end{itemize}

         \underline{Sources:} 
         \rowcolors{3}{white}{light-gray} 
         {\scriptsize 
         \begin{longtable}{p{0.2\textwidth}p{0.8\textwidth}} 
         \hiderowcolors% 
         \toprule 
         Source & Excerpt \\ 
         \midrule 
         \showrowcolors% 
         \endfirsthead% 
         \hiderowcolors% 
         \toprule 
         Source & Excerpt \\ 
         \midrule 
         \showrowcolors% 
         \endhead% 
         \midrule 
         \endfoot% 
         \hiderowcolors% 
         \bottomrule 
         \showrowcolors% 
         \endlastfoot% 
         
Gibler, Douglas M. 2018. International Conflicts, 1816–2010: Militarized Interstate Dispute Narratives. Vol. 2 Rowman \& Littlefield. & A series of conflicts between Vietnam and Cambodia occurred almost immediately after North Vietnam eliminated South Vietnam and after the Khmer Rouge had come to power in Cambodia. The two former allies were quickly at odds. The origins of the conflict largely lay in the legacy of French colonial rule, which drew arbitrary borders between the new states the French left behind. Vietnamese enclaves became part of Cambodia, and Cambodian enclaves persisted in Vietnam. Both governments wanted not just to rectify this arbitrary line between them, but maximize their influence as well. Cambodia moved first and attacked Vietnamese positions in the island Phu Quoc in May 1975. Conflict simmered at a low level for the next two years before both sides dispatched regular military personnel to the conflict areas in lieu of local militias. A Cambodian attack in the Tay Ninh border province dated September 24, 1977, inaugurated war in earnest between the two states. The border issue served to start the war, but the general ambition for Vietnam became clear: Pol Pot had to go. Pol Pot was aware of his precarious situation as well. He needed Chinese support—Vietnam was pro-Soviet—because he began facing increasing opposition domestically as the conflict progressed. Vietnam struggled through 1978 but had a breakthrough by year's end. A new offensive launched on December 26, 1978, led to the successful occupation of Phnom Penh on January 7, 1979. Pol Pot and his confidantes fled the capital. Vietnam installed Heng Samrin as head of state, ending the conflict.
\\
Abuza, Zachary. 1995. “The Khmer Rouge and the Crisis of Vietnamese Settlers in Cambodia.” Contemporary Southeast Asia. & The first explanation is perhaps that the KR leadership truly believes its own propaganda. Vietnam, to them, is expansionist and seeks to annex Cambodia. Vietnam invaded and occupied the country once, and will attempt to do it again. This deep-seated hostility towards, and fear of, the Vietnamese is grounded in history. [...] The second reason may be that regimes in trouble often use foreign crises in order to bolster their domestic position. [...] The third explanation is that the issue of Vietnamese settlers is a rallying point for the Khmer population. [...] The issue of the Vietnamese settlers gives the Khmer Rouge a raison d'etre and increases their legitimacy [...] The fifth explanation is that since the election, targeting the Vietnamese has been a means of splitting the RGC. Clearly, the Khmer Rouge did not believe that the PRK would support FUNCINPEC's stance on the Vietnamese settlers.
\\
Chandler, David. 2018. A History of Cambodia. 4th ed. Routledge. & The third phase of the DK era, between the political crisis of September 1976 and a speech by Pol Pot twelve months later in which he announced the existence of the CPK, was marked by waves of purges and by a shift toward blaming Cambodia's difficulties and counterrevolutionary activity to an increasing extent on Vietnam. Open conflict with Vietnam had been a possibility ever since April 1975, when Cambodian forces had attacked several Vietnamese-held islands in the Gulf of Thailand with the hope of making territorial gains in the final stages of the Vietnam War. The Cambodian forces had been driven back and differences between the two Communist regimes had more or less been papered over, but DK's distrust of Vietnamese territorial intentions was very deep. So were Pol Pot's suspicions of the Vietnamese Communist Party, whose leaders had been patronizing toward their Cambodian counterparts for many years and had allowed Cambodia's armed struggle to flourish only in the shadow of Vietnam's. Pol Pot's suspicions deepened in July 1977 when Vietnam signed a treaty of cooperation with Laos, a move that he interpreted as part of Vietnam's plan to encircle Cambodia and to reconstitute and control what had once been French Indochina. Realizing the relative strengths of the two countries, however, Pol Pot tried at first to maintain correct relations and was unwilling to expand DK's armed forces to defend eastern Cambodia against possible Vietnamese incursions. The Vietnamese, recovering from almost thirty years of fighting, were also cautious. (...). Cambodians claimed parts of the Gulf of Thailand, where they hoped to profit from partially explored but unexploited offshore oil deposits, but these claims were rejected by the Vietnamese, who had similar hopes. (...). As so often in modern Cambodian history, what Cambodians interpreted as an internal affair or a quarrel between neighbors had unforeseen international repercussions. For several months after the death of Mao Zedong and the arrest of his radical subordinates known as the Gang of Four, the Chinese regime was in disarray. Although the four radicals were soon arrested, the new ruler, Hua Guofeng, tried to maintain Mao's momentum by opposing the Soviet Union, praising Mao's ideas and supporting third world revolutions like DK's. Many Chinese officials, including Hua's successor, Deng Xiaoping, perceived Vietnam as a pro-Soviet threat along their southern border—much as the United States at the time saw Cuba. (...). Vietnam saw the DK-Beijing alliance that was strengthened during Pol Pot's visit to China as a provocation, and in mid-December 1977 Vietnam mounted a military offensive against Cambodia. Fourteen divisions were involved, and Vietnamese troops penetrated up to thirty-two kilometers (twenty miles) into Cambodia in some areas. In the first week of 1978, after DK had broken off diplomatic relations with Vietnam, most of the Vietnamese troops went home, taking along thousands of Cambodian villagers as hostages. The Vietnamese soon began grooming some of these hostages as a government in exile; others were given military training. One of the exiles, a DK regimental commander named Hun Sen, who had fled Cambodia in 1977, emerged as the premier of the People's Republic of Kampuchea (PRK) in 1985. (...). On Christmas Day 1978, Vietnamese forces numbering over one hundred thousand attacked DK on several fronts. Because DK forces were crowded into the eastern and southwestern zones, Vietnamese attacks in the northeast encountered little resistance, and by the end of the year several major roads to Phnom Penh were in Vietnamese hands. At this point the Vietnamese altered their strategy, which had been to occupy the eastern half of the country, and decided to capture the capital itself. (pp. 269-276)
\\
Parker, Geoffrey, ed. 2020. The Cambridge History of Warfare. 2nd ed. Cambridge: Cambridge University Press. & Buoyed by their success in reuniting Vietnam under their rule in 1975, the Vietnamese communists involved themselves in both Cambodia and Laos - in effect, taking over the lands of peoples they regarded as inferior. Major Chinese attacks in the North, however, underlined to the Vietnamese that their megalomania would not remain unchallenged.
\\
\end{longtable}}
\clearpage
         \textbf{Ugandian-Tanzanian} 

          \underline{Onset:} 1978         

\underline{Reasons:}
\begin{itemize}
\item Nationalism
\item Border Clashes
\item Revenge/Retribution
\end{itemize}

         \underline{Sources:} 
         \rowcolors{3}{white}{light-gray} 
         {\scriptsize 
         \begin{longtable}{p{0.2\textwidth}p{0.8\textwidth}} 
         \hiderowcolors% 
         \toprule 
         Source & Excerpt \\ 
         \midrule 
         \showrowcolors% 
         \endfirsthead% 
         \hiderowcolors% 
         \toprule 
         Source & Excerpt \\ 
         \midrule 
         \showrowcolors% 
         \endhead% 
         \midrule 
         \endfoot% 
         \hiderowcolors% 
         \bottomrule 
         \showrowcolors% 
         \endlastfoot% 
         
Gibler, Douglas M. 2018. International Conflicts, 1816–2010: Militarized Interstate Dispute Narratives. Vol. 1 Rowman \& Littlefield. & In January 1971, Uganda's President Obote was ousted in a military coup and replaced by General Idi Amin. Obote sought refuge in Tanzania, where the country's president offered asylum and support. Soon after the military coup, though, Amin asserted that Tanzania was preparing for an attack to reinstate Obote and instructed Uganda's military forces to begin preparations for the attack. Tanzania's president denied these allegations but also warned citizens in his country to prepare to defend themselves. Meanwhile, the Tanzanian government refused to acknowledge General Amin as head of Uganda and continued to support President Obote as the legitimate leader. Military clashes began and escalated in August 1971, although both militaries were mobilized earlier and Amin began making threats of force as early as January 1971. Thanks to mediation by Kenya both two countries reached an agreement to keep economic ties through the East African Community. While this agreement appears to have ended the fighting, it did not resolve the larger issue of Tanzania's acceptance of Amin as Uganda's legitimate head of state. As of December 1971, Tanzania still did not recognize Amin's government. (...). 
The Ugandan-Tanzanian War of 1978–1979 ousted Idi Amin from his position as head of the Ugandan state. (...). Tanzania did not trust Amin and undertook several measures through the 1970s to influence his removal as Ugandan dictator. Meanwhile, Uganda had ambitions to annex the Kagera area in northwest Tanzania. Uganda made the first move, initiating the war with an October 1978 invasion of Tanzania. Libyan head of state Muammar al-Qaddafi came to the assistance of Idi Amin in this war, sending Libyan troops to fight alongside Ugandan troops in Tanzania. They were not successful. Tanzania, with assistance from Ugandan rebels who had become enemies of Amin through the course of the 1970s, captured the Ugandan capital of Kampala. Amin fled, ultimately taking sanctuary in Libya. Tanzania installed the former Ugandan head of state, Milton Obote, to his position. (p. 454-457)
\\
Thomas, Charles. 2022. “Uganda–Tanzania War.” Oxford Research Encyclopedia of African History. Accessed August 20, 2023. https://oxfordre.com/africanhistory/display/10.1093/acrefore/9780190277734.001.0001/acrefore-9780190277734-e-1040. & The Uganda–Tanzania War was a military conflict between Idi Amin's regime in Uganda and Julius Nyerere's Tanzania. The roots of the conflict can be traced back to Amin's seizure of power in Uganda in a coup against Nyerere's ally Milton Obote in 1971. Their mutual animosity was then cemented by Nyerere's approval of an invasion of Uganda by armed exiles the following year. The two countries continued to have strained relationships, leading to border disputes and faltering regional relations. Finally, in October 1978, Amin's military invaded Tanzania and declared the annexation of all territory north of the Kagera River.
\\
Roberts, George. 2014. “The Uganda–Tanzania War, the Fall of Idi Amin, and the Failure of African Diplomacy, 1978–1979.” Journal of Eastern African Studies, 8(4): 692–709. & Distinct among contemporaneous African conflicts for its noticeable lack of a Cold War context, the war demonstrated the shortcomings of the OAU in resolving African conflicts. (...). The Uganda–Tanzania War of 1978–1979 was a landmark event in post-colonial East African history. In response to Idi Amin's annexation of the Kagera Salient in northwestern Tanzania in November 1978, Julius Nyerere launched a controversial counterattack that routed Amin's forces and swept him from power in April 1979. Rooted in a deep rivalry between Amin and Nyerere, the conflict provoked bitter exchanges at the Organisation of African Unity (OAU), contributed to the failure of ujamaa in Tanzania, and brought an end to eight years of brutal dictatorship in Uganda. (...). The roots of the war lay in the coup d'état that brought General Amin to power in January 1971. Nyerere had enjoyed close relations with deposed President Milton Obote, having backed his socialist ‘Move to the Left' policies. Nyerere refused to recognise the new regime in Kampala and offered Obote and many of his supporters exile in Tanzania. Thus began a bitter rivalry between the two presidents. (...). By 1978, political tensions in Uganda were running high. Amin had become increasingly uncompromising towards any sign of dissent. In August 1976, up to 100 student protestors were shot dead at Makerere University. In February 1977, Amin caused outrage around the world when the Archbishop Luwum was murdered (the regime disguising it as a ‘car accident'), following allegations that he was plotting with Obote, days after the Archbishop had openly criticised the regime.
\\
Mambo, Andrew, and Julian Schofield. 2007. “Military Diversion in the 1978 Uganda–Tanzania War.” Journal of Political \& Military Sociology, 35(2): 299–321. & In May of 1978, Uganda's President and Field Marshal Idi Amin Dada protested on Radio Uganda of Tanzanian military incursions to the Ugandan towns of Kikagarti and Mutukula. This was not the first time. Tanzania had been accused of conducting raids into Uganda in July of 1974, September 1975, and of troop build-ups in preparation for war in March 1973, August 1975 and February-May 1977. On October 12, 1978, Radio Uganda announced fierce resistance against a Tanzanian invasion fifteen kilometers inside Uganda. On October 26, it was announced that Tanzania was repelled, and on October 31 Uganda counter-attacked, seizing control of the eighteen-hundred square kilometer Kagera salient in Northern Tanzania, from where the invasion jumped-off (Southall 1980:634). On November 1, Idi Amin declared the annexation of Kagera and offered to negotiate a peace with Tanzania. During the battle, there was not the least bit of evidence of a rally-around-the-flag effect in Uganda. Had Amin misread his people? The reality was that no Tanzanian attack took place. There was just one under-strength company of Tanzanian soldiers nearly thirty kilometers from the Ugandan border. In fact, Uganda was involved in a violent military purge in mid- 1978, which led to a series of small Ugandan incursions as early as May of 1978. The most probable link is that the Amin loyalist factions fell beyond his control during the purges. Then, on October 9, a substantial force of Ugandan soldiers crossed into Kagera, followed on October 22 by 3,000 troops, who, besides their occupation, engaged in widespread looting (Avirgan and Honey 1982:51,54, 57-61). This was the fatal mistake of the Amin regime. Tanzania deployed 7,000 soldiers opposite Kagera, prompting Amin to withdraw his army on November 11. In February 1979, Tanzanian President Julius Nyerere invaded Uganda assisted by a small army of exiles. By May, the Amin regime had been destroyed (Clodfelter 2002:627).
\\
\end{longtable}}
\clearpage
         \textbf{Sino-Vietnamese Punitive} 

          \underline{Onset:} 1979         

\underline{Reasons:}
\begin{itemize}
\item Power Transition or Security Dilemma
\item Border Clashes
\item Domestic Politics
\item Revenge/Retribution
\end{itemize}

         \underline{Sources:} 
         \rowcolors{3}{white}{light-gray} 
         {\scriptsize 
         \begin{longtable}{p{0.2\textwidth}p{0.8\textwidth}} 
         \hiderowcolors% 
         \toprule 
         Source & Excerpt \\ 
         \midrule 
         \showrowcolors% 
         \endfirsthead% 
         \hiderowcolors% 
         \toprule 
         Source & Excerpt \\ 
         \midrule 
         \showrowcolors% 
         \endhead% 
         \midrule 
         \endfoot% 
         \hiderowcolors% 
         \bottomrule 
         \showrowcolors% 
         \endlastfoot% 
         
Gibler, Douglas M. 2018. International Conflicts, 1816–2010: Militarized Interstate Dispute Narratives. Vol. 2 Rowman \& Littlefield. & Sino-Vietnamese relations had worsened following the Vietnam War for several reasons. China perceived the Soviet Union's friendship and support for the Vietnamese as an attempt to encircle China. China was also concerned with the Vietnamese government's treatment of ethnic Chinese in Vietnam, especially after almost 170,000 ethnic Chinese migrated to China in the spring of 1978. Land and maritime border claims also remained unresolved. On June 29, 1978, Vietnam joined COMECON, the Soviet bloc's economic community, but China was alarmed when Vietnam invaded Cambodia and quickly overthrew the Chinese-supported Pol Pot regime in December 1978. On February 17, 1979, Chinese forces crossed into Vietnam and continued fighting until March 5, when China announced that Vietnam had been sufficiently “punished” and that Chinese forces would soon withdraw. Vietnam stopped military operations against Chinese troops and agreed to talks following the Chinese withdrawal. The two sides held talks throughout the rest of the year but made no progress. Normalization would not be achieved until November 1991. (p. 827)
\\
Britannica. 2023. “20th Century International Relations - American Uncertainty.” Accessed August 20, 2023. https://www.britannica.com/topic/20th-century-international-relations-2085155/American-uncertainty\#ref305042. & After their 1975 victory the North Vietnamese showed a natural strategic preference for the distant U.S.S.R. and fell out with their historic enemy, neighbouring China. In quick succession Vietnam expelled Chinese merchants, opened Cam Ranh Bay to the Soviet navy, and signed a treaty of friendship with Moscow. Vietnamese troops had also invaded Cambodia to oust the pro-Peking Khmer Rouge. Soon after Deng Xiaoping's celebrated visit to the United States, Peking announced its intention to punish the Vietnamese, and, in February 1979, its forces invaded Vietnam in strength.
\\
Gompert, David C., Hans Binnendijk, and Bonny Lin. 2014. “China's Punitive War Against Vietnam, 1979.” In Blinders, Blunders, and Wars: What America and China Can Learn, 117–128. RAND Corporation. & China's Vietnam War was a mitigated strategic blunder. The Chinese leader Deng Xiaoping made the decision to attack Vietnam under complex domestic and international circumstances. Though Deng was angry at Vietnam, underestimated Vietnam's military strength, and overestimated the Chinese People's Liberation Army's (PLA's) capabilities, he did not completely discount the potential for failure. He had planned for the best-case scenario of quick victory, but he also considered the possibility of a less-than-successful military campaign. Most important, Deng correctly assessed that a limited war could bring China neither great victory nor greater loss. When operations on the battlefront were less than desirable, Deng kept his original plan of a short war and did not let China become bogged down in Vietnam. Despite Beijing's dismal military performance in early 1979, Deng was able to use the war to consolidate his domestic control. The war also undermined Vietnamese illusions of strong Soviet support and made Vietnam question whether cooperating with Moscow to encircle China was in its best interests. Over the long term, the war caused Vietnam to overextend itself militarily and economically. China's attack on Vietnam further bolstered Beijing's image in Southeast Asia by standing up against an aggressive Vietnam that was threatening its Southeast Asian neighbors. (pp. 117-118)
\\
Yu, Miles Maochun. 2022. “The 1979 Sino-Vietnamese War and Its Consequences.” Military History in the News, Hoover Institution, December 20. & In a grand struggle with the Soviet Union for the leadership role of the global communist movement, the Chinese Communist Party (CCP) waged a full-scale war of aggression against communist Vietnam in February and March 1979. Vietnam had abandoned Beijing and joined Moscow as a mutual defense-treaty ally and invaded and toppled China's Maoist puppet government, the Khmer Rouge in Cambodia. The war was also triggered by an internal CCP power struggle: Deng Xiaoping wanted to consolidate his control over the People's Liberation Army (PLA) to finally force the CCP General Secretary Hua Guofeng, Mao's outmaneuvered chosen successor, to cede supreme power to him.
\\
Parker, Geoffrey, ed. 2020. The Cambridge History of Warfare. 2nd ed. Cambridge: Cambridge University Press. & Buoyed by their success in reuniting Vietnam under their rule in 1975, the Vietnamese communists involved themselves in both Cambodia and Laos - in effect, taking over the lands of peoples they regarded as inferior. Major Chinese attacks in the North, however, underlined to the Vietnamese that their megalomania would not remain unchallenged.
\\
\end{longtable}}
\clearpage
         \textbf{Iran-Iraq} 

          \underline{Onset:} 1980         

\underline{Reasons:}
\begin{itemize}
\item Power Transition or Security Dilemma
\item Religion or Ideology
\item Border Clashes
\item Economic, Long-Run
\end{itemize}

         \underline{Sources:} 
         \rowcolors{3}{white}{light-gray} 
         {\scriptsize 
         \begin{longtable}{p{0.2\textwidth}p{0.8\textwidth}} 
         \hiderowcolors% 
         \toprule 
         Source & Excerpt \\ 
         \midrule 
         \showrowcolors% 
         \endfirsthead% 
         \hiderowcolors% 
         \toprule 
         Source & Excerpt \\ 
         \midrule 
         \showrowcolors% 
         \endhead% 
         \midrule 
         \endfoot% 
         \hiderowcolors% 
         \bottomrule 
         \showrowcolors% 
         \endlastfoot% 
         
Gibler, Douglas M. 2018. International Conflicts, 1816–2010: Militarized Interstate Dispute Narratives. Vol. 2 Rowman \& Littlefield. & Iran and Iraq returned to Cold War relations after the fall of the Shah in Iran. While Iraq initially tried to make friendly diplomatic gestures toward the new regime, relations between the two states quickly deteriorated in the spring of 1979. Both governments accused the other of 500–600 territorial violations between February 1979 and September 1980, and a war of words erupted between both sides shortly after the establishment of the Islamic Republic in Iran. Iraq was providing substantial arms to Arab rebels inside Iran who were fighting the Revolutionary Guard. Border clashes intensified in early 1980. The war was fought over contested territory, specifically the Shatt al-Arab and Khuzestan bounded by the Khorramshahr-Ahvaz-Susangerd-Messian line. On September 22, 1980, Iraq invaded Iran, and Iran bombed Iraqi air bases the next day. The United Nations offered its first resolution calling for a ceasefire on September 28 (479). However, Iraqi forces gained momentum and pushed forward. Iraq fortified the Iranian city of Khorramshahr, but by June 9, 1982, Iran had seized control. Iraq subsequently implemented a unilateral ceasefire and withdrew from 5,500 square kilometers of Iranian territory. Iran wanted Saddam Hussein ousted, so it continued fighting, but their forces were bogged down trying to cut Basra off. In 1983 Iran began to use human waves to attack Iraq, and Iraq used mustard gas against Iranian troops. In 1984, both sides started targeting oil tankers in the Persian Gulf both to cut off supplies to the other state and also affect the other's economy. In February 1986 Iran captured al-Faw and held it for two years. Iran lost the territory it had gained by the summer. On July 20, 1988, Iran and Iraq accepted a ceasefire as demanded by UN Resolution 598. The ceasefire went into effect on August 20. More than 500,000 soldiers died in the conflict.
\\
Britannica. 2023. “Iraq War.” Accessed August 19, 2023. https://www.britannica.com/event/Iraq-War. & Long-standing territorial disagreements between Iran and Iraq were reignited and remained a source of tension throughout the 1970s. With Iraq's new Ba'ath regime facing instability at home, its de facto ruler, Saddam Hussein, conceded some of the country's claims in 1975 in exchange for the cessation of Iranian meddling in Iraq. The primary concession was the eastern bank of the Shaṭṭ Al-'Arab (Persian: Arvand Rūd), a river formed by the confluence of the Tigris and Euphrates rivers that had high strategic value as a principal waterway for the countries' maritime commerce. [...] The Iranian Revolution of 1978–79 brought the 1975 agreements into questionable standing. Border clashes began occurring from time to time while signs of Iranian interference in Iraq were pronounced. Leaders of Iran's Islamic republic indicated their desire to “export” the revolution's concept of Islamic governance guided by Shi'i clergy. Ruhollah Khomeini, the ideological architect of the revolution, had already found an audience for his ideas in Iraq when he was living there during his exile. In April 1980 the deputy prime minister of Iraq, Tariq Aziz, survived an assassination attempt by Iraqi Shi'is sympathetic to the Iranian Revolution. [...] Meanwhile, Iran's new government—then at loggerheads with the United States over the seizure of the U.S. embassy in Tehrān by Iranian militants—was in apparent isolation and disorder, and its regular armed forces were gutted and demoralized. Saddam, who had considered the 1975 agreements lopsided, saw in this an opportune moment to reassert territorial claims he had previously conceded, notably the Shaṭṭ Al-'Arab. He made additional demands as well, including some form of Arab self-determination in Khūzestān, a rich oil-producing border region in Iran that was inhabited largely by ethnic
\\
Johnson, David. 2011. “Iraq–Iran War (1980–1988).” In The Encyclopedia of War, edited by G. Martel. & This eight-year war began as a series of miscalculations and spiraled out of control as new tactics brought forth unexpected reactions (...). (...). The centrality of oil made the war a matter of serious interest to all industrialized nations (...). Wars do not usually begin over long-standing animosities alone. In the case of Iran and Iraq those animosities go back at least three hundred years and some would say even more.The rise of the secular Ba'ath Party in Iraq and somewhat later the emergence of the Persian Islamic revival under the Ayatollah Khomeini accentuated the long-standing animosities between the two nations. The secularist Ba'ath political program resembled the firm centralized control and strict internal security of the East Germans who served as advisors. (...).
The real reason was that a window of opportunity opened for Iraq and it stepped through it. The cal-culus had to have been approximately as follows: the Ayatollah Khomeini came to power on the strength of Islamic renewal and had been seeking the overthrow of the Ba'ath Party and Saddam Hussein as apostates. Khomeini had been attempting to incite the large Iraqi Shi'a population in southern Iraq to join the revolution and overthrow the apostate Hussein. Saddam seems to have judged the verbal vitriol from Tehran to be dangerous enough possibly to upset the Shi'a and thus create a serious threat to his position, the issue that ultimately mattered most.
\\
Parker, Geoffrey, ed. 2020. The Cambridge History of Warfare. 2nd ed. Cambridge: Cambridge University Press. & When Iran collapsed into seeming anarchy after religious ideologues seized power in 1979, the Iraqi leader invaded his neighbour in order to take advantage of the situation. p. 420
\\
\end{longtable}}
\clearpage
         \textbf{War over Lebanon} 

          \underline{Onset:} 1982         

\underline{Reasons:}
\begin{itemize}
\item Nationalism
\item Religion or Ideology
\item Border Clashes
\end{itemize}

         \underline{Sources:} 
         \rowcolors{3}{white}{light-gray} 
         {\scriptsize 
         \begin{longtable}{p{0.2\textwidth}p{0.8\textwidth}} 
         \hiderowcolors% 
         \toprule 
         Source & Excerpt \\ 
         \midrule 
         \showrowcolors% 
         \endfirsthead% 
         \hiderowcolors% 
         \toprule 
         Source & Excerpt \\ 
         \midrule 
         \showrowcolors% 
         \endhead% 
         \midrule 
         \endfoot% 
         \hiderowcolors% 
         \bottomrule 
         \showrowcolors% 
         \endlastfoot% 
         
Gibler, Douglas M. 2018. International Conflicts, 1816–2010: Militarized Interstate Dispute Narratives. Vol. 2 Rowman \& Littlefield. & This dispute describes what was colloquially known as the First Lebanon War. In late March 1982, after completing the last of the withdrawals from the Sinai, Israel turned its attention to Lebanon, which they believed could contain Soviet missiles placed there by Syria. On April 11, Israel moved troops to the Lebanese border but did not break the truce and cross. The conflict stayed in this state with a few smaller skirmishes until Israeli diplomat Shlomo Argov was paralyzed in an assassination attempt in London on June 3. While the group that organized the attempt was not a part of the Palestine Liberation Organization (PLO), the PLO was suspected, blamed, and Israel decided it was time to attack. On June 6, Israeli Defense Forces invaded southern Lebanon. The primary objective of this attack was to remove the PLO and its influence in southern Lebanon and remove any weapons staged by Syria in that territory. Not long after the fighting began, the Syrian military joined the effort on the Lebanese side but was of little help. Most of the heavy fighting was over by September 1982 with the PLO removing its troops from southern Lebanon. On May 17, 1983, with the assistance of the United States, Israel and Lebanon signed an accord officially ending the war and staging Israeli withdrawal from Lebanon. This withdrawal took over two years to complete with Israel only pulling back small numbers of troops at a time so as to leave a security force in Lebanon to help the Israeli-backed militia, the Southern Lebanon Army. On June 5, 1985, Israel completed its troop withdrawal to the security zone. (p. 695)
\\
Britannica. 2023. “Lebanese Civil War.” Accessed August 20, 2023. https://www.britannica.com/event/Lebanese-Civil-War. & The most significant Israeli intervention during the course of the Lebanese Civil War, however, was the invasion that began on June 6, 1982. Although the stated goal of Israel was only to secure the territory north of its border with Lebanon so as to stop PLO raids, Israeli forces quickly progressed as far as Beirut's suburbs and laid siege to the capital, particularly to West Beirut. The invasion resulted in the eventual removal of PLO militia from Lebanon under the supervision of a multinational peacekeeping force, the transfer of the PLO headquarters to Tunis, Tunisia, and the temporary withdrawal of Syrian forces back to al-Biqā'. Galvanized by the Israeli invasion, a number of Shi'i groups subsequently emerged, including Hezbollah.
\\
Bregman, Ahron. 2011. “Arab—Israeli Conflict.” In The Encyclopedia of War, edited by Gordon Martel. & From the day of his appointment, Sharon's attention was firmly focused on Lebanon, where he identified two main problems. The first was the presence of Syrian troops and their ground-to-air missile system in the Beka'a Valley, which hindered the IAF's freedom to fly over Lebanon; the second was the presence of the Palestinian Liberation Organization (PLO), led by Yasser Arafat, whom Sharon suspected of wanting to take over Lebanon and turn it into a base to attack Israel. Sharon wished to strike at both the PLO and the Syrians in Lebanon. (...) on June 3, 1982 (...) gunmen of a dissident Palestinian faction (...) shot the Israeli ambassador to London and seriously injured him. (...) [Such] was the mood in Israel following the attempt on the life of the ambassador that (...) Israel needed to attack the PLO. (...).
[The invasion of Lebanon] gave the IDF the mission of “freeing all the Galilee settlements from the range of fire of terrorists” and instructed that “the Syrian army [stationed in Lebanon] should not be attacked unless it attacks our forces.” (...) [The] operation's objective was to remove the PLO from firing range of Israel's northern border, “approxi-mately 45 kilometers” (Resolution 676 of the Israeli cabinet). (...). 
Sharon ordered [the IDF] to proceed up to Lebanon's capital Beirut in order to hunt down PLO leader Yasser Arafat. (...).
(...) [Some] of his aims: the IDF pushed the PLO and Syrian forces out of Beirut and the IAF destroyed the Syrian ground-to-air missiles in eastern Lebanon.
\\
Zisser, Eyal. 2011. “Lebanese Civil War (1975–1990).” In The Encyclopedia of War, edited by Gordon Martel. & From April 1975 to October 1990 Lebanon suffered a bloody civil war. At its peak it threatened to bring about the complete collapse and disintegration of the Lebanese state. During the fighting nearly 150,000 people were killed and another 200,000 wounded. Six hundred thousand people were made homeless, and nearly 250,000 left Lebanon, most of whom never returned. All this took place among a population that numbered merely 3.1 million at the beginning of the war. (...). It can be seen infact as a conglomeration of internal Lebanese conflicts, joined by regional confrontations (...). (...). Thus, during the civil war we find Maronites fighting Sunnis, Druze fighting Maronites, Druze fighting Shi'ites, and Shi'ites fighting Sunnis, as well as Maronites, Druze, and Shi'ites fighting other members of their own community. (...). Three foreign actors in particular played a role: Syria, Israel, and the PLO (...). (...) [The] demographic composition between Christians and Muslims (...) were totally different from those which had prevailed in the Ottoman district of Mount Lebanon. (...). Almost 120,000 Palestinian refugees had come to Lebanon in the wake of the 1948 [Arab-Israeli] war. (...) [The] PLO became an armed presence on Lebanese soil. (...) [This] (...) led to the border with Israel heating up (...). (...). The leaders of the Christian militias (...) linked up with Israel in an effort to expel the PLO, and even the Syrians, from Lebanon. (...) The struggle between Israel and Syria for dominance in Lebanon developed finally into all-out warfare when Israel launched Operation Peace for the Galilee on June 6, 1982.
\\
Fraser, T. G. 2025. “Lebanese Civil Wars (1958, 1975–90).” In The Encyclopedia of Diplomacy, edited by Gordon Martel. & There was poverty amongst the Shias, who had never really shared in the Maronite–Sunni duumvirate, and who were growing increasingly restive over their marginal status. They found a ready home in the Movement of the Disinherited, founded by the Shia Musa al-Sadr, born in Iran of Lebanese origin, who proved to be a dynamic spokesman. There were also the Palestinians who were in Lebanon but not of it. The 1967 war saw a new influx of Palestinians, while the following year the reconstructed Palestine Liberation Organization (PLO), with the Fatah leader Yasir Arafat (1929–2004) as chairman, was increasingly seen as operating independently of the structures of the Lebanese state. After King Hussein of Jordan (1935–99) acted against the Palestinian guerrillas in September 1970, Lebanon became the PLO's main base and area of operations. There was no lack of recruits from amongst the refugee community. In a portent of what was to come, its activities provoked Israeli retaliatory raids into Lebanese territory. The Palestinians attracted the particular attention of the two main rightist Christian parties, the National Liberal Party of Camille Chamoun, who had never renounced his role in politics, and the Kataeb or Phalange Party founded in 1936 by Pierre Gemayel (1905–84). Opposing them was a left-wing alliance, the National Movement, formed in 1969 by Kamal Jumblatt, which was broadly sympathetic to the PLO. A surface gloss concealed a state likely to combust, which in April 1975 it did.
\\
\end{longtable}}
\clearpage
         \textbf{Falkland Islands} 

          \underline{Onset:} 1982         

\underline{Reasons:}
\begin{itemize}
\item Nationalism
\item Border Clashes
\item Economic, Short-Run
\end{itemize}

         \underline{Sources:} 
         \rowcolors{3}{white}{light-gray} 
         {\scriptsize 
         \begin{longtable}{p{0.2\textwidth}p{0.8\textwidth}} 
         \hiderowcolors% 
         \toprule 
         Source & Excerpt \\ 
         \midrule 
         \showrowcolors% 
         \endfirsthead% 
         \hiderowcolors% 
         \toprule 
         Source & Excerpt \\ 
         \midrule 
         \showrowcolors% 
         \endhead% 
         \midrule 
         \endfoot% 
         \hiderowcolors% 
         \bottomrule 
         \showrowcolors% 
         \endlastfoot% 
         
Gibler, Douglas M. 2018. International Conflicts, 1816–2010: Militarized Interstate Dispute Narratives. Vol. 1. Rowman \& Littlefield. & This dispute describes the Falklands War or Guerra de las Islas Malvinas between the United Kingdom and Argentina over control of Falkland Islands and accompanying South Georgia and South Sandwich Islands just off the coast of Argentina. The islands were owned and occupied by the British since 1833, but popular sentiment in Argentina had long held those islands to be unredeemed Argentine territory that became British only after a period of Spanish, Argentine rule. The military junta of Leopoldo Galtieri in Argentina had underwhelming public support and tried to encourage greater legitimacy at home by seizing the islands from the British. (...). The dispute began when a group of Argentine marines posed as a group of civilians on board a naval transport en route to the Falkland Islands. When they arrived, they scouted the island in preparation for a potential invasion. While this was concealed from the British at the time, the Argentines had created a diplomatic incident by ignoring protocol for visits to the island and, importantly, raised the Argentine flag on the islands. They were asked to leave for a wanton violation of protocol—including slaughtering protected reindeer on the island and defacing British flags—and tensions between both states increased. Thinking that the British would try to send reinforcements to the islands, Argentina made the decision for an invasion on April 2. War followed, culminating in a British victory as codified in a June 14 surrender from Argentine general Mario Menendez to British general Jeremy Moore. The South Georgia, South Sandwich, and Falkland Islands were once more British. The military junta in Argentina, which had gambled all on this invasion, collapsed following the war. (p. 198)
\\
Britannica. 2023. “Falkland Islands War.” Encyclopædia Britannica. Accessed August 19, 2023. https://www.britannica.com/event/Falkland-Islands-War. & Brief but undeclared war between Argentina and Great Britain over control of the Falkland Islands (Islas Malvinas) and associated island dependencies. Both countries long had claimed sovereignty over the Falklands and had been in protracted negotiations over them. [...] The decision to invade was chiefly political: the junta, which was being criticized for economic mismanagement and human rights abuses, believed that the “recovery” of the islands would unite Argentines behind the government in a patriotic fervour.
\\
Turner, Blair. 2011. “Falklands War (1982).” In The Encyclopedia of War, edited by Gordon Martel. & (...). There had been long and drawn-out discussions in the United Nations (UN) and other venues for decades about the contentious issue of sovereign control of the islands. Multiple conflicting claims dating back to the Seven Years' War of the eighteenth century ensured that both sides had established rather rigid positions: Britain in favor of “self-determination” for the few settlers on the islands (Scots and Welsh) and the Argentines in favor of exclusive sovereignty over islands on the Argentine continental shelf, with guarantees of civil rights for the settlers as Argentine citizens. (...). The Argentine plan was based on forcing a fait accompli. That the British would undertake the costly and huge operation necessary to recover the islands some 8,500 miles distant in the coming South Atlantic winter was deemed too unlikely; a negotiated settlement would be reached. Argentine planners understood, however, that defeat was a distinct probability if the British did fight.
\\
Ashton, Nigel. 2025. “Falklands War (1982).” In The Encyclopedia of Diplomacy, edited by Gordon Martel. & The Falklands War of 1982 resulted from a sovereignty dispute between Britain and Argentina over the Falklands (Malvinas) Islands which dates back to the early nineteenth century. The immediate causes of the April 2, 1982, Argentinian invasion of the islands lay in the domestic difficulties experienced by the military regime in Buenos Aires and in its belief that Britain would not seek to reverse any invasion by force. The British government of prime minister Margaret Thatcher responded by dispatching a naval task force to the South Atlantic, which first recovered the Falklands dependency of South Georgia and then moved on to land forces on the Falklands themselves. By June 14, 1982, the Argentinian forces defending the islands had surrendered. In the aftermath of the conflict Britain refused to discuss the question of sovereignty with Argentina and proceeded to enhance its capabilities to defend the islands. (...). 
The military junta which ruled in Buenos Aires, led by General Leopoldo Galtieri, found itself under significant domestic pressure through a combination of its human rights abuses and economic mismanagement. A successful operation to seize the Malvinas promised to harness Argentine nationalism and unite the people behind the regime. Meanwhile in Britain the 1981 Defence Review had taken place against the backdrop of significant cuts in public expenditure caused by economic problems and included the decision to withdraw the ice patrol vessel HMS Endurance from the South Atlantic.
\\
Parker, Geoffrey, ed. 2020. The Cambridge History of Warfare. 2nd ed. Cambridge: Cambridge University Press. & On 2 April 1982, at the height of the Cold War, the Argentine government launched an amphibious expedition to seize the Falkland Islands. The Argentinians had claimed those cold, desolate islands in the South Atlantic since the nineteenth century, but few of its citizens had ever lived on the archipelago. A small group of British settlers had inhabited the islands since the early nineteenth century. Nevertheless, the Argentinians had created a myth as to the importance of those islands to their national identity. Given the pusillanimous diplomatic approach that the British Foreign Office had taken to negotiating with the Argentinians, the military junta that governed Argentina had little expectation that the British would use force to respond to their invasion. In fact, most Britons felt outrage, while the nation's prime minister, Margaret Thatcher, proved a formidable leader with the ruthlessness necessary to see the war through to its conclusion. Neither country was ready for the kind of war that resulted. The focus of British military preparations had been on supporting western Europe against Soviet invasion, and the outbreak of the war found the Ministry of Defense in the process of scaling down its presence in other areas. Nevertheless, the British still possessed residual amphibious resources, including two anti-submarine carriers able to support Harrier jump-jets with significant air-to-air capabilities. Their opponents, the Argentine military, had not engaged in serious military operations in the twentieth century, other than low-level counter-insurgency. (pp. 414-415)
\\
Morillo, Stephen, Jeremy Black, and Paul Lococo. 2008. War in World History: Society, Technology, and War from Ancient Times to the Present. New York: McGraw-Hill Professional. & Decolonization also ensured that earlier disputes that would have been limited by the strength and influence of the imperial powers, whether as participants or as arbiters, became more urgent. This led, for example, to war between Britain and Argentina over the Falkland Islands in 1982, with an Argentine invasion eventually totally defeated (...).
\\
\end{longtable}}
\clearpage
         \textbf{War over the Aouzou Strip} 

          \underline{Onset:} 1986         

\underline{Reasons:}
\begin{itemize}
\item Nationalism
\item Economic, Long-Run
\end{itemize}

         \underline{Sources:} 
         \rowcolors{3}{white}{light-gray} 
         {\scriptsize 
         \begin{longtable}{p{0.2\textwidth}p{0.8\textwidth}} 
         \hiderowcolors% 
         \toprule 
         Source & Excerpt \\ 
         \midrule 
         \showrowcolors% 
         \endfirsthead% 
         \hiderowcolors% 
         \toprule 
         Source & Excerpt \\ 
         \midrule 
         \showrowcolors% 
         \endhead% 
         \midrule 
         \endfoot% 
         \hiderowcolors% 
         \bottomrule 
         \showrowcolors% 
         \endlastfoot% 
         
Gibler, Douglas M. 2018. International Conflicts, 1816–2010: Militarized Interstate Dispute Narratives. Vol. 2 Rowman \& Littlefield. & This dispute describes what is fundamentally a Libya-Chad conflict over the Aouzou Strip that buffers their respective borders. In an opening session of the Organization of African Unity (OAU), Chad charged that Libya had capitalized on the civil war pitting the Chadian government against the rebel group Frolinat and had seized part of the Aouzou Strip. Chad further alleged that Libya was providing full support for the rebel group in order to undermine the Chadian government. Chad momentarily suspended relations with Libya on February 6, 1978, though this was followed by multiparty negotiations in Sabha that led to a ceasefire on March 28. Near the end of the negotiations, the French intervened in Chad under the auspices of a defense pact signed with the Chadian government. The ceasefire agreements were nullified in mid-April as a result and clashes ensued through the rest of 1978, pitting the Chadian government and French forces against the Chadian rebels and the Libyans. The conflict took on another dimension when a schism between President Malloum and Prime Minister Habre resulted in fighting between armed forces loyal to both leaders. On March 20, 1979, France announced that it had decided to incrementally withdraw its military personnel from Chad under these circumstances. While the French planned withdrawal and coordinate some type of settlement for all parties involved in Chad, Frolinat forces and the Libyan government continued its push farther south into Chad. Internal dissension in Chad continued throughout 1979. A reconciliation plan signed in Lagos, Nigeria, on August 21, led to a transitional government on November 10 that ultimately collapsed in March 1980. (...). 
This dispute describes continued hostilities between Libya and Chad (see disputes MID\#3631 and MID\#3624). This particular outbreak began on February 10, 1986, when Libyan forces attacked Chadian army troops in central Chad while fighting in support of former Chad President Oueddei. The Chad military mobilized in response. On February 15, French forces were put on alert over fears the incidents would reignite the civil war. Limited actions followed throughout the rest of the year. Then, on January 5, 1987, Libya bombed a French protected area of Chad. Fighting continued into the summer until Chad and Libya agreed to a ceasefire on September 11. (pp. 586-589)
\\
Britannica. 2011. “Aozou Strip.” Accessed August 20, 2023. https://www.britannica.com/place/Aozou-Strip. & The Aozou Strip has an area of about 44,000 square miles (114,000 square km) and consists almost entirely of the Sahara desert. The Tibesti Mountains interrupt the desert in the northwest, near the only town, Aozou. The population consists of scattered livestock herders and a few subsistence farmers, but interest in the strip intensified in the 1970s with the discovery that the area might be rich in uranium deposits. The Aozou Strip became the object of a fierce sovereignty dispute after Libya occupied the region in 1973 and unilaterally annexed it in 1975. Over the next 15 years, armed conflicts periodically erupted between Libya and Chad as each nation tried to assert its control over the strip. In 1988, however, the two countries agreed to settle the dispute peacefully, and in 1990 they submitted the dispute to the International Court of Justice (The Hague, Neth.).
\\
Naldi, Gino J. 2009. “The Aouzou Strip Dispute — A Legal Analysis.” Journal of African Law. & The outbreak of further hostilities between Chad and Libya in August 1987 was occasioned by a dispute concerning sovereignty over the so-called Aouzou Strip in northern Chad. The extent of Libyan involvement in Chad is motivated to a large degree by this territorial claim. This dispute must be distinguished from Libya's wider ambitions for Arab unity or its involvement in Chad's civil war, although it would appear true to say that Libya was thereby seeking to consolidate its hold on the Aouzou Strip.
\\
\end{longtable}}
\clearpage
         \textbf{Sino-Vietnamese Border War} 

          \underline{Onset:} 1987         

\underline{Reasons:}
\begin{itemize}
\item Power Transition or Security Dilemma
\item Domestic Politics
\end{itemize}

         \underline{Sources:} 
         \rowcolors{3}{white}{light-gray} 
         {\scriptsize 
         \begin{longtable}{p{0.2\textwidth}p{0.8\textwidth}} 
         \hiderowcolors% 
         \toprule 
         Source & Excerpt \\ 
         \midrule 
         \showrowcolors% 
         \endfirsthead% 
         \hiderowcolors% 
         \toprule 
         Source & Excerpt \\ 
         \midrule 
         \showrowcolors% 
         \endhead% 
         \midrule 
         \endfoot% 
         \hiderowcolors% 
         \bottomrule 
         \showrowcolors% 
         \endlastfoot% 
         
Gibler, Douglas M. 2018. International Conflicts, 1816–2010: Militarized Interstate Dispute Narratives. Vol. 2 Rowman \& Littlefield. & China and Vietnam had a brief conflict beginning in October 1986, again over the role Vietnam had in Cambodian affairs. This episode's fighting centered on Ha Tuyen province in Vietnam and the Yunnan province in China. China fired approximately 35,000 shells into Vietnam and conducted three separate division-sized raids into Vietnam. Fighting intensified further in January 1987 as Taiwan launched 15 separate division-sized raids. Losses were high for both, with Vietnam claiming more than 1,500 Chinese killed, while China provided numbers about one-third that size. The Vietnamese proposed talks in January, but China said that it refused to negotiate while Vietnamese forces continued to occupy Cambodia. The conflict continued the next year. (pp. 829-830)
\\
Britannica. 2023. “20th Century International Relations - American Uncertainty.” Accessed August 20, 2023. https://www.britannica.com/topic/20th-century-international-relations-2085155/American-uncertainty\#ref305042. & After their 1975 victory the North Vietnamese showed a natural strategic preference for the distant U.S.S.R. and fell out with their historic enemy, neighbouring China. In quick succession Vietnam expelled Chinese merchants, opened Cam Ranh Bay to the Soviet navy, and signed a treaty of friendship with Moscow. Vietnamese troops had also invaded Cambodia to oust the pro-Peking Khmer Rouge. Soon after Deng Xiaoping's celebrated visit to the United States, Peking announced its intention to punish the Vietnamese, and, in February 1979, its forces invaded Vietnam in strength.
\\
Yu, Miles M. 2022. “The 1979 Sino-Vietnamese War and Its Consequences.” Hoover Institution. & The war had a profound impact. China failed to achieve its strategic objective of ousting Vietnam from Cambodia, as Vietnam's occupation would continue until 1989, but with the United States' inexplicable assistance, China successfully prevented Vietnam's mutual defense- treaty ally the Soviet Union from an invasion from the North. The war deepened Vietnam's hostility toward China, and the two communist countries would be engaged in a series of intermittent, brutal small wars for the next 12 years until peaceful negotiation finally took place in 1991, when the Soviet Union under Mikhail Gorbachev was collapsing. While the war facilitated Deng Xiaoping's rise to supreme power inside the CCP, the PLA's abysmal performance during the conflict stimulated a Chinese miliary modernization drive that has made the PLA a formidable modern miliary force today, threatening the United States, the very country that lent a helping hand to save the CCP in 1979 from a two-front, unwinnable internecine communist war in an international power struggle for communist ideological correctness.
\\
\end{longtable}}
\clearpage
         \textbf{Gulf War} 

          \underline{Onset:} 1990         

\underline{Reasons:}
\begin{itemize}
\item Power Transition or Security Dilemma
\item Border Clashes
\item Economic, Long-Run
\end{itemize}

         \underline{Sources:} 
         \rowcolors{3}{white}{light-gray} 
         {\scriptsize 
         \begin{longtable}{p{0.2\textwidth}p{0.8\textwidth}} 
         \hiderowcolors% 
         \toprule 
         Source & Excerpt \\ 
         \midrule 
         \showrowcolors% 
         \endfirsthead% 
         \hiderowcolors% 
         \toprule 
         Source & Excerpt \\ 
         \midrule 
         \showrowcolors% 
         \endhead% 
         \midrule 
         \endfoot% 
         \hiderowcolors% 
         \bottomrule 
         \showrowcolors% 
         \endlastfoot% 
         
Gibler, Douglas M. 2018. International Conflicts, 1816–2010: Militarized Interstate Dispute Narratives. Vol. 2 Rowman \& Littlefield. & Iraq invaded and occupied Kuwait, and an allied coalition largely spearheaded by the United States was formed to expel Saddam Hussein's Iraqi troops from Kuwait. This began the Gulf War. There are two important sources for this dispute. First, Kuwait, and Saudi Arabia took an active role in financing Saddam Hussein's war with neighboring Iran through the 1980s. When Kuwait and Saudi Arabia began pressing Hussein for repayment of the loans, Hussein responded with belligerence and threats to Kuwaiti and Saudi sovereignty. Kuwait was particularly vulnerable. Disputes over oil drilling mounted between both sides and Iraq increasingly felt that Kuwait was independent from Iraq only as an artifact of British imperialism. Hussein's Iraqi forces invaded on August 2, 1990. This forced the Saudi government to plead with the United States for intervention. The United States responded with Operation Desert Shield to deter further Iraqi aggression. Shortly thereafter, the United States decided to expel the Iraqi military from Kuwait. The United States put together a far-reaching coalition of states to assist in the effort. Hussein was given an ultimatum via a UN resolution that authorized the use of force against Iraq should Iraqi troops still occupy after January 15, 1991. Hussein refused and war followed. The efforts of the coalitions were more than effective, resulting in a ceasefire on February 28 and an Iraqi surrender on March 3. UN Resolution 687, passed on April 3, formalized an end to the conflict. (pp. 648-649)
\\
Britannica. 2023. “Persian Gulf War.” Encyclopædia Britannica. Accessed August 19, 2023. https://www.britannica.com/event/Persian-Gulf-War. & Iraq's leader, Saddam Hussein, ordered the invasion and occupation of Kuwait with the apparent aim of acquiring that nation's large oil reserves, canceling a large debt Iraq owed Kuwait, and expanding Iraqi power in the region.
\\
Estes, Kenneth W. 2011. “Gulf Wars (1990–1991, 2003–present).” In The Encyclopedia of War, edited by G. Martel. & Known in the region as the Second and Third Gulf Wars, with the Iraq–Iran War (1980–1988) as the first of the series, these conflicts marked the apogee and nadir of the regime of the Iraqi dictator Saddam Hussein. In the case of the subject wars, each consisted of a coalition war organized by the United States to first limit and then overthrow the Iraqi regime, with the aim of improving peace and stability in the petroleum-rich region so vital to the economies of most of the rest of the world. (...) [The] Gulf War of 1990-1991 (...) decisively reversed the issue of Iraq's seizure of Kuwait on August 2, 1991. (...). The 1990–1991 Gulf War drew some origins from the Iraq–Iran War (...). In July 1988 both sides ceased hostilities. Less than two years later, Iraq then seized Kuwait, partly out of the need to recover financially from the ravages of the First Gulf War. Kuwait had loaned Saddam Hussein's government fourteen billion dollars in the interval and was loathe to remit any of it, despite Iraqi requests. Given incipient quarrels with Kuwait over borders and slant drilling into the Rumelia oil fields, Hussein determined to settle the issue by outright conquest.
\\
Johnson, David. 2011. “Iraq–Iran War (1980–1988).” In The Encyclopedia of War, edited by G. Martel. & By [the Iran-Iraq War's] end, both sides were exhausted, but Iraq had access to reserves Iran did not possess and could have continued the war for a while longer; however, its debt settlement issues led to the invasion of Kuwait in 1989 and the coalition counterstroke, Operation Desert Storm.
\\
Parker, Geoffrey, ed. 2020. The Cambridge History of Warfare. 2nd ed. Cambridge: Cambridge University Press. & To Saddam Hussein, 'victory' over Iran opened up the possibility of economic and political control of the Middle East. As leader of the Sunni Ba'ath party, he aimed to redress the ancient wrongs infliced on the Arab and Islamic worlds by Western interlopers over the previous five centuries. Burdened with debt from the war with Iran and believing it inconceivable that the Americans would use their power, Saddam Hussein struck the small but oil-rich adjoining state of Kuwait in summer 1990. (p. 420)
\\
\end{longtable}}
\clearpage
         \textbf{Bosnian Independence} 

          \underline{Onset:} 1992         

\underline{Reasons:}
\begin{itemize}
\item Nationalism
\item Economic, Long-Run
\item Domestic Politics
\end{itemize}

         \underline{Sources:} 
         \rowcolors{3}{white}{light-gray} 
         {\scriptsize 
         \begin{longtable}{p{0.2\textwidth}p{0.8\textwidth}} 
         \hiderowcolors% 
         \toprule 
         Source & Excerpt \\ 
         \midrule 
         \showrowcolors% 
         \endfirsthead% 
         \hiderowcolors% 
         \toprule 
         Source & Excerpt \\ 
         \midrule 
         \showrowcolors% 
         \endhead% 
         \midrule 
         \endfoot% 
         \hiderowcolors% 
         \bottomrule 
         \showrowcolors% 
         \endlastfoot% 
         
Gibler, Douglas M. 2018. International Conflicts, 1816–2010: Militarized Interstate Dispute Narratives. Vol. 1. Rowman \& Littlefield. & Croatia and Bosnia enjoyed a pragmatic alliance against Serbian forces. That changed as Croatia made an agreement with Serbia to partition Bosnia and fighting broke out on October 20, 1992. The fighting lasted a little over one week and ended with an immediate ceasefire on October 27. Much of Bosnia had been partitioned. (...).
In 1992 Bosnia held a referendum on independence. The Bosnian Serbs opposed independence and boycotted the election, so the voters who did turn out voted overwhelmingly for independence. Fighting broke out as soon as the election results were announced, and the Bosnian government moved toward independence. Over the protest of the Bosnian Serbs, several countries recognized Bosnia in April, and the United Nations admitted Bosnia in May. The Bosnian Serbs soon declared their own independent state, although it did not receive recognition from the international community. Croatia also had interests in the fate of Bosnia. The Croats wanted Bosnian territory and sometimes worked with the Serbs. However, the Croats also wanted to weaken the Bosnian Serbs, including removing a substantial portion of the Bosnian Serb lands. (...).
Bosnia-Herzegovina declared independence on March 3, 1992. Serb nationalists then began their siege of Sarajevo on April 6. As Bosnia began receiving international recognition through April and May 1992, the fighting in Bosnia intensified. Serbia and Montenegro announce a truncated Yugoslavia on April 27 and said it had no territorial claims to Bosnia, though that proclamation was made under American pressure. Bosnia responded by formally demanding that Yugoslavia withdraw its troops from the new state. Serbian nationalists (backed by the new Yugoslavia) continued their shelling and bombing of cities in Bosnia, with particularly heavy fighting in Sarajevo. Fighting intensified in October, November, and December 1992. While Bosnian Muslims and Croats had been allied with Croatia against the Serbs, this alliance began to falter in October 1992 (pp. 331-336)
\\
Britannica. 2023. “Bosnian War.” Encyclopædia Britannica. Accessed August 19, 2023. https://www.britannica.com/event/Bosnian-War. & Bosnian War, ethnically rooted war (1992–95) in Bosnia and Herzegovina, a former republic of Yugoslavia with a multiethnic population comprising Bosniaks (Bosnian Muslims), Serbs, and Croats. [...] In 1991 several self-styled “Serb Autonomous Regions” were declared in areas of Bosnia and Herzegovina with large Serb populations. Evidence emerged that the Yugoslav People's Army was being used to send secret arms deliveries to the Bosnian Serbs from Belgrade (Serbia). In August the Serb Democratic Party began boycotting the Bosnian presidency meetings, and in October it removed its deputies from the Bosnian assembly and set up a “Serb National Assembly” in Banja Luka. By then full-scale war had broken out in Croatia, and the breakup of Yugoslavia was under way. Bosnia and Herzegovina's position became highly vulnerable. The possibility of partitioning Bosnia and Herzegovina had been discussed during talks between the Croatian president, Franjo Tudjman, and the Serbian president, Slobodan Milošević, earlier in the year, and two Croat “communities” in northern and southwestern Bosnia and Herzegovina, similar in some ways to the “Serb Autonomous Regions,” were proclaimed in November 1991.
\\
Ashbrook, John E. 2011. “Yugoslav Succession, Wars of (1990–1999).” In The Encyclopedia of War, edited by Gordon Martel. & The Wars of Yugoslav Succession were four major conflicts in Slovenia, Croatia, Bosnia-Hercegovina, and Serbia-Kosovo. Of these, the largest occurred in Croatia and Bosnia-Hercegovina in the early 1990s. The wars and dissolution resulted from a number of unresolved historical issues; the rise of opportunistic, nationalist politicians; a cumbersome federal government; and a long-term economic crisis. The wars themselves occurred in no small part from long-standing tensions between ethnic groups in Yugoslavia.(...). During the 1960s the country's economy, seemingly strong, reduced ethnic friction, but with the economic downturn in the 1970s and Tito's death in 1980, tensions reemerged without the leader's charm or his ability to assuage them. A number of politicians and academics began to revive historical national grievances and to disparage the multi-national experiment in communism for political and economic reasons, eventually throwing the country into the wars of dissolution. (...). The Yugoslav People's Army (JNA) supplied Serbs with weapons to resist secessionist drives in Croatia and Bosnia. (...). Unlike the other republics, there was no majority ethnic group in Bosnia-Hercegovina (BiH). In 1991 Bosnian Muslims (Bosˇnjaks) constituted 43 percent of the population, Serbs 31 percent, and Croats 17 percent. In general, most Bosnians were ethnically tolerant, especially in urban areas. However, three new political parties, each national in character, emerged and increasingly agitated against parties seeking to keep the country united. After the first elections the Serbian and Croatian parties radicalized, due in part to outside pressures from Serbia and Croatia. For Milosevic and Tudman a functioning, multi-ethnic BiH was anathema to their nationalist projects. Both also desired an unstable Bosnia to focus internal discontent on an external crisis.
\\
Parker, Geoffrey, ed. 2020. The Cambridge History of Warfare. 2nd ed. Cambridge: Cambridge University Press. & Shortly after the collapse of the USSR, deep ethnic divisions resurfaced in Yugoslavia, a federation of six republics (Serbia, Croatia, Macedonia, Bosnia-Herzegovina, Slovenia, and Montenegro) held together by an authoritarian communist regime in Belgrade, the capital of Serbia as well as of the federation. Internal divisions appeared in the 1980s and, burdened by international debts, inflation,  and rising unemployment, the economy declined. In 1990 multi-party elections in the constituent republics brought nationalists to the fore in Slovenia and Croatia, paving the way for their simultaneous declaration of independence in June 1991. The central government of Yugoslavia, dominated by the Serb leader Slobodan Milosevic, attacked until in January 1992, after much bloodshed and destruction, United Nations troops and American diplomats forced him to recognize the secession of the two states. In Bosnia-Herzegovina, the most ethnically diverse of the former Yugoslav republics, Muslims and Croats pushed for independene as the Bosnian Serbs sought to remain in a union with the republic of Serbia. From April 1992 each of the three ethnic groups sought to 'cleanse' the areas under its control of all opponents and the Serbs, with support from Belgrade, conducted a particularly savage campaign in which non-Serb men were routinely killed and non-Serb women systematically raped. Within a yera Serb forces controlled around 70 per cent of Bosnia.  (pp. 425-427)
\\
\end{longtable}}
\clearpage
         \textbf{Azeri-Armenian} 

          \underline{Onset:} 1993         

\underline{Reasons:}
\begin{itemize}
\item Nationalism
\item Religion or Ideology
\item Border Clashes
\end{itemize}

         \underline{Sources:} 
         \rowcolors{3}{white}{light-gray} 
         {\scriptsize 
         \begin{longtable}{p{0.2\textwidth}p{0.8\textwidth}} 
         \hiderowcolors% 
         \toprule 
         Source & Excerpt \\ 
         \midrule 
         \showrowcolors% 
         \endfirsthead% 
         \hiderowcolors% 
         \toprule 
         Source & Excerpt \\ 
         \midrule 
         \showrowcolors% 
         \endhead% 
         \midrule 
         \endfoot% 
         \hiderowcolors% 
         \bottomrule 
         \showrowcolors% 
         \endlastfoot% 
         
Gibler, Douglas M. 2018. International Conflicts, 1816–2010: Militarized Interstate Dispute Narratives. Vol. 1 Rowman \& Littlefield. & The Nagorno-Karabakh War was fought principally by Armenia and Azerbaijan over control of the disputed territorial enclave, with both citing the prevalence of co-ethnics to justify their claims. By December 30, 1991, thousands of Azerbaijan troops were grouped at the border of the enclave, ready to strike at Nagorno-Karabakh as ethnic fighting continued in the region. On January 27, 1992, there was open fighting between Armenian and Azeri troops, with 60 fatalities recorded in the first clashes. Initial actions in the conflict focused on Armenia trying to secure access to the enclave—Armenia was dependent on a narrow mountain pass that could be reached only by helicopter. The war reached a critical point in 1993 when Iran and Turkey both protested that its border security was compromised by the intensified fighting between Armenia and Azerbaijan. The fighting ultimately ended in a stalemate in May 1994 that was brokered by the Commonwealth of Independent States, and minor incidents followed until January 1995. The issue remains unresolved to this day (pp. 399-400)
\\
Melander, Erik. 2001. “The Nagorno-Karabakh Conflict Revisited.” Journal of Cold War Studies. & The escalation of the Nagorno-Karabakh conºict in 1988–1992 from a non-violent political struggle to a full-scale ethnic war was a momentous development in the ªnal phase of the Cold War. Some analysts have suggested that this local conºict gravely weakened the Soviet Union and thus directly contributed to the end of the Cold War. The Nagorno-Karabakh conºict has also been depicted as part of a general wave of ethnic violence that arose once the Cold War was largely over. Previous scholarly work has tended to portray the escalation of the Nagorno-Karabakh conºict as the result of a surge of nationalist sentiment unleashed under the more relaxed regime of Mikhail Gorbachev
\\
Cornell, Svante E. 1999. “The Nagorno-Karabakh Conflict.” Report No. 46, Department of East European Studies, Uppsala University. & The conflict over Nagorno-Karabakh clearly possesses an intra-state dimension, that of the struggle for independence on the part of the Armenian population of Nagorno-Karabakh. However, since the beginning of 1992 the conflict also possesses an inter-state dimension in the sense that it involved two sovereign states as belligerents: Armenia and Azerbaijan. The existence of three parties to the conflict, that is the governments of the two sovereign states as well as that of the unrecognized ‘Republic of Nagorno-Karabakh' is a factor which has made a solution to the conflict all the more difficult. The conflict has developed into one of the most intractable disputes in the international arena. By virtue of being the only one among the various Caucasian ethnopolitical conflicts that involve two internationally recognized states as parties, it is also the conflict of the region that carries the largest geopolitical significance. (...). By early 1992, the power vacuum created by the dissolution of the Soviet Union led to the loss of the last factor containing the conflict. Thus with the imminent withdrawal of the formerly Soviet forces, Karabakh became the scene of what gradually increased to a full-scale war. The Armenian side, having prepared itself to solve the conflict through military means, did not loose any time to act. From early February onwards, the Azeri villages of Malybeili, Karadagly, and Agdaban were conquered an their population evicted, leading to at least 99 civilian deaths and 140 wounded.133 After two days of artillery fire Armenian forces on 27 February, according to many impartial observers supported by the 366th CIS (formerly Soviet) motor rifle regiment, seized the small but strategically placed town of Khojaly, on the Agdam-Stepanakert road.134 This conquest was the first step in a series of atrocities to follow during the subsequent Armenian conquest of Karabakh and its surrounding areas. (p. 1, 31)
\\
Mooradian, Moorad, and Daniel Druckman. 1999. “Hurting Stalemate or Mediation? The Conflict over Nagorno-Karabakh, 1990–95.” Journal of Peace Research 36 (6): 709–727. & The conflict over Nagorno-Karabakh has a long history. In 1921, Joseph Stalin and Vladimir Lenin, working through the Caucasian Bureau of the Russian Communist Party, pacified Mustafa Kemal, the demanding leader of the Turkish Nationalist Army, by assigning the dis- puted territory of Nagorno-Karabakh to Azerbaijan. This decision angered the Armenian leaders who, having been forced into the Soviet Union, learned that Stalin had little patience for their complaints, holding out the prospect of severe recrimi- nations if actions were taken by them. The conflict thus remained latent during the Stalin era, leading to a perpetuation of the myth of brotherly cooperation between the neighboring Soviet Republics. The leadership transition from Stalin to Khrushchev in 1953 provided opportunities for the Armenian leaders to express their dissatisfaction with the status quo. Moscow was overwhelmed with protests from the Armenians and petitions from thousands of Armenians living in Nagorno-Karabakh favoring annexation to Armenia. Although the protests continued to be made to Khrushchev's successors, they fell on deaf ears. Bolstered by Gorbachev's open policies, the protests took on a strong nationalist flavor, preparing the way for an independent Armenia. In 1990, just before the Soviet Union was about to be dissolved, the conflict between Armenia, the Nagorno-Karabakh Armenians, and Azerbaijan over Nagorno-Karabakh escalated dramatically, becoming the first and arguably the most violent conflict between post-Soviet republics.
\\
\end{longtable}}
\clearpage
         \textbf{Cenepa Valley} 

          \underline{Onset:} 1995         

\underline{Reasons:}
\begin{itemize}
\item Border Clashes
\item Economic, Long-Run
\item Revenge/Retribution
\end{itemize}

         \underline{Sources:} 
         \rowcolors{3}{white}{light-gray} 
         {\scriptsize 
         \begin{longtable}{p{0.2\textwidth}p{0.8\textwidth}} 
         \hiderowcolors% 
         \toprule 
         Source & Excerpt \\ 
         \midrule 
         \showrowcolors% 
         \endfirsthead% 
         \hiderowcolors% 
         \toprule 
         Source & Excerpt \\ 
         \midrule 
         \showrowcolors% 
         \endhead% 
         \midrule 
         \endfoot% 
         \hiderowcolors% 
         \bottomrule 
         \showrowcolors% 
         \endlastfoot% 
         
Gibler, Douglas M. 2018. International Conflicts, 1816–2010: Militarized Interstate Dispute Narratives. Vol. 1. Rowman \& Littlefield. & This dispute describes the Cenepa War, another dispute over the contested area of Cordillera del Condor, more than 125,000 square miles resting between Ecuador and Peru. These states had contested Cordillera del Condor since the end of Spanish rule in South America, including militarized disputes in 1932, 1939–1942, 1981, and 1991. In December 1994 the Peruvian military discovered Ecuadorian outposts in the disputed territory. On January 9, 1995, Ecuadorian troops captured four Peruvian soldiers. By the third week of January, Peru launched air and ground strikes around Cenepa and the meeting point of the Santiago and Yaupi rivers. More than 5,000 troops were involved and hundreds were killed. On February 17, 1995, Ecuador and Peru signed the Itamaraty Peace Declaration, which called for a cessation of hostilities, demobilization of forces, and talks to resolve the border dispute. Argentina, Brazil, Chile, and the United States also agreed to establish the Military Observer Mission Ecuador-Peru (MOMEP). However, violations of the agreement occurred for nearly a month. On March 10, the parties signed the definition of procedures, and MOMEP headquarters were established within two days. By May 1995 MOMEP had completed most of its preparations for the creation of a demilitarized zone. Ecuador and Peru signed a demilitarization agreement on July 25 in Lima. The parties signed the Brasilia Agreements on October 26, 1998, which among other things settled the Cordillera del Condor border disputes between them. (pp. 144-145)
\\
The Economist. 1998. “Peace in the Andes.” The Economist, October 29, 1998. Accessed August 19, 2023. https://www.economist.com/the-americas/1998/10/29/peace-in-the-andes. & The line of the resultant frontier was mostly plain, save for some 80 kilometres (50 miles). Over this—and after decades of unhappiness in Ecuador, leading it in 1960 to denounce the Rio Protocol itself—shooting broke out again in 1981 and 1984, and in early 1995 a minor war. [...] Ecuador—though not winning the sovereign access it wanted—gets navigation rights on the Amazon and its tributaries within Peru, and can set up two trading centres there.
\\
Radcliffe, Sarah A. 1998. “Frontiers and Popular Nationhood: Geographies of Identity in the 1995 Ecuador–Peru Border Dispute.” Political Geography, 17(3): 273–293. & In late January 1995, a series of incidences occurred along one section of the international boundary between the South American countries of Ecuador and Peru. (...) [The] disputed frontier lies in an area of wood, petrol, rubber and gold resources (...). There were immedaite political and military interests involved in the eruption of the dispute at this particular time, due to declining public support for both national presidents. While the motivations of the major decision-makers in this dispute remain sketchy, these political actors were able to mobilize the population around issues of territory, nationhood and the Amazon (...). (...). According to Ecuadorean official sources, from autumn 1994 and particularlyfrom early Januray 1995, there were an increasing number of incursions by Peruvian troops into an Ecuadorean area of the undelimited zone of the Cordillera del Condor (...). (...). Following a major dispute in 1941, the United States, Brazil, Chile, and Argentina had called the warring nations together in Rio de Janeiro to sign an agreement on the location of the new frontier (...). (...). 
(...) [Survey] material, examined here in light of the 1995 dispute with Peru, found multifaceted and complex affiliations to places within and beyond the nation, in addition to national identifications, which suggests that geographies of identity are not fixed, but neither are they, as some postmodern approaches might suggest, infinitely in play and flux. 
\\
Mares, David R., and David Scott Palmer. 2012. Power, Institutions, and Leadership in War and Peace: Lessons from Peru and Ecuador, 1995–1998. New York: University of Texas Press. & In January 1995 fighting broke out between Ecuadorean and Peruvian military forces in a remote section of the Amazon that ultimately cost hundreds and perhaps even more than a thousand lives. Ecuador refused to abandon outposts constructed in territory it disputed with Peru. (...). In the specific case of Ecuador and Peru, the return and routinization of democratic practice were insufficient to overcome the long history of a festering border dispute between the two countries. Ecuador's political leaders and citizenry alike came to believe that war was a legitimate means to produce a “just” settlement. As fighting broke out in the midst of an electoral campaign for president in Peru, the leading opposition candidate, former United Nations (UN) secretary general Javier Pérez de Cuellar, advocated a much tougher response than the economic carrot and personal diplomacy approach that had been pursued up to that point by incumbent candidate Alberto Fujimori. (...). It is plausible to argue that Alberto Fujimori hijacked Peruvian politics in 1992-1993 with his autogolpe (self-coup) before the country returned to democratic forms in 1993-1994; however, his international strategy was not an aggressive military one. (pp. 1-28)

The institutional structure of the democracies created in the late 1970s in both countries was such that it produced a situation in which it was not in any politician's interest to agree on settlement terms. For Ecuador, with no reelection, legislators might have taken a moderately unpopular stance that promised important benefits in the medium term and voted for a settlement. But the economic benefi ts were likely to be large only for the lightly populated border regions rather than for the nation as a whole. In addition, with Peru insisting on the sanctity of a treaty denounced in Ecuador by democratic and authoritarian leaders alike over decades, no Ecuadorean president who cared about a place in history or legislator hoping to become president or run for other public office in the new democracy could accept a treaty under such terms. For Peru, where reelection was permitted, no legislator or president could accept a resolution adjusting any terms of the 1942 treaty without expecting punishment at the hands of the electorate, given multiple governments' historic position on the border issue. (pp. 126-136)

\\
Simmons, Beth A. 1999. “Territorial Disputes and Their Resolution: The Case of Ecuador and Peru.” United States Institute of Peace, 31(27). & [The border conflict] survived World War II, outlasted the Cold War, and most recently was the locus of military conflict between Peruvian and Ecuadoran forces in 1995. Yet the natural resources of the shared Condor Mountain range-gold, uranium, and oil-do not appear to be significant, and the territory, especially on the Peruvian side, is difficult to access and has only minimal infrastructure. (...). Until negotiations reopened in 1995, Peru had denied that a territorial dispute between the two countries existed. (...). By the early 1990s, prospects for settling the border dispute seemed to improve. In Peru, the newly elected Fujimori government faced serious economic problems, including cumulative inflation of two million percent over the five-year course of the previous Garcia administration (1985-90), widespread impoverishment, and continuing problems servicing external debt, as well as internal unrest fomented by the Maoist guerrilla group Shining Path. Beset by a number of urgent domestic problems, the Fujimori government judged that Peru was in no position to fight a border war with its neighbor. Consequently, his administration tried to improve relations with all of Peru's neighbors, including Ecuador. In 1992, President Fujimori wrote a letter to his Ecuadoran counterpart, Rodrigo Borja, proposing to complete the demarcation of two small sections of the common border in exchange for an agreement to grant navigation rights to and through the Amazon River so that Ecuador would have an outlet to the Atlantic Ocean, in accordance with Article VI of the Rio Protocol. Peruvian proposals offering the use of port facilities on the Amazon and its tributaries in return for final border demarcation were reiterated in various forms during 1992-93, but the Ecuadoran government continued in its insistence that the Rio Protocol was invalid and demanded its revision to ensure Ecuador's sovereign territorial access to the Amazon. Minor military clashes broke out near the Sabanilla River in August 1993, but the dispute remained stalemated until violent conflict erupted in early 1995.  (pp. 10-12)
\\
Marcella, Gabriel. 1995. War and Peace in the Amazon: Strategic Implications for the United States and Latin America of the 1995 Ecuador-Perú War. DIANE Publishing. & Ecuadorean foreign policy for 30 years actively pursued the nullification of the Rio Protocol, arguing further that an unjust settlement was imposed in 1941-42 by the force of a Peruvian occupation army acting in defiance of international law and of civilian control in Lima. Lately, it has advanced the concept that the Rio Protocol is not executable in the 78 kilometers. In domestic politics the Amazon issue has become a national crusade. The January 29 annual commemoration of the Rio Protocol is an emotional event for Ecuadoreans. Each January is a sensitive time along the disputed border, with occasional skirmishes between the two sides, as occurred on January 9 and 11, 1995. These were a prelude to the more serious fighting of that ensued on January 26 and in February.
\\
\end{longtable}}
\clearpage
         \textbf{Badme Border} 

          \underline{Onset:} 1998         

\underline{Reasons:}
\begin{itemize}
\item Nationalism
\item Border Clashes
\item Economic, Long-Run
\item Revenge/Retribution
\end{itemize}

         \underline{Sources:} 
         \rowcolors{3}{white}{light-gray} 
         {\scriptsize 
         \begin{longtable}{p{0.2\textwidth}p{0.8\textwidth}} 
         \hiderowcolors% 
         \toprule 
         Source & Excerpt \\ 
         \midrule 
         \showrowcolors% 
         \endfirsthead% 
         \hiderowcolors% 
         \toprule 
         Source & Excerpt \\ 
         \midrule 
         \showrowcolors% 
         \endhead% 
         \midrule 
         \endfoot% 
         \hiderowcolors% 
         \bottomrule 
         \showrowcolors% 
         \endlastfoot% 
         
Gibler, Douglas M. 2018. International Conflicts, 1816–2010: Militarized Interstate Dispute Narratives. Vol. 1. Rowman \& Littlefield. & Ethiopia and Eritrea contested several territories in a two-year war, and the most important was Badme. On May 6, 1998, Eritrean soldiers occupied Sheraro and Badme. Two weeks later Eritrea reported that Ethiopian troops occupied Sorona and Badda. Ethiopia launched air strikes against Asmara, the Eritrean capital, and Eritrean planes bombed Ethiopian towns Adigrat and Mekele. Each country spent hundreds of millions on new military equipment. After tens of thousands of deaths and hundreds of thousands of displacements, Eritrea and Ethiopia ultimately signed a ceasefire on June 18, 2000. On December 12, the two countries signed a peace agreement in Algiers that was brokered by both the United States and the Organization of African Unity. The agreement included explicit provisions for settling the border dispute, and on April 13, 2002, the Permanent Court of Arbitration at The Hague delimited most of the border. (p. 477)
\\
Britannica. 2023. “Independent Eritrea.” Encyclopædia Britannica. Accessed August 19, 2023. https://www.britannica.com/place/Eritrea/Independent-Eritrea. & In 1998 relations deteriorated rapidly when a border dispute, centred around the hamlet of Badme, exploded into violence.
\\
Phillips, Gervase. 2011. “Eritrea–Ethiopia Conflict (1998–2000).” In The Encyclopedia of War, edited by Gordon Martel. & The war of 1998–2000 waged between Eritrea and Ethiopia had a peculiar irony: the respective governments of the rival nations, the Ethiopian People's Revolutionary Democratic Front (EPRDF) of Prime Minister Meles Zenawi, and the Eritrean People's Liberation Front (EPLF) led by President Isaias Afewerki, were former allies, who had driven the dictator Mengistu Haile Mariam from Addis Abeba in 1991. Yet, even during that conflict, the condescending attitude of the veteran Eritrean fighters towards the less experienced Tigrean People's Liberation Front (who had gone on to establish the EPRDF) had often caused friction. So too did the Eritrean vision of their own future as an independent nation, confirmed by popular vote in 1993, that ran counter to Ethiopian hopes for confederation. By 1997, landlocked Ethiopia smarted over the costs now incurred in accessing the Eritrean ports of Assab and Mitsiwa. Meanwhile, in the Eritrean capital of Asmara, there was concern as the EPRDF strove to regulate the entire regional economy. The Ethiopians refused to accept parity between their currency, the birr, and the new Eritrean nakfa. There was another element to the growing tension: their border was poorly demarcated and sovereignty over much territory remained disputed. When war finally came, the United Nations (UN), the European Union (EU), and the Organization of African Unity (OAU) would all prove humiliatingly impotent as two of the poorest countries in the world squandered millions of dollars on expensive military equipment. For the Eritreans and Ethiopians themselves, the conflict had one more terrible irony: it killed tens of thousands; it displaced hundreds of thousands; it settled nothing.


\\
Bereketeab, Redie. 2010. “The Complex Roots of the Second Eritrea-Ethiopia War: Re-Examining the Causes.” African Journal of International Affairs, 13(1–3): 15–60. & [The] Eritrea-Ethiopia conflict will be found to revolve around the status of Eritrean independence. (...). The second Eritrean-Ethiopian war was depicted as one between brothers (cf. Negash and Tronvoll 2000). It was also described as a conflict between two bald men's fight for a comb, a senseless war (cf. Mengisteab 1999; Mengisteab and Yohannes 2005:231). Others described it as prestige war, border war, Eritrea's economic non-viability (Abbay 2001; Tadesse 1999). Yet, others were convinced that it has to do with Ethiopia's failure to accept Eritrean sovereignty (Fessehatzion 2002). Ethiopia's obsession with gaining an outlet to the sea, collusion of Eritrea-Tigray identities (Abbay 1998; Reid 2003), unresolved pre-liberation problematic relations of the liberation fronts (Lata 2003; Bereketeab 2009) were also given as plausible explanations. To majority of Eritreans, the second war was not a senseless war as many observers and the media have tried to present it. In the perception of many Eritreans the conflict had to do with Eritrean sovereignty. (...).
The issuance of the Eritrean currency in 1997 seemed to have shattered the thick cloud of illusions and misplaced expectations on the Ethiopian side. Ethiopia realised that the economic ‘free rider' chance given to Eritrea failed to engender its long-term political objectives of ‘political union' between both countries. Hence, the Ethiopian government unilaterally abrogated the 1993 Friendship and Cooperation Agreement (FCA). Ethiopia, in what could completely rearrange and reverse the FCA if implemented, proposed that all trade transactions between the two countries be carried out in foreign currency. Eritrea rejected the idea and proposed that the two national currencies trade on an equal basis. The two governments were unable to resolve this policy difference, which led Ethiopia to boycott Eritrean seaports and redirect its trade through Djibouti. After this a number of incidents followed that eventually triggered the outbreak of the war in May 1998.
\\
Pratt, Martin. 2006. “A Terminal Crisis? Examining the Breakdown of the Eritrea-Ethiopia Boundary Dispute Resolution Process.” Conflict Management and Peace Science, 23(4): 329-341. & there is no doubt that disagreement over the alignment of the boundary between the two countries was an important element in the conflict [...] However, as it became clear that the boundary identified by the commission placed the village of Badme inside Eritrea, satisfaction gave way to triumphalism within Eritrea and dismay within Ethiopia. Despite its tiny size and lack of any apparent strategic or economic value, Badme—the location of the spark that ignited the conflagration—had taken on great symbolic significance over the course of the war; previous research (Hensel and Mitchell, 2005) indicates that symbolically valued territory may be especially prone to violent conflict. For many people in both countries, Badme's fate became the primary indicator of whether the enormous loss of life during the fighting had been justified.
\\
\end{longtable}}
\clearpage
         \textbf{War for Kosovo} 

          \underline{Onset:} 1999         

\underline{Reasons:}
\begin{itemize}
\item Nationalism
\item Religion or Ideology
\end{itemize}

         \underline{Sources:} 
         \rowcolors{3}{white}{light-gray} 
         {\scriptsize 
         \begin{longtable}{p{0.2\textwidth}p{0.8\textwidth}} 
         \hiderowcolors% 
         \toprule 
         Source & Excerpt \\ 
         \midrule 
         \showrowcolors% 
         \endfirsthead% 
         \hiderowcolors% 
         \toprule 
         Source & Excerpt \\ 
         \midrule 
         \showrowcolors% 
         \endhead% 
         \midrule 
         \endfoot% 
         \hiderowcolors% 
         \bottomrule 
         \showrowcolors% 
         \endlastfoot% 
         
Gibler, Douglas M. 2018. International Conflicts, 1816–2010: Militarized Interstate Dispute Narratives. Vol. 1 Rowman \& Littlefield. & This dispute describes the military response by Albania and other nations to Yugoslav military activity in the Kosovo region. Albania took a particular interest in Kosovo because its population is comprised of primarily of ethnic Albanians. Yugoslavia had put an end to Kosovo's autonomy by force in 1989 and had ruled the region with a strong military and police presence thereafter. In late February 1998, protests by tens of thousands of ethnic Albanians in Kosovo turned violent, increasing tensions that could ignite broader conflict. Albania responded by fortifying their border. In the beginning of March, other countries began to consider economic sanctions and possible military action against Yugoslavia. On April 1, the UN Security Council placed an arms embargo on Yugoslavia for its actions in Kosovo. The embargo would only be lifted after the Yugoslav Prime Minister Slobodan Milosevic's government started making substantive talks about increased autonomy for Kosovo's Albanian population. Albania again fortified its border with Yugoslavia, and engaged in numerous subsequent incidents with Yugoslavia over its stance toward ethnic Albanians in Kosovo. Yugoslavia claimed that the Albanians were training, supplying, and supporting the ethnic Albanian extremists that were terrorizing the Kosovo region. Tensions began increasing dramatically as Yugoslavia threatened Albania with war due to its involvement with the ethnic Albanians fighting for autonomy in Kosovo. Although the United States was able to broker talks between Milosevic and the leader of the Kosovo Liberation Army in late May, no progress was made between the two groups and war was eminent. (pp. 209-210)
\\
Larson, Eric V., and Bogdan Savych. 1999. “Operation Allied Force (Kosovo, 1999).” In Misfortunes of War. RAND Corporation. & Operation Allied Force in Kosovo was the NATO air war to halt a Federal Republic of Yugoslavia (FRY) campaign of ethnic cleansing of Albanian Kosovars that already had killed an estimated 1,000 civilians in 1998 and resumed again in early 1999.1
\\
Ashbrook, John E. 2011. “Yugoslav Succession, Wars of (1990–1999).” In The Encyclopedia of War, edited by Gordon Martel. & Tito kept ethnic tensions between Kosovar Serbs and Albanians, who made up nearly 90 percent of Kosovo's population, relatively quiet. They boiled over shortly after his death as Serbian nationalists brought up the question of the province's autonomy. This led to Milošević 's meteoric rise to power. He purged the Kosovar Communist party and reduced the province's autonomy in 1989, sparking increased Albanian resistance. Ibrahim Rugova led a separatist movement employing non-violent means, while another, the Kosovo Liberation Army (KLA), used terrorism. Rugova kept the region relatively calm during the wars in Croatia and Bosnia, but his ineffective policies led to the growth of the KLA. It began ethnically cleansing some rural areas by 1997. In response, Serbian security forces used brutal means to crush the uprising. Thousands of Albanians abandoned their homes, pouring into hastily erected refugee camps in neighboring Albania and Macedonia. The international community, sick of the violence and perceiving the Serbs as its root cause, condemned Belgrade's actions.
\\
Parker, Geoffrey, ed. 2020. The Cambridge History of Warfare. 2nd ed. Cambridge: Cambridge University Press. & The ICTY [International Criminal Tribunal for the former Yugoslavia] also brought charges agaist Milosevic arising from another war. In 1990 he had unilaterally ended the semi-autonomy within Serbia enjoyed by Kosovo, a territory central to Orthodox Christian Serb national identity, although it had a mostly Muslim population. After six years of mainly passive resistance, a group of militant Muslims formed the Kosovo Liberation Army (KLA), dedicated to gaining independence. In 1998, in response to KLA attacks on local Serbs, Milosevic sanctioned reprisals that forced hundreds of thousands of Muslims to flee Kosovo, a process he proudly termed 'ethnic cleansing'. The threat of armed intervention by NATO led to peace discussions over the winter of 1998-9, but, although the Kosovar leaders declared themselves willing to accept autonomy within Serbia as a temporary compromise, with a later referendum on independence, Milosevic rejected the deal and ordered ethnic cleansing to resume. NATO warplanes flew some 10,000 sorties in spring 1999, targeting both Serbia's main cities and Serb units in Kosovo in the hope of forcing Belgrade to accept the peace agreement. The mission failed. (pp. 429-430)
\\
\end{longtable}}
\clearpage
         \textbf{Kargil War} 

          \underline{Onset:} 1999         

\underline{Reasons:}
\begin{itemize}
\item Nationalism
\item Power Transition or Security Dilemma
\item Border Clashes
\item Revenge/Retribution
\end{itemize}

         \underline{Sources:} 
         \rowcolors{3}{white}{light-gray} 
         {\scriptsize 
         \begin{longtable}{p{0.2\textwidth}p{0.8\textwidth}} 
         \hiderowcolors% 
         \toprule 
         Source & Excerpt \\ 
         \midrule 
         \showrowcolors% 
         \endfirsthead% 
         \hiderowcolors% 
         \toprule 
         Source & Excerpt \\ 
         \midrule 
         \showrowcolors% 
         \endhead% 
         \midrule 
         \endfoot% 
         \hiderowcolors% 
         \bottomrule 
         \showrowcolors% 
         \endlastfoot% 
         
Gibler, Douglas M. 2018. International Conflicts, 1816–2010: Militarized Interstate Dispute Narratives. Vol. 2 Rowman \& Littlefield. & On September 17, 1993, soldiers from India and Pakistan clashed, killing three. Indian officials claimed the fighting occurred as Pakistani soldiers fired on Indian soldiers in an attempt to sneak guerrillas into India. This brief clash started almost continual fighting that lasted for almost six years, causing more than one thousand deaths. The conflict ended in stalemate when Pakistani troops finally exited Kashmir (on September 9, 1993), and the majority of fighting ceased. (...).
This dispute describes border tensions between India and Pakistan beginning in July 1999 with a clash of forces in Kashmir. Both sides were involved in border and airspace violations, and both initiated troop and air clashes in a total of almost 10 militarized incidents. The conflict continued even after the October 1999 coup in Pakistan. The last incident in the conflict was a series of Indian military exercises in October 2000. (p. 880)
\\
Britannica. 2023. “Kargil War.” Accessed August 20, 2023. https://www.britannica.com/event/Kargil-War. & The sector has often been the site of border skirmishes between the two countries, and the Kargil War was the largest and deadliest of these clashes. The conflict began in early May when the Indian military learned that Pakistani fighters had infiltrated the Indian-administered territory. After detecting the infiltration, India ordered its army and air force to push back the intruders, who included regulars of the Pakistani army. The bitter fighting took place in harsh terrain 5,000 metres (16,400 feet) above sea level while intensive diplomatic activity took place elsewhere.
\\
Gill, John H. 2011. “Kargil War (1999).” In The Encyclopedia of War, edited by Gordon Martel. & The Kargil War was a limited conflict fought between India and Pakistan along the Line of Control (LOC) in Kashmir during the spring/summer of 1999. The conflict takes its name from the principal town in the combat zone on the Indian side of the LOC. Although observers debate whether Kargil constituted the “Fourth India–Pakistan War,” the conflict represented the first open fighting between India and Pakistan as declared nuclear-weapons states and had significant repercussions for both. The conflict took place in an extremely rugged and remote section of the LOC that divides the Indian and Pakistani portions of Kashmir, a region claimed by both countries and the scene of fighting in previous wars. It occurred against the backdrop of a violent separatist insurgency that had flared – with Pakistani support – in Indian Kashmir since 1989. Frequent exchanges of small arms and artillery fire across the LOC compounded the mutual suspicion generated by the insurgency and competing claims to the region. Nuclear tests carried out by both countries in May 1998 added further tension and complexity. In February 1999, however, Indian Prime Minister Atal Bihari Vajpayee visited Lahore in Pakistan for a bilateral summit to reduce tensions and lay the foundations for more normalized relations. In this atmosphere of deep mistrust and potential reconciliation, Pakistani army soldiers and paramilitary troops crossed the LOC during the spring of 1999 after the Lahore summit. (...). Pakistani motivations are also obscure, but likely included threatening a vulnerable Indian highway, reinvigorating the flagging insurgency, raising the prominence of the Kashmir issue internationally, and redeeming the army's reputation after defeat in 1971 and India's occupation of a glaciated, undemarcated portion of Kashmir in 1984. Some commentators, however, maintain that it was merely a tactical operation that suffered “mission creep” and unintentionally acquired strategic proportions.
\\
Constable, Philip. 2025. “Kashmir Dispute since 1947.” In The Encyclopedia of Diplomacy, edited by Gordon Martel. &  Within the context of ongoing Kashmiri insurgency from 1987 and testing of nuclear arsenals (India from 1974 and Pakistan in 1998), the Indian prime minister A. B. Vajpayee visited Pakistan's prime minister Nawaz Sharif in February 1999 to discuss the Kashmir issue. The resultant Lahore Declaration promised bilateral cooperation to resolve border disputes, but a few months later the Pakistani army occupied Indian positions across the line of control in the Kargil Mountains, which Indian troops vacated during winter months. Pakistani occupation of the Kargil range allowed them to block the strategic road between the Kashmir Valley and Ladakh as a bargaining position for Indian withdrawal from the Siachen Glacier and international discussion of the Kashmir issue. Pakistani occupation resulted in sustained Indian military action to eject Pakistani forces. As in the 1965 and 1971 conflicts, the PRC proved unresponsive to Pakistan's appeal for support, the United States pressured Pakistan to withdraw, and the UN refused to reopen debate on a plebiscite. The Kargil War from May to July 1999 therefore failed as did the 1965 war to internationalize Pakistani claims or debate within Kashmir on its right to self-determination.
\\
Tellis, Ashley J., C. Christine Fair, and Jamison Jo Medby. 2001. Limited Conflicts Under the Nuclear Umbrella: Indian and Pakistani Lessons from the Kargil Crisis. 1st ed. Santa Monica, CA: RAND Corporation. & Although the “real” reasons for Pakistan's prosecution of the Kargil war cannot be discerned with any certainty right now, a variety of Pakistani writings and public statements suggest that Islamabad likely had several motivations: a desire to redeem itself after its humiliating defeat in the 1971 war with India; India's occupation of the Siachen glacier; a desire to punish India for its periodic shelling of the Neelum Valley road and its other “provocations” along the line of control (LOC) in Kashmir; a desire to energize what at that point appeared to be a flagging insurgency in the Kashmir valley; and, finally, a desire to exploit its newly confirmed nuclear capabilities to achieve those lasting political changes in Kashmir that had hitherto eluded Islamabad.
\\
\end{longtable}}
\clearpage
         \textbf{Invasion of Afghanistan} 

          \underline{Onset:} 2001         

\underline{Reasons:}
\begin{itemize}
\item Power Transition or Security Dilemma
\item Religion or Ideology
\item Revenge/Retribution
\end{itemize}

         \underline{Sources:} 
         \rowcolors{3}{white}{light-gray} 
         {\scriptsize 
         \begin{longtable}{p{0.2\textwidth}p{0.8\textwidth}} 
         \hiderowcolors% 
         \toprule 
         Source & Excerpt \\ 
         \midrule 
         \showrowcolors% 
         \endfirsthead% 
         \hiderowcolors% 
         \toprule 
         Source & Excerpt \\ 
         \midrule 
         \showrowcolors% 
         \endhead% 
         \midrule 
         \endfoot% 
         \hiderowcolors% 
         \bottomrule 
         \showrowcolors% 
         \endlastfoot% 
         
Gibler, Douglas M. 2018. International Conflicts, 1816–2010: Militarized Interstate Dispute Narratives. Vol. 2 Rowman \& Littlefield. & This dispute describes American retaliation on Taliban-controlled Afghanistan following the September 11, 2001, attack on New York City and Washington, DC. The United States, Britain, and France deployed naval forces to the vicinity, with the United States tasking two aircraft carrier groups. Afghani forces went on alert in response to these moves. By October 7, 2001, American strikes on targets in Afghanistan had begun. The war that followed was quick, though the target of the war effort—Osama bin Laden and the Al-Qaeda leadership—had escaped through the mountain territories to Pakistan. American forces remain in Afghanistan today. (pp. 710-711)
\\
Britannica. 2023. “Afghanistan War.” Accessed August 19, 2023. https://www.britannica.com/event/Afghanistan-War. & Afghanistan War, international conflict in Afghanistan beginning in 2001 that was triggered by the September 11 attacks and consisted of three phases. The first phase—toppling the Taliban (the ultraconservative political and religious faction that ruled Afghanistan and provided sanctuary for al-Qaeda, perpetrators of the September 11 attacks)—was brief, lasting just two months.
\\
Miller, Seumas. 2011. “Terrorism, War against.” In The Encyclopedia of War, edited by Gordon Martel. & [The idea of a war against terrorism] has come into vogue primarily, it seems, as a consequence of the September 11, 2001 attacks on the World Trade Center in New York and the Pentagon in Washington, DC by al-Qaeda operatives. The person most famously associated with prosecuting what he called a “war against terrorism” was US President George Bush in the aftermath of 9/11 (Coady and O'Keefe 2002).
\\
Estes, Kenneth W. 2011. “Gulf Wars (1990–1991, 2003–present).” In The Encyclopedia of War, edited by G. Martel. & As extraordinary as it would have seemed years ago, nobody could have predicted that the drama of the 2001 terrorist attack on New York and Washington would have been upstaged by any other event in the presidential administration of George W. Bush. (...). The US government had already demonstrated its desire to settle direct and latent threats to security in the aftermath of the September 11, 2001, attacks on US soil. (...). Any doubts in the Bush administration on how best to put an end to the seemingly endless defense drain posed by the Saddam Hussein regime quickly gelled with the destruction of the World Trade Center and the resulting need to provide a strong response to reassert US power in the region and destroy the amorphous band that could be held responsible for the attacks. Even as the quick reprisal campaign against Afghanistan to destroy terrorist bases of operations took shape, US deployments and war planning efforts against Iraq could be discerned. The US Army V Corps headquarters, in Heidelberg, Germany, received its assignment to begin planning operations in “Southwest Asia” during the first week of November 2001.
\\
Parker, Geoffrey, ed. 2020. The Cambridge History of Warfare. 2nd ed. Cambridge: Cambridge University Press. & With the exception of Palestine, Iraq, and a few other countries that felt little sympathy for the United States, the '9/11 attacks' (...) outraged world opinion. (...). Within a week, fifty-eight countries pledged assistance (...) for a US-led invasion of Afghanistan. On September 20 2001 President Bush issued an ultimatum: The Taliban must act, and act immediately, he announced in a televised address to Congress. They will hand over the terrorists or they will share their fate. He also approved plans for air strikes and attacks on al-Qaeda and Taliban targets by Northern Alliance forces aided by a modest contingent US ground forces. Since the Taliban declined to surrender Bin Laden to President Bush, an aerial bombardment began on 7 October 2001 (...). (p. 438)
\\
\end{longtable}}
\clearpage
         \textbf{Invasion of Iraq} 

          \underline{Onset:} 2003         

\underline{Reasons:}
\begin{itemize}
\item Power Transition or Security Dilemma
\item Religion or Ideology
\item Economic, Long-Run
\end{itemize}

         \underline{Sources:} 
         \rowcolors{3}{white}{light-gray} 
         {\scriptsize 
         \begin{longtable}{p{0.2\textwidth}p{0.8\textwidth}} 
         \hiderowcolors% 
         \toprule 
         Source & Excerpt \\ 
         \midrule 
         \showrowcolors% 
         \endfirsthead% 
         \hiderowcolors% 
         \toprule 
         Source & Excerpt \\ 
         \midrule 
         \showrowcolors% 
         \endhead% 
         \midrule 
         \endfoot% 
         \hiderowcolors% 
         \bottomrule 
         \showrowcolors% 
         \endlastfoot% 
         
Gibler, Douglas M. 2018. International Conflicts, 1816–2010: Militarized Interstate Dispute Narratives. Vol. 2 Rowman \& Littlefield. & This prolonged dispute concerned various attempts by the United States and coalition allies to force Iraq's Saddam Hussein to comply with weapons inspections. The incidents mostly comprised shows of force by both sides with threats to use force made by Iraq against American and British warplanes flying patrols in their imposed no-fly zone. Saudi Arabia and Kuwait were also involved with clashes along their respective borders. Circumstances changed in 2003 when the United States demanded compliance from Saddam Hussein and began a buildup of forces, preparing for an attack. Coalition forces—mostly comprised by US military—attacked and quickly overran Iraq, while Saddam Hussein went into hiding. Occupation of Iraq followed, and American ground forces remain in the country as of 2017. (pp. 506-507)
\\
Britannica. 2023. “Iraq War.” Accessed August 19, 2023. https://www.britannica.com/event/Iraq-War. & In 2002 the new U.S. president, George W. Bush, argued that the vulnerability of the United States following the September 11 attacks of 2001, combined with Iraq's alleged continued possession and manufacture of weapons of mass destruction (an accusation that was later proved erroneous) and its support for terrorist groups—which, according to the Bush administration, included al-Qaeda, the perpetrators of the September 11 attacks—made disarming Iraq a renewed priority. UN Security Council Resolution 1441, passed on November 8, 2002, demanded that Iraq readmit inspectors and that it comply with all previous resolutions. Iraq appeared to comply with the resolution, but in early 2003 President Bush and British Prime Minister Tony Blair declared that Iraq was actually continuing to hinder UN inspections and that it still retained proscribed weapons.
\\
Estes, Kenneth W. 2011. “Gulf Wars (1990–1991, 2003–present).” In The Encyclopedia of War, edited by G. Martel. & As extraordinary as it would have seemedyears ago, nobody could have predicted thatthe drama of the 2001 terrorist attack on New York and Washington would have been upstaged by any other event in the presidential administration of George W. Bush. However, this became the case with the unilaterally initiated US invasion of Iraq, which produced vigorous debate in an already polarized diplomatic and US political landscape. The Iraq campaign of 2003 (inaugurating the Third Gulf War) nevertheless met the needs of the US political leadership to settle its issues with Iraq. During 2002, indicators continued to build signaling impending US military action against Iraq. The US government had already demonstrated its desire to settle direct and latent threats to security in the aftermath of the September 11, 2001, attacks on US soil. Thus, even before the campaign against Afghanistan reached its culminating point in the establishment of a friendly interim government there, planning continued for removing the onerous government of Iraq and removing any future threat it presented to the region and US interests. (...). (...) Paul Wolfowitz, the US deputy secretary of defense, revealed publicly that the emphasis on Iraq's rumored possession of weapons of mass destruction was brandished most vigorously because the US public would grasp it as a threat, making war justification more easily than the other possibilities. (...). The immediate objective aimed at overthrowing Saddam Hussein and installing a more friendly government that would ease tensions, isolate other opponents [Iran], and permit greater US influence in the region, including military basing rights.
\\
Gompert, David C., Hans Binnendijk, and Bonny Lin. 2014. “The U.S. Invasion of Iraq, 2003.” In Blinders, Blunders, and Wars: What America and China Can Learn, 161–174. RAND Corporation. & President George W. Bush's decision to invade Iraq on March 20, 2003, was not a blunder on the scale of those of Napoleon, Hitler, and Tojo. There was a case to be made on several grounds for operations against Saddam Hussein. The initial phase of combat was highly successful, and some still argue that the American investment was worth the cost of toppling the Saddam regime. Bush was reelected in November of 2004 as much because of as despite his invasion of Iraq. His subsequent 2007 decision to launch the “surge” did limit some of the damage. The main premise for the war was that Saddam had weapons of mass destruction (WMDs) and that these were at risk of falling into the hands of terrorists. In the end, however, there were no such weapons, and Saddam's links to al Qaeda were unproven. This robbed the invasion of legitimacy. The insurgency that ensued after initial combat operation robbed the invasion of success. Today, the United States has less influence in Baghdad than Iran does. Iraq is a Shia-dominated state with an alienated Sunni minority, rampant violence, and virtually no control over the Kurdish north. At least 134,000 Iraqis died as a direct result of the American invasion, and the violence there continues.
\\
Parker, Geoffrey, ed. 2020. The Cambridge History of Warfare. 2nd ed. Cambridge: Cambridge University Press. & In 2002, Iraqi air defenders attacked Coalition aircraft on 500 occasions, promting 90 coalition air strikes in response; patrolling the contested northern and southern no-fly zones eventually involved up to 200 aircraft on a daily basis. In addition, al-Qaeda's oft-expressed desire to acquire weapons of mass destruction focused American attention on those who already poosessed them, or were thought to have them, and fostered the desire to engage potential threats proactively. Both developments led Bush and his advisers to start planning the invasion of Iraq. (p. 440)
\\
\end{longtable}}
\clearpage
         \textbf{Invasion of Ukraine} 

          \underline{Onset:} 2022         

\underline{Reasons:}
\begin{itemize}
\item Nationalism
\item Power Transition or Security Dilemma
\end{itemize}

         \underline{Sources:} 
         \rowcolors{3}{white}{light-gray} 
         {\scriptsize 
         \begin{longtable}{p{0.2\textwidth}p{0.8\textwidth}} 
         \hiderowcolors% 
         \toprule 
         Source & Excerpt \\ 
         \midrule 
         \showrowcolors% 
         \endfirsthead% 
         \hiderowcolors% 
         \toprule 
         Source & Excerpt \\ 
         \midrule 
         \showrowcolors% 
         \endhead% 
         \midrule 
         \endfoot% 
         \hiderowcolors% 
         \bottomrule 
         \showrowcolors% 
         \endlastfoot% 
         
Atlantic Council. 2023. “Putin’s Dreams of a New Russian Empire Are Unraveling in Ukraine.”  & Putin saw the invasion of Ukraine as a key step toward rebuilding the Russian Empire.
\\
The Economist. 2022. “John Mearsheimer on Why the West Is Principally Responsible for the Ukrainian Crisis.” March 19, 2022.  & The trouble over Ukraine actually started at NATO's Bucharest summit in April 2008, when George W. Bush's administration pushed the alliance to announce that Ukraine and Georgia “will become members”. [...] These efforts eventually sparked hostilities in February 2014, after an uprising (which was supported by America) caused Ukraine's pro-Russian president, Viktor Yanukovych, to flee the country. In response, Russia took Crimea from Ukraine and helped fuel a civil war that broke out in the Donbas region of eastern Ukraine. [...] The next major confrontation came in December 2021 and led directly to the current war. The main cause was that Ukraine was becoming a de facto member of NATO. The process started in December 2017, when the Trump administration decided to sell Kyiv “defensive weapons”. What counts as “defensive” is hardly clear-cut, however, and these weapons certainly looked offensive to Moscow and its allies in the Donbas region. Other NATO countries got in on the act, shipping weapons to Ukraine, training its armed forces and allowing it to participate in joint air and naval exercises. In July 2021, Ukraine and America co-hosted a major naval exercise in the Black Sea region involving navies from 32 countries. Operation Sea Breeze almost provoked Russia to fire at a British naval destroyer that deliberately entered what Russia considers its territorial waters.
\\
\end{longtable}}
\clearpage
